\begin{python}
from scripts.calc_11 import *
\end{python}

\section{Расчёт реакции при периодическом воздействии}

Расчёт реакции в общем виде осуществляется на основе спектра воздействия и АФХ цепи $H(j\omega)$ (см. \cref{form:Ainp,form:Phiinp} для спектров; \cref{form:AChH,form:FChH} для частотных характеристик): 
\begin{equation}
    A_{\text{вых}}(k) = A_{\text{вх}}(k) \cdot A(k\omega_1);
\end{equation}
\begin{equation}
    \begin{aligned}
        A_{\text{вых}}(k) 
        &= 
        \frac{2}{T} \left|
            \frac{4I_m}{\omega} \sin^2\left(
                \frac{\omega T}{8}
            \right)
        \right|_{\omega=k\omega_{1}} \times \\
        &\times
        \left.
            \frac{
                \pv[.2]{tf['K']}
                \left|\omega^2 - \pv[.2]{tf['z1_im']}^2\right|
            }{
                \sqrt{\omega^2 + \pv[.2]{fact['alpha1']}^2}
                \sqrt{(\omega - \pv[.2]{fact['beta2']})^2 + \pv[.2]{fact['alpha2']}^2}
                \sqrt{(\omega + \pv[.2]{fact['beta2']})^2 + \pv[.2]{fact['alpha2']}^2}
            }
        \right|_{\omega=k\omega_{1}};
    \end{aligned}
\end{equation}

\begin{equation}
    \Phi_{\text{вых}}(k) = \Phi_{\text{вх}}(k) + \Phi(k\omega_1);
\end{equation}
\begin{equation}
    \begin{aligned}
        \Phi_{\text{вых}}(k)
        &= 
        \left.\left(
            90^\circ - \frac{\omega T}{4}
        \right)\right|_{\omega=k\omega_{1}} + \left[
            \arg{( \pv[.2]{tf['z1_im']}^2 - \omega^2 )} -
            \arctan{ \frac{ \omega }{ \pv[.2]{fact['alpha1']} } } 
        \right. \\
        &- \left. 
            \arctan{ \frac{ \omega + \pv[.2]{fact['beta2']} }{ \pv[.2]{fact['alpha2']} } } -
            \arctan{ \frac{ \omega - \pv[.2]{fact['beta2']} }{ \pv[.2]{fact['alpha2']} } }
        \right]_{\omega=k\omega_{1}},
    \end{aligned}
\end{equation}
где $\omega_1=2\pi/T$ — основная частота.

\subsection{Первый сигнал}
\begin{equation}
    \omega = \omega_1 k; \quad \omega_{1(1)} = \pv[.3]{ex10['w1_1']}
\end{equation}

Амплитудный и фазовый спектры реакции приведены соответственно на \cref{fig:plot_disc_spec_A_out_1,fig:plot_disc_spec_Phi_out_1}. Пунктирной линией обозначены огибающие непрерывного спектра, вертикальными линиями — дискретные гармоники.

\begin{figure}[H]
    \centering
    %% Creator: Matplotlib, PGF backend
%%
%% To include the figure in your LaTeX document, write
%%   \input{<filename>.pgf}
%%
%% Make sure the required packages are loaded in your preamble
%%   \usepackage{pgf}
%%
%% Also ensure that all the required font packages are loaded; for instance,
%% the lmodern package is sometimes necessary when using math font.
%%   \usepackage{lmodern}
%%
%% Figures using additional raster images can only be included by \input if
%% they are in the same directory as the main LaTeX file. For loading figures
%% from other directories you can use the `import` package
%%   \usepackage{import}
%%
%% and then include the figures with
%%   \import{<path to file>}{<filename>.pgf}
%%
%% Matplotlib used the following preamble
%%   \def\mathdefault#1{#1}
%%   \everymath=\expandafter{\the\everymath\displaystyle}
%%   \IfFileExists{scrextend.sty}{
%%     \usepackage[fontsize=12.000000pt]{scrextend}
%%   }{
%%     \renewcommand{\normalsize}{\fontsize{12.000000}{14.400000}\selectfont}
%%     \normalsize
%%   }
%%   \usepackage{fontspec}\setmainfont{Times New Roman}\usepackage[english,russian]{babel}\usepackage{amsmath}
%%   \ifdefined\pdftexversion\else  % non-pdftex case.
%%     \usepackage{fontspec}
%%     \setmainfont{times.ttf}[Path=\detokenize{/usr/local/share/fonts/truetype/times-new-roman/}]
%%     \setsansfont{DejaVuSans.ttf}[Path=\detokenize{/usr/share/matplotlib/mpl-data/fonts/ttf/}]
%%     \setmonofont{DejaVuSansMono.ttf}[Path=\detokenize{/usr/share/matplotlib/mpl-data/fonts/ttf/}]
%%   \fi
%%   \makeatletter\@ifpackageloaded{underscore}{}{\usepackage[strings]{underscore}}\makeatother
%%
\begingroup%
\makeatletter%
\begin{pgfpicture}%
\pgfpathrectangle{\pgfpointorigin}{\pgfqpoint{6.490000in}{3.740000in}}%
\pgfusepath{use as bounding box, clip}%
\begin{pgfscope}%
\pgfsetbuttcap%
\pgfsetmiterjoin%
\definecolor{currentfill}{rgb}{1.000000,1.000000,1.000000}%
\pgfsetfillcolor{currentfill}%
\pgfsetlinewidth{0.000000pt}%
\definecolor{currentstroke}{rgb}{1.000000,1.000000,1.000000}%
\pgfsetstrokecolor{currentstroke}%
\pgfsetdash{}{0pt}%
\pgfpathmoveto{\pgfqpoint{0.000000in}{0.000000in}}%
\pgfpathlineto{\pgfqpoint{6.490000in}{0.000000in}}%
\pgfpathlineto{\pgfqpoint{6.490000in}{3.740000in}}%
\pgfpathlineto{\pgfqpoint{0.000000in}{3.740000in}}%
\pgfpathlineto{\pgfqpoint{0.000000in}{0.000000in}}%
\pgfpathclose%
\pgfusepath{fill}%
\end{pgfscope}%
\begin{pgfscope}%
\pgfsetbuttcap%
\pgfsetmiterjoin%
\definecolor{currentfill}{rgb}{1.000000,1.000000,1.000000}%
\pgfsetfillcolor{currentfill}%
\pgfsetlinewidth{0.000000pt}%
\definecolor{currentstroke}{rgb}{0.000000,0.000000,0.000000}%
\pgfsetstrokecolor{currentstroke}%
\pgfsetstrokeopacity{0.000000}%
\pgfsetdash{}{0pt}%
\pgfpathmoveto{\pgfqpoint{0.571024in}{0.456901in}}%
\pgfpathlineto{\pgfqpoint{6.440000in}{0.456901in}}%
\pgfpathlineto{\pgfqpoint{6.440000in}{3.273333in}}%
\pgfpathlineto{\pgfqpoint{0.571024in}{3.273333in}}%
\pgfpathlineto{\pgfqpoint{0.571024in}{0.456901in}}%
\pgfpathclose%
\pgfusepath{fill}%
\end{pgfscope}%
\begin{pgfscope}%
\pgfsetbuttcap%
\pgfsetroundjoin%
\definecolor{currentfill}{rgb}{1.000000,1.000000,1.000000}%
\pgfsetfillcolor{currentfill}%
\pgfsetlinewidth{0.803000pt}%
\definecolor{currentstroke}{rgb}{0.000000,0.000000,0.000000}%
\pgfsetstrokecolor{currentstroke}%
\pgfsetdash{}{0pt}%
\pgfsys@defobject{currentmarker}{\pgfqpoint{0.000000in}{0.000000in}}{\pgfqpoint{0.000000in}{0.048611in}}{%
\pgfpathmoveto{\pgfqpoint{0.000000in}{0.000000in}}%
\pgfpathlineto{\pgfqpoint{0.000000in}{0.048611in}}%
\pgfusepath{stroke,fill}%
}%
\begin{pgfscope}%
\pgfsys@transformshift{0.719532in}{0.456901in}%
\pgfsys@useobject{currentmarker}{}%
\end{pgfscope}%
\end{pgfscope}%
\begin{pgfscope}%
\definecolor{textcolor}{rgb}{0.000000,0.000000,0.000000}%
\pgfsetstrokecolor{textcolor}%
\pgfsetfillcolor{textcolor}%
\pgftext[x=0.719532in,y=0.408290in,,top]{\color{textcolor}{\rmfamily\fontsize{12.000000}{14.400000}\selectfont\catcode`\^=\active\def^{\ifmmode\sp\else\^{}\fi}\catcode`\%=\active\def%{\%}$0$}}%
\end{pgfscope}%
\begin{pgfscope}%
\pgfsetbuttcap%
\pgfsetroundjoin%
\definecolor{currentfill}{rgb}{1.000000,1.000000,1.000000}%
\pgfsetfillcolor{currentfill}%
\pgfsetlinewidth{0.803000pt}%
\definecolor{currentstroke}{rgb}{0.000000,0.000000,0.000000}%
\pgfsetstrokecolor{currentstroke}%
\pgfsetdash{}{0pt}%
\pgfsys@defobject{currentmarker}{\pgfqpoint{0.000000in}{0.000000in}}{\pgfqpoint{0.000000in}{0.048611in}}{%
\pgfpathmoveto{\pgfqpoint{0.000000in}{0.000000in}}%
\pgfpathlineto{\pgfqpoint{0.000000in}{0.048611in}}%
\pgfusepath{stroke,fill}%
}%
\begin{pgfscope}%
\pgfsys@transformshift{1.710089in}{0.456901in}%
\pgfsys@useobject{currentmarker}{}%
\end{pgfscope}%
\end{pgfscope}%
\begin{pgfscope}%
\definecolor{textcolor}{rgb}{0.000000,0.000000,0.000000}%
\pgfsetstrokecolor{textcolor}%
\pgfsetfillcolor{textcolor}%
\pgftext[x=1.710089in,y=0.408290in,,top]{\color{textcolor}{\rmfamily\fontsize{12.000000}{14.400000}\selectfont\catcode`\^=\active\def^{\ifmmode\sp\else\^{}\fi}\catcode`\%=\active\def%{\%}$1\omega_1$}}%
\end{pgfscope}%
\begin{pgfscope}%
\pgfsetbuttcap%
\pgfsetroundjoin%
\definecolor{currentfill}{rgb}{1.000000,1.000000,1.000000}%
\pgfsetfillcolor{currentfill}%
\pgfsetlinewidth{0.803000pt}%
\definecolor{currentstroke}{rgb}{0.000000,0.000000,0.000000}%
\pgfsetstrokecolor{currentstroke}%
\pgfsetdash{}{0pt}%
\pgfsys@defobject{currentmarker}{\pgfqpoint{0.000000in}{0.000000in}}{\pgfqpoint{0.000000in}{0.048611in}}{%
\pgfpathmoveto{\pgfqpoint{0.000000in}{0.000000in}}%
\pgfpathlineto{\pgfqpoint{0.000000in}{0.048611in}}%
\pgfusepath{stroke,fill}%
}%
\begin{pgfscope}%
\pgfsys@transformshift{2.700647in}{0.456901in}%
\pgfsys@useobject{currentmarker}{}%
\end{pgfscope}%
\end{pgfscope}%
\begin{pgfscope}%
\definecolor{textcolor}{rgb}{0.000000,0.000000,0.000000}%
\pgfsetstrokecolor{textcolor}%
\pgfsetfillcolor{textcolor}%
\pgftext[x=2.700647in,y=0.408290in,,top]{\color{textcolor}{\rmfamily\fontsize{12.000000}{14.400000}\selectfont\catcode`\^=\active\def^{\ifmmode\sp\else\^{}\fi}\catcode`\%=\active\def%{\%}$2\omega_1$}}%
\end{pgfscope}%
\begin{pgfscope}%
\pgfsetbuttcap%
\pgfsetroundjoin%
\definecolor{currentfill}{rgb}{1.000000,1.000000,1.000000}%
\pgfsetfillcolor{currentfill}%
\pgfsetlinewidth{0.803000pt}%
\definecolor{currentstroke}{rgb}{0.000000,0.000000,0.000000}%
\pgfsetstrokecolor{currentstroke}%
\pgfsetdash{}{0pt}%
\pgfsys@defobject{currentmarker}{\pgfqpoint{0.000000in}{0.000000in}}{\pgfqpoint{0.000000in}{0.048611in}}{%
\pgfpathmoveto{\pgfqpoint{0.000000in}{0.000000in}}%
\pgfpathlineto{\pgfqpoint{0.000000in}{0.048611in}}%
\pgfusepath{stroke,fill}%
}%
\begin{pgfscope}%
\pgfsys@transformshift{3.691204in}{0.456901in}%
\pgfsys@useobject{currentmarker}{}%
\end{pgfscope}%
\end{pgfscope}%
\begin{pgfscope}%
\definecolor{textcolor}{rgb}{0.000000,0.000000,0.000000}%
\pgfsetstrokecolor{textcolor}%
\pgfsetfillcolor{textcolor}%
\pgftext[x=3.691204in,y=0.408290in,,top]{\color{textcolor}{\rmfamily\fontsize{12.000000}{14.400000}\selectfont\catcode`\^=\active\def^{\ifmmode\sp\else\^{}\fi}\catcode`\%=\active\def%{\%}$3\omega_1$}}%
\end{pgfscope}%
\begin{pgfscope}%
\pgfsetbuttcap%
\pgfsetroundjoin%
\definecolor{currentfill}{rgb}{1.000000,1.000000,1.000000}%
\pgfsetfillcolor{currentfill}%
\pgfsetlinewidth{0.803000pt}%
\definecolor{currentstroke}{rgb}{0.000000,0.000000,0.000000}%
\pgfsetstrokecolor{currentstroke}%
\pgfsetdash{}{0pt}%
\pgfsys@defobject{currentmarker}{\pgfqpoint{0.000000in}{0.000000in}}{\pgfqpoint{0.000000in}{0.048611in}}{%
\pgfpathmoveto{\pgfqpoint{0.000000in}{0.000000in}}%
\pgfpathlineto{\pgfqpoint{0.000000in}{0.048611in}}%
\pgfusepath{stroke,fill}%
}%
\begin{pgfscope}%
\pgfsys@transformshift{4.681761in}{0.456901in}%
\pgfsys@useobject{currentmarker}{}%
\end{pgfscope}%
\end{pgfscope}%
\begin{pgfscope}%
\definecolor{textcolor}{rgb}{0.000000,0.000000,0.000000}%
\pgfsetstrokecolor{textcolor}%
\pgfsetfillcolor{textcolor}%
\pgftext[x=4.681761in,y=0.408290in,,top]{\color{textcolor}{\rmfamily\fontsize{12.000000}{14.400000}\selectfont\catcode`\^=\active\def^{\ifmmode\sp\else\^{}\fi}\catcode`\%=\active\def%{\%}$4\omega_1$}}%
\end{pgfscope}%
\begin{pgfscope}%
\pgfsetbuttcap%
\pgfsetroundjoin%
\definecolor{currentfill}{rgb}{1.000000,1.000000,1.000000}%
\pgfsetfillcolor{currentfill}%
\pgfsetlinewidth{0.803000pt}%
\definecolor{currentstroke}{rgb}{0.000000,0.000000,0.000000}%
\pgfsetstrokecolor{currentstroke}%
\pgfsetdash{}{0pt}%
\pgfsys@defobject{currentmarker}{\pgfqpoint{0.000000in}{0.000000in}}{\pgfqpoint{0.000000in}{0.048611in}}{%
\pgfpathmoveto{\pgfqpoint{0.000000in}{0.000000in}}%
\pgfpathlineto{\pgfqpoint{0.000000in}{0.048611in}}%
\pgfusepath{stroke,fill}%
}%
\begin{pgfscope}%
\pgfsys@transformshift{5.672318in}{0.456901in}%
\pgfsys@useobject{currentmarker}{}%
\end{pgfscope}%
\end{pgfscope}%
\begin{pgfscope}%
\definecolor{textcolor}{rgb}{0.000000,0.000000,0.000000}%
\pgfsetstrokecolor{textcolor}%
\pgfsetfillcolor{textcolor}%
\pgftext[x=5.672318in,y=0.408290in,,top]{\color{textcolor}{\rmfamily\fontsize{12.000000}{14.400000}\selectfont\catcode`\^=\active\def^{\ifmmode\sp\else\^{}\fi}\catcode`\%=\active\def%{\%}$5\omega_1$}}%
\end{pgfscope}%
\begin{pgfscope}%
\pgfpathrectangle{\pgfqpoint{0.571024in}{0.456901in}}{\pgfqpoint{5.868976in}{2.816432in}}%
\pgfusepath{clip}%
\pgfsetbuttcap%
\pgfsetroundjoin%
\pgfsetlinewidth{0.501875pt}%
\definecolor{currentstroke}{rgb}{0.501961,0.501961,0.501961}%
\pgfsetstrokecolor{currentstroke}%
\pgfsetstrokeopacity{0.400000}%
\pgfsetdash{{1.850000pt}{0.800000pt}}{0.000000pt}%
\pgfpathmoveto{\pgfqpoint{0.967172in}{0.456901in}}%
\pgfpathlineto{\pgfqpoint{0.967172in}{3.273333in}}%
\pgfusepath{stroke}%
\end{pgfscope}%
\begin{pgfscope}%
\pgfsetbuttcap%
\pgfsetroundjoin%
\definecolor{currentfill}{rgb}{1.000000,1.000000,1.000000}%
\pgfsetfillcolor{currentfill}%
\pgfsetlinewidth{0.602250pt}%
\definecolor{currentstroke}{rgb}{0.000000,0.000000,0.000000}%
\pgfsetstrokecolor{currentstroke}%
\pgfsetdash{}{0pt}%
\pgfsys@defobject{currentmarker}{\pgfqpoint{0.000000in}{0.000000in}}{\pgfqpoint{0.000000in}{0.027778in}}{%
\pgfpathmoveto{\pgfqpoint{0.000000in}{0.000000in}}%
\pgfpathlineto{\pgfqpoint{0.000000in}{0.027778in}}%
\pgfusepath{stroke,fill}%
}%
\begin{pgfscope}%
\pgfsys@transformshift{0.967172in}{0.456901in}%
\pgfsys@useobject{currentmarker}{}%
\end{pgfscope}%
\end{pgfscope}%
\begin{pgfscope}%
\pgfpathrectangle{\pgfqpoint{0.571024in}{0.456901in}}{\pgfqpoint{5.868976in}{2.816432in}}%
\pgfusepath{clip}%
\pgfsetbuttcap%
\pgfsetroundjoin%
\pgfsetlinewidth{0.501875pt}%
\definecolor{currentstroke}{rgb}{0.501961,0.501961,0.501961}%
\pgfsetstrokecolor{currentstroke}%
\pgfsetstrokeopacity{0.400000}%
\pgfsetdash{{1.850000pt}{0.800000pt}}{0.000000pt}%
\pgfpathmoveto{\pgfqpoint{1.214811in}{0.456901in}}%
\pgfpathlineto{\pgfqpoint{1.214811in}{3.273333in}}%
\pgfusepath{stroke}%
\end{pgfscope}%
\begin{pgfscope}%
\pgfsetbuttcap%
\pgfsetroundjoin%
\definecolor{currentfill}{rgb}{1.000000,1.000000,1.000000}%
\pgfsetfillcolor{currentfill}%
\pgfsetlinewidth{0.602250pt}%
\definecolor{currentstroke}{rgb}{0.000000,0.000000,0.000000}%
\pgfsetstrokecolor{currentstroke}%
\pgfsetdash{}{0pt}%
\pgfsys@defobject{currentmarker}{\pgfqpoint{0.000000in}{0.000000in}}{\pgfqpoint{0.000000in}{0.027778in}}{%
\pgfpathmoveto{\pgfqpoint{0.000000in}{0.000000in}}%
\pgfpathlineto{\pgfqpoint{0.000000in}{0.027778in}}%
\pgfusepath{stroke,fill}%
}%
\begin{pgfscope}%
\pgfsys@transformshift{1.214811in}{0.456901in}%
\pgfsys@useobject{currentmarker}{}%
\end{pgfscope}%
\end{pgfscope}%
\begin{pgfscope}%
\pgfpathrectangle{\pgfqpoint{0.571024in}{0.456901in}}{\pgfqpoint{5.868976in}{2.816432in}}%
\pgfusepath{clip}%
\pgfsetbuttcap%
\pgfsetroundjoin%
\pgfsetlinewidth{0.501875pt}%
\definecolor{currentstroke}{rgb}{0.501961,0.501961,0.501961}%
\pgfsetstrokecolor{currentstroke}%
\pgfsetstrokeopacity{0.400000}%
\pgfsetdash{{1.850000pt}{0.800000pt}}{0.000000pt}%
\pgfpathmoveto{\pgfqpoint{1.462450in}{0.456901in}}%
\pgfpathlineto{\pgfqpoint{1.462450in}{3.273333in}}%
\pgfusepath{stroke}%
\end{pgfscope}%
\begin{pgfscope}%
\pgfsetbuttcap%
\pgfsetroundjoin%
\definecolor{currentfill}{rgb}{1.000000,1.000000,1.000000}%
\pgfsetfillcolor{currentfill}%
\pgfsetlinewidth{0.602250pt}%
\definecolor{currentstroke}{rgb}{0.000000,0.000000,0.000000}%
\pgfsetstrokecolor{currentstroke}%
\pgfsetdash{}{0pt}%
\pgfsys@defobject{currentmarker}{\pgfqpoint{0.000000in}{0.000000in}}{\pgfqpoint{0.000000in}{0.027778in}}{%
\pgfpathmoveto{\pgfqpoint{0.000000in}{0.000000in}}%
\pgfpathlineto{\pgfqpoint{0.000000in}{0.027778in}}%
\pgfusepath{stroke,fill}%
}%
\begin{pgfscope}%
\pgfsys@transformshift{1.462450in}{0.456901in}%
\pgfsys@useobject{currentmarker}{}%
\end{pgfscope}%
\end{pgfscope}%
\begin{pgfscope}%
\pgfpathrectangle{\pgfqpoint{0.571024in}{0.456901in}}{\pgfqpoint{5.868976in}{2.816432in}}%
\pgfusepath{clip}%
\pgfsetbuttcap%
\pgfsetroundjoin%
\pgfsetlinewidth{0.501875pt}%
\definecolor{currentstroke}{rgb}{0.501961,0.501961,0.501961}%
\pgfsetstrokecolor{currentstroke}%
\pgfsetstrokeopacity{0.400000}%
\pgfsetdash{{1.850000pt}{0.800000pt}}{0.000000pt}%
\pgfpathmoveto{\pgfqpoint{1.957729in}{0.456901in}}%
\pgfpathlineto{\pgfqpoint{1.957729in}{3.273333in}}%
\pgfusepath{stroke}%
\end{pgfscope}%
\begin{pgfscope}%
\pgfsetbuttcap%
\pgfsetroundjoin%
\definecolor{currentfill}{rgb}{1.000000,1.000000,1.000000}%
\pgfsetfillcolor{currentfill}%
\pgfsetlinewidth{0.602250pt}%
\definecolor{currentstroke}{rgb}{0.000000,0.000000,0.000000}%
\pgfsetstrokecolor{currentstroke}%
\pgfsetdash{}{0pt}%
\pgfsys@defobject{currentmarker}{\pgfqpoint{0.000000in}{0.000000in}}{\pgfqpoint{0.000000in}{0.027778in}}{%
\pgfpathmoveto{\pgfqpoint{0.000000in}{0.000000in}}%
\pgfpathlineto{\pgfqpoint{0.000000in}{0.027778in}}%
\pgfusepath{stroke,fill}%
}%
\begin{pgfscope}%
\pgfsys@transformshift{1.957729in}{0.456901in}%
\pgfsys@useobject{currentmarker}{}%
\end{pgfscope}%
\end{pgfscope}%
\begin{pgfscope}%
\pgfpathrectangle{\pgfqpoint{0.571024in}{0.456901in}}{\pgfqpoint{5.868976in}{2.816432in}}%
\pgfusepath{clip}%
\pgfsetbuttcap%
\pgfsetroundjoin%
\pgfsetlinewidth{0.501875pt}%
\definecolor{currentstroke}{rgb}{0.501961,0.501961,0.501961}%
\pgfsetstrokecolor{currentstroke}%
\pgfsetstrokeopacity{0.400000}%
\pgfsetdash{{1.850000pt}{0.800000pt}}{0.000000pt}%
\pgfpathmoveto{\pgfqpoint{2.205368in}{0.456901in}}%
\pgfpathlineto{\pgfqpoint{2.205368in}{3.273333in}}%
\pgfusepath{stroke}%
\end{pgfscope}%
\begin{pgfscope}%
\pgfsetbuttcap%
\pgfsetroundjoin%
\definecolor{currentfill}{rgb}{1.000000,1.000000,1.000000}%
\pgfsetfillcolor{currentfill}%
\pgfsetlinewidth{0.602250pt}%
\definecolor{currentstroke}{rgb}{0.000000,0.000000,0.000000}%
\pgfsetstrokecolor{currentstroke}%
\pgfsetdash{}{0pt}%
\pgfsys@defobject{currentmarker}{\pgfqpoint{0.000000in}{0.000000in}}{\pgfqpoint{0.000000in}{0.027778in}}{%
\pgfpathmoveto{\pgfqpoint{0.000000in}{0.000000in}}%
\pgfpathlineto{\pgfqpoint{0.000000in}{0.027778in}}%
\pgfusepath{stroke,fill}%
}%
\begin{pgfscope}%
\pgfsys@transformshift{2.205368in}{0.456901in}%
\pgfsys@useobject{currentmarker}{}%
\end{pgfscope}%
\end{pgfscope}%
\begin{pgfscope}%
\pgfpathrectangle{\pgfqpoint{0.571024in}{0.456901in}}{\pgfqpoint{5.868976in}{2.816432in}}%
\pgfusepath{clip}%
\pgfsetbuttcap%
\pgfsetroundjoin%
\pgfsetlinewidth{0.501875pt}%
\definecolor{currentstroke}{rgb}{0.501961,0.501961,0.501961}%
\pgfsetstrokecolor{currentstroke}%
\pgfsetstrokeopacity{0.400000}%
\pgfsetdash{{1.850000pt}{0.800000pt}}{0.000000pt}%
\pgfpathmoveto{\pgfqpoint{2.453007in}{0.456901in}}%
\pgfpathlineto{\pgfqpoint{2.453007in}{3.273333in}}%
\pgfusepath{stroke}%
\end{pgfscope}%
\begin{pgfscope}%
\pgfsetbuttcap%
\pgfsetroundjoin%
\definecolor{currentfill}{rgb}{1.000000,1.000000,1.000000}%
\pgfsetfillcolor{currentfill}%
\pgfsetlinewidth{0.602250pt}%
\definecolor{currentstroke}{rgb}{0.000000,0.000000,0.000000}%
\pgfsetstrokecolor{currentstroke}%
\pgfsetdash{}{0pt}%
\pgfsys@defobject{currentmarker}{\pgfqpoint{0.000000in}{0.000000in}}{\pgfqpoint{0.000000in}{0.027778in}}{%
\pgfpathmoveto{\pgfqpoint{0.000000in}{0.000000in}}%
\pgfpathlineto{\pgfqpoint{0.000000in}{0.027778in}}%
\pgfusepath{stroke,fill}%
}%
\begin{pgfscope}%
\pgfsys@transformshift{2.453007in}{0.456901in}%
\pgfsys@useobject{currentmarker}{}%
\end{pgfscope}%
\end{pgfscope}%
\begin{pgfscope}%
\pgfpathrectangle{\pgfqpoint{0.571024in}{0.456901in}}{\pgfqpoint{5.868976in}{2.816432in}}%
\pgfusepath{clip}%
\pgfsetbuttcap%
\pgfsetroundjoin%
\pgfsetlinewidth{0.501875pt}%
\definecolor{currentstroke}{rgb}{0.501961,0.501961,0.501961}%
\pgfsetstrokecolor{currentstroke}%
\pgfsetstrokeopacity{0.400000}%
\pgfsetdash{{1.850000pt}{0.800000pt}}{0.000000pt}%
\pgfpathmoveto{\pgfqpoint{2.948286in}{0.456901in}}%
\pgfpathlineto{\pgfqpoint{2.948286in}{3.273333in}}%
\pgfusepath{stroke}%
\end{pgfscope}%
\begin{pgfscope}%
\pgfsetbuttcap%
\pgfsetroundjoin%
\definecolor{currentfill}{rgb}{1.000000,1.000000,1.000000}%
\pgfsetfillcolor{currentfill}%
\pgfsetlinewidth{0.602250pt}%
\definecolor{currentstroke}{rgb}{0.000000,0.000000,0.000000}%
\pgfsetstrokecolor{currentstroke}%
\pgfsetdash{}{0pt}%
\pgfsys@defobject{currentmarker}{\pgfqpoint{0.000000in}{0.000000in}}{\pgfqpoint{0.000000in}{0.027778in}}{%
\pgfpathmoveto{\pgfqpoint{0.000000in}{0.000000in}}%
\pgfpathlineto{\pgfqpoint{0.000000in}{0.027778in}}%
\pgfusepath{stroke,fill}%
}%
\begin{pgfscope}%
\pgfsys@transformshift{2.948286in}{0.456901in}%
\pgfsys@useobject{currentmarker}{}%
\end{pgfscope}%
\end{pgfscope}%
\begin{pgfscope}%
\pgfpathrectangle{\pgfqpoint{0.571024in}{0.456901in}}{\pgfqpoint{5.868976in}{2.816432in}}%
\pgfusepath{clip}%
\pgfsetbuttcap%
\pgfsetroundjoin%
\pgfsetlinewidth{0.501875pt}%
\definecolor{currentstroke}{rgb}{0.501961,0.501961,0.501961}%
\pgfsetstrokecolor{currentstroke}%
\pgfsetstrokeopacity{0.400000}%
\pgfsetdash{{1.850000pt}{0.800000pt}}{0.000000pt}%
\pgfpathmoveto{\pgfqpoint{3.195925in}{0.456901in}}%
\pgfpathlineto{\pgfqpoint{3.195925in}{3.273333in}}%
\pgfusepath{stroke}%
\end{pgfscope}%
\begin{pgfscope}%
\pgfsetbuttcap%
\pgfsetroundjoin%
\definecolor{currentfill}{rgb}{1.000000,1.000000,1.000000}%
\pgfsetfillcolor{currentfill}%
\pgfsetlinewidth{0.602250pt}%
\definecolor{currentstroke}{rgb}{0.000000,0.000000,0.000000}%
\pgfsetstrokecolor{currentstroke}%
\pgfsetdash{}{0pt}%
\pgfsys@defobject{currentmarker}{\pgfqpoint{0.000000in}{0.000000in}}{\pgfqpoint{0.000000in}{0.027778in}}{%
\pgfpathmoveto{\pgfqpoint{0.000000in}{0.000000in}}%
\pgfpathlineto{\pgfqpoint{0.000000in}{0.027778in}}%
\pgfusepath{stroke,fill}%
}%
\begin{pgfscope}%
\pgfsys@transformshift{3.195925in}{0.456901in}%
\pgfsys@useobject{currentmarker}{}%
\end{pgfscope}%
\end{pgfscope}%
\begin{pgfscope}%
\pgfpathrectangle{\pgfqpoint{0.571024in}{0.456901in}}{\pgfqpoint{5.868976in}{2.816432in}}%
\pgfusepath{clip}%
\pgfsetbuttcap%
\pgfsetroundjoin%
\pgfsetlinewidth{0.501875pt}%
\definecolor{currentstroke}{rgb}{0.501961,0.501961,0.501961}%
\pgfsetstrokecolor{currentstroke}%
\pgfsetstrokeopacity{0.400000}%
\pgfsetdash{{1.850000pt}{0.800000pt}}{0.000000pt}%
\pgfpathmoveto{\pgfqpoint{3.443565in}{0.456901in}}%
\pgfpathlineto{\pgfqpoint{3.443565in}{3.273333in}}%
\pgfusepath{stroke}%
\end{pgfscope}%
\begin{pgfscope}%
\pgfsetbuttcap%
\pgfsetroundjoin%
\definecolor{currentfill}{rgb}{1.000000,1.000000,1.000000}%
\pgfsetfillcolor{currentfill}%
\pgfsetlinewidth{0.602250pt}%
\definecolor{currentstroke}{rgb}{0.000000,0.000000,0.000000}%
\pgfsetstrokecolor{currentstroke}%
\pgfsetdash{}{0pt}%
\pgfsys@defobject{currentmarker}{\pgfqpoint{0.000000in}{0.000000in}}{\pgfqpoint{0.000000in}{0.027778in}}{%
\pgfpathmoveto{\pgfqpoint{0.000000in}{0.000000in}}%
\pgfpathlineto{\pgfqpoint{0.000000in}{0.027778in}}%
\pgfusepath{stroke,fill}%
}%
\begin{pgfscope}%
\pgfsys@transformshift{3.443565in}{0.456901in}%
\pgfsys@useobject{currentmarker}{}%
\end{pgfscope}%
\end{pgfscope}%
\begin{pgfscope}%
\pgfpathrectangle{\pgfqpoint{0.571024in}{0.456901in}}{\pgfqpoint{5.868976in}{2.816432in}}%
\pgfusepath{clip}%
\pgfsetbuttcap%
\pgfsetroundjoin%
\pgfsetlinewidth{0.501875pt}%
\definecolor{currentstroke}{rgb}{0.501961,0.501961,0.501961}%
\pgfsetstrokecolor{currentstroke}%
\pgfsetstrokeopacity{0.400000}%
\pgfsetdash{{1.850000pt}{0.800000pt}}{0.000000pt}%
\pgfpathmoveto{\pgfqpoint{3.938843in}{0.456901in}}%
\pgfpathlineto{\pgfqpoint{3.938843in}{3.273333in}}%
\pgfusepath{stroke}%
\end{pgfscope}%
\begin{pgfscope}%
\pgfsetbuttcap%
\pgfsetroundjoin%
\definecolor{currentfill}{rgb}{1.000000,1.000000,1.000000}%
\pgfsetfillcolor{currentfill}%
\pgfsetlinewidth{0.602250pt}%
\definecolor{currentstroke}{rgb}{0.000000,0.000000,0.000000}%
\pgfsetstrokecolor{currentstroke}%
\pgfsetdash{}{0pt}%
\pgfsys@defobject{currentmarker}{\pgfqpoint{0.000000in}{0.000000in}}{\pgfqpoint{0.000000in}{0.027778in}}{%
\pgfpathmoveto{\pgfqpoint{0.000000in}{0.000000in}}%
\pgfpathlineto{\pgfqpoint{0.000000in}{0.027778in}}%
\pgfusepath{stroke,fill}%
}%
\begin{pgfscope}%
\pgfsys@transformshift{3.938843in}{0.456901in}%
\pgfsys@useobject{currentmarker}{}%
\end{pgfscope}%
\end{pgfscope}%
\begin{pgfscope}%
\pgfpathrectangle{\pgfqpoint{0.571024in}{0.456901in}}{\pgfqpoint{5.868976in}{2.816432in}}%
\pgfusepath{clip}%
\pgfsetbuttcap%
\pgfsetroundjoin%
\pgfsetlinewidth{0.501875pt}%
\definecolor{currentstroke}{rgb}{0.501961,0.501961,0.501961}%
\pgfsetstrokecolor{currentstroke}%
\pgfsetstrokeopacity{0.400000}%
\pgfsetdash{{1.850000pt}{0.800000pt}}{0.000000pt}%
\pgfpathmoveto{\pgfqpoint{4.186482in}{0.456901in}}%
\pgfpathlineto{\pgfqpoint{4.186482in}{3.273333in}}%
\pgfusepath{stroke}%
\end{pgfscope}%
\begin{pgfscope}%
\pgfsetbuttcap%
\pgfsetroundjoin%
\definecolor{currentfill}{rgb}{1.000000,1.000000,1.000000}%
\pgfsetfillcolor{currentfill}%
\pgfsetlinewidth{0.602250pt}%
\definecolor{currentstroke}{rgb}{0.000000,0.000000,0.000000}%
\pgfsetstrokecolor{currentstroke}%
\pgfsetdash{}{0pt}%
\pgfsys@defobject{currentmarker}{\pgfqpoint{0.000000in}{0.000000in}}{\pgfqpoint{0.000000in}{0.027778in}}{%
\pgfpathmoveto{\pgfqpoint{0.000000in}{0.000000in}}%
\pgfpathlineto{\pgfqpoint{0.000000in}{0.027778in}}%
\pgfusepath{stroke,fill}%
}%
\begin{pgfscope}%
\pgfsys@transformshift{4.186482in}{0.456901in}%
\pgfsys@useobject{currentmarker}{}%
\end{pgfscope}%
\end{pgfscope}%
\begin{pgfscope}%
\pgfpathrectangle{\pgfqpoint{0.571024in}{0.456901in}}{\pgfqpoint{5.868976in}{2.816432in}}%
\pgfusepath{clip}%
\pgfsetbuttcap%
\pgfsetroundjoin%
\pgfsetlinewidth{0.501875pt}%
\definecolor{currentstroke}{rgb}{0.501961,0.501961,0.501961}%
\pgfsetstrokecolor{currentstroke}%
\pgfsetstrokeopacity{0.400000}%
\pgfsetdash{{1.850000pt}{0.800000pt}}{0.000000pt}%
\pgfpathmoveto{\pgfqpoint{4.434122in}{0.456901in}}%
\pgfpathlineto{\pgfqpoint{4.434122in}{3.273333in}}%
\pgfusepath{stroke}%
\end{pgfscope}%
\begin{pgfscope}%
\pgfsetbuttcap%
\pgfsetroundjoin%
\definecolor{currentfill}{rgb}{1.000000,1.000000,1.000000}%
\pgfsetfillcolor{currentfill}%
\pgfsetlinewidth{0.602250pt}%
\definecolor{currentstroke}{rgb}{0.000000,0.000000,0.000000}%
\pgfsetstrokecolor{currentstroke}%
\pgfsetdash{}{0pt}%
\pgfsys@defobject{currentmarker}{\pgfqpoint{0.000000in}{0.000000in}}{\pgfqpoint{0.000000in}{0.027778in}}{%
\pgfpathmoveto{\pgfqpoint{0.000000in}{0.000000in}}%
\pgfpathlineto{\pgfqpoint{0.000000in}{0.027778in}}%
\pgfusepath{stroke,fill}%
}%
\begin{pgfscope}%
\pgfsys@transformshift{4.434122in}{0.456901in}%
\pgfsys@useobject{currentmarker}{}%
\end{pgfscope}%
\end{pgfscope}%
\begin{pgfscope}%
\pgfpathrectangle{\pgfqpoint{0.571024in}{0.456901in}}{\pgfqpoint{5.868976in}{2.816432in}}%
\pgfusepath{clip}%
\pgfsetbuttcap%
\pgfsetroundjoin%
\pgfsetlinewidth{0.501875pt}%
\definecolor{currentstroke}{rgb}{0.501961,0.501961,0.501961}%
\pgfsetstrokecolor{currentstroke}%
\pgfsetstrokeopacity{0.400000}%
\pgfsetdash{{1.850000pt}{0.800000pt}}{0.000000pt}%
\pgfpathmoveto{\pgfqpoint{4.929400in}{0.456901in}}%
\pgfpathlineto{\pgfqpoint{4.929400in}{3.273333in}}%
\pgfusepath{stroke}%
\end{pgfscope}%
\begin{pgfscope}%
\pgfsetbuttcap%
\pgfsetroundjoin%
\definecolor{currentfill}{rgb}{1.000000,1.000000,1.000000}%
\pgfsetfillcolor{currentfill}%
\pgfsetlinewidth{0.602250pt}%
\definecolor{currentstroke}{rgb}{0.000000,0.000000,0.000000}%
\pgfsetstrokecolor{currentstroke}%
\pgfsetdash{}{0pt}%
\pgfsys@defobject{currentmarker}{\pgfqpoint{0.000000in}{0.000000in}}{\pgfqpoint{0.000000in}{0.027778in}}{%
\pgfpathmoveto{\pgfqpoint{0.000000in}{0.000000in}}%
\pgfpathlineto{\pgfqpoint{0.000000in}{0.027778in}}%
\pgfusepath{stroke,fill}%
}%
\begin{pgfscope}%
\pgfsys@transformshift{4.929400in}{0.456901in}%
\pgfsys@useobject{currentmarker}{}%
\end{pgfscope}%
\end{pgfscope}%
\begin{pgfscope}%
\pgfpathrectangle{\pgfqpoint{0.571024in}{0.456901in}}{\pgfqpoint{5.868976in}{2.816432in}}%
\pgfusepath{clip}%
\pgfsetbuttcap%
\pgfsetroundjoin%
\pgfsetlinewidth{0.501875pt}%
\definecolor{currentstroke}{rgb}{0.501961,0.501961,0.501961}%
\pgfsetstrokecolor{currentstroke}%
\pgfsetstrokeopacity{0.400000}%
\pgfsetdash{{1.850000pt}{0.800000pt}}{0.000000pt}%
\pgfpathmoveto{\pgfqpoint{5.177040in}{0.456901in}}%
\pgfpathlineto{\pgfqpoint{5.177040in}{3.273333in}}%
\pgfusepath{stroke}%
\end{pgfscope}%
\begin{pgfscope}%
\pgfsetbuttcap%
\pgfsetroundjoin%
\definecolor{currentfill}{rgb}{1.000000,1.000000,1.000000}%
\pgfsetfillcolor{currentfill}%
\pgfsetlinewidth{0.602250pt}%
\definecolor{currentstroke}{rgb}{0.000000,0.000000,0.000000}%
\pgfsetstrokecolor{currentstroke}%
\pgfsetdash{}{0pt}%
\pgfsys@defobject{currentmarker}{\pgfqpoint{0.000000in}{0.000000in}}{\pgfqpoint{0.000000in}{0.027778in}}{%
\pgfpathmoveto{\pgfqpoint{0.000000in}{0.000000in}}%
\pgfpathlineto{\pgfqpoint{0.000000in}{0.027778in}}%
\pgfusepath{stroke,fill}%
}%
\begin{pgfscope}%
\pgfsys@transformshift{5.177040in}{0.456901in}%
\pgfsys@useobject{currentmarker}{}%
\end{pgfscope}%
\end{pgfscope}%
\begin{pgfscope}%
\pgfpathrectangle{\pgfqpoint{0.571024in}{0.456901in}}{\pgfqpoint{5.868976in}{2.816432in}}%
\pgfusepath{clip}%
\pgfsetbuttcap%
\pgfsetroundjoin%
\pgfsetlinewidth{0.501875pt}%
\definecolor{currentstroke}{rgb}{0.501961,0.501961,0.501961}%
\pgfsetstrokecolor{currentstroke}%
\pgfsetstrokeopacity{0.400000}%
\pgfsetdash{{1.850000pt}{0.800000pt}}{0.000000pt}%
\pgfpathmoveto{\pgfqpoint{5.424679in}{0.456901in}}%
\pgfpathlineto{\pgfqpoint{5.424679in}{3.273333in}}%
\pgfusepath{stroke}%
\end{pgfscope}%
\begin{pgfscope}%
\pgfsetbuttcap%
\pgfsetroundjoin%
\definecolor{currentfill}{rgb}{1.000000,1.000000,1.000000}%
\pgfsetfillcolor{currentfill}%
\pgfsetlinewidth{0.602250pt}%
\definecolor{currentstroke}{rgb}{0.000000,0.000000,0.000000}%
\pgfsetstrokecolor{currentstroke}%
\pgfsetdash{}{0pt}%
\pgfsys@defobject{currentmarker}{\pgfqpoint{0.000000in}{0.000000in}}{\pgfqpoint{0.000000in}{0.027778in}}{%
\pgfpathmoveto{\pgfqpoint{0.000000in}{0.000000in}}%
\pgfpathlineto{\pgfqpoint{0.000000in}{0.027778in}}%
\pgfusepath{stroke,fill}%
}%
\begin{pgfscope}%
\pgfsys@transformshift{5.424679in}{0.456901in}%
\pgfsys@useobject{currentmarker}{}%
\end{pgfscope}%
\end{pgfscope}%
\begin{pgfscope}%
\pgfpathrectangle{\pgfqpoint{0.571024in}{0.456901in}}{\pgfqpoint{5.868976in}{2.816432in}}%
\pgfusepath{clip}%
\pgfsetbuttcap%
\pgfsetroundjoin%
\pgfsetlinewidth{0.501875pt}%
\definecolor{currentstroke}{rgb}{0.501961,0.501961,0.501961}%
\pgfsetstrokecolor{currentstroke}%
\pgfsetstrokeopacity{0.400000}%
\pgfsetdash{{1.850000pt}{0.800000pt}}{0.000000pt}%
\pgfpathmoveto{\pgfqpoint{5.919957in}{0.456901in}}%
\pgfpathlineto{\pgfqpoint{5.919957in}{3.273333in}}%
\pgfusepath{stroke}%
\end{pgfscope}%
\begin{pgfscope}%
\pgfsetbuttcap%
\pgfsetroundjoin%
\definecolor{currentfill}{rgb}{1.000000,1.000000,1.000000}%
\pgfsetfillcolor{currentfill}%
\pgfsetlinewidth{0.602250pt}%
\definecolor{currentstroke}{rgb}{0.000000,0.000000,0.000000}%
\pgfsetstrokecolor{currentstroke}%
\pgfsetdash{}{0pt}%
\pgfsys@defobject{currentmarker}{\pgfqpoint{0.000000in}{0.000000in}}{\pgfqpoint{0.000000in}{0.027778in}}{%
\pgfpathmoveto{\pgfqpoint{0.000000in}{0.000000in}}%
\pgfpathlineto{\pgfqpoint{0.000000in}{0.027778in}}%
\pgfusepath{stroke,fill}%
}%
\begin{pgfscope}%
\pgfsys@transformshift{5.919957in}{0.456901in}%
\pgfsys@useobject{currentmarker}{}%
\end{pgfscope}%
\end{pgfscope}%
\begin{pgfscope}%
\pgfpathrectangle{\pgfqpoint{0.571024in}{0.456901in}}{\pgfqpoint{5.868976in}{2.816432in}}%
\pgfusepath{clip}%
\pgfsetbuttcap%
\pgfsetroundjoin%
\pgfsetlinewidth{0.501875pt}%
\definecolor{currentstroke}{rgb}{0.501961,0.501961,0.501961}%
\pgfsetstrokecolor{currentstroke}%
\pgfsetstrokeopacity{0.400000}%
\pgfsetdash{{1.850000pt}{0.800000pt}}{0.000000pt}%
\pgfpathmoveto{\pgfqpoint{6.167597in}{0.456901in}}%
\pgfpathlineto{\pgfqpoint{6.167597in}{3.273333in}}%
\pgfusepath{stroke}%
\end{pgfscope}%
\begin{pgfscope}%
\pgfsetbuttcap%
\pgfsetroundjoin%
\definecolor{currentfill}{rgb}{1.000000,1.000000,1.000000}%
\pgfsetfillcolor{currentfill}%
\pgfsetlinewidth{0.602250pt}%
\definecolor{currentstroke}{rgb}{0.000000,0.000000,0.000000}%
\pgfsetstrokecolor{currentstroke}%
\pgfsetdash{}{0pt}%
\pgfsys@defobject{currentmarker}{\pgfqpoint{0.000000in}{0.000000in}}{\pgfqpoint{0.000000in}{0.027778in}}{%
\pgfpathmoveto{\pgfqpoint{0.000000in}{0.000000in}}%
\pgfpathlineto{\pgfqpoint{0.000000in}{0.027778in}}%
\pgfusepath{stroke,fill}%
}%
\begin{pgfscope}%
\pgfsys@transformshift{6.167597in}{0.456901in}%
\pgfsys@useobject{currentmarker}{}%
\end{pgfscope}%
\end{pgfscope}%
\begin{pgfscope}%
\pgfpathrectangle{\pgfqpoint{0.571024in}{0.456901in}}{\pgfqpoint{5.868976in}{2.816432in}}%
\pgfusepath{clip}%
\pgfsetbuttcap%
\pgfsetroundjoin%
\pgfsetlinewidth{0.501875pt}%
\definecolor{currentstroke}{rgb}{0.501961,0.501961,0.501961}%
\pgfsetstrokecolor{currentstroke}%
\pgfsetstrokeopacity{0.400000}%
\pgfsetdash{{1.850000pt}{0.800000pt}}{0.000000pt}%
\pgfpathmoveto{\pgfqpoint{6.415236in}{0.456901in}}%
\pgfpathlineto{\pgfqpoint{6.415236in}{3.273333in}}%
\pgfusepath{stroke}%
\end{pgfscope}%
\begin{pgfscope}%
\pgfsetbuttcap%
\pgfsetroundjoin%
\definecolor{currentfill}{rgb}{1.000000,1.000000,1.000000}%
\pgfsetfillcolor{currentfill}%
\pgfsetlinewidth{0.602250pt}%
\definecolor{currentstroke}{rgb}{0.000000,0.000000,0.000000}%
\pgfsetstrokecolor{currentstroke}%
\pgfsetdash{}{0pt}%
\pgfsys@defobject{currentmarker}{\pgfqpoint{0.000000in}{0.000000in}}{\pgfqpoint{0.000000in}{0.027778in}}{%
\pgfpathmoveto{\pgfqpoint{0.000000in}{0.000000in}}%
\pgfpathlineto{\pgfqpoint{0.000000in}{0.027778in}}%
\pgfusepath{stroke,fill}%
}%
\begin{pgfscope}%
\pgfsys@transformshift{6.415236in}{0.456901in}%
\pgfsys@useobject{currentmarker}{}%
\end{pgfscope}%
\end{pgfscope}%
\begin{pgfscope}%
\definecolor{textcolor}{rgb}{0.000000,0.000000,0.000000}%
\pgfsetstrokecolor{textcolor}%
\pgfsetfillcolor{textcolor}%
\pgftext[x=3.505512in,y=0.201367in,,top]{\color{textcolor}{\rmfamily\fontsize{12.000000}{14.400000}\selectfont\catcode`\^=\active\def^{\ifmmode\sp\else\^{}\fi}\catcode`\%=\active\def%{\%}$\omega$}}%
\end{pgfscope}%
\begin{pgfscope}%
\pgfsetbuttcap%
\pgfsetroundjoin%
\definecolor{currentfill}{rgb}{1.000000,1.000000,1.000000}%
\pgfsetfillcolor{currentfill}%
\pgfsetlinewidth{0.803000pt}%
\definecolor{currentstroke}{rgb}{0.000000,0.000000,0.000000}%
\pgfsetstrokecolor{currentstroke}%
\pgfsetdash{}{0pt}%
\pgfsys@defobject{currentmarker}{\pgfqpoint{0.000000in}{0.000000in}}{\pgfqpoint{0.048611in}{0.000000in}}{%
\pgfpathmoveto{\pgfqpoint{0.000000in}{0.000000in}}%
\pgfpathlineto{\pgfqpoint{0.048611in}{0.000000in}}%
\pgfusepath{stroke,fill}%
}%
\begin{pgfscope}%
\pgfsys@transformshift{0.571024in}{0.456901in}%
\pgfsys@useobject{currentmarker}{}%
\end{pgfscope}%
\end{pgfscope}%
\begin{pgfscope}%
\definecolor{textcolor}{rgb}{0.000000,0.000000,0.000000}%
\pgfsetstrokecolor{textcolor}%
\pgfsetfillcolor{textcolor}%
\pgftext[x=0.272222in, y=0.399040in, left, base]{\color{textcolor}{\rmfamily\fontsize{12.000000}{14.400000}\selectfont\catcode`\^=\active\def^{\ifmmode\sp\else\^{}\fi}\catcode`\%=\active\def%{\%}$\mathdefault{0.0}$}}%
\end{pgfscope}%
\begin{pgfscope}%
\pgfsetbuttcap%
\pgfsetroundjoin%
\definecolor{currentfill}{rgb}{1.000000,1.000000,1.000000}%
\pgfsetfillcolor{currentfill}%
\pgfsetlinewidth{0.803000pt}%
\definecolor{currentstroke}{rgb}{0.000000,0.000000,0.000000}%
\pgfsetstrokecolor{currentstroke}%
\pgfsetdash{}{0pt}%
\pgfsys@defobject{currentmarker}{\pgfqpoint{0.000000in}{0.000000in}}{\pgfqpoint{0.048611in}{0.000000in}}{%
\pgfpathmoveto{\pgfqpoint{0.000000in}{0.000000in}}%
\pgfpathlineto{\pgfqpoint{0.048611in}{0.000000in}}%
\pgfusepath{stroke,fill}%
}%
\begin{pgfscope}%
\pgfsys@transformshift{0.571024in}{0.827084in}%
\pgfsys@useobject{currentmarker}{}%
\end{pgfscope}%
\end{pgfscope}%
\begin{pgfscope}%
\definecolor{textcolor}{rgb}{0.000000,0.000000,0.000000}%
\pgfsetstrokecolor{textcolor}%
\pgfsetfillcolor{textcolor}%
\pgftext[x=0.272222in, y=0.769222in, left, base]{\color{textcolor}{\rmfamily\fontsize{12.000000}{14.400000}\selectfont\catcode`\^=\active\def^{\ifmmode\sp\else\^{}\fi}\catcode`\%=\active\def%{\%}$\mathdefault{0.2}$}}%
\end{pgfscope}%
\begin{pgfscope}%
\pgfsetbuttcap%
\pgfsetroundjoin%
\definecolor{currentfill}{rgb}{1.000000,1.000000,1.000000}%
\pgfsetfillcolor{currentfill}%
\pgfsetlinewidth{0.803000pt}%
\definecolor{currentstroke}{rgb}{0.000000,0.000000,0.000000}%
\pgfsetstrokecolor{currentstroke}%
\pgfsetdash{}{0pt}%
\pgfsys@defobject{currentmarker}{\pgfqpoint{0.000000in}{0.000000in}}{\pgfqpoint{0.048611in}{0.000000in}}{%
\pgfpathmoveto{\pgfqpoint{0.000000in}{0.000000in}}%
\pgfpathlineto{\pgfqpoint{0.048611in}{0.000000in}}%
\pgfusepath{stroke,fill}%
}%
\begin{pgfscope}%
\pgfsys@transformshift{0.571024in}{1.197266in}%
\pgfsys@useobject{currentmarker}{}%
\end{pgfscope}%
\end{pgfscope}%
\begin{pgfscope}%
\definecolor{textcolor}{rgb}{0.000000,0.000000,0.000000}%
\pgfsetstrokecolor{textcolor}%
\pgfsetfillcolor{textcolor}%
\pgftext[x=0.272222in, y=1.139405in, left, base]{\color{textcolor}{\rmfamily\fontsize{12.000000}{14.400000}\selectfont\catcode`\^=\active\def^{\ifmmode\sp\else\^{}\fi}\catcode`\%=\active\def%{\%}$\mathdefault{0.4}$}}%
\end{pgfscope}%
\begin{pgfscope}%
\pgfsetbuttcap%
\pgfsetroundjoin%
\definecolor{currentfill}{rgb}{1.000000,1.000000,1.000000}%
\pgfsetfillcolor{currentfill}%
\pgfsetlinewidth{0.803000pt}%
\definecolor{currentstroke}{rgb}{0.000000,0.000000,0.000000}%
\pgfsetstrokecolor{currentstroke}%
\pgfsetdash{}{0pt}%
\pgfsys@defobject{currentmarker}{\pgfqpoint{0.000000in}{0.000000in}}{\pgfqpoint{0.048611in}{0.000000in}}{%
\pgfpathmoveto{\pgfqpoint{0.000000in}{0.000000in}}%
\pgfpathlineto{\pgfqpoint{0.048611in}{0.000000in}}%
\pgfusepath{stroke,fill}%
}%
\begin{pgfscope}%
\pgfsys@transformshift{0.571024in}{1.567449in}%
\pgfsys@useobject{currentmarker}{}%
\end{pgfscope}%
\end{pgfscope}%
\begin{pgfscope}%
\definecolor{textcolor}{rgb}{0.000000,0.000000,0.000000}%
\pgfsetstrokecolor{textcolor}%
\pgfsetfillcolor{textcolor}%
\pgftext[x=0.272222in, y=1.509588in, left, base]{\color{textcolor}{\rmfamily\fontsize{12.000000}{14.400000}\selectfont\catcode`\^=\active\def^{\ifmmode\sp\else\^{}\fi}\catcode`\%=\active\def%{\%}$\mathdefault{0.6}$}}%
\end{pgfscope}%
\begin{pgfscope}%
\pgfsetbuttcap%
\pgfsetroundjoin%
\definecolor{currentfill}{rgb}{1.000000,1.000000,1.000000}%
\pgfsetfillcolor{currentfill}%
\pgfsetlinewidth{0.803000pt}%
\definecolor{currentstroke}{rgb}{0.000000,0.000000,0.000000}%
\pgfsetstrokecolor{currentstroke}%
\pgfsetdash{}{0pt}%
\pgfsys@defobject{currentmarker}{\pgfqpoint{0.000000in}{0.000000in}}{\pgfqpoint{0.048611in}{0.000000in}}{%
\pgfpathmoveto{\pgfqpoint{0.000000in}{0.000000in}}%
\pgfpathlineto{\pgfqpoint{0.048611in}{0.000000in}}%
\pgfusepath{stroke,fill}%
}%
\begin{pgfscope}%
\pgfsys@transformshift{0.571024in}{1.937632in}%
\pgfsys@useobject{currentmarker}{}%
\end{pgfscope}%
\end{pgfscope}%
\begin{pgfscope}%
\definecolor{textcolor}{rgb}{0.000000,0.000000,0.000000}%
\pgfsetstrokecolor{textcolor}%
\pgfsetfillcolor{textcolor}%
\pgftext[x=0.272222in, y=1.879770in, left, base]{\color{textcolor}{\rmfamily\fontsize{12.000000}{14.400000}\selectfont\catcode`\^=\active\def^{\ifmmode\sp\else\^{}\fi}\catcode`\%=\active\def%{\%}$\mathdefault{0.8}$}}%
\end{pgfscope}%
\begin{pgfscope}%
\pgfsetbuttcap%
\pgfsetroundjoin%
\definecolor{currentfill}{rgb}{1.000000,1.000000,1.000000}%
\pgfsetfillcolor{currentfill}%
\pgfsetlinewidth{0.803000pt}%
\definecolor{currentstroke}{rgb}{0.000000,0.000000,0.000000}%
\pgfsetstrokecolor{currentstroke}%
\pgfsetdash{}{0pt}%
\pgfsys@defobject{currentmarker}{\pgfqpoint{0.000000in}{0.000000in}}{\pgfqpoint{0.048611in}{0.000000in}}{%
\pgfpathmoveto{\pgfqpoint{0.000000in}{0.000000in}}%
\pgfpathlineto{\pgfqpoint{0.048611in}{0.000000in}}%
\pgfusepath{stroke,fill}%
}%
\begin{pgfscope}%
\pgfsys@transformshift{0.571024in}{2.307814in}%
\pgfsys@useobject{currentmarker}{}%
\end{pgfscope}%
\end{pgfscope}%
\begin{pgfscope}%
\definecolor{textcolor}{rgb}{0.000000,0.000000,0.000000}%
\pgfsetstrokecolor{textcolor}%
\pgfsetfillcolor{textcolor}%
\pgftext[x=0.272222in, y=2.249953in, left, base]{\color{textcolor}{\rmfamily\fontsize{12.000000}{14.400000}\selectfont\catcode`\^=\active\def^{\ifmmode\sp\else\^{}\fi}\catcode`\%=\active\def%{\%}$\mathdefault{1.0}$}}%
\end{pgfscope}%
\begin{pgfscope}%
\pgfsetbuttcap%
\pgfsetroundjoin%
\definecolor{currentfill}{rgb}{1.000000,1.000000,1.000000}%
\pgfsetfillcolor{currentfill}%
\pgfsetlinewidth{0.803000pt}%
\definecolor{currentstroke}{rgb}{0.000000,0.000000,0.000000}%
\pgfsetstrokecolor{currentstroke}%
\pgfsetdash{}{0pt}%
\pgfsys@defobject{currentmarker}{\pgfqpoint{0.000000in}{0.000000in}}{\pgfqpoint{0.048611in}{0.000000in}}{%
\pgfpathmoveto{\pgfqpoint{0.000000in}{0.000000in}}%
\pgfpathlineto{\pgfqpoint{0.048611in}{0.000000in}}%
\pgfusepath{stroke,fill}%
}%
\begin{pgfscope}%
\pgfsys@transformshift{0.571024in}{2.677997in}%
\pgfsys@useobject{currentmarker}{}%
\end{pgfscope}%
\end{pgfscope}%
\begin{pgfscope}%
\definecolor{textcolor}{rgb}{0.000000,0.000000,0.000000}%
\pgfsetstrokecolor{textcolor}%
\pgfsetfillcolor{textcolor}%
\pgftext[x=0.272222in, y=2.620136in, left, base]{\color{textcolor}{\rmfamily\fontsize{12.000000}{14.400000}\selectfont\catcode`\^=\active\def^{\ifmmode\sp\else\^{}\fi}\catcode`\%=\active\def%{\%}$\mathdefault{1.2}$}}%
\end{pgfscope}%
\begin{pgfscope}%
\pgfsetbuttcap%
\pgfsetroundjoin%
\definecolor{currentfill}{rgb}{1.000000,1.000000,1.000000}%
\pgfsetfillcolor{currentfill}%
\pgfsetlinewidth{0.803000pt}%
\definecolor{currentstroke}{rgb}{0.000000,0.000000,0.000000}%
\pgfsetstrokecolor{currentstroke}%
\pgfsetdash{}{0pt}%
\pgfsys@defobject{currentmarker}{\pgfqpoint{0.000000in}{0.000000in}}{\pgfqpoint{0.048611in}{0.000000in}}{%
\pgfpathmoveto{\pgfqpoint{0.000000in}{0.000000in}}%
\pgfpathlineto{\pgfqpoint{0.048611in}{0.000000in}}%
\pgfusepath{stroke,fill}%
}%
\begin{pgfscope}%
\pgfsys@transformshift{0.571024in}{3.048180in}%
\pgfsys@useobject{currentmarker}{}%
\end{pgfscope}%
\end{pgfscope}%
\begin{pgfscope}%
\definecolor{textcolor}{rgb}{0.000000,0.000000,0.000000}%
\pgfsetstrokecolor{textcolor}%
\pgfsetfillcolor{textcolor}%
\pgftext[x=0.272222in, y=2.990318in, left, base]{\color{textcolor}{\rmfamily\fontsize{12.000000}{14.400000}\selectfont\catcode`\^=\active\def^{\ifmmode\sp\else\^{}\fi}\catcode`\%=\active\def%{\%}$\mathdefault{1.4}$}}%
\end{pgfscope}%
\begin{pgfscope}%
\definecolor{textcolor}{rgb}{0.000000,0.000000,0.000000}%
\pgfsetstrokecolor{textcolor}%
\pgfsetfillcolor{textcolor}%
\pgftext[x=0.216667in,y=1.865117in,,bottom,rotate=90.000000]{\color{textcolor}{\rmfamily\fontsize{12.000000}{14.400000}\selectfont\catcode`\^=\active\def^{\ifmmode\sp\else\^{}\fi}\catcode`\%=\active\def%{\%}$A_{\mathrm{вых}}(\omega)$}}%
\end{pgfscope}%
\begin{pgfscope}%
\pgfpathrectangle{\pgfqpoint{0.571024in}{0.456901in}}{\pgfqpoint{5.868976in}{2.816432in}}%
\pgfusepath{clip}%
\pgfsetbuttcap%
\pgfsetroundjoin%
\pgfsetlinewidth{2.007500pt}%
\definecolor{currentstroke}{rgb}{1.000000,0.000000,0.000000}%
\pgfsetstrokecolor{currentstroke}%
\pgfsetdash{{7.400000pt}{3.200000pt}}{0.000000pt}%
\pgfpathmoveto{\pgfqpoint{0.719532in}{0.456901in}}%
\pgfpathlineto{\pgfqpoint{0.828712in}{1.091434in}}%
\pgfpathlineto{\pgfqpoint{0.883302in}{1.396849in}}%
\pgfpathlineto{\pgfqpoint{0.926974in}{1.631338in}}%
\pgfpathlineto{\pgfqpoint{0.970645in}{1.854711in}}%
\pgfpathlineto{\pgfqpoint{1.014317in}{2.064931in}}%
\pgfpathlineto{\pgfqpoint{1.047071in}{2.212833in}}%
\pgfpathlineto{\pgfqpoint{1.079825in}{2.351560in}}%
\pgfpathlineto{\pgfqpoint{1.112579in}{2.480458in}}%
\pgfpathlineto{\pgfqpoint{1.145333in}{2.598941in}}%
\pgfpathlineto{\pgfqpoint{1.178087in}{2.706497in}}%
\pgfpathlineto{\pgfqpoint{1.199923in}{2.771910in}}%
\pgfpathlineto{\pgfqpoint{1.221759in}{2.832162in}}%
\pgfpathlineto{\pgfqpoint{1.243595in}{2.887157in}}%
\pgfpathlineto{\pgfqpoint{1.265431in}{2.936818in}}%
\pgfpathlineto{\pgfqpoint{1.287266in}{2.981084in}}%
\pgfpathlineto{\pgfqpoint{1.309102in}{3.019908in}}%
\pgfpathlineto{\pgfqpoint{1.330938in}{3.053263in}}%
\pgfpathlineto{\pgfqpoint{1.352774in}{3.081136in}}%
\pgfpathlineto{\pgfqpoint{1.374610in}{3.103532in}}%
\pgfpathlineto{\pgfqpoint{1.396446in}{3.120470in}}%
\pgfpathlineto{\pgfqpoint{1.407364in}{3.126903in}}%
\pgfpathlineto{\pgfqpoint{1.418282in}{3.131987in}}%
\pgfpathlineto{\pgfqpoint{1.429200in}{3.135729in}}%
\pgfpathlineto{\pgfqpoint{1.440118in}{3.138136in}}%
\pgfpathlineto{\pgfqpoint{1.451036in}{3.139217in}}%
\pgfpathlineto{\pgfqpoint{1.461954in}{3.138984in}}%
\pgfpathlineto{\pgfqpoint{1.472872in}{3.137445in}}%
\pgfpathlineto{\pgfqpoint{1.483790in}{3.134614in}}%
\pgfpathlineto{\pgfqpoint{1.494708in}{3.130503in}}%
\pgfpathlineto{\pgfqpoint{1.505626in}{3.125125in}}%
\pgfpathlineto{\pgfqpoint{1.527462in}{3.110628in}}%
\pgfpathlineto{\pgfqpoint{1.549298in}{3.091251in}}%
\pgfpathlineto{\pgfqpoint{1.571134in}{3.067132in}}%
\pgfpathlineto{\pgfqpoint{1.592969in}{3.038425in}}%
\pgfpathlineto{\pgfqpoint{1.614805in}{3.005294in}}%
\pgfpathlineto{\pgfqpoint{1.636641in}{2.967916in}}%
\pgfpathlineto{\pgfqpoint{1.658477in}{2.926478in}}%
\pgfpathlineto{\pgfqpoint{1.680313in}{2.881177in}}%
\pgfpathlineto{\pgfqpoint{1.713067in}{2.806438in}}%
\pgfpathlineto{\pgfqpoint{1.745821in}{2.724208in}}%
\pgfpathlineto{\pgfqpoint{1.778575in}{2.635258in}}%
\pgfpathlineto{\pgfqpoint{1.811329in}{2.540390in}}%
\pgfpathlineto{\pgfqpoint{1.855001in}{2.406123in}}%
\pgfpathlineto{\pgfqpoint{1.898672in}{2.264810in}}%
\pgfpathlineto{\pgfqpoint{1.964180in}{2.044073in}}%
\pgfpathlineto{\pgfqpoint{2.117032in}{1.523045in}}%
\pgfpathlineto{\pgfqpoint{2.171622in}{1.346568in}}%
\pgfpathlineto{\pgfqpoint{2.215293in}{1.212674in}}%
\pgfpathlineto{\pgfqpoint{2.258965in}{1.086833in}}%
\pgfpathlineto{\pgfqpoint{2.291719in}{0.998514in}}%
\pgfpathlineto{\pgfqpoint{2.324473in}{0.915920in}}%
\pgfpathlineto{\pgfqpoint{2.357227in}{0.839464in}}%
\pgfpathlineto{\pgfqpoint{2.389981in}{0.769495in}}%
\pgfpathlineto{\pgfqpoint{2.422735in}{0.706301in}}%
\pgfpathlineto{\pgfqpoint{2.444571in}{0.668048in}}%
\pgfpathlineto{\pgfqpoint{2.466407in}{0.632954in}}%
\pgfpathlineto{\pgfqpoint{2.488243in}{0.601057in}}%
\pgfpathlineto{\pgfqpoint{2.510079in}{0.572380in}}%
\pgfpathlineto{\pgfqpoint{2.531914in}{0.546935in}}%
\pgfpathlineto{\pgfqpoint{2.553750in}{0.524719in}}%
\pgfpathlineto{\pgfqpoint{2.575586in}{0.505721in}}%
\pgfpathlineto{\pgfqpoint{2.597422in}{0.489912in}}%
\pgfpathlineto{\pgfqpoint{2.619258in}{0.477256in}}%
\pgfpathlineto{\pgfqpoint{2.641094in}{0.467703in}}%
\pgfpathlineto{\pgfqpoint{2.662930in}{0.461193in}}%
\pgfpathlineto{\pgfqpoint{2.684766in}{0.457654in}}%
\pgfpathlineto{\pgfqpoint{2.706602in}{0.457006in}}%
\pgfpathlineto{\pgfqpoint{2.728438in}{0.459156in}}%
\pgfpathlineto{\pgfqpoint{2.750274in}{0.464005in}}%
\pgfpathlineto{\pgfqpoint{2.772110in}{0.471444in}}%
\pgfpathlineto{\pgfqpoint{2.793946in}{0.481354in}}%
\pgfpathlineto{\pgfqpoint{2.815782in}{0.493611in}}%
\pgfpathlineto{\pgfqpoint{2.837617in}{0.508082in}}%
\pgfpathlineto{\pgfqpoint{2.859453in}{0.524631in}}%
\pgfpathlineto{\pgfqpoint{2.892207in}{0.553031in}}%
\pgfpathlineto{\pgfqpoint{2.924961in}{0.585272in}}%
\pgfpathlineto{\pgfqpoint{2.957715in}{0.620827in}}%
\pgfpathlineto{\pgfqpoint{3.001387in}{0.672454in}}%
\pgfpathlineto{\pgfqpoint{3.045059in}{0.727702in}}%
\pgfpathlineto{\pgfqpoint{3.110567in}{0.814494in}}%
\pgfpathlineto{\pgfqpoint{3.208828in}{0.944940in}}%
\pgfpathlineto{\pgfqpoint{3.263418in}{1.013154in}}%
\pgfpathlineto{\pgfqpoint{3.307090in}{1.063698in}}%
\pgfpathlineto{\pgfqpoint{3.339844in}{1.098627in}}%
\pgfpathlineto{\pgfqpoint{3.372598in}{1.130602in}}%
\pgfpathlineto{\pgfqpoint{3.405352in}{1.159307in}}%
\pgfpathlineto{\pgfqpoint{3.438106in}{1.184467in}}%
\pgfpathlineto{\pgfqpoint{3.470859in}{1.205850in}}%
\pgfpathlineto{\pgfqpoint{3.492695in}{1.217913in}}%
\pgfpathlineto{\pgfqpoint{3.514531in}{1.228169in}}%
\pgfpathlineto{\pgfqpoint{3.536367in}{1.236584in}}%
\pgfpathlineto{\pgfqpoint{3.558203in}{1.243130in}}%
\pgfpathlineto{\pgfqpoint{3.580039in}{1.247793in}}%
\pgfpathlineto{\pgfqpoint{3.601875in}{1.250562in}}%
\pgfpathlineto{\pgfqpoint{3.623711in}{1.251440in}}%
\pgfpathlineto{\pgfqpoint{3.645547in}{1.250437in}}%
\pgfpathlineto{\pgfqpoint{3.667383in}{1.247570in}}%
\pgfpathlineto{\pgfqpoint{3.689219in}{1.242866in}}%
\pgfpathlineto{\pgfqpoint{3.711055in}{1.236361in}}%
\pgfpathlineto{\pgfqpoint{3.732891in}{1.228097in}}%
\pgfpathlineto{\pgfqpoint{3.754727in}{1.218124in}}%
\pgfpathlineto{\pgfqpoint{3.787480in}{1.200090in}}%
\pgfpathlineto{\pgfqpoint{3.820234in}{1.178567in}}%
\pgfpathlineto{\pgfqpoint{3.852988in}{1.153810in}}%
\pgfpathlineto{\pgfqpoint{3.885742in}{1.126103in}}%
\pgfpathlineto{\pgfqpoint{3.918496in}{1.095756in}}%
\pgfpathlineto{\pgfqpoint{3.962168in}{1.051761in}}%
\pgfpathlineto{\pgfqpoint{4.005840in}{1.004495in}}%
\pgfpathlineto{\pgfqpoint{4.060430in}{0.942144in}}%
\pgfpathlineto{\pgfqpoint{4.256953in}{0.713656in}}%
\pgfpathlineto{\pgfqpoint{4.300625in}{0.667345in}}%
\pgfpathlineto{\pgfqpoint{4.344297in}{0.624334in}}%
\pgfpathlineto{\pgfqpoint{4.377051in}{0.594635in}}%
\pgfpathlineto{\pgfqpoint{4.409804in}{0.567401in}}%
\pgfpathlineto{\pgfqpoint{4.442558in}{0.542845in}}%
\pgfpathlineto{\pgfqpoint{4.475312in}{0.521149in}}%
\pgfpathlineto{\pgfqpoint{4.508066in}{0.502462in}}%
\pgfpathlineto{\pgfqpoint{4.540820in}{0.486900in}}%
\pgfpathlineto{\pgfqpoint{4.573574in}{0.474547in}}%
\pgfpathlineto{\pgfqpoint{4.606328in}{0.465450in}}%
\pgfpathlineto{\pgfqpoint{4.639082in}{0.459624in}}%
\pgfpathlineto{\pgfqpoint{4.671836in}{0.457047in}}%
\pgfpathlineto{\pgfqpoint{4.704589in}{0.457668in}}%
\pgfpathlineto{\pgfqpoint{4.737343in}{0.461401in}}%
\pgfpathlineto{\pgfqpoint{4.770097in}{0.468131in}}%
\pgfpathlineto{\pgfqpoint{4.802851in}{0.477712in}}%
\pgfpathlineto{\pgfqpoint{4.835605in}{0.489974in}}%
\pgfpathlineto{\pgfqpoint{4.868359in}{0.504720in}}%
\pgfpathlineto{\pgfqpoint{4.901113in}{0.521731in}}%
\pgfpathlineto{\pgfqpoint{4.944785in}{0.547519in}}%
\pgfpathlineto{\pgfqpoint{4.988457in}{0.576294in}}%
\pgfpathlineto{\pgfqpoint{5.043046in}{0.615464in}}%
\pgfpathlineto{\pgfqpoint{5.119472in}{0.673867in}}%
\pgfpathlineto{\pgfqpoint{5.228652in}{0.757713in}}%
\pgfpathlineto{\pgfqpoint{5.283242in}{0.796917in}}%
\pgfpathlineto{\pgfqpoint{5.326913in}{0.825866in}}%
\pgfpathlineto{\pgfqpoint{5.370585in}{0.852071in}}%
\pgfpathlineto{\pgfqpoint{5.403339in}{0.869629in}}%
\pgfpathlineto{\pgfqpoint{5.436093in}{0.885189in}}%
\pgfpathlineto{\pgfqpoint{5.468847in}{0.898594in}}%
\pgfpathlineto{\pgfqpoint{5.501601in}{0.909715in}}%
\pgfpathlineto{\pgfqpoint{5.534355in}{0.918445in}}%
\pgfpathlineto{\pgfqpoint{5.567109in}{0.924707in}}%
\pgfpathlineto{\pgfqpoint{5.599863in}{0.928451in}}%
\pgfpathlineto{\pgfqpoint{5.632616in}{0.929653in}}%
\pgfpathlineto{\pgfqpoint{5.665370in}{0.928317in}}%
\pgfpathlineto{\pgfqpoint{5.698124in}{0.924474in}}%
\pgfpathlineto{\pgfqpoint{5.730878in}{0.918181in}}%
\pgfpathlineto{\pgfqpoint{5.763632in}{0.909520in}}%
\pgfpathlineto{\pgfqpoint{5.796386in}{0.898599in}}%
\pgfpathlineto{\pgfqpoint{5.829140in}{0.885546in}}%
\pgfpathlineto{\pgfqpoint{5.861894in}{0.870513in}}%
\pgfpathlineto{\pgfqpoint{5.905566in}{0.847681in}}%
\pgfpathlineto{\pgfqpoint{5.949237in}{0.822079in}}%
\pgfpathlineto{\pgfqpoint{6.003827in}{0.786936in}}%
\pgfpathlineto{\pgfqpoint{6.069335in}{0.741535in}}%
\pgfpathlineto{\pgfqpoint{6.167597in}{0.670785in}}%
\pgfpathlineto{\pgfqpoint{6.167597in}{0.670785in}}%
\pgfusepath{stroke}%
\end{pgfscope}%
\begin{pgfscope}%
\pgfpathrectangle{\pgfqpoint{0.571024in}{0.456901in}}{\pgfqpoint{5.868976in}{2.816432in}}%
\pgfusepath{clip}%
\pgfsetbuttcap%
\pgfsetroundjoin%
\pgfsetlinewidth{2.007500pt}%
\definecolor{currentstroke}{rgb}{0.000000,0.000000,1.000000}%
\pgfsetstrokecolor{currentstroke}%
\pgfsetdash{}{0pt}%
\pgfpathmoveto{\pgfqpoint{0.719532in}{0.456901in}}%
\pgfpathlineto{\pgfqpoint{0.719532in}{0.456901in}}%
\pgfusepath{stroke}%
\end{pgfscope}%
\begin{pgfscope}%
\pgfpathrectangle{\pgfqpoint{0.571024in}{0.456901in}}{\pgfqpoint{5.868976in}{2.816432in}}%
\pgfusepath{clip}%
\pgfsetbuttcap%
\pgfsetroundjoin%
\pgfsetlinewidth{2.007500pt}%
\definecolor{currentstroke}{rgb}{0.000000,0.000000,1.000000}%
\pgfsetstrokecolor{currentstroke}%
\pgfsetdash{}{0pt}%
\pgfpathmoveto{\pgfqpoint{1.710089in}{0.456901in}}%
\pgfpathlineto{\pgfqpoint{1.710089in}{2.813553in}}%
\pgfusepath{stroke}%
\end{pgfscope}%
\begin{pgfscope}%
\pgfpathrectangle{\pgfqpoint{0.571024in}{0.456901in}}{\pgfqpoint{5.868976in}{2.816432in}}%
\pgfusepath{clip}%
\pgfsetbuttcap%
\pgfsetroundjoin%
\pgfsetlinewidth{2.007500pt}%
\definecolor{currentstroke}{rgb}{0.000000,0.000000,1.000000}%
\pgfsetstrokecolor{currentstroke}%
\pgfsetdash{}{0pt}%
\pgfpathmoveto{\pgfqpoint{2.700647in}{0.456901in}}%
\pgfpathlineto{\pgfqpoint{2.700647in}{0.456901in}}%
\pgfusepath{stroke}%
\end{pgfscope}%
\begin{pgfscope}%
\pgfpathrectangle{\pgfqpoint{0.571024in}{0.456901in}}{\pgfqpoint{5.868976in}{2.816432in}}%
\pgfusepath{clip}%
\pgfsetbuttcap%
\pgfsetroundjoin%
\pgfsetlinewidth{2.007500pt}%
\definecolor{currentstroke}{rgb}{0.000000,0.000000,1.000000}%
\pgfsetstrokecolor{currentstroke}%
\pgfsetdash{}{0pt}%
\pgfpathmoveto{\pgfqpoint{3.691204in}{0.456901in}}%
\pgfpathlineto{\pgfqpoint{3.691204in}{1.242348in}}%
\pgfusepath{stroke}%
\end{pgfscope}%
\begin{pgfscope}%
\pgfpathrectangle{\pgfqpoint{0.571024in}{0.456901in}}{\pgfqpoint{5.868976in}{2.816432in}}%
\pgfusepath{clip}%
\pgfsetbuttcap%
\pgfsetroundjoin%
\pgfsetlinewidth{2.007500pt}%
\definecolor{currentstroke}{rgb}{0.000000,0.000000,1.000000}%
\pgfsetstrokecolor{currentstroke}%
\pgfsetdash{}{0pt}%
\pgfpathmoveto{\pgfqpoint{4.681761in}{0.456901in}}%
\pgfpathlineto{\pgfqpoint{4.681761in}{0.456901in}}%
\pgfusepath{stroke}%
\end{pgfscope}%
\begin{pgfscope}%
\pgfpathrectangle{\pgfqpoint{0.571024in}{0.456901in}}{\pgfqpoint{5.868976in}{2.816432in}}%
\pgfusepath{clip}%
\pgfsetbuttcap%
\pgfsetroundjoin%
\pgfsetlinewidth{2.007500pt}%
\definecolor{currentstroke}{rgb}{0.000000,0.000000,1.000000}%
\pgfsetstrokecolor{currentstroke}%
\pgfsetdash{}{0pt}%
\pgfpathmoveto{\pgfqpoint{5.672318in}{0.456901in}}%
\pgfpathlineto{\pgfqpoint{5.672318in}{0.927710in}}%
\pgfusepath{stroke}%
\end{pgfscope}%
\begin{pgfscope}%
\pgfpathrectangle{\pgfqpoint{0.571024in}{0.456901in}}{\pgfqpoint{5.868976in}{2.816432in}}%
\pgfusepath{clip}%
\pgfsetbuttcap%
\pgfsetroundjoin%
\definecolor{currentfill}{rgb}{1.000000,1.000000,1.000000}%
\pgfsetfillcolor{currentfill}%
\pgfsetlinewidth{1.505625pt}%
\definecolor{currentstroke}{rgb}{0.000000,0.000000,1.000000}%
\pgfsetstrokecolor{currentstroke}%
\pgfsetdash{}{0pt}%
\pgfsys@defobject{currentmarker}{\pgfqpoint{-0.027778in}{-0.027778in}}{\pgfqpoint{0.027778in}{0.027778in}}{%
\pgfpathmoveto{\pgfqpoint{0.000000in}{-0.027778in}}%
\pgfpathcurveto{\pgfqpoint{0.007367in}{-0.027778in}}{\pgfqpoint{0.014433in}{-0.024851in}}{\pgfqpoint{0.019642in}{-0.019642in}}%
\pgfpathcurveto{\pgfqpoint{0.024851in}{-0.014433in}}{\pgfqpoint{0.027778in}{-0.007367in}}{\pgfqpoint{0.027778in}{0.000000in}}%
\pgfpathcurveto{\pgfqpoint{0.027778in}{0.007367in}}{\pgfqpoint{0.024851in}{0.014433in}}{\pgfqpoint{0.019642in}{0.019642in}}%
\pgfpathcurveto{\pgfqpoint{0.014433in}{0.024851in}}{\pgfqpoint{0.007367in}{0.027778in}}{\pgfqpoint{0.000000in}{0.027778in}}%
\pgfpathcurveto{\pgfqpoint{-0.007367in}{0.027778in}}{\pgfqpoint{-0.014433in}{0.024851in}}{\pgfqpoint{-0.019642in}{0.019642in}}%
\pgfpathcurveto{\pgfqpoint{-0.024851in}{0.014433in}}{\pgfqpoint{-0.027778in}{0.007367in}}{\pgfqpoint{-0.027778in}{0.000000in}}%
\pgfpathcurveto{\pgfqpoint{-0.027778in}{-0.007367in}}{\pgfqpoint{-0.024851in}{-0.014433in}}{\pgfqpoint{-0.019642in}{-0.019642in}}%
\pgfpathcurveto{\pgfqpoint{-0.014433in}{-0.024851in}}{\pgfqpoint{-0.007367in}{-0.027778in}}{\pgfqpoint{0.000000in}{-0.027778in}}%
\pgfpathlineto{\pgfqpoint{0.000000in}{-0.027778in}}%
\pgfpathclose%
\pgfusepath{stroke,fill}%
}%
\begin{pgfscope}%
\pgfsys@transformshift{0.719532in}{0.456901in}%
\pgfsys@useobject{currentmarker}{}%
\end{pgfscope}%
\begin{pgfscope}%
\pgfsys@transformshift{1.710089in}{2.813553in}%
\pgfsys@useobject{currentmarker}{}%
\end{pgfscope}%
\begin{pgfscope}%
\pgfsys@transformshift{2.700647in}{0.456901in}%
\pgfsys@useobject{currentmarker}{}%
\end{pgfscope}%
\begin{pgfscope}%
\pgfsys@transformshift{3.691204in}{1.242348in}%
\pgfsys@useobject{currentmarker}{}%
\end{pgfscope}%
\begin{pgfscope}%
\pgfsys@transformshift{4.681761in}{0.456901in}%
\pgfsys@useobject{currentmarker}{}%
\end{pgfscope}%
\begin{pgfscope}%
\pgfsys@transformshift{5.672318in}{0.927710in}%
\pgfsys@useobject{currentmarker}{}%
\end{pgfscope}%
\end{pgfscope}%
\begin{pgfscope}%
\pgfpathrectangle{\pgfqpoint{0.571024in}{0.456901in}}{\pgfqpoint{5.868976in}{2.816432in}}%
\pgfusepath{clip}%
\pgfsetrectcap%
\pgfsetroundjoin%
\pgfsetlinewidth{2.007500pt}%
\definecolor{currentstroke}{rgb}{0.000000,0.000000,0.000000}%
\pgfsetstrokecolor{currentstroke}%
\pgfsetdash{}{0pt}%
\pgfpathmoveto{\pgfqpoint{0.719532in}{0.456901in}}%
\pgfpathlineto{\pgfqpoint{5.672318in}{0.456901in}}%
\pgfusepath{stroke}%
\end{pgfscope}%
\begin{pgfscope}%
\pgfsetrectcap%
\pgfsetmiterjoin%
\pgfsetlinewidth{0.602250pt}%
\definecolor{currentstroke}{rgb}{0.000000,0.000000,0.000000}%
\pgfsetstrokecolor{currentstroke}%
\pgfsetdash{}{0pt}%
\pgfpathmoveto{\pgfqpoint{0.571024in}{0.456901in}}%
\pgfpathlineto{\pgfqpoint{0.571024in}{3.273333in}}%
\pgfusepath{stroke}%
\end{pgfscope}%
\begin{pgfscope}%
\pgfsetrectcap%
\pgfsetmiterjoin%
\pgfsetlinewidth{0.602250pt}%
\definecolor{currentstroke}{rgb}{0.000000,0.000000,0.000000}%
\pgfsetstrokecolor{currentstroke}%
\pgfsetdash{}{0pt}%
\pgfpathmoveto{\pgfqpoint{0.571024in}{0.456901in}}%
\pgfpathlineto{\pgfqpoint{6.440000in}{0.456901in}}%
\pgfusepath{stroke}%
\end{pgfscope}%
\begin{pgfscope}%
\pgfsetbuttcap%
\pgfsetroundjoin%
\pgfsetlinewidth{2.007500pt}%
\definecolor{currentstroke}{rgb}{1.000000,0.000000,0.000000}%
\pgfsetstrokecolor{currentstroke}%
\pgfsetdash{{7.400000pt}{3.200000pt}}{0.000000pt}%
\pgfpathmoveto{\pgfqpoint{2.474688in}{3.556667in}}%
\pgfpathlineto{\pgfqpoint{2.724688in}{3.556667in}}%
\pgfpathlineto{\pgfqpoint{2.974688in}{3.556667in}}%
\pgfusepath{stroke}%
\end{pgfscope}%
\begin{pgfscope}%
\definecolor{textcolor}{rgb}{0.000000,0.000000,0.000000}%
\pgfsetstrokecolor{textcolor}%
\pgfsetfillcolor{textcolor}%
\pgftext[x=3.108021in,y=3.498333in,left,base]{\color{textcolor}{\rmfamily\fontsize{12.000000}{14.400000}\selectfont\catcode`\^=\active\def^{\ifmmode\sp\else\^{}\fi}\catcode`\%=\active\def%{\%}Огибающая $A_{\mathrm{вых}}(\omega)$}}%
\end{pgfscope}%
\end{pgfpicture}%
\makeatother%
\endgroup%

    \caption{Амплитудный дискретный спектр реакции на первый входной сигнал}
    \label{fig:plot_disc_spec_A_out_1}
\end{figure}

\begin{figure}[H]
    \centering
    %% Creator: Matplotlib, PGF backend
%%
%% To include the figure in your LaTeX document, write
%%   \input{<filename>.pgf}
%%
%% Make sure the required packages are loaded in your preamble
%%   \usepackage{pgf}
%%
%% Also ensure that all the required font packages are loaded; for instance,
%% the lmodern package is sometimes necessary when using math font.
%%   \usepackage{lmodern}
%%
%% Figures using additional raster images can only be included by \input if
%% they are in the same directory as the main LaTeX file. For loading figures
%% from other directories you can use the `import` package
%%   \usepackage{import}
%%
%% and then include the figures with
%%   \import{<path to file>}{<filename>.pgf}
%%
%% Matplotlib used the following preamble
%%   \def\mathdefault#1{#1}
%%   \everymath=\expandafter{\the\everymath\displaystyle}
%%   \IfFileExists{scrextend.sty}{
%%     \usepackage[fontsize=12.000000pt]{scrextend}
%%   }{
%%     \renewcommand{\normalsize}{\fontsize{12.000000}{14.400000}\selectfont}
%%     \normalsize
%%   }
%%   \usepackage{fontspec}\setmainfont{Times New Roman}\usepackage[english,russian]{babel}\usepackage{amsmath}
%%   \ifdefined\pdftexversion\else  % non-pdftex case.
%%     \usepackage{fontspec}
%%     \setmainfont{times.ttf}[Path=\detokenize{/usr/local/share/fonts/truetype/times-new-roman/}]
%%     \setsansfont{DejaVuSans.ttf}[Path=\detokenize{/usr/share/matplotlib/mpl-data/fonts/ttf/}]
%%     \setmonofont{DejaVuSansMono.ttf}[Path=\detokenize{/usr/share/matplotlib/mpl-data/fonts/ttf/}]
%%   \fi
%%   \makeatletter\@ifpackageloaded{underscore}{}{\usepackage[strings]{underscore}}\makeatother
%%
\begingroup%
\makeatletter%
\begin{pgfpicture}%
\pgfpathrectangle{\pgfpointorigin}{\pgfqpoint{6.490000in}{3.740000in}}%
\pgfusepath{use as bounding box, clip}%
\begin{pgfscope}%
\pgfsetbuttcap%
\pgfsetmiterjoin%
\definecolor{currentfill}{rgb}{1.000000,1.000000,1.000000}%
\pgfsetfillcolor{currentfill}%
\pgfsetlinewidth{0.000000pt}%
\definecolor{currentstroke}{rgb}{1.000000,1.000000,1.000000}%
\pgfsetstrokecolor{currentstroke}%
\pgfsetdash{}{0pt}%
\pgfpathmoveto{\pgfqpoint{0.000000in}{0.000000in}}%
\pgfpathlineto{\pgfqpoint{6.490000in}{0.000000in}}%
\pgfpathlineto{\pgfqpoint{6.490000in}{3.740000in}}%
\pgfpathlineto{\pgfqpoint{0.000000in}{3.740000in}}%
\pgfpathlineto{\pgfqpoint{0.000000in}{0.000000in}}%
\pgfpathclose%
\pgfusepath{fill}%
\end{pgfscope}%
\begin{pgfscope}%
\pgfsetbuttcap%
\pgfsetmiterjoin%
\definecolor{currentfill}{rgb}{1.000000,1.000000,1.000000}%
\pgfsetfillcolor{currentfill}%
\pgfsetlinewidth{0.000000pt}%
\definecolor{currentstroke}{rgb}{0.000000,0.000000,0.000000}%
\pgfsetstrokecolor{currentstroke}%
\pgfsetstrokeopacity{0.000000}%
\pgfsetdash{}{0pt}%
\pgfpathmoveto{\pgfqpoint{0.695253in}{0.456901in}}%
\pgfpathlineto{\pgfqpoint{6.440000in}{0.456901in}}%
\pgfpathlineto{\pgfqpoint{6.440000in}{3.273333in}}%
\pgfpathlineto{\pgfqpoint{0.695253in}{3.273333in}}%
\pgfpathlineto{\pgfqpoint{0.695253in}{0.456901in}}%
\pgfpathclose%
\pgfusepath{fill}%
\end{pgfscope}%
\begin{pgfscope}%
\pgfsetbuttcap%
\pgfsetroundjoin%
\definecolor{currentfill}{rgb}{1.000000,1.000000,1.000000}%
\pgfsetfillcolor{currentfill}%
\pgfsetlinewidth{0.803000pt}%
\definecolor{currentstroke}{rgb}{0.000000,0.000000,0.000000}%
\pgfsetstrokecolor{currentstroke}%
\pgfsetdash{}{0pt}%
\pgfsys@defobject{currentmarker}{\pgfqpoint{0.000000in}{0.000000in}}{\pgfqpoint{0.000000in}{0.048611in}}{%
\pgfpathmoveto{\pgfqpoint{0.000000in}{0.000000in}}%
\pgfpathlineto{\pgfqpoint{0.000000in}{0.048611in}}%
\pgfusepath{stroke,fill}%
}%
\begin{pgfscope}%
\pgfsys@transformshift{0.768866in}{0.456901in}%
\pgfsys@useobject{currentmarker}{}%
\end{pgfscope}%
\end{pgfscope}%
\begin{pgfscope}%
\definecolor{textcolor}{rgb}{0.000000,0.000000,0.000000}%
\pgfsetstrokecolor{textcolor}%
\pgfsetfillcolor{textcolor}%
\pgftext[x=0.768866in,y=0.408290in,,top]{\color{textcolor}{\rmfamily\fontsize{12.000000}{14.400000}\selectfont\catcode`\^=\active\def^{\ifmmode\sp\else\^{}\fi}\catcode`\%=\active\def%{\%}$0$}}%
\end{pgfscope}%
\begin{pgfscope}%
\pgfsetbuttcap%
\pgfsetroundjoin%
\definecolor{currentfill}{rgb}{1.000000,1.000000,1.000000}%
\pgfsetfillcolor{currentfill}%
\pgfsetlinewidth{0.803000pt}%
\definecolor{currentstroke}{rgb}{0.000000,0.000000,0.000000}%
\pgfsetstrokecolor{currentstroke}%
\pgfsetdash{}{0pt}%
\pgfsys@defobject{currentmarker}{\pgfqpoint{0.000000in}{0.000000in}}{\pgfqpoint{0.000000in}{0.048611in}}{%
\pgfpathmoveto{\pgfqpoint{0.000000in}{0.000000in}}%
\pgfpathlineto{\pgfqpoint{0.000000in}{0.048611in}}%
\pgfusepath{stroke,fill}%
}%
\begin{pgfscope}%
\pgfsys@transformshift{1.750881in}{0.456901in}%
\pgfsys@useobject{currentmarker}{}%
\end{pgfscope}%
\end{pgfscope}%
\begin{pgfscope}%
\definecolor{textcolor}{rgb}{0.000000,0.000000,0.000000}%
\pgfsetstrokecolor{textcolor}%
\pgfsetfillcolor{textcolor}%
\pgftext[x=1.750881in,y=0.408290in,,top]{\color{textcolor}{\rmfamily\fontsize{12.000000}{14.400000}\selectfont\catcode`\^=\active\def^{\ifmmode\sp\else\^{}\fi}\catcode`\%=\active\def%{\%}$1\omega_1$}}%
\end{pgfscope}%
\begin{pgfscope}%
\pgfsetbuttcap%
\pgfsetroundjoin%
\definecolor{currentfill}{rgb}{1.000000,1.000000,1.000000}%
\pgfsetfillcolor{currentfill}%
\pgfsetlinewidth{0.803000pt}%
\definecolor{currentstroke}{rgb}{0.000000,0.000000,0.000000}%
\pgfsetstrokecolor{currentstroke}%
\pgfsetdash{}{0pt}%
\pgfsys@defobject{currentmarker}{\pgfqpoint{0.000000in}{0.000000in}}{\pgfqpoint{0.000000in}{0.048611in}}{%
\pgfpathmoveto{\pgfqpoint{0.000000in}{0.000000in}}%
\pgfpathlineto{\pgfqpoint{0.000000in}{0.048611in}}%
\pgfusepath{stroke,fill}%
}%
\begin{pgfscope}%
\pgfsys@transformshift{2.732895in}{0.456901in}%
\pgfsys@useobject{currentmarker}{}%
\end{pgfscope}%
\end{pgfscope}%
\begin{pgfscope}%
\definecolor{textcolor}{rgb}{0.000000,0.000000,0.000000}%
\pgfsetstrokecolor{textcolor}%
\pgfsetfillcolor{textcolor}%
\pgftext[x=2.732895in,y=0.408290in,,top]{\color{textcolor}{\rmfamily\fontsize{12.000000}{14.400000}\selectfont\catcode`\^=\active\def^{\ifmmode\sp\else\^{}\fi}\catcode`\%=\active\def%{\%}$2\omega_1$}}%
\end{pgfscope}%
\begin{pgfscope}%
\pgfsetbuttcap%
\pgfsetroundjoin%
\definecolor{currentfill}{rgb}{1.000000,1.000000,1.000000}%
\pgfsetfillcolor{currentfill}%
\pgfsetlinewidth{0.803000pt}%
\definecolor{currentstroke}{rgb}{0.000000,0.000000,0.000000}%
\pgfsetstrokecolor{currentstroke}%
\pgfsetdash{}{0pt}%
\pgfsys@defobject{currentmarker}{\pgfqpoint{0.000000in}{0.000000in}}{\pgfqpoint{0.000000in}{0.048611in}}{%
\pgfpathmoveto{\pgfqpoint{0.000000in}{0.000000in}}%
\pgfpathlineto{\pgfqpoint{0.000000in}{0.048611in}}%
\pgfusepath{stroke,fill}%
}%
\begin{pgfscope}%
\pgfsys@transformshift{3.714910in}{0.456901in}%
\pgfsys@useobject{currentmarker}{}%
\end{pgfscope}%
\end{pgfscope}%
\begin{pgfscope}%
\definecolor{textcolor}{rgb}{0.000000,0.000000,0.000000}%
\pgfsetstrokecolor{textcolor}%
\pgfsetfillcolor{textcolor}%
\pgftext[x=3.714910in,y=0.408290in,,top]{\color{textcolor}{\rmfamily\fontsize{12.000000}{14.400000}\selectfont\catcode`\^=\active\def^{\ifmmode\sp\else\^{}\fi}\catcode`\%=\active\def%{\%}$3\omega_1$}}%
\end{pgfscope}%
\begin{pgfscope}%
\pgfsetbuttcap%
\pgfsetroundjoin%
\definecolor{currentfill}{rgb}{1.000000,1.000000,1.000000}%
\pgfsetfillcolor{currentfill}%
\pgfsetlinewidth{0.803000pt}%
\definecolor{currentstroke}{rgb}{0.000000,0.000000,0.000000}%
\pgfsetstrokecolor{currentstroke}%
\pgfsetdash{}{0pt}%
\pgfsys@defobject{currentmarker}{\pgfqpoint{0.000000in}{0.000000in}}{\pgfqpoint{0.000000in}{0.048611in}}{%
\pgfpathmoveto{\pgfqpoint{0.000000in}{0.000000in}}%
\pgfpathlineto{\pgfqpoint{0.000000in}{0.048611in}}%
\pgfusepath{stroke,fill}%
}%
\begin{pgfscope}%
\pgfsys@transformshift{4.696924in}{0.456901in}%
\pgfsys@useobject{currentmarker}{}%
\end{pgfscope}%
\end{pgfscope}%
\begin{pgfscope}%
\definecolor{textcolor}{rgb}{0.000000,0.000000,0.000000}%
\pgfsetstrokecolor{textcolor}%
\pgfsetfillcolor{textcolor}%
\pgftext[x=4.696924in,y=0.408290in,,top]{\color{textcolor}{\rmfamily\fontsize{12.000000}{14.400000}\selectfont\catcode`\^=\active\def^{\ifmmode\sp\else\^{}\fi}\catcode`\%=\active\def%{\%}$4\omega_1$}}%
\end{pgfscope}%
\begin{pgfscope}%
\pgfsetbuttcap%
\pgfsetroundjoin%
\definecolor{currentfill}{rgb}{1.000000,1.000000,1.000000}%
\pgfsetfillcolor{currentfill}%
\pgfsetlinewidth{0.803000pt}%
\definecolor{currentstroke}{rgb}{0.000000,0.000000,0.000000}%
\pgfsetstrokecolor{currentstroke}%
\pgfsetdash{}{0pt}%
\pgfsys@defobject{currentmarker}{\pgfqpoint{0.000000in}{0.000000in}}{\pgfqpoint{0.000000in}{0.048611in}}{%
\pgfpathmoveto{\pgfqpoint{0.000000in}{0.000000in}}%
\pgfpathlineto{\pgfqpoint{0.000000in}{0.048611in}}%
\pgfusepath{stroke,fill}%
}%
\begin{pgfscope}%
\pgfsys@transformshift{5.678939in}{0.456901in}%
\pgfsys@useobject{currentmarker}{}%
\end{pgfscope}%
\end{pgfscope}%
\begin{pgfscope}%
\definecolor{textcolor}{rgb}{0.000000,0.000000,0.000000}%
\pgfsetstrokecolor{textcolor}%
\pgfsetfillcolor{textcolor}%
\pgftext[x=5.678939in,y=0.408290in,,top]{\color{textcolor}{\rmfamily\fontsize{12.000000}{14.400000}\selectfont\catcode`\^=\active\def^{\ifmmode\sp\else\^{}\fi}\catcode`\%=\active\def%{\%}$5\omega_1$}}%
\end{pgfscope}%
\begin{pgfscope}%
\pgfpathrectangle{\pgfqpoint{0.695253in}{0.456901in}}{\pgfqpoint{5.744747in}{2.816432in}}%
\pgfusepath{clip}%
\pgfsetbuttcap%
\pgfsetroundjoin%
\pgfsetlinewidth{0.501875pt}%
\definecolor{currentstroke}{rgb}{0.501961,0.501961,0.501961}%
\pgfsetstrokecolor{currentstroke}%
\pgfsetstrokeopacity{0.400000}%
\pgfsetdash{{1.850000pt}{0.800000pt}}{0.000000pt}%
\pgfpathmoveto{\pgfqpoint{1.014370in}{0.456901in}}%
\pgfpathlineto{\pgfqpoint{1.014370in}{3.273333in}}%
\pgfusepath{stroke}%
\end{pgfscope}%
\begin{pgfscope}%
\pgfsetbuttcap%
\pgfsetroundjoin%
\definecolor{currentfill}{rgb}{1.000000,1.000000,1.000000}%
\pgfsetfillcolor{currentfill}%
\pgfsetlinewidth{0.602250pt}%
\definecolor{currentstroke}{rgb}{0.000000,0.000000,0.000000}%
\pgfsetstrokecolor{currentstroke}%
\pgfsetdash{}{0pt}%
\pgfsys@defobject{currentmarker}{\pgfqpoint{0.000000in}{0.000000in}}{\pgfqpoint{0.000000in}{0.027778in}}{%
\pgfpathmoveto{\pgfqpoint{0.000000in}{0.000000in}}%
\pgfpathlineto{\pgfqpoint{0.000000in}{0.027778in}}%
\pgfusepath{stroke,fill}%
}%
\begin{pgfscope}%
\pgfsys@transformshift{1.014370in}{0.456901in}%
\pgfsys@useobject{currentmarker}{}%
\end{pgfscope}%
\end{pgfscope}%
\begin{pgfscope}%
\pgfpathrectangle{\pgfqpoint{0.695253in}{0.456901in}}{\pgfqpoint{5.744747in}{2.816432in}}%
\pgfusepath{clip}%
\pgfsetbuttcap%
\pgfsetroundjoin%
\pgfsetlinewidth{0.501875pt}%
\definecolor{currentstroke}{rgb}{0.501961,0.501961,0.501961}%
\pgfsetstrokecolor{currentstroke}%
\pgfsetstrokeopacity{0.400000}%
\pgfsetdash{{1.850000pt}{0.800000pt}}{0.000000pt}%
\pgfpathmoveto{\pgfqpoint{1.259874in}{0.456901in}}%
\pgfpathlineto{\pgfqpoint{1.259874in}{3.273333in}}%
\pgfusepath{stroke}%
\end{pgfscope}%
\begin{pgfscope}%
\pgfsetbuttcap%
\pgfsetroundjoin%
\definecolor{currentfill}{rgb}{1.000000,1.000000,1.000000}%
\pgfsetfillcolor{currentfill}%
\pgfsetlinewidth{0.602250pt}%
\definecolor{currentstroke}{rgb}{0.000000,0.000000,0.000000}%
\pgfsetstrokecolor{currentstroke}%
\pgfsetdash{}{0pt}%
\pgfsys@defobject{currentmarker}{\pgfqpoint{0.000000in}{0.000000in}}{\pgfqpoint{0.000000in}{0.027778in}}{%
\pgfpathmoveto{\pgfqpoint{0.000000in}{0.000000in}}%
\pgfpathlineto{\pgfqpoint{0.000000in}{0.027778in}}%
\pgfusepath{stroke,fill}%
}%
\begin{pgfscope}%
\pgfsys@transformshift{1.259874in}{0.456901in}%
\pgfsys@useobject{currentmarker}{}%
\end{pgfscope}%
\end{pgfscope}%
\begin{pgfscope}%
\pgfpathrectangle{\pgfqpoint{0.695253in}{0.456901in}}{\pgfqpoint{5.744747in}{2.816432in}}%
\pgfusepath{clip}%
\pgfsetbuttcap%
\pgfsetroundjoin%
\pgfsetlinewidth{0.501875pt}%
\definecolor{currentstroke}{rgb}{0.501961,0.501961,0.501961}%
\pgfsetstrokecolor{currentstroke}%
\pgfsetstrokeopacity{0.400000}%
\pgfsetdash{{1.850000pt}{0.800000pt}}{0.000000pt}%
\pgfpathmoveto{\pgfqpoint{1.505377in}{0.456901in}}%
\pgfpathlineto{\pgfqpoint{1.505377in}{3.273333in}}%
\pgfusepath{stroke}%
\end{pgfscope}%
\begin{pgfscope}%
\pgfsetbuttcap%
\pgfsetroundjoin%
\definecolor{currentfill}{rgb}{1.000000,1.000000,1.000000}%
\pgfsetfillcolor{currentfill}%
\pgfsetlinewidth{0.602250pt}%
\definecolor{currentstroke}{rgb}{0.000000,0.000000,0.000000}%
\pgfsetstrokecolor{currentstroke}%
\pgfsetdash{}{0pt}%
\pgfsys@defobject{currentmarker}{\pgfqpoint{0.000000in}{0.000000in}}{\pgfqpoint{0.000000in}{0.027778in}}{%
\pgfpathmoveto{\pgfqpoint{0.000000in}{0.000000in}}%
\pgfpathlineto{\pgfqpoint{0.000000in}{0.027778in}}%
\pgfusepath{stroke,fill}%
}%
\begin{pgfscope}%
\pgfsys@transformshift{1.505377in}{0.456901in}%
\pgfsys@useobject{currentmarker}{}%
\end{pgfscope}%
\end{pgfscope}%
\begin{pgfscope}%
\pgfpathrectangle{\pgfqpoint{0.695253in}{0.456901in}}{\pgfqpoint{5.744747in}{2.816432in}}%
\pgfusepath{clip}%
\pgfsetbuttcap%
\pgfsetroundjoin%
\pgfsetlinewidth{0.501875pt}%
\definecolor{currentstroke}{rgb}{0.501961,0.501961,0.501961}%
\pgfsetstrokecolor{currentstroke}%
\pgfsetstrokeopacity{0.400000}%
\pgfsetdash{{1.850000pt}{0.800000pt}}{0.000000pt}%
\pgfpathmoveto{\pgfqpoint{1.996384in}{0.456901in}}%
\pgfpathlineto{\pgfqpoint{1.996384in}{3.273333in}}%
\pgfusepath{stroke}%
\end{pgfscope}%
\begin{pgfscope}%
\pgfsetbuttcap%
\pgfsetroundjoin%
\definecolor{currentfill}{rgb}{1.000000,1.000000,1.000000}%
\pgfsetfillcolor{currentfill}%
\pgfsetlinewidth{0.602250pt}%
\definecolor{currentstroke}{rgb}{0.000000,0.000000,0.000000}%
\pgfsetstrokecolor{currentstroke}%
\pgfsetdash{}{0pt}%
\pgfsys@defobject{currentmarker}{\pgfqpoint{0.000000in}{0.000000in}}{\pgfqpoint{0.000000in}{0.027778in}}{%
\pgfpathmoveto{\pgfqpoint{0.000000in}{0.000000in}}%
\pgfpathlineto{\pgfqpoint{0.000000in}{0.027778in}}%
\pgfusepath{stroke,fill}%
}%
\begin{pgfscope}%
\pgfsys@transformshift{1.996384in}{0.456901in}%
\pgfsys@useobject{currentmarker}{}%
\end{pgfscope}%
\end{pgfscope}%
\begin{pgfscope}%
\pgfpathrectangle{\pgfqpoint{0.695253in}{0.456901in}}{\pgfqpoint{5.744747in}{2.816432in}}%
\pgfusepath{clip}%
\pgfsetbuttcap%
\pgfsetroundjoin%
\pgfsetlinewidth{0.501875pt}%
\definecolor{currentstroke}{rgb}{0.501961,0.501961,0.501961}%
\pgfsetstrokecolor{currentstroke}%
\pgfsetstrokeopacity{0.400000}%
\pgfsetdash{{1.850000pt}{0.800000pt}}{0.000000pt}%
\pgfpathmoveto{\pgfqpoint{2.241888in}{0.456901in}}%
\pgfpathlineto{\pgfqpoint{2.241888in}{3.273333in}}%
\pgfusepath{stroke}%
\end{pgfscope}%
\begin{pgfscope}%
\pgfsetbuttcap%
\pgfsetroundjoin%
\definecolor{currentfill}{rgb}{1.000000,1.000000,1.000000}%
\pgfsetfillcolor{currentfill}%
\pgfsetlinewidth{0.602250pt}%
\definecolor{currentstroke}{rgb}{0.000000,0.000000,0.000000}%
\pgfsetstrokecolor{currentstroke}%
\pgfsetdash{}{0pt}%
\pgfsys@defobject{currentmarker}{\pgfqpoint{0.000000in}{0.000000in}}{\pgfqpoint{0.000000in}{0.027778in}}{%
\pgfpathmoveto{\pgfqpoint{0.000000in}{0.000000in}}%
\pgfpathlineto{\pgfqpoint{0.000000in}{0.027778in}}%
\pgfusepath{stroke,fill}%
}%
\begin{pgfscope}%
\pgfsys@transformshift{2.241888in}{0.456901in}%
\pgfsys@useobject{currentmarker}{}%
\end{pgfscope}%
\end{pgfscope}%
\begin{pgfscope}%
\pgfpathrectangle{\pgfqpoint{0.695253in}{0.456901in}}{\pgfqpoint{5.744747in}{2.816432in}}%
\pgfusepath{clip}%
\pgfsetbuttcap%
\pgfsetroundjoin%
\pgfsetlinewidth{0.501875pt}%
\definecolor{currentstroke}{rgb}{0.501961,0.501961,0.501961}%
\pgfsetstrokecolor{currentstroke}%
\pgfsetstrokeopacity{0.400000}%
\pgfsetdash{{1.850000pt}{0.800000pt}}{0.000000pt}%
\pgfpathmoveto{\pgfqpoint{2.487392in}{0.456901in}}%
\pgfpathlineto{\pgfqpoint{2.487392in}{3.273333in}}%
\pgfusepath{stroke}%
\end{pgfscope}%
\begin{pgfscope}%
\pgfsetbuttcap%
\pgfsetroundjoin%
\definecolor{currentfill}{rgb}{1.000000,1.000000,1.000000}%
\pgfsetfillcolor{currentfill}%
\pgfsetlinewidth{0.602250pt}%
\definecolor{currentstroke}{rgb}{0.000000,0.000000,0.000000}%
\pgfsetstrokecolor{currentstroke}%
\pgfsetdash{}{0pt}%
\pgfsys@defobject{currentmarker}{\pgfqpoint{0.000000in}{0.000000in}}{\pgfqpoint{0.000000in}{0.027778in}}{%
\pgfpathmoveto{\pgfqpoint{0.000000in}{0.000000in}}%
\pgfpathlineto{\pgfqpoint{0.000000in}{0.027778in}}%
\pgfusepath{stroke,fill}%
}%
\begin{pgfscope}%
\pgfsys@transformshift{2.487392in}{0.456901in}%
\pgfsys@useobject{currentmarker}{}%
\end{pgfscope}%
\end{pgfscope}%
\begin{pgfscope}%
\pgfpathrectangle{\pgfqpoint{0.695253in}{0.456901in}}{\pgfqpoint{5.744747in}{2.816432in}}%
\pgfusepath{clip}%
\pgfsetbuttcap%
\pgfsetroundjoin%
\pgfsetlinewidth{0.501875pt}%
\definecolor{currentstroke}{rgb}{0.501961,0.501961,0.501961}%
\pgfsetstrokecolor{currentstroke}%
\pgfsetstrokeopacity{0.400000}%
\pgfsetdash{{1.850000pt}{0.800000pt}}{0.000000pt}%
\pgfpathmoveto{\pgfqpoint{2.978399in}{0.456901in}}%
\pgfpathlineto{\pgfqpoint{2.978399in}{3.273333in}}%
\pgfusepath{stroke}%
\end{pgfscope}%
\begin{pgfscope}%
\pgfsetbuttcap%
\pgfsetroundjoin%
\definecolor{currentfill}{rgb}{1.000000,1.000000,1.000000}%
\pgfsetfillcolor{currentfill}%
\pgfsetlinewidth{0.602250pt}%
\definecolor{currentstroke}{rgb}{0.000000,0.000000,0.000000}%
\pgfsetstrokecolor{currentstroke}%
\pgfsetdash{}{0pt}%
\pgfsys@defobject{currentmarker}{\pgfqpoint{0.000000in}{0.000000in}}{\pgfqpoint{0.000000in}{0.027778in}}{%
\pgfpathmoveto{\pgfqpoint{0.000000in}{0.000000in}}%
\pgfpathlineto{\pgfqpoint{0.000000in}{0.027778in}}%
\pgfusepath{stroke,fill}%
}%
\begin{pgfscope}%
\pgfsys@transformshift{2.978399in}{0.456901in}%
\pgfsys@useobject{currentmarker}{}%
\end{pgfscope}%
\end{pgfscope}%
\begin{pgfscope}%
\pgfpathrectangle{\pgfqpoint{0.695253in}{0.456901in}}{\pgfqpoint{5.744747in}{2.816432in}}%
\pgfusepath{clip}%
\pgfsetbuttcap%
\pgfsetroundjoin%
\pgfsetlinewidth{0.501875pt}%
\definecolor{currentstroke}{rgb}{0.501961,0.501961,0.501961}%
\pgfsetstrokecolor{currentstroke}%
\pgfsetstrokeopacity{0.400000}%
\pgfsetdash{{1.850000pt}{0.800000pt}}{0.000000pt}%
\pgfpathmoveto{\pgfqpoint{3.223903in}{0.456901in}}%
\pgfpathlineto{\pgfqpoint{3.223903in}{3.273333in}}%
\pgfusepath{stroke}%
\end{pgfscope}%
\begin{pgfscope}%
\pgfsetbuttcap%
\pgfsetroundjoin%
\definecolor{currentfill}{rgb}{1.000000,1.000000,1.000000}%
\pgfsetfillcolor{currentfill}%
\pgfsetlinewidth{0.602250pt}%
\definecolor{currentstroke}{rgb}{0.000000,0.000000,0.000000}%
\pgfsetstrokecolor{currentstroke}%
\pgfsetdash{}{0pt}%
\pgfsys@defobject{currentmarker}{\pgfqpoint{0.000000in}{0.000000in}}{\pgfqpoint{0.000000in}{0.027778in}}{%
\pgfpathmoveto{\pgfqpoint{0.000000in}{0.000000in}}%
\pgfpathlineto{\pgfqpoint{0.000000in}{0.027778in}}%
\pgfusepath{stroke,fill}%
}%
\begin{pgfscope}%
\pgfsys@transformshift{3.223903in}{0.456901in}%
\pgfsys@useobject{currentmarker}{}%
\end{pgfscope}%
\end{pgfscope}%
\begin{pgfscope}%
\pgfpathrectangle{\pgfqpoint{0.695253in}{0.456901in}}{\pgfqpoint{5.744747in}{2.816432in}}%
\pgfusepath{clip}%
\pgfsetbuttcap%
\pgfsetroundjoin%
\pgfsetlinewidth{0.501875pt}%
\definecolor{currentstroke}{rgb}{0.501961,0.501961,0.501961}%
\pgfsetstrokecolor{currentstroke}%
\pgfsetstrokeopacity{0.400000}%
\pgfsetdash{{1.850000pt}{0.800000pt}}{0.000000pt}%
\pgfpathmoveto{\pgfqpoint{3.469406in}{0.456901in}}%
\pgfpathlineto{\pgfqpoint{3.469406in}{3.273333in}}%
\pgfusepath{stroke}%
\end{pgfscope}%
\begin{pgfscope}%
\pgfsetbuttcap%
\pgfsetroundjoin%
\definecolor{currentfill}{rgb}{1.000000,1.000000,1.000000}%
\pgfsetfillcolor{currentfill}%
\pgfsetlinewidth{0.602250pt}%
\definecolor{currentstroke}{rgb}{0.000000,0.000000,0.000000}%
\pgfsetstrokecolor{currentstroke}%
\pgfsetdash{}{0pt}%
\pgfsys@defobject{currentmarker}{\pgfqpoint{0.000000in}{0.000000in}}{\pgfqpoint{0.000000in}{0.027778in}}{%
\pgfpathmoveto{\pgfqpoint{0.000000in}{0.000000in}}%
\pgfpathlineto{\pgfqpoint{0.000000in}{0.027778in}}%
\pgfusepath{stroke,fill}%
}%
\begin{pgfscope}%
\pgfsys@transformshift{3.469406in}{0.456901in}%
\pgfsys@useobject{currentmarker}{}%
\end{pgfscope}%
\end{pgfscope}%
\begin{pgfscope}%
\pgfpathrectangle{\pgfqpoint{0.695253in}{0.456901in}}{\pgfqpoint{5.744747in}{2.816432in}}%
\pgfusepath{clip}%
\pgfsetbuttcap%
\pgfsetroundjoin%
\pgfsetlinewidth{0.501875pt}%
\definecolor{currentstroke}{rgb}{0.501961,0.501961,0.501961}%
\pgfsetstrokecolor{currentstroke}%
\pgfsetstrokeopacity{0.400000}%
\pgfsetdash{{1.850000pt}{0.800000pt}}{0.000000pt}%
\pgfpathmoveto{\pgfqpoint{3.960413in}{0.456901in}}%
\pgfpathlineto{\pgfqpoint{3.960413in}{3.273333in}}%
\pgfusepath{stroke}%
\end{pgfscope}%
\begin{pgfscope}%
\pgfsetbuttcap%
\pgfsetroundjoin%
\definecolor{currentfill}{rgb}{1.000000,1.000000,1.000000}%
\pgfsetfillcolor{currentfill}%
\pgfsetlinewidth{0.602250pt}%
\definecolor{currentstroke}{rgb}{0.000000,0.000000,0.000000}%
\pgfsetstrokecolor{currentstroke}%
\pgfsetdash{}{0pt}%
\pgfsys@defobject{currentmarker}{\pgfqpoint{0.000000in}{0.000000in}}{\pgfqpoint{0.000000in}{0.027778in}}{%
\pgfpathmoveto{\pgfqpoint{0.000000in}{0.000000in}}%
\pgfpathlineto{\pgfqpoint{0.000000in}{0.027778in}}%
\pgfusepath{stroke,fill}%
}%
\begin{pgfscope}%
\pgfsys@transformshift{3.960413in}{0.456901in}%
\pgfsys@useobject{currentmarker}{}%
\end{pgfscope}%
\end{pgfscope}%
\begin{pgfscope}%
\pgfpathrectangle{\pgfqpoint{0.695253in}{0.456901in}}{\pgfqpoint{5.744747in}{2.816432in}}%
\pgfusepath{clip}%
\pgfsetbuttcap%
\pgfsetroundjoin%
\pgfsetlinewidth{0.501875pt}%
\definecolor{currentstroke}{rgb}{0.501961,0.501961,0.501961}%
\pgfsetstrokecolor{currentstroke}%
\pgfsetstrokeopacity{0.400000}%
\pgfsetdash{{1.850000pt}{0.800000pt}}{0.000000pt}%
\pgfpathmoveto{\pgfqpoint{4.205917in}{0.456901in}}%
\pgfpathlineto{\pgfqpoint{4.205917in}{3.273333in}}%
\pgfusepath{stroke}%
\end{pgfscope}%
\begin{pgfscope}%
\pgfsetbuttcap%
\pgfsetroundjoin%
\definecolor{currentfill}{rgb}{1.000000,1.000000,1.000000}%
\pgfsetfillcolor{currentfill}%
\pgfsetlinewidth{0.602250pt}%
\definecolor{currentstroke}{rgb}{0.000000,0.000000,0.000000}%
\pgfsetstrokecolor{currentstroke}%
\pgfsetdash{}{0pt}%
\pgfsys@defobject{currentmarker}{\pgfqpoint{0.000000in}{0.000000in}}{\pgfqpoint{0.000000in}{0.027778in}}{%
\pgfpathmoveto{\pgfqpoint{0.000000in}{0.000000in}}%
\pgfpathlineto{\pgfqpoint{0.000000in}{0.027778in}}%
\pgfusepath{stroke,fill}%
}%
\begin{pgfscope}%
\pgfsys@transformshift{4.205917in}{0.456901in}%
\pgfsys@useobject{currentmarker}{}%
\end{pgfscope}%
\end{pgfscope}%
\begin{pgfscope}%
\pgfpathrectangle{\pgfqpoint{0.695253in}{0.456901in}}{\pgfqpoint{5.744747in}{2.816432in}}%
\pgfusepath{clip}%
\pgfsetbuttcap%
\pgfsetroundjoin%
\pgfsetlinewidth{0.501875pt}%
\definecolor{currentstroke}{rgb}{0.501961,0.501961,0.501961}%
\pgfsetstrokecolor{currentstroke}%
\pgfsetstrokeopacity{0.400000}%
\pgfsetdash{{1.850000pt}{0.800000pt}}{0.000000pt}%
\pgfpathmoveto{\pgfqpoint{4.451421in}{0.456901in}}%
\pgfpathlineto{\pgfqpoint{4.451421in}{3.273333in}}%
\pgfusepath{stroke}%
\end{pgfscope}%
\begin{pgfscope}%
\pgfsetbuttcap%
\pgfsetroundjoin%
\definecolor{currentfill}{rgb}{1.000000,1.000000,1.000000}%
\pgfsetfillcolor{currentfill}%
\pgfsetlinewidth{0.602250pt}%
\definecolor{currentstroke}{rgb}{0.000000,0.000000,0.000000}%
\pgfsetstrokecolor{currentstroke}%
\pgfsetdash{}{0pt}%
\pgfsys@defobject{currentmarker}{\pgfqpoint{0.000000in}{0.000000in}}{\pgfqpoint{0.000000in}{0.027778in}}{%
\pgfpathmoveto{\pgfqpoint{0.000000in}{0.000000in}}%
\pgfpathlineto{\pgfqpoint{0.000000in}{0.027778in}}%
\pgfusepath{stroke,fill}%
}%
\begin{pgfscope}%
\pgfsys@transformshift{4.451421in}{0.456901in}%
\pgfsys@useobject{currentmarker}{}%
\end{pgfscope}%
\end{pgfscope}%
\begin{pgfscope}%
\pgfpathrectangle{\pgfqpoint{0.695253in}{0.456901in}}{\pgfqpoint{5.744747in}{2.816432in}}%
\pgfusepath{clip}%
\pgfsetbuttcap%
\pgfsetroundjoin%
\pgfsetlinewidth{0.501875pt}%
\definecolor{currentstroke}{rgb}{0.501961,0.501961,0.501961}%
\pgfsetstrokecolor{currentstroke}%
\pgfsetstrokeopacity{0.400000}%
\pgfsetdash{{1.850000pt}{0.800000pt}}{0.000000pt}%
\pgfpathmoveto{\pgfqpoint{4.942428in}{0.456901in}}%
\pgfpathlineto{\pgfqpoint{4.942428in}{3.273333in}}%
\pgfusepath{stroke}%
\end{pgfscope}%
\begin{pgfscope}%
\pgfsetbuttcap%
\pgfsetroundjoin%
\definecolor{currentfill}{rgb}{1.000000,1.000000,1.000000}%
\pgfsetfillcolor{currentfill}%
\pgfsetlinewidth{0.602250pt}%
\definecolor{currentstroke}{rgb}{0.000000,0.000000,0.000000}%
\pgfsetstrokecolor{currentstroke}%
\pgfsetdash{}{0pt}%
\pgfsys@defobject{currentmarker}{\pgfqpoint{0.000000in}{0.000000in}}{\pgfqpoint{0.000000in}{0.027778in}}{%
\pgfpathmoveto{\pgfqpoint{0.000000in}{0.000000in}}%
\pgfpathlineto{\pgfqpoint{0.000000in}{0.027778in}}%
\pgfusepath{stroke,fill}%
}%
\begin{pgfscope}%
\pgfsys@transformshift{4.942428in}{0.456901in}%
\pgfsys@useobject{currentmarker}{}%
\end{pgfscope}%
\end{pgfscope}%
\begin{pgfscope}%
\pgfpathrectangle{\pgfqpoint{0.695253in}{0.456901in}}{\pgfqpoint{5.744747in}{2.816432in}}%
\pgfusepath{clip}%
\pgfsetbuttcap%
\pgfsetroundjoin%
\pgfsetlinewidth{0.501875pt}%
\definecolor{currentstroke}{rgb}{0.501961,0.501961,0.501961}%
\pgfsetstrokecolor{currentstroke}%
\pgfsetstrokeopacity{0.400000}%
\pgfsetdash{{1.850000pt}{0.800000pt}}{0.000000pt}%
\pgfpathmoveto{\pgfqpoint{5.187932in}{0.456901in}}%
\pgfpathlineto{\pgfqpoint{5.187932in}{3.273333in}}%
\pgfusepath{stroke}%
\end{pgfscope}%
\begin{pgfscope}%
\pgfsetbuttcap%
\pgfsetroundjoin%
\definecolor{currentfill}{rgb}{1.000000,1.000000,1.000000}%
\pgfsetfillcolor{currentfill}%
\pgfsetlinewidth{0.602250pt}%
\definecolor{currentstroke}{rgb}{0.000000,0.000000,0.000000}%
\pgfsetstrokecolor{currentstroke}%
\pgfsetdash{}{0pt}%
\pgfsys@defobject{currentmarker}{\pgfqpoint{0.000000in}{0.000000in}}{\pgfqpoint{0.000000in}{0.027778in}}{%
\pgfpathmoveto{\pgfqpoint{0.000000in}{0.000000in}}%
\pgfpathlineto{\pgfqpoint{0.000000in}{0.027778in}}%
\pgfusepath{stroke,fill}%
}%
\begin{pgfscope}%
\pgfsys@transformshift{5.187932in}{0.456901in}%
\pgfsys@useobject{currentmarker}{}%
\end{pgfscope}%
\end{pgfscope}%
\begin{pgfscope}%
\pgfpathrectangle{\pgfqpoint{0.695253in}{0.456901in}}{\pgfqpoint{5.744747in}{2.816432in}}%
\pgfusepath{clip}%
\pgfsetbuttcap%
\pgfsetroundjoin%
\pgfsetlinewidth{0.501875pt}%
\definecolor{currentstroke}{rgb}{0.501961,0.501961,0.501961}%
\pgfsetstrokecolor{currentstroke}%
\pgfsetstrokeopacity{0.400000}%
\pgfsetdash{{1.850000pt}{0.800000pt}}{0.000000pt}%
\pgfpathmoveto{\pgfqpoint{5.433435in}{0.456901in}}%
\pgfpathlineto{\pgfqpoint{5.433435in}{3.273333in}}%
\pgfusepath{stroke}%
\end{pgfscope}%
\begin{pgfscope}%
\pgfsetbuttcap%
\pgfsetroundjoin%
\definecolor{currentfill}{rgb}{1.000000,1.000000,1.000000}%
\pgfsetfillcolor{currentfill}%
\pgfsetlinewidth{0.602250pt}%
\definecolor{currentstroke}{rgb}{0.000000,0.000000,0.000000}%
\pgfsetstrokecolor{currentstroke}%
\pgfsetdash{}{0pt}%
\pgfsys@defobject{currentmarker}{\pgfqpoint{0.000000in}{0.000000in}}{\pgfqpoint{0.000000in}{0.027778in}}{%
\pgfpathmoveto{\pgfqpoint{0.000000in}{0.000000in}}%
\pgfpathlineto{\pgfqpoint{0.000000in}{0.027778in}}%
\pgfusepath{stroke,fill}%
}%
\begin{pgfscope}%
\pgfsys@transformshift{5.433435in}{0.456901in}%
\pgfsys@useobject{currentmarker}{}%
\end{pgfscope}%
\end{pgfscope}%
\begin{pgfscope}%
\pgfpathrectangle{\pgfqpoint{0.695253in}{0.456901in}}{\pgfqpoint{5.744747in}{2.816432in}}%
\pgfusepath{clip}%
\pgfsetbuttcap%
\pgfsetroundjoin%
\pgfsetlinewidth{0.501875pt}%
\definecolor{currentstroke}{rgb}{0.501961,0.501961,0.501961}%
\pgfsetstrokecolor{currentstroke}%
\pgfsetstrokeopacity{0.400000}%
\pgfsetdash{{1.850000pt}{0.800000pt}}{0.000000pt}%
\pgfpathmoveto{\pgfqpoint{5.924442in}{0.456901in}}%
\pgfpathlineto{\pgfqpoint{5.924442in}{3.273333in}}%
\pgfusepath{stroke}%
\end{pgfscope}%
\begin{pgfscope}%
\pgfsetbuttcap%
\pgfsetroundjoin%
\definecolor{currentfill}{rgb}{1.000000,1.000000,1.000000}%
\pgfsetfillcolor{currentfill}%
\pgfsetlinewidth{0.602250pt}%
\definecolor{currentstroke}{rgb}{0.000000,0.000000,0.000000}%
\pgfsetstrokecolor{currentstroke}%
\pgfsetdash{}{0pt}%
\pgfsys@defobject{currentmarker}{\pgfqpoint{0.000000in}{0.000000in}}{\pgfqpoint{0.000000in}{0.027778in}}{%
\pgfpathmoveto{\pgfqpoint{0.000000in}{0.000000in}}%
\pgfpathlineto{\pgfqpoint{0.000000in}{0.027778in}}%
\pgfusepath{stroke,fill}%
}%
\begin{pgfscope}%
\pgfsys@transformshift{5.924442in}{0.456901in}%
\pgfsys@useobject{currentmarker}{}%
\end{pgfscope}%
\end{pgfscope}%
\begin{pgfscope}%
\pgfpathrectangle{\pgfqpoint{0.695253in}{0.456901in}}{\pgfqpoint{5.744747in}{2.816432in}}%
\pgfusepath{clip}%
\pgfsetbuttcap%
\pgfsetroundjoin%
\pgfsetlinewidth{0.501875pt}%
\definecolor{currentstroke}{rgb}{0.501961,0.501961,0.501961}%
\pgfsetstrokecolor{currentstroke}%
\pgfsetstrokeopacity{0.400000}%
\pgfsetdash{{1.850000pt}{0.800000pt}}{0.000000pt}%
\pgfpathmoveto{\pgfqpoint{6.169946in}{0.456901in}}%
\pgfpathlineto{\pgfqpoint{6.169946in}{3.273333in}}%
\pgfusepath{stroke}%
\end{pgfscope}%
\begin{pgfscope}%
\pgfsetbuttcap%
\pgfsetroundjoin%
\definecolor{currentfill}{rgb}{1.000000,1.000000,1.000000}%
\pgfsetfillcolor{currentfill}%
\pgfsetlinewidth{0.602250pt}%
\definecolor{currentstroke}{rgb}{0.000000,0.000000,0.000000}%
\pgfsetstrokecolor{currentstroke}%
\pgfsetdash{}{0pt}%
\pgfsys@defobject{currentmarker}{\pgfqpoint{0.000000in}{0.000000in}}{\pgfqpoint{0.000000in}{0.027778in}}{%
\pgfpathmoveto{\pgfqpoint{0.000000in}{0.000000in}}%
\pgfpathlineto{\pgfqpoint{0.000000in}{0.027778in}}%
\pgfusepath{stroke,fill}%
}%
\begin{pgfscope}%
\pgfsys@transformshift{6.169946in}{0.456901in}%
\pgfsys@useobject{currentmarker}{}%
\end{pgfscope}%
\end{pgfscope}%
\begin{pgfscope}%
\pgfpathrectangle{\pgfqpoint{0.695253in}{0.456901in}}{\pgfqpoint{5.744747in}{2.816432in}}%
\pgfusepath{clip}%
\pgfsetbuttcap%
\pgfsetroundjoin%
\pgfsetlinewidth{0.501875pt}%
\definecolor{currentstroke}{rgb}{0.501961,0.501961,0.501961}%
\pgfsetstrokecolor{currentstroke}%
\pgfsetstrokeopacity{0.400000}%
\pgfsetdash{{1.850000pt}{0.800000pt}}{0.000000pt}%
\pgfpathmoveto{\pgfqpoint{6.415450in}{0.456901in}}%
\pgfpathlineto{\pgfqpoint{6.415450in}{3.273333in}}%
\pgfusepath{stroke}%
\end{pgfscope}%
\begin{pgfscope}%
\pgfsetbuttcap%
\pgfsetroundjoin%
\definecolor{currentfill}{rgb}{1.000000,1.000000,1.000000}%
\pgfsetfillcolor{currentfill}%
\pgfsetlinewidth{0.602250pt}%
\definecolor{currentstroke}{rgb}{0.000000,0.000000,0.000000}%
\pgfsetstrokecolor{currentstroke}%
\pgfsetdash{}{0pt}%
\pgfsys@defobject{currentmarker}{\pgfqpoint{0.000000in}{0.000000in}}{\pgfqpoint{0.000000in}{0.027778in}}{%
\pgfpathmoveto{\pgfqpoint{0.000000in}{0.000000in}}%
\pgfpathlineto{\pgfqpoint{0.000000in}{0.027778in}}%
\pgfusepath{stroke,fill}%
}%
\begin{pgfscope}%
\pgfsys@transformshift{6.415450in}{0.456901in}%
\pgfsys@useobject{currentmarker}{}%
\end{pgfscope}%
\end{pgfscope}%
\begin{pgfscope}%
\definecolor{textcolor}{rgb}{0.000000,0.000000,0.000000}%
\pgfsetstrokecolor{textcolor}%
\pgfsetfillcolor{textcolor}%
\pgftext[x=3.567626in,y=0.201367in,,top]{\color{textcolor}{\rmfamily\fontsize{12.000000}{14.400000}\selectfont\catcode`\^=\active\def^{\ifmmode\sp\else\^{}\fi}\catcode`\%=\active\def%{\%}$\omega$}}%
\end{pgfscope}%
\begin{pgfscope}%
\pgfsetbuttcap%
\pgfsetroundjoin%
\definecolor{currentfill}{rgb}{1.000000,1.000000,1.000000}%
\pgfsetfillcolor{currentfill}%
\pgfsetlinewidth{0.803000pt}%
\definecolor{currentstroke}{rgb}{0.000000,0.000000,0.000000}%
\pgfsetstrokecolor{currentstroke}%
\pgfsetdash{}{0pt}%
\pgfsys@defobject{currentmarker}{\pgfqpoint{0.000000in}{0.000000in}}{\pgfqpoint{0.048611in}{0.000000in}}{%
\pgfpathmoveto{\pgfqpoint{0.000000in}{0.000000in}}%
\pgfpathlineto{\pgfqpoint{0.048611in}{0.000000in}}%
\pgfusepath{stroke,fill}%
}%
\begin{pgfscope}%
\pgfsys@transformshift{0.695253in}{0.797667in}%
\pgfsys@useobject{currentmarker}{}%
\end{pgfscope}%
\end{pgfscope}%
\begin{pgfscope}%
\definecolor{textcolor}{rgb}{0.000000,0.000000,0.000000}%
\pgfsetstrokecolor{textcolor}%
\pgfsetfillcolor{textcolor}%
\pgftext[x=0.272222in, y=0.739805in, left, base]{\color{textcolor}{\rmfamily\fontsize{12.000000}{14.400000}\selectfont\catcode`\^=\active\def^{\ifmmode\sp\else\^{}\fi}\catcode`\%=\active\def%{\%}$\mathdefault{\ensuremath{-}150}$}}%
\end{pgfscope}%
\begin{pgfscope}%
\pgfsetbuttcap%
\pgfsetroundjoin%
\definecolor{currentfill}{rgb}{1.000000,1.000000,1.000000}%
\pgfsetfillcolor{currentfill}%
\pgfsetlinewidth{0.803000pt}%
\definecolor{currentstroke}{rgb}{0.000000,0.000000,0.000000}%
\pgfsetstrokecolor{currentstroke}%
\pgfsetdash{}{0pt}%
\pgfsys@defobject{currentmarker}{\pgfqpoint{0.000000in}{0.000000in}}{\pgfqpoint{0.048611in}{0.000000in}}{%
\pgfpathmoveto{\pgfqpoint{0.000000in}{0.000000in}}%
\pgfpathlineto{\pgfqpoint{0.048611in}{0.000000in}}%
\pgfusepath{stroke,fill}%
}%
\begin{pgfscope}%
\pgfsys@transformshift{0.695253in}{1.153397in}%
\pgfsys@useobject{currentmarker}{}%
\end{pgfscope}%
\end{pgfscope}%
\begin{pgfscope}%
\definecolor{textcolor}{rgb}{0.000000,0.000000,0.000000}%
\pgfsetstrokecolor{textcolor}%
\pgfsetfillcolor{textcolor}%
\pgftext[x=0.272222in, y=1.095536in, left, base]{\color{textcolor}{\rmfamily\fontsize{12.000000}{14.400000}\selectfont\catcode`\^=\active\def^{\ifmmode\sp\else\^{}\fi}\catcode`\%=\active\def%{\%}$\mathdefault{\ensuremath{-}100}$}}%
\end{pgfscope}%
\begin{pgfscope}%
\pgfsetbuttcap%
\pgfsetroundjoin%
\definecolor{currentfill}{rgb}{1.000000,1.000000,1.000000}%
\pgfsetfillcolor{currentfill}%
\pgfsetlinewidth{0.803000pt}%
\definecolor{currentstroke}{rgb}{0.000000,0.000000,0.000000}%
\pgfsetstrokecolor{currentstroke}%
\pgfsetdash{}{0pt}%
\pgfsys@defobject{currentmarker}{\pgfqpoint{0.000000in}{0.000000in}}{\pgfqpoint{0.048611in}{0.000000in}}{%
\pgfpathmoveto{\pgfqpoint{0.000000in}{0.000000in}}%
\pgfpathlineto{\pgfqpoint{0.048611in}{0.000000in}}%
\pgfusepath{stroke,fill}%
}%
\begin{pgfscope}%
\pgfsys@transformshift{0.695253in}{1.509128in}%
\pgfsys@useobject{currentmarker}{}%
\end{pgfscope}%
\end{pgfscope}%
\begin{pgfscope}%
\definecolor{textcolor}{rgb}{0.000000,0.000000,0.000000}%
\pgfsetstrokecolor{textcolor}%
\pgfsetfillcolor{textcolor}%
\pgftext[x=0.353819in, y=1.451266in, left, base]{\color{textcolor}{\rmfamily\fontsize{12.000000}{14.400000}\selectfont\catcode`\^=\active\def^{\ifmmode\sp\else\^{}\fi}\catcode`\%=\active\def%{\%}$\mathdefault{\ensuremath{-}50}$}}%
\end{pgfscope}%
\begin{pgfscope}%
\pgfsetbuttcap%
\pgfsetroundjoin%
\definecolor{currentfill}{rgb}{1.000000,1.000000,1.000000}%
\pgfsetfillcolor{currentfill}%
\pgfsetlinewidth{0.803000pt}%
\definecolor{currentstroke}{rgb}{0.000000,0.000000,0.000000}%
\pgfsetstrokecolor{currentstroke}%
\pgfsetdash{}{0pt}%
\pgfsys@defobject{currentmarker}{\pgfqpoint{0.000000in}{0.000000in}}{\pgfqpoint{0.048611in}{0.000000in}}{%
\pgfpathmoveto{\pgfqpoint{0.000000in}{0.000000in}}%
\pgfpathlineto{\pgfqpoint{0.048611in}{0.000000in}}%
\pgfusepath{stroke,fill}%
}%
\begin{pgfscope}%
\pgfsys@transformshift{0.695253in}{1.864858in}%
\pgfsys@useobject{currentmarker}{}%
\end{pgfscope}%
\end{pgfscope}%
\begin{pgfscope}%
\definecolor{textcolor}{rgb}{0.000000,0.000000,0.000000}%
\pgfsetstrokecolor{textcolor}%
\pgfsetfillcolor{textcolor}%
\pgftext[x=0.565045in, y=1.806997in, left, base]{\color{textcolor}{\rmfamily\fontsize{12.000000}{14.400000}\selectfont\catcode`\^=\active\def^{\ifmmode\sp\else\^{}\fi}\catcode`\%=\active\def%{\%}$\mathdefault{0}$}}%
\end{pgfscope}%
\begin{pgfscope}%
\pgfsetbuttcap%
\pgfsetroundjoin%
\definecolor{currentfill}{rgb}{1.000000,1.000000,1.000000}%
\pgfsetfillcolor{currentfill}%
\pgfsetlinewidth{0.803000pt}%
\definecolor{currentstroke}{rgb}{0.000000,0.000000,0.000000}%
\pgfsetstrokecolor{currentstroke}%
\pgfsetdash{}{0pt}%
\pgfsys@defobject{currentmarker}{\pgfqpoint{0.000000in}{0.000000in}}{\pgfqpoint{0.048611in}{0.000000in}}{%
\pgfpathmoveto{\pgfqpoint{0.000000in}{0.000000in}}%
\pgfpathlineto{\pgfqpoint{0.048611in}{0.000000in}}%
\pgfusepath{stroke,fill}%
}%
\begin{pgfscope}%
\pgfsys@transformshift{0.695253in}{2.220589in}%
\pgfsys@useobject{currentmarker}{}%
\end{pgfscope}%
\end{pgfscope}%
\begin{pgfscope}%
\definecolor{textcolor}{rgb}{0.000000,0.000000,0.000000}%
\pgfsetstrokecolor{textcolor}%
\pgfsetfillcolor{textcolor}%
\pgftext[x=0.483449in, y=2.162727in, left, base]{\color{textcolor}{\rmfamily\fontsize{12.000000}{14.400000}\selectfont\catcode`\^=\active\def^{\ifmmode\sp\else\^{}\fi}\catcode`\%=\active\def%{\%}$\mathdefault{50}$}}%
\end{pgfscope}%
\begin{pgfscope}%
\pgfsetbuttcap%
\pgfsetroundjoin%
\definecolor{currentfill}{rgb}{1.000000,1.000000,1.000000}%
\pgfsetfillcolor{currentfill}%
\pgfsetlinewidth{0.803000pt}%
\definecolor{currentstroke}{rgb}{0.000000,0.000000,0.000000}%
\pgfsetstrokecolor{currentstroke}%
\pgfsetdash{}{0pt}%
\pgfsys@defobject{currentmarker}{\pgfqpoint{0.000000in}{0.000000in}}{\pgfqpoint{0.048611in}{0.000000in}}{%
\pgfpathmoveto{\pgfqpoint{0.000000in}{0.000000in}}%
\pgfpathlineto{\pgfqpoint{0.048611in}{0.000000in}}%
\pgfusepath{stroke,fill}%
}%
\begin{pgfscope}%
\pgfsys@transformshift{0.695253in}{2.576319in}%
\pgfsys@useobject{currentmarker}{}%
\end{pgfscope}%
\end{pgfscope}%
\begin{pgfscope}%
\definecolor{textcolor}{rgb}{0.000000,0.000000,0.000000}%
\pgfsetstrokecolor{textcolor}%
\pgfsetfillcolor{textcolor}%
\pgftext[x=0.401852in, y=2.518458in, left, base]{\color{textcolor}{\rmfamily\fontsize{12.000000}{14.400000}\selectfont\catcode`\^=\active\def^{\ifmmode\sp\else\^{}\fi}\catcode`\%=\active\def%{\%}$\mathdefault{100}$}}%
\end{pgfscope}%
\begin{pgfscope}%
\pgfsetbuttcap%
\pgfsetroundjoin%
\definecolor{currentfill}{rgb}{1.000000,1.000000,1.000000}%
\pgfsetfillcolor{currentfill}%
\pgfsetlinewidth{0.803000pt}%
\definecolor{currentstroke}{rgb}{0.000000,0.000000,0.000000}%
\pgfsetstrokecolor{currentstroke}%
\pgfsetdash{}{0pt}%
\pgfsys@defobject{currentmarker}{\pgfqpoint{0.000000in}{0.000000in}}{\pgfqpoint{0.048611in}{0.000000in}}{%
\pgfpathmoveto{\pgfqpoint{0.000000in}{0.000000in}}%
\pgfpathlineto{\pgfqpoint{0.048611in}{0.000000in}}%
\pgfusepath{stroke,fill}%
}%
\begin{pgfscope}%
\pgfsys@transformshift{0.695253in}{2.932050in}%
\pgfsys@useobject{currentmarker}{}%
\end{pgfscope}%
\end{pgfscope}%
\begin{pgfscope}%
\definecolor{textcolor}{rgb}{0.000000,0.000000,0.000000}%
\pgfsetstrokecolor{textcolor}%
\pgfsetfillcolor{textcolor}%
\pgftext[x=0.401852in, y=2.874188in, left, base]{\color{textcolor}{\rmfamily\fontsize{12.000000}{14.400000}\selectfont\catcode`\^=\active\def^{\ifmmode\sp\else\^{}\fi}\catcode`\%=\active\def%{\%}$\mathdefault{150}$}}%
\end{pgfscope}%
\begin{pgfscope}%
\definecolor{textcolor}{rgb}{0.000000,0.000000,0.000000}%
\pgfsetstrokecolor{textcolor}%
\pgfsetfillcolor{textcolor}%
\pgftext[x=0.216667in,y=1.865117in,,bottom,rotate=90.000000]{\color{textcolor}{\rmfamily\fontsize{12.000000}{14.400000}\selectfont\catcode`\^=\active\def^{\ifmmode\sp\else\^{}\fi}\catcode`\%=\active\def%{\%}$\Phi_{\mathrm{вых}}(\omega), ^\circ$}}%
\end{pgfscope}%
\begin{pgfscope}%
\pgfpathrectangle{\pgfqpoint{0.695253in}{0.456901in}}{\pgfqpoint{5.744747in}{2.816432in}}%
\pgfusepath{clip}%
\pgfsetbuttcap%
\pgfsetroundjoin%
\pgfsetlinewidth{2.007500pt}%
\definecolor{currentstroke}{rgb}{1.000000,0.000000,0.000000}%
\pgfsetstrokecolor{currentstroke}%
\pgfsetdash{{7.400000pt}{3.200000pt}}{0.000000pt}%
\pgfpathmoveto{\pgfqpoint{0.768866in}{2.505173in}}%
\pgfpathlineto{\pgfqpoint{2.186785in}{0.584920in}}%
\pgfpathlineto{\pgfqpoint{2.197609in}{3.131507in}}%
\pgfpathlineto{\pgfqpoint{3.193399in}{1.779716in}}%
\pgfpathlineto{\pgfqpoint{3.961889in}{0.732877in}}%
\pgfpathlineto{\pgfqpoint{4.070127in}{0.585089in}}%
\pgfpathlineto{\pgfqpoint{4.080951in}{3.131565in}}%
\pgfpathlineto{\pgfqpoint{4.708732in}{2.272213in}}%
\pgfpathlineto{\pgfqpoint{5.260746in}{1.512950in}}%
\pgfpathlineto{\pgfqpoint{5.758641in}{0.824470in}}%
\pgfpathlineto{\pgfqpoint{5.920998in}{0.599095in}}%
\pgfpathlineto{\pgfqpoint{5.931822in}{3.145314in}}%
\pgfpathlineto{\pgfqpoint{6.169946in}{2.813857in}}%
\pgfpathlineto{\pgfqpoint{6.169946in}{2.813857in}}%
\pgfusepath{stroke}%
\end{pgfscope}%
\begin{pgfscope}%
\pgfpathrectangle{\pgfqpoint{0.695253in}{0.456901in}}{\pgfqpoint{5.744747in}{2.816432in}}%
\pgfusepath{clip}%
\pgfsetbuttcap%
\pgfsetroundjoin%
\pgfsetlinewidth{2.007500pt}%
\definecolor{currentstroke}{rgb}{0.000000,0.000000,1.000000}%
\pgfsetstrokecolor{currentstroke}%
\pgfsetdash{}{0pt}%
\pgfpathmoveto{\pgfqpoint{0.768866in}{1.864858in}}%
\pgfpathlineto{\pgfqpoint{0.768866in}{2.505173in}}%
\pgfusepath{stroke}%
\end{pgfscope}%
\begin{pgfscope}%
\pgfpathrectangle{\pgfqpoint{0.695253in}{0.456901in}}{\pgfqpoint{5.744747in}{2.816432in}}%
\pgfusepath{clip}%
\pgfsetbuttcap%
\pgfsetroundjoin%
\pgfsetlinewidth{2.007500pt}%
\definecolor{currentstroke}{rgb}{0.000000,0.000000,1.000000}%
\pgfsetstrokecolor{currentstroke}%
\pgfsetdash{}{0pt}%
\pgfpathmoveto{\pgfqpoint{1.750881in}{1.864858in}}%
\pgfpathlineto{\pgfqpoint{1.750881in}{1.175595in}}%
\pgfusepath{stroke}%
\end{pgfscope}%
\begin{pgfscope}%
\pgfpathrectangle{\pgfqpoint{0.695253in}{0.456901in}}{\pgfqpoint{5.744747in}{2.816432in}}%
\pgfusepath{clip}%
\pgfsetbuttcap%
\pgfsetroundjoin%
\pgfsetlinewidth{2.007500pt}%
\definecolor{currentstroke}{rgb}{0.000000,0.000000,1.000000}%
\pgfsetstrokecolor{currentstroke}%
\pgfsetdash{}{0pt}%
\pgfpathmoveto{\pgfqpoint{2.732895in}{1.864858in}}%
\pgfpathlineto{\pgfqpoint{2.732895in}{2.405367in}}%
\pgfusepath{stroke}%
\end{pgfscope}%
\begin{pgfscope}%
\pgfpathrectangle{\pgfqpoint{0.695253in}{0.456901in}}{\pgfqpoint{5.744747in}{2.816432in}}%
\pgfusepath{clip}%
\pgfsetbuttcap%
\pgfsetroundjoin%
\pgfsetlinewidth{2.007500pt}%
\definecolor{currentstroke}{rgb}{0.000000,0.000000,1.000000}%
\pgfsetstrokecolor{currentstroke}%
\pgfsetdash{}{0pt}%
\pgfpathmoveto{\pgfqpoint{3.714910in}{1.864858in}}%
\pgfpathlineto{\pgfqpoint{3.714910in}{1.069755in}}%
\pgfusepath{stroke}%
\end{pgfscope}%
\begin{pgfscope}%
\pgfpathrectangle{\pgfqpoint{0.695253in}{0.456901in}}{\pgfqpoint{5.744747in}{2.816432in}}%
\pgfusepath{clip}%
\pgfsetbuttcap%
\pgfsetroundjoin%
\pgfsetlinewidth{2.007500pt}%
\definecolor{currentstroke}{rgb}{0.000000,0.000000,1.000000}%
\pgfsetstrokecolor{currentstroke}%
\pgfsetdash{}{0pt}%
\pgfpathmoveto{\pgfqpoint{4.696924in}{1.864858in}}%
\pgfpathlineto{\pgfqpoint{4.696924in}{2.288414in}}%
\pgfusepath{stroke}%
\end{pgfscope}%
\begin{pgfscope}%
\pgfpathrectangle{\pgfqpoint{0.695253in}{0.456901in}}{\pgfqpoint{5.744747in}{2.816432in}}%
\pgfusepath{clip}%
\pgfsetbuttcap%
\pgfsetroundjoin%
\pgfsetlinewidth{2.007500pt}%
\definecolor{currentstroke}{rgb}{0.000000,0.000000,1.000000}%
\pgfsetstrokecolor{currentstroke}%
\pgfsetdash{}{0pt}%
\pgfpathmoveto{\pgfqpoint{5.678939in}{1.864858in}}%
\pgfpathlineto{\pgfqpoint{5.678939in}{0.934944in}}%
\pgfusepath{stroke}%
\end{pgfscope}%
\begin{pgfscope}%
\pgfpathrectangle{\pgfqpoint{0.695253in}{0.456901in}}{\pgfqpoint{5.744747in}{2.816432in}}%
\pgfusepath{clip}%
\pgfsetbuttcap%
\pgfsetroundjoin%
\definecolor{currentfill}{rgb}{1.000000,1.000000,1.000000}%
\pgfsetfillcolor{currentfill}%
\pgfsetlinewidth{1.505625pt}%
\definecolor{currentstroke}{rgb}{0.000000,0.000000,1.000000}%
\pgfsetstrokecolor{currentstroke}%
\pgfsetdash{}{0pt}%
\pgfsys@defobject{currentmarker}{\pgfqpoint{-0.027778in}{-0.027778in}}{\pgfqpoint{0.027778in}{0.027778in}}{%
\pgfpathmoveto{\pgfqpoint{0.000000in}{-0.027778in}}%
\pgfpathcurveto{\pgfqpoint{0.007367in}{-0.027778in}}{\pgfqpoint{0.014433in}{-0.024851in}}{\pgfqpoint{0.019642in}{-0.019642in}}%
\pgfpathcurveto{\pgfqpoint{0.024851in}{-0.014433in}}{\pgfqpoint{0.027778in}{-0.007367in}}{\pgfqpoint{0.027778in}{0.000000in}}%
\pgfpathcurveto{\pgfqpoint{0.027778in}{0.007367in}}{\pgfqpoint{0.024851in}{0.014433in}}{\pgfqpoint{0.019642in}{0.019642in}}%
\pgfpathcurveto{\pgfqpoint{0.014433in}{0.024851in}}{\pgfqpoint{0.007367in}{0.027778in}}{\pgfqpoint{0.000000in}{0.027778in}}%
\pgfpathcurveto{\pgfqpoint{-0.007367in}{0.027778in}}{\pgfqpoint{-0.014433in}{0.024851in}}{\pgfqpoint{-0.019642in}{0.019642in}}%
\pgfpathcurveto{\pgfqpoint{-0.024851in}{0.014433in}}{\pgfqpoint{-0.027778in}{0.007367in}}{\pgfqpoint{-0.027778in}{0.000000in}}%
\pgfpathcurveto{\pgfqpoint{-0.027778in}{-0.007367in}}{\pgfqpoint{-0.024851in}{-0.014433in}}{\pgfqpoint{-0.019642in}{-0.019642in}}%
\pgfpathcurveto{\pgfqpoint{-0.014433in}{-0.024851in}}{\pgfqpoint{-0.007367in}{-0.027778in}}{\pgfqpoint{0.000000in}{-0.027778in}}%
\pgfpathlineto{\pgfqpoint{0.000000in}{-0.027778in}}%
\pgfpathclose%
\pgfusepath{stroke,fill}%
}%
\begin{pgfscope}%
\pgfsys@transformshift{0.768866in}{2.505173in}%
\pgfsys@useobject{currentmarker}{}%
\end{pgfscope}%
\begin{pgfscope}%
\pgfsys@transformshift{1.750881in}{1.175595in}%
\pgfsys@useobject{currentmarker}{}%
\end{pgfscope}%
\begin{pgfscope}%
\pgfsys@transformshift{2.732895in}{2.405367in}%
\pgfsys@useobject{currentmarker}{}%
\end{pgfscope}%
\begin{pgfscope}%
\pgfsys@transformshift{3.714910in}{1.069755in}%
\pgfsys@useobject{currentmarker}{}%
\end{pgfscope}%
\begin{pgfscope}%
\pgfsys@transformshift{4.696924in}{2.288414in}%
\pgfsys@useobject{currentmarker}{}%
\end{pgfscope}%
\begin{pgfscope}%
\pgfsys@transformshift{5.678939in}{0.934944in}%
\pgfsys@useobject{currentmarker}{}%
\end{pgfscope}%
\end{pgfscope}%
\begin{pgfscope}%
\pgfpathrectangle{\pgfqpoint{0.695253in}{0.456901in}}{\pgfqpoint{5.744747in}{2.816432in}}%
\pgfusepath{clip}%
\pgfsetrectcap%
\pgfsetroundjoin%
\pgfsetlinewidth{2.007500pt}%
\definecolor{currentstroke}{rgb}{0.000000,0.000000,0.000000}%
\pgfsetstrokecolor{currentstroke}%
\pgfsetdash{}{0pt}%
\pgfpathmoveto{\pgfqpoint{0.768866in}{1.864858in}}%
\pgfpathlineto{\pgfqpoint{5.678939in}{1.864858in}}%
\pgfusepath{stroke}%
\end{pgfscope}%
\begin{pgfscope}%
\pgfsetrectcap%
\pgfsetmiterjoin%
\pgfsetlinewidth{0.602250pt}%
\definecolor{currentstroke}{rgb}{0.000000,0.000000,0.000000}%
\pgfsetstrokecolor{currentstroke}%
\pgfsetdash{}{0pt}%
\pgfpathmoveto{\pgfqpoint{0.695253in}{0.456901in}}%
\pgfpathlineto{\pgfqpoint{0.695253in}{3.273333in}}%
\pgfusepath{stroke}%
\end{pgfscope}%
\begin{pgfscope}%
\pgfsetrectcap%
\pgfsetmiterjoin%
\pgfsetlinewidth{0.602250pt}%
\definecolor{currentstroke}{rgb}{0.000000,0.000000,0.000000}%
\pgfsetstrokecolor{currentstroke}%
\pgfsetdash{}{0pt}%
\pgfpathmoveto{\pgfqpoint{0.695253in}{0.456901in}}%
\pgfpathlineto{\pgfqpoint{6.440000in}{0.456901in}}%
\pgfusepath{stroke}%
\end{pgfscope}%
\begin{pgfscope}%
\pgfsetbuttcap%
\pgfsetroundjoin%
\pgfsetlinewidth{2.007500pt}%
\definecolor{currentstroke}{rgb}{1.000000,0.000000,0.000000}%
\pgfsetstrokecolor{currentstroke}%
\pgfsetdash{{7.400000pt}{3.200000pt}}{0.000000pt}%
\pgfpathmoveto{\pgfqpoint{2.589040in}{3.556667in}}%
\pgfpathlineto{\pgfqpoint{2.989040in}{3.556667in}}%
\pgfusepath{stroke}%
\end{pgfscope}%
\begin{pgfscope}%
\definecolor{textcolor}{rgb}{0.000000,0.000000,0.000000}%
\pgfsetstrokecolor{textcolor}%
\pgfsetfillcolor{textcolor}%
\pgftext[x=3.172373in,y=3.498333in,left,base]{\color{textcolor}{\rmfamily\fontsize{12.000000}{14.400000}\selectfont\catcode`\^=\active\def^{\ifmmode\sp\else\^{}\fi}\catcode`\%=\active\def%{\%}Огибающая $\Phi_{\mathrm{вых}}(\omega)$}}%
\end{pgfscope}%
\end{pgfpicture}%
\makeatother%
\endgroup%

    \caption{Фазовый дискретный спектр реакции на первый входной сигнал}
    \label{fig:plot_disc_spec_Phi_out_1}
\end{figure}

\subsection{Второй сигнал}

\begin{equation}
    \omega = \omega_1 k; \quad \omega_{1(2)} = \pv[.3]{ex10['w1_2']}
\end{equation}

Амплитудный и фазовый спектры реакции на второй сигнал приведены на \cref{fig:plot_disc_spec_A_out_2,fig:plot_disc_spec_Phi_out_2}.

\begin{figure}[H]
    \centering
    %% Creator: Matplotlib, PGF backend
%%
%% To include the figure in your LaTeX document, write
%%   \input{<filename>.pgf}
%%
%% Make sure the required packages are loaded in your preamble
%%   \usepackage{pgf}
%%
%% Also ensure that all the required font packages are loaded; for instance,
%% the lmodern package is sometimes necessary when using math font.
%%   \usepackage{lmodern}
%%
%% Figures using additional raster images can only be included by \input if
%% they are in the same directory as the main LaTeX file. For loading figures
%% from other directories you can use the `import` package
%%   \usepackage{import}
%%
%% and then include the figures with
%%   \import{<path to file>}{<filename>.pgf}
%%
%% Matplotlib used the following preamble
%%   \def\mathdefault#1{#1}
%%   \everymath=\expandafter{\the\everymath\displaystyle}
%%   \IfFileExists{scrextend.sty}{
%%     \usepackage[fontsize=12.000000pt]{scrextend}
%%   }{
%%     \renewcommand{\normalsize}{\fontsize{12.000000}{14.400000}\selectfont}
%%     \normalsize
%%   }
%%   \usepackage{fontspec}\setmainfont{Times New Roman}\usepackage[english,russian]{babel}\usepackage{amsmath}
%%   \ifdefined\pdftexversion\else  % non-pdftex case.
%%     \usepackage{fontspec}
%%     \setmainfont{times.ttf}[Path=\detokenize{/usr/local/share/fonts/truetype/times-new-roman/}]
%%     \setsansfont{DejaVuSans.ttf}[Path=\detokenize{/usr/share/matplotlib/mpl-data/fonts/ttf/}]
%%     \setmonofont{DejaVuSansMono.ttf}[Path=\detokenize{/usr/share/matplotlib/mpl-data/fonts/ttf/}]
%%   \fi
%%   \makeatletter\@ifpackageloaded{underscore}{}{\usepackage[strings]{underscore}}\makeatother
%%
\begingroup%
\makeatletter%
\begin{pgfpicture}%
\pgfpathrectangle{\pgfpointorigin}{\pgfqpoint{6.490000in}{3.740000in}}%
\pgfusepath{use as bounding box, clip}%
\begin{pgfscope}%
\pgfsetbuttcap%
\pgfsetmiterjoin%
\definecolor{currentfill}{rgb}{1.000000,1.000000,1.000000}%
\pgfsetfillcolor{currentfill}%
\pgfsetlinewidth{0.000000pt}%
\definecolor{currentstroke}{rgb}{1.000000,1.000000,1.000000}%
\pgfsetstrokecolor{currentstroke}%
\pgfsetdash{}{0pt}%
\pgfpathmoveto{\pgfqpoint{0.000000in}{0.000000in}}%
\pgfpathlineto{\pgfqpoint{6.490000in}{0.000000in}}%
\pgfpathlineto{\pgfqpoint{6.490000in}{3.740000in}}%
\pgfpathlineto{\pgfqpoint{0.000000in}{3.740000in}}%
\pgfpathlineto{\pgfqpoint{0.000000in}{0.000000in}}%
\pgfpathclose%
\pgfusepath{fill}%
\end{pgfscope}%
\begin{pgfscope}%
\pgfsetbuttcap%
\pgfsetmiterjoin%
\definecolor{currentfill}{rgb}{1.000000,1.000000,1.000000}%
\pgfsetfillcolor{currentfill}%
\pgfsetlinewidth{0.000000pt}%
\definecolor{currentstroke}{rgb}{0.000000,0.000000,0.000000}%
\pgfsetstrokecolor{currentstroke}%
\pgfsetstrokeopacity{0.000000}%
\pgfsetdash{}{0pt}%
\pgfpathmoveto{\pgfqpoint{0.571024in}{0.456901in}}%
\pgfpathlineto{\pgfqpoint{6.440000in}{0.456901in}}%
\pgfpathlineto{\pgfqpoint{6.440000in}{3.273333in}}%
\pgfpathlineto{\pgfqpoint{0.571024in}{3.273333in}}%
\pgfpathlineto{\pgfqpoint{0.571024in}{0.456901in}}%
\pgfpathclose%
\pgfusepath{fill}%
\end{pgfscope}%
\begin{pgfscope}%
\pgfsetbuttcap%
\pgfsetroundjoin%
\definecolor{currentfill}{rgb}{1.000000,1.000000,1.000000}%
\pgfsetfillcolor{currentfill}%
\pgfsetlinewidth{0.803000pt}%
\definecolor{currentstroke}{rgb}{0.000000,0.000000,0.000000}%
\pgfsetstrokecolor{currentstroke}%
\pgfsetdash{}{0pt}%
\pgfsys@defobject{currentmarker}{\pgfqpoint{0.000000in}{0.000000in}}{\pgfqpoint{0.000000in}{0.048611in}}{%
\pgfpathmoveto{\pgfqpoint{0.000000in}{0.000000in}}%
\pgfpathlineto{\pgfqpoint{0.000000in}{0.048611in}}%
\pgfusepath{stroke,fill}%
}%
\begin{pgfscope}%
\pgfsys@transformshift{0.601339in}{0.456901in}%
\pgfsys@useobject{currentmarker}{}%
\end{pgfscope}%
\end{pgfscope}%
\begin{pgfscope}%
\definecolor{textcolor}{rgb}{0.000000,0.000000,0.000000}%
\pgfsetstrokecolor{textcolor}%
\pgfsetfillcolor{textcolor}%
\pgftext[x=0.601339in,y=0.408290in,,top]{\color{textcolor}{\rmfamily\fontsize{12.000000}{14.400000}\selectfont\catcode`\^=\active\def^{\ifmmode\sp\else\^{}\fi}\catcode`\%=\active\def%{\%}$0$}}%
\end{pgfscope}%
\begin{pgfscope}%
\pgfsetbuttcap%
\pgfsetroundjoin%
\definecolor{currentfill}{rgb}{1.000000,1.000000,1.000000}%
\pgfsetfillcolor{currentfill}%
\pgfsetlinewidth{0.803000pt}%
\definecolor{currentstroke}{rgb}{0.000000,0.000000,0.000000}%
\pgfsetstrokecolor{currentstroke}%
\pgfsetdash{}{0pt}%
\pgfsys@defobject{currentmarker}{\pgfqpoint{0.000000in}{0.000000in}}{\pgfqpoint{0.000000in}{0.048611in}}{%
\pgfpathmoveto{\pgfqpoint{0.000000in}{0.000000in}}%
\pgfpathlineto{\pgfqpoint{0.000000in}{0.048611in}}%
\pgfusepath{stroke,fill}%
}%
\begin{pgfscope}%
\pgfsys@transformshift{1.612363in}{0.456901in}%
\pgfsys@useobject{currentmarker}{}%
\end{pgfscope}%
\end{pgfscope}%
\begin{pgfscope}%
\definecolor{textcolor}{rgb}{0.000000,0.000000,0.000000}%
\pgfsetstrokecolor{textcolor}%
\pgfsetfillcolor{textcolor}%
\pgftext[x=1.612363in,y=0.408290in,,top]{\color{textcolor}{\rmfamily\fontsize{12.000000}{14.400000}\selectfont\catcode`\^=\active\def^{\ifmmode\sp\else\^{}\fi}\catcode`\%=\active\def%{\%}$1\omega_1$}}%
\end{pgfscope}%
\begin{pgfscope}%
\pgfsetbuttcap%
\pgfsetroundjoin%
\definecolor{currentfill}{rgb}{1.000000,1.000000,1.000000}%
\pgfsetfillcolor{currentfill}%
\pgfsetlinewidth{0.803000pt}%
\definecolor{currentstroke}{rgb}{0.000000,0.000000,0.000000}%
\pgfsetstrokecolor{currentstroke}%
\pgfsetdash{}{0pt}%
\pgfsys@defobject{currentmarker}{\pgfqpoint{0.000000in}{0.000000in}}{\pgfqpoint{0.000000in}{0.048611in}}{%
\pgfpathmoveto{\pgfqpoint{0.000000in}{0.000000in}}%
\pgfpathlineto{\pgfqpoint{0.000000in}{0.048611in}}%
\pgfusepath{stroke,fill}%
}%
\begin{pgfscope}%
\pgfsys@transformshift{2.623386in}{0.456901in}%
\pgfsys@useobject{currentmarker}{}%
\end{pgfscope}%
\end{pgfscope}%
\begin{pgfscope}%
\definecolor{textcolor}{rgb}{0.000000,0.000000,0.000000}%
\pgfsetstrokecolor{textcolor}%
\pgfsetfillcolor{textcolor}%
\pgftext[x=2.623386in,y=0.408290in,,top]{\color{textcolor}{\rmfamily\fontsize{12.000000}{14.400000}\selectfont\catcode`\^=\active\def^{\ifmmode\sp\else\^{}\fi}\catcode`\%=\active\def%{\%}$2\omega_1$}}%
\end{pgfscope}%
\begin{pgfscope}%
\pgfsetbuttcap%
\pgfsetroundjoin%
\definecolor{currentfill}{rgb}{1.000000,1.000000,1.000000}%
\pgfsetfillcolor{currentfill}%
\pgfsetlinewidth{0.803000pt}%
\definecolor{currentstroke}{rgb}{0.000000,0.000000,0.000000}%
\pgfsetstrokecolor{currentstroke}%
\pgfsetdash{}{0pt}%
\pgfsys@defobject{currentmarker}{\pgfqpoint{0.000000in}{0.000000in}}{\pgfqpoint{0.000000in}{0.048611in}}{%
\pgfpathmoveto{\pgfqpoint{0.000000in}{0.000000in}}%
\pgfpathlineto{\pgfqpoint{0.000000in}{0.048611in}}%
\pgfusepath{stroke,fill}%
}%
\begin{pgfscope}%
\pgfsys@transformshift{3.634410in}{0.456901in}%
\pgfsys@useobject{currentmarker}{}%
\end{pgfscope}%
\end{pgfscope}%
\begin{pgfscope}%
\definecolor{textcolor}{rgb}{0.000000,0.000000,0.000000}%
\pgfsetstrokecolor{textcolor}%
\pgfsetfillcolor{textcolor}%
\pgftext[x=3.634410in,y=0.408290in,,top]{\color{textcolor}{\rmfamily\fontsize{12.000000}{14.400000}\selectfont\catcode`\^=\active\def^{\ifmmode\sp\else\^{}\fi}\catcode`\%=\active\def%{\%}$3\omega_1$}}%
\end{pgfscope}%
\begin{pgfscope}%
\pgfsetbuttcap%
\pgfsetroundjoin%
\definecolor{currentfill}{rgb}{1.000000,1.000000,1.000000}%
\pgfsetfillcolor{currentfill}%
\pgfsetlinewidth{0.803000pt}%
\definecolor{currentstroke}{rgb}{0.000000,0.000000,0.000000}%
\pgfsetstrokecolor{currentstroke}%
\pgfsetdash{}{0pt}%
\pgfsys@defobject{currentmarker}{\pgfqpoint{0.000000in}{0.000000in}}{\pgfqpoint{0.000000in}{0.048611in}}{%
\pgfpathmoveto{\pgfqpoint{0.000000in}{0.000000in}}%
\pgfpathlineto{\pgfqpoint{0.000000in}{0.048611in}}%
\pgfusepath{stroke,fill}%
}%
\begin{pgfscope}%
\pgfsys@transformshift{4.645433in}{0.456901in}%
\pgfsys@useobject{currentmarker}{}%
\end{pgfscope}%
\end{pgfscope}%
\begin{pgfscope}%
\definecolor{textcolor}{rgb}{0.000000,0.000000,0.000000}%
\pgfsetstrokecolor{textcolor}%
\pgfsetfillcolor{textcolor}%
\pgftext[x=4.645433in,y=0.408290in,,top]{\color{textcolor}{\rmfamily\fontsize{12.000000}{14.400000}\selectfont\catcode`\^=\active\def^{\ifmmode\sp\else\^{}\fi}\catcode`\%=\active\def%{\%}$4\omega_1$}}%
\end{pgfscope}%
\begin{pgfscope}%
\pgfsetbuttcap%
\pgfsetroundjoin%
\definecolor{currentfill}{rgb}{1.000000,1.000000,1.000000}%
\pgfsetfillcolor{currentfill}%
\pgfsetlinewidth{0.803000pt}%
\definecolor{currentstroke}{rgb}{0.000000,0.000000,0.000000}%
\pgfsetstrokecolor{currentstroke}%
\pgfsetdash{}{0pt}%
\pgfsys@defobject{currentmarker}{\pgfqpoint{0.000000in}{0.000000in}}{\pgfqpoint{0.000000in}{0.048611in}}{%
\pgfpathmoveto{\pgfqpoint{0.000000in}{0.000000in}}%
\pgfpathlineto{\pgfqpoint{0.000000in}{0.048611in}}%
\pgfusepath{stroke,fill}%
}%
\begin{pgfscope}%
\pgfsys@transformshift{5.656457in}{0.456901in}%
\pgfsys@useobject{currentmarker}{}%
\end{pgfscope}%
\end{pgfscope}%
\begin{pgfscope}%
\definecolor{textcolor}{rgb}{0.000000,0.000000,0.000000}%
\pgfsetstrokecolor{textcolor}%
\pgfsetfillcolor{textcolor}%
\pgftext[x=5.656457in,y=0.408290in,,top]{\color{textcolor}{\rmfamily\fontsize{12.000000}{14.400000}\selectfont\catcode`\^=\active\def^{\ifmmode\sp\else\^{}\fi}\catcode`\%=\active\def%{\%}$5\omega_1$}}%
\end{pgfscope}%
\begin{pgfscope}%
\pgfpathrectangle{\pgfqpoint{0.571024in}{0.456901in}}{\pgfqpoint{5.868976in}{2.816432in}}%
\pgfusepath{clip}%
\pgfsetbuttcap%
\pgfsetroundjoin%
\pgfsetlinewidth{0.501875pt}%
\definecolor{currentstroke}{rgb}{0.501961,0.501961,0.501961}%
\pgfsetstrokecolor{currentstroke}%
\pgfsetstrokeopacity{0.400000}%
\pgfsetdash{{1.850000pt}{0.800000pt}}{0.000000pt}%
\pgfpathmoveto{\pgfqpoint{0.854095in}{0.456901in}}%
\pgfpathlineto{\pgfqpoint{0.854095in}{3.273333in}}%
\pgfusepath{stroke}%
\end{pgfscope}%
\begin{pgfscope}%
\pgfsetbuttcap%
\pgfsetroundjoin%
\definecolor{currentfill}{rgb}{1.000000,1.000000,1.000000}%
\pgfsetfillcolor{currentfill}%
\pgfsetlinewidth{0.602250pt}%
\definecolor{currentstroke}{rgb}{0.000000,0.000000,0.000000}%
\pgfsetstrokecolor{currentstroke}%
\pgfsetdash{}{0pt}%
\pgfsys@defobject{currentmarker}{\pgfqpoint{0.000000in}{0.000000in}}{\pgfqpoint{0.000000in}{0.027778in}}{%
\pgfpathmoveto{\pgfqpoint{0.000000in}{0.000000in}}%
\pgfpathlineto{\pgfqpoint{0.000000in}{0.027778in}}%
\pgfusepath{stroke,fill}%
}%
\begin{pgfscope}%
\pgfsys@transformshift{0.854095in}{0.456901in}%
\pgfsys@useobject{currentmarker}{}%
\end{pgfscope}%
\end{pgfscope}%
\begin{pgfscope}%
\pgfpathrectangle{\pgfqpoint{0.571024in}{0.456901in}}{\pgfqpoint{5.868976in}{2.816432in}}%
\pgfusepath{clip}%
\pgfsetbuttcap%
\pgfsetroundjoin%
\pgfsetlinewidth{0.501875pt}%
\definecolor{currentstroke}{rgb}{0.501961,0.501961,0.501961}%
\pgfsetstrokecolor{currentstroke}%
\pgfsetstrokeopacity{0.400000}%
\pgfsetdash{{1.850000pt}{0.800000pt}}{0.000000pt}%
\pgfpathmoveto{\pgfqpoint{1.106851in}{0.456901in}}%
\pgfpathlineto{\pgfqpoint{1.106851in}{3.273333in}}%
\pgfusepath{stroke}%
\end{pgfscope}%
\begin{pgfscope}%
\pgfsetbuttcap%
\pgfsetroundjoin%
\definecolor{currentfill}{rgb}{1.000000,1.000000,1.000000}%
\pgfsetfillcolor{currentfill}%
\pgfsetlinewidth{0.602250pt}%
\definecolor{currentstroke}{rgb}{0.000000,0.000000,0.000000}%
\pgfsetstrokecolor{currentstroke}%
\pgfsetdash{}{0pt}%
\pgfsys@defobject{currentmarker}{\pgfqpoint{0.000000in}{0.000000in}}{\pgfqpoint{0.000000in}{0.027778in}}{%
\pgfpathmoveto{\pgfqpoint{0.000000in}{0.000000in}}%
\pgfpathlineto{\pgfqpoint{0.000000in}{0.027778in}}%
\pgfusepath{stroke,fill}%
}%
\begin{pgfscope}%
\pgfsys@transformshift{1.106851in}{0.456901in}%
\pgfsys@useobject{currentmarker}{}%
\end{pgfscope}%
\end{pgfscope}%
\begin{pgfscope}%
\pgfpathrectangle{\pgfqpoint{0.571024in}{0.456901in}}{\pgfqpoint{5.868976in}{2.816432in}}%
\pgfusepath{clip}%
\pgfsetbuttcap%
\pgfsetroundjoin%
\pgfsetlinewidth{0.501875pt}%
\definecolor{currentstroke}{rgb}{0.501961,0.501961,0.501961}%
\pgfsetstrokecolor{currentstroke}%
\pgfsetstrokeopacity{0.400000}%
\pgfsetdash{{1.850000pt}{0.800000pt}}{0.000000pt}%
\pgfpathmoveto{\pgfqpoint{1.359607in}{0.456901in}}%
\pgfpathlineto{\pgfqpoint{1.359607in}{3.273333in}}%
\pgfusepath{stroke}%
\end{pgfscope}%
\begin{pgfscope}%
\pgfsetbuttcap%
\pgfsetroundjoin%
\definecolor{currentfill}{rgb}{1.000000,1.000000,1.000000}%
\pgfsetfillcolor{currentfill}%
\pgfsetlinewidth{0.602250pt}%
\definecolor{currentstroke}{rgb}{0.000000,0.000000,0.000000}%
\pgfsetstrokecolor{currentstroke}%
\pgfsetdash{}{0pt}%
\pgfsys@defobject{currentmarker}{\pgfqpoint{0.000000in}{0.000000in}}{\pgfqpoint{0.000000in}{0.027778in}}{%
\pgfpathmoveto{\pgfqpoint{0.000000in}{0.000000in}}%
\pgfpathlineto{\pgfqpoint{0.000000in}{0.027778in}}%
\pgfusepath{stroke,fill}%
}%
\begin{pgfscope}%
\pgfsys@transformshift{1.359607in}{0.456901in}%
\pgfsys@useobject{currentmarker}{}%
\end{pgfscope}%
\end{pgfscope}%
\begin{pgfscope}%
\pgfpathrectangle{\pgfqpoint{0.571024in}{0.456901in}}{\pgfqpoint{5.868976in}{2.816432in}}%
\pgfusepath{clip}%
\pgfsetbuttcap%
\pgfsetroundjoin%
\pgfsetlinewidth{0.501875pt}%
\definecolor{currentstroke}{rgb}{0.501961,0.501961,0.501961}%
\pgfsetstrokecolor{currentstroke}%
\pgfsetstrokeopacity{0.400000}%
\pgfsetdash{{1.850000pt}{0.800000pt}}{0.000000pt}%
\pgfpathmoveto{\pgfqpoint{1.865119in}{0.456901in}}%
\pgfpathlineto{\pgfqpoint{1.865119in}{3.273333in}}%
\pgfusepath{stroke}%
\end{pgfscope}%
\begin{pgfscope}%
\pgfsetbuttcap%
\pgfsetroundjoin%
\definecolor{currentfill}{rgb}{1.000000,1.000000,1.000000}%
\pgfsetfillcolor{currentfill}%
\pgfsetlinewidth{0.602250pt}%
\definecolor{currentstroke}{rgb}{0.000000,0.000000,0.000000}%
\pgfsetstrokecolor{currentstroke}%
\pgfsetdash{}{0pt}%
\pgfsys@defobject{currentmarker}{\pgfqpoint{0.000000in}{0.000000in}}{\pgfqpoint{0.000000in}{0.027778in}}{%
\pgfpathmoveto{\pgfqpoint{0.000000in}{0.000000in}}%
\pgfpathlineto{\pgfqpoint{0.000000in}{0.027778in}}%
\pgfusepath{stroke,fill}%
}%
\begin{pgfscope}%
\pgfsys@transformshift{1.865119in}{0.456901in}%
\pgfsys@useobject{currentmarker}{}%
\end{pgfscope}%
\end{pgfscope}%
\begin{pgfscope}%
\pgfpathrectangle{\pgfqpoint{0.571024in}{0.456901in}}{\pgfqpoint{5.868976in}{2.816432in}}%
\pgfusepath{clip}%
\pgfsetbuttcap%
\pgfsetroundjoin%
\pgfsetlinewidth{0.501875pt}%
\definecolor{currentstroke}{rgb}{0.501961,0.501961,0.501961}%
\pgfsetstrokecolor{currentstroke}%
\pgfsetstrokeopacity{0.400000}%
\pgfsetdash{{1.850000pt}{0.800000pt}}{0.000000pt}%
\pgfpathmoveto{\pgfqpoint{2.117875in}{0.456901in}}%
\pgfpathlineto{\pgfqpoint{2.117875in}{3.273333in}}%
\pgfusepath{stroke}%
\end{pgfscope}%
\begin{pgfscope}%
\pgfsetbuttcap%
\pgfsetroundjoin%
\definecolor{currentfill}{rgb}{1.000000,1.000000,1.000000}%
\pgfsetfillcolor{currentfill}%
\pgfsetlinewidth{0.602250pt}%
\definecolor{currentstroke}{rgb}{0.000000,0.000000,0.000000}%
\pgfsetstrokecolor{currentstroke}%
\pgfsetdash{}{0pt}%
\pgfsys@defobject{currentmarker}{\pgfqpoint{0.000000in}{0.000000in}}{\pgfqpoint{0.000000in}{0.027778in}}{%
\pgfpathmoveto{\pgfqpoint{0.000000in}{0.000000in}}%
\pgfpathlineto{\pgfqpoint{0.000000in}{0.027778in}}%
\pgfusepath{stroke,fill}%
}%
\begin{pgfscope}%
\pgfsys@transformshift{2.117875in}{0.456901in}%
\pgfsys@useobject{currentmarker}{}%
\end{pgfscope}%
\end{pgfscope}%
\begin{pgfscope}%
\pgfpathrectangle{\pgfqpoint{0.571024in}{0.456901in}}{\pgfqpoint{5.868976in}{2.816432in}}%
\pgfusepath{clip}%
\pgfsetbuttcap%
\pgfsetroundjoin%
\pgfsetlinewidth{0.501875pt}%
\definecolor{currentstroke}{rgb}{0.501961,0.501961,0.501961}%
\pgfsetstrokecolor{currentstroke}%
\pgfsetstrokeopacity{0.400000}%
\pgfsetdash{{1.850000pt}{0.800000pt}}{0.000000pt}%
\pgfpathmoveto{\pgfqpoint{2.370630in}{0.456901in}}%
\pgfpathlineto{\pgfqpoint{2.370630in}{3.273333in}}%
\pgfusepath{stroke}%
\end{pgfscope}%
\begin{pgfscope}%
\pgfsetbuttcap%
\pgfsetroundjoin%
\definecolor{currentfill}{rgb}{1.000000,1.000000,1.000000}%
\pgfsetfillcolor{currentfill}%
\pgfsetlinewidth{0.602250pt}%
\definecolor{currentstroke}{rgb}{0.000000,0.000000,0.000000}%
\pgfsetstrokecolor{currentstroke}%
\pgfsetdash{}{0pt}%
\pgfsys@defobject{currentmarker}{\pgfqpoint{0.000000in}{0.000000in}}{\pgfqpoint{0.000000in}{0.027778in}}{%
\pgfpathmoveto{\pgfqpoint{0.000000in}{0.000000in}}%
\pgfpathlineto{\pgfqpoint{0.000000in}{0.027778in}}%
\pgfusepath{stroke,fill}%
}%
\begin{pgfscope}%
\pgfsys@transformshift{2.370630in}{0.456901in}%
\pgfsys@useobject{currentmarker}{}%
\end{pgfscope}%
\end{pgfscope}%
\begin{pgfscope}%
\pgfpathrectangle{\pgfqpoint{0.571024in}{0.456901in}}{\pgfqpoint{5.868976in}{2.816432in}}%
\pgfusepath{clip}%
\pgfsetbuttcap%
\pgfsetroundjoin%
\pgfsetlinewidth{0.501875pt}%
\definecolor{currentstroke}{rgb}{0.501961,0.501961,0.501961}%
\pgfsetstrokecolor{currentstroke}%
\pgfsetstrokeopacity{0.400000}%
\pgfsetdash{{1.850000pt}{0.800000pt}}{0.000000pt}%
\pgfpathmoveto{\pgfqpoint{2.876142in}{0.456901in}}%
\pgfpathlineto{\pgfqpoint{2.876142in}{3.273333in}}%
\pgfusepath{stroke}%
\end{pgfscope}%
\begin{pgfscope}%
\pgfsetbuttcap%
\pgfsetroundjoin%
\definecolor{currentfill}{rgb}{1.000000,1.000000,1.000000}%
\pgfsetfillcolor{currentfill}%
\pgfsetlinewidth{0.602250pt}%
\definecolor{currentstroke}{rgb}{0.000000,0.000000,0.000000}%
\pgfsetstrokecolor{currentstroke}%
\pgfsetdash{}{0pt}%
\pgfsys@defobject{currentmarker}{\pgfqpoint{0.000000in}{0.000000in}}{\pgfqpoint{0.000000in}{0.027778in}}{%
\pgfpathmoveto{\pgfqpoint{0.000000in}{0.000000in}}%
\pgfpathlineto{\pgfqpoint{0.000000in}{0.027778in}}%
\pgfusepath{stroke,fill}%
}%
\begin{pgfscope}%
\pgfsys@transformshift{2.876142in}{0.456901in}%
\pgfsys@useobject{currentmarker}{}%
\end{pgfscope}%
\end{pgfscope}%
\begin{pgfscope}%
\pgfpathrectangle{\pgfqpoint{0.571024in}{0.456901in}}{\pgfqpoint{5.868976in}{2.816432in}}%
\pgfusepath{clip}%
\pgfsetbuttcap%
\pgfsetroundjoin%
\pgfsetlinewidth{0.501875pt}%
\definecolor{currentstroke}{rgb}{0.501961,0.501961,0.501961}%
\pgfsetstrokecolor{currentstroke}%
\pgfsetstrokeopacity{0.400000}%
\pgfsetdash{{1.850000pt}{0.800000pt}}{0.000000pt}%
\pgfpathmoveto{\pgfqpoint{3.128898in}{0.456901in}}%
\pgfpathlineto{\pgfqpoint{3.128898in}{3.273333in}}%
\pgfusepath{stroke}%
\end{pgfscope}%
\begin{pgfscope}%
\pgfsetbuttcap%
\pgfsetroundjoin%
\definecolor{currentfill}{rgb}{1.000000,1.000000,1.000000}%
\pgfsetfillcolor{currentfill}%
\pgfsetlinewidth{0.602250pt}%
\definecolor{currentstroke}{rgb}{0.000000,0.000000,0.000000}%
\pgfsetstrokecolor{currentstroke}%
\pgfsetdash{}{0pt}%
\pgfsys@defobject{currentmarker}{\pgfqpoint{0.000000in}{0.000000in}}{\pgfqpoint{0.000000in}{0.027778in}}{%
\pgfpathmoveto{\pgfqpoint{0.000000in}{0.000000in}}%
\pgfpathlineto{\pgfqpoint{0.000000in}{0.027778in}}%
\pgfusepath{stroke,fill}%
}%
\begin{pgfscope}%
\pgfsys@transformshift{3.128898in}{0.456901in}%
\pgfsys@useobject{currentmarker}{}%
\end{pgfscope}%
\end{pgfscope}%
\begin{pgfscope}%
\pgfpathrectangle{\pgfqpoint{0.571024in}{0.456901in}}{\pgfqpoint{5.868976in}{2.816432in}}%
\pgfusepath{clip}%
\pgfsetbuttcap%
\pgfsetroundjoin%
\pgfsetlinewidth{0.501875pt}%
\definecolor{currentstroke}{rgb}{0.501961,0.501961,0.501961}%
\pgfsetstrokecolor{currentstroke}%
\pgfsetstrokeopacity{0.400000}%
\pgfsetdash{{1.850000pt}{0.800000pt}}{0.000000pt}%
\pgfpathmoveto{\pgfqpoint{3.381654in}{0.456901in}}%
\pgfpathlineto{\pgfqpoint{3.381654in}{3.273333in}}%
\pgfusepath{stroke}%
\end{pgfscope}%
\begin{pgfscope}%
\pgfsetbuttcap%
\pgfsetroundjoin%
\definecolor{currentfill}{rgb}{1.000000,1.000000,1.000000}%
\pgfsetfillcolor{currentfill}%
\pgfsetlinewidth{0.602250pt}%
\definecolor{currentstroke}{rgb}{0.000000,0.000000,0.000000}%
\pgfsetstrokecolor{currentstroke}%
\pgfsetdash{}{0pt}%
\pgfsys@defobject{currentmarker}{\pgfqpoint{0.000000in}{0.000000in}}{\pgfqpoint{0.000000in}{0.027778in}}{%
\pgfpathmoveto{\pgfqpoint{0.000000in}{0.000000in}}%
\pgfpathlineto{\pgfqpoint{0.000000in}{0.027778in}}%
\pgfusepath{stroke,fill}%
}%
\begin{pgfscope}%
\pgfsys@transformshift{3.381654in}{0.456901in}%
\pgfsys@useobject{currentmarker}{}%
\end{pgfscope}%
\end{pgfscope}%
\begin{pgfscope}%
\pgfpathrectangle{\pgfqpoint{0.571024in}{0.456901in}}{\pgfqpoint{5.868976in}{2.816432in}}%
\pgfusepath{clip}%
\pgfsetbuttcap%
\pgfsetroundjoin%
\pgfsetlinewidth{0.501875pt}%
\definecolor{currentstroke}{rgb}{0.501961,0.501961,0.501961}%
\pgfsetstrokecolor{currentstroke}%
\pgfsetstrokeopacity{0.400000}%
\pgfsetdash{{1.850000pt}{0.800000pt}}{0.000000pt}%
\pgfpathmoveto{\pgfqpoint{3.887166in}{0.456901in}}%
\pgfpathlineto{\pgfqpoint{3.887166in}{3.273333in}}%
\pgfusepath{stroke}%
\end{pgfscope}%
\begin{pgfscope}%
\pgfsetbuttcap%
\pgfsetroundjoin%
\definecolor{currentfill}{rgb}{1.000000,1.000000,1.000000}%
\pgfsetfillcolor{currentfill}%
\pgfsetlinewidth{0.602250pt}%
\definecolor{currentstroke}{rgb}{0.000000,0.000000,0.000000}%
\pgfsetstrokecolor{currentstroke}%
\pgfsetdash{}{0pt}%
\pgfsys@defobject{currentmarker}{\pgfqpoint{0.000000in}{0.000000in}}{\pgfqpoint{0.000000in}{0.027778in}}{%
\pgfpathmoveto{\pgfqpoint{0.000000in}{0.000000in}}%
\pgfpathlineto{\pgfqpoint{0.000000in}{0.027778in}}%
\pgfusepath{stroke,fill}%
}%
\begin{pgfscope}%
\pgfsys@transformshift{3.887166in}{0.456901in}%
\pgfsys@useobject{currentmarker}{}%
\end{pgfscope}%
\end{pgfscope}%
\begin{pgfscope}%
\pgfpathrectangle{\pgfqpoint{0.571024in}{0.456901in}}{\pgfqpoint{5.868976in}{2.816432in}}%
\pgfusepath{clip}%
\pgfsetbuttcap%
\pgfsetroundjoin%
\pgfsetlinewidth{0.501875pt}%
\definecolor{currentstroke}{rgb}{0.501961,0.501961,0.501961}%
\pgfsetstrokecolor{currentstroke}%
\pgfsetstrokeopacity{0.400000}%
\pgfsetdash{{1.850000pt}{0.800000pt}}{0.000000pt}%
\pgfpathmoveto{\pgfqpoint{4.139922in}{0.456901in}}%
\pgfpathlineto{\pgfqpoint{4.139922in}{3.273333in}}%
\pgfusepath{stroke}%
\end{pgfscope}%
\begin{pgfscope}%
\pgfsetbuttcap%
\pgfsetroundjoin%
\definecolor{currentfill}{rgb}{1.000000,1.000000,1.000000}%
\pgfsetfillcolor{currentfill}%
\pgfsetlinewidth{0.602250pt}%
\definecolor{currentstroke}{rgb}{0.000000,0.000000,0.000000}%
\pgfsetstrokecolor{currentstroke}%
\pgfsetdash{}{0pt}%
\pgfsys@defobject{currentmarker}{\pgfqpoint{0.000000in}{0.000000in}}{\pgfqpoint{0.000000in}{0.027778in}}{%
\pgfpathmoveto{\pgfqpoint{0.000000in}{0.000000in}}%
\pgfpathlineto{\pgfqpoint{0.000000in}{0.027778in}}%
\pgfusepath{stroke,fill}%
}%
\begin{pgfscope}%
\pgfsys@transformshift{4.139922in}{0.456901in}%
\pgfsys@useobject{currentmarker}{}%
\end{pgfscope}%
\end{pgfscope}%
\begin{pgfscope}%
\pgfpathrectangle{\pgfqpoint{0.571024in}{0.456901in}}{\pgfqpoint{5.868976in}{2.816432in}}%
\pgfusepath{clip}%
\pgfsetbuttcap%
\pgfsetroundjoin%
\pgfsetlinewidth{0.501875pt}%
\definecolor{currentstroke}{rgb}{0.501961,0.501961,0.501961}%
\pgfsetstrokecolor{currentstroke}%
\pgfsetstrokeopacity{0.400000}%
\pgfsetdash{{1.850000pt}{0.800000pt}}{0.000000pt}%
\pgfpathmoveto{\pgfqpoint{4.392677in}{0.456901in}}%
\pgfpathlineto{\pgfqpoint{4.392677in}{3.273333in}}%
\pgfusepath{stroke}%
\end{pgfscope}%
\begin{pgfscope}%
\pgfsetbuttcap%
\pgfsetroundjoin%
\definecolor{currentfill}{rgb}{1.000000,1.000000,1.000000}%
\pgfsetfillcolor{currentfill}%
\pgfsetlinewidth{0.602250pt}%
\definecolor{currentstroke}{rgb}{0.000000,0.000000,0.000000}%
\pgfsetstrokecolor{currentstroke}%
\pgfsetdash{}{0pt}%
\pgfsys@defobject{currentmarker}{\pgfqpoint{0.000000in}{0.000000in}}{\pgfqpoint{0.000000in}{0.027778in}}{%
\pgfpathmoveto{\pgfqpoint{0.000000in}{0.000000in}}%
\pgfpathlineto{\pgfqpoint{0.000000in}{0.027778in}}%
\pgfusepath{stroke,fill}%
}%
\begin{pgfscope}%
\pgfsys@transformshift{4.392677in}{0.456901in}%
\pgfsys@useobject{currentmarker}{}%
\end{pgfscope}%
\end{pgfscope}%
\begin{pgfscope}%
\pgfpathrectangle{\pgfqpoint{0.571024in}{0.456901in}}{\pgfqpoint{5.868976in}{2.816432in}}%
\pgfusepath{clip}%
\pgfsetbuttcap%
\pgfsetroundjoin%
\pgfsetlinewidth{0.501875pt}%
\definecolor{currentstroke}{rgb}{0.501961,0.501961,0.501961}%
\pgfsetstrokecolor{currentstroke}%
\pgfsetstrokeopacity{0.400000}%
\pgfsetdash{{1.850000pt}{0.800000pt}}{0.000000pt}%
\pgfpathmoveto{\pgfqpoint{4.898189in}{0.456901in}}%
\pgfpathlineto{\pgfqpoint{4.898189in}{3.273333in}}%
\pgfusepath{stroke}%
\end{pgfscope}%
\begin{pgfscope}%
\pgfsetbuttcap%
\pgfsetroundjoin%
\definecolor{currentfill}{rgb}{1.000000,1.000000,1.000000}%
\pgfsetfillcolor{currentfill}%
\pgfsetlinewidth{0.602250pt}%
\definecolor{currentstroke}{rgb}{0.000000,0.000000,0.000000}%
\pgfsetstrokecolor{currentstroke}%
\pgfsetdash{}{0pt}%
\pgfsys@defobject{currentmarker}{\pgfqpoint{0.000000in}{0.000000in}}{\pgfqpoint{0.000000in}{0.027778in}}{%
\pgfpathmoveto{\pgfqpoint{0.000000in}{0.000000in}}%
\pgfpathlineto{\pgfqpoint{0.000000in}{0.027778in}}%
\pgfusepath{stroke,fill}%
}%
\begin{pgfscope}%
\pgfsys@transformshift{4.898189in}{0.456901in}%
\pgfsys@useobject{currentmarker}{}%
\end{pgfscope}%
\end{pgfscope}%
\begin{pgfscope}%
\pgfpathrectangle{\pgfqpoint{0.571024in}{0.456901in}}{\pgfqpoint{5.868976in}{2.816432in}}%
\pgfusepath{clip}%
\pgfsetbuttcap%
\pgfsetroundjoin%
\pgfsetlinewidth{0.501875pt}%
\definecolor{currentstroke}{rgb}{0.501961,0.501961,0.501961}%
\pgfsetstrokecolor{currentstroke}%
\pgfsetstrokeopacity{0.400000}%
\pgfsetdash{{1.850000pt}{0.800000pt}}{0.000000pt}%
\pgfpathmoveto{\pgfqpoint{5.150945in}{0.456901in}}%
\pgfpathlineto{\pgfqpoint{5.150945in}{3.273333in}}%
\pgfusepath{stroke}%
\end{pgfscope}%
\begin{pgfscope}%
\pgfsetbuttcap%
\pgfsetroundjoin%
\definecolor{currentfill}{rgb}{1.000000,1.000000,1.000000}%
\pgfsetfillcolor{currentfill}%
\pgfsetlinewidth{0.602250pt}%
\definecolor{currentstroke}{rgb}{0.000000,0.000000,0.000000}%
\pgfsetstrokecolor{currentstroke}%
\pgfsetdash{}{0pt}%
\pgfsys@defobject{currentmarker}{\pgfqpoint{0.000000in}{0.000000in}}{\pgfqpoint{0.000000in}{0.027778in}}{%
\pgfpathmoveto{\pgfqpoint{0.000000in}{0.000000in}}%
\pgfpathlineto{\pgfqpoint{0.000000in}{0.027778in}}%
\pgfusepath{stroke,fill}%
}%
\begin{pgfscope}%
\pgfsys@transformshift{5.150945in}{0.456901in}%
\pgfsys@useobject{currentmarker}{}%
\end{pgfscope}%
\end{pgfscope}%
\begin{pgfscope}%
\pgfpathrectangle{\pgfqpoint{0.571024in}{0.456901in}}{\pgfqpoint{5.868976in}{2.816432in}}%
\pgfusepath{clip}%
\pgfsetbuttcap%
\pgfsetroundjoin%
\pgfsetlinewidth{0.501875pt}%
\definecolor{currentstroke}{rgb}{0.501961,0.501961,0.501961}%
\pgfsetstrokecolor{currentstroke}%
\pgfsetstrokeopacity{0.400000}%
\pgfsetdash{{1.850000pt}{0.800000pt}}{0.000000pt}%
\pgfpathmoveto{\pgfqpoint{5.403701in}{0.456901in}}%
\pgfpathlineto{\pgfqpoint{5.403701in}{3.273333in}}%
\pgfusepath{stroke}%
\end{pgfscope}%
\begin{pgfscope}%
\pgfsetbuttcap%
\pgfsetroundjoin%
\definecolor{currentfill}{rgb}{1.000000,1.000000,1.000000}%
\pgfsetfillcolor{currentfill}%
\pgfsetlinewidth{0.602250pt}%
\definecolor{currentstroke}{rgb}{0.000000,0.000000,0.000000}%
\pgfsetstrokecolor{currentstroke}%
\pgfsetdash{}{0pt}%
\pgfsys@defobject{currentmarker}{\pgfqpoint{0.000000in}{0.000000in}}{\pgfqpoint{0.000000in}{0.027778in}}{%
\pgfpathmoveto{\pgfqpoint{0.000000in}{0.000000in}}%
\pgfpathlineto{\pgfqpoint{0.000000in}{0.027778in}}%
\pgfusepath{stroke,fill}%
}%
\begin{pgfscope}%
\pgfsys@transformshift{5.403701in}{0.456901in}%
\pgfsys@useobject{currentmarker}{}%
\end{pgfscope}%
\end{pgfscope}%
\begin{pgfscope}%
\pgfpathrectangle{\pgfqpoint{0.571024in}{0.456901in}}{\pgfqpoint{5.868976in}{2.816432in}}%
\pgfusepath{clip}%
\pgfsetbuttcap%
\pgfsetroundjoin%
\pgfsetlinewidth{0.501875pt}%
\definecolor{currentstroke}{rgb}{0.501961,0.501961,0.501961}%
\pgfsetstrokecolor{currentstroke}%
\pgfsetstrokeopacity{0.400000}%
\pgfsetdash{{1.850000pt}{0.800000pt}}{0.000000pt}%
\pgfpathmoveto{\pgfqpoint{5.909213in}{0.456901in}}%
\pgfpathlineto{\pgfqpoint{5.909213in}{3.273333in}}%
\pgfusepath{stroke}%
\end{pgfscope}%
\begin{pgfscope}%
\pgfsetbuttcap%
\pgfsetroundjoin%
\definecolor{currentfill}{rgb}{1.000000,1.000000,1.000000}%
\pgfsetfillcolor{currentfill}%
\pgfsetlinewidth{0.602250pt}%
\definecolor{currentstroke}{rgb}{0.000000,0.000000,0.000000}%
\pgfsetstrokecolor{currentstroke}%
\pgfsetdash{}{0pt}%
\pgfsys@defobject{currentmarker}{\pgfqpoint{0.000000in}{0.000000in}}{\pgfqpoint{0.000000in}{0.027778in}}{%
\pgfpathmoveto{\pgfqpoint{0.000000in}{0.000000in}}%
\pgfpathlineto{\pgfqpoint{0.000000in}{0.027778in}}%
\pgfusepath{stroke,fill}%
}%
\begin{pgfscope}%
\pgfsys@transformshift{5.909213in}{0.456901in}%
\pgfsys@useobject{currentmarker}{}%
\end{pgfscope}%
\end{pgfscope}%
\begin{pgfscope}%
\pgfpathrectangle{\pgfqpoint{0.571024in}{0.456901in}}{\pgfqpoint{5.868976in}{2.816432in}}%
\pgfusepath{clip}%
\pgfsetbuttcap%
\pgfsetroundjoin%
\pgfsetlinewidth{0.501875pt}%
\definecolor{currentstroke}{rgb}{0.501961,0.501961,0.501961}%
\pgfsetstrokecolor{currentstroke}%
\pgfsetstrokeopacity{0.400000}%
\pgfsetdash{{1.850000pt}{0.800000pt}}{0.000000pt}%
\pgfpathmoveto{\pgfqpoint{6.161969in}{0.456901in}}%
\pgfpathlineto{\pgfqpoint{6.161969in}{3.273333in}}%
\pgfusepath{stroke}%
\end{pgfscope}%
\begin{pgfscope}%
\pgfsetbuttcap%
\pgfsetroundjoin%
\definecolor{currentfill}{rgb}{1.000000,1.000000,1.000000}%
\pgfsetfillcolor{currentfill}%
\pgfsetlinewidth{0.602250pt}%
\definecolor{currentstroke}{rgb}{0.000000,0.000000,0.000000}%
\pgfsetstrokecolor{currentstroke}%
\pgfsetdash{}{0pt}%
\pgfsys@defobject{currentmarker}{\pgfqpoint{0.000000in}{0.000000in}}{\pgfqpoint{0.000000in}{0.027778in}}{%
\pgfpathmoveto{\pgfqpoint{0.000000in}{0.000000in}}%
\pgfpathlineto{\pgfqpoint{0.000000in}{0.027778in}}%
\pgfusepath{stroke,fill}%
}%
\begin{pgfscope}%
\pgfsys@transformshift{6.161969in}{0.456901in}%
\pgfsys@useobject{currentmarker}{}%
\end{pgfscope}%
\end{pgfscope}%
\begin{pgfscope}%
\pgfpathrectangle{\pgfqpoint{0.571024in}{0.456901in}}{\pgfqpoint{5.868976in}{2.816432in}}%
\pgfusepath{clip}%
\pgfsetbuttcap%
\pgfsetroundjoin%
\pgfsetlinewidth{0.501875pt}%
\definecolor{currentstroke}{rgb}{0.501961,0.501961,0.501961}%
\pgfsetstrokecolor{currentstroke}%
\pgfsetstrokeopacity{0.400000}%
\pgfsetdash{{1.850000pt}{0.800000pt}}{0.000000pt}%
\pgfpathmoveto{\pgfqpoint{6.414724in}{0.456901in}}%
\pgfpathlineto{\pgfqpoint{6.414724in}{3.273333in}}%
\pgfusepath{stroke}%
\end{pgfscope}%
\begin{pgfscope}%
\pgfsetbuttcap%
\pgfsetroundjoin%
\definecolor{currentfill}{rgb}{1.000000,1.000000,1.000000}%
\pgfsetfillcolor{currentfill}%
\pgfsetlinewidth{0.602250pt}%
\definecolor{currentstroke}{rgb}{0.000000,0.000000,0.000000}%
\pgfsetstrokecolor{currentstroke}%
\pgfsetdash{}{0pt}%
\pgfsys@defobject{currentmarker}{\pgfqpoint{0.000000in}{0.000000in}}{\pgfqpoint{0.000000in}{0.027778in}}{%
\pgfpathmoveto{\pgfqpoint{0.000000in}{0.000000in}}%
\pgfpathlineto{\pgfqpoint{0.000000in}{0.027778in}}%
\pgfusepath{stroke,fill}%
}%
\begin{pgfscope}%
\pgfsys@transformshift{6.414724in}{0.456901in}%
\pgfsys@useobject{currentmarker}{}%
\end{pgfscope}%
\end{pgfscope}%
\begin{pgfscope}%
\definecolor{textcolor}{rgb}{0.000000,0.000000,0.000000}%
\pgfsetstrokecolor{textcolor}%
\pgfsetfillcolor{textcolor}%
\pgftext[x=3.505512in,y=0.201367in,,top]{\color{textcolor}{\rmfamily\fontsize{12.000000}{14.400000}\selectfont\catcode`\^=\active\def^{\ifmmode\sp\else\^{}\fi}\catcode`\%=\active\def%{\%}$\omega$}}%
\end{pgfscope}%
\begin{pgfscope}%
\pgfsetbuttcap%
\pgfsetroundjoin%
\definecolor{currentfill}{rgb}{1.000000,1.000000,1.000000}%
\pgfsetfillcolor{currentfill}%
\pgfsetlinewidth{0.803000pt}%
\definecolor{currentstroke}{rgb}{0.000000,0.000000,0.000000}%
\pgfsetstrokecolor{currentstroke}%
\pgfsetdash{}{0pt}%
\pgfsys@defobject{currentmarker}{\pgfqpoint{0.000000in}{0.000000in}}{\pgfqpoint{0.048611in}{0.000000in}}{%
\pgfpathmoveto{\pgfqpoint{0.000000in}{0.000000in}}%
\pgfpathlineto{\pgfqpoint{0.048611in}{0.000000in}}%
\pgfusepath{stroke,fill}%
}%
\begin{pgfscope}%
\pgfsys@transformshift{0.571024in}{0.456901in}%
\pgfsys@useobject{currentmarker}{}%
\end{pgfscope}%
\end{pgfscope}%
\begin{pgfscope}%
\definecolor{textcolor}{rgb}{0.000000,0.000000,0.000000}%
\pgfsetstrokecolor{textcolor}%
\pgfsetfillcolor{textcolor}%
\pgftext[x=0.272222in, y=0.399040in, left, base]{\color{textcolor}{\rmfamily\fontsize{12.000000}{14.400000}\selectfont\catcode`\^=\active\def^{\ifmmode\sp\else\^{}\fi}\catcode`\%=\active\def%{\%}$\mathdefault{0.0}$}}%
\end{pgfscope}%
\begin{pgfscope}%
\pgfsetbuttcap%
\pgfsetroundjoin%
\definecolor{currentfill}{rgb}{1.000000,1.000000,1.000000}%
\pgfsetfillcolor{currentfill}%
\pgfsetlinewidth{0.803000pt}%
\definecolor{currentstroke}{rgb}{0.000000,0.000000,0.000000}%
\pgfsetstrokecolor{currentstroke}%
\pgfsetdash{}{0pt}%
\pgfsys@defobject{currentmarker}{\pgfqpoint{0.000000in}{0.000000in}}{\pgfqpoint{0.048611in}{0.000000in}}{%
\pgfpathmoveto{\pgfqpoint{0.000000in}{0.000000in}}%
\pgfpathlineto{\pgfqpoint{0.048611in}{0.000000in}}%
\pgfusepath{stroke,fill}%
}%
\begin{pgfscope}%
\pgfsys@transformshift{0.571024in}{0.829341in}%
\pgfsys@useobject{currentmarker}{}%
\end{pgfscope}%
\end{pgfscope}%
\begin{pgfscope}%
\definecolor{textcolor}{rgb}{0.000000,0.000000,0.000000}%
\pgfsetstrokecolor{textcolor}%
\pgfsetfillcolor{textcolor}%
\pgftext[x=0.272222in, y=0.771480in, left, base]{\color{textcolor}{\rmfamily\fontsize{12.000000}{14.400000}\selectfont\catcode`\^=\active\def^{\ifmmode\sp\else\^{}\fi}\catcode`\%=\active\def%{\%}$\mathdefault{0.1}$}}%
\end{pgfscope}%
\begin{pgfscope}%
\pgfsetbuttcap%
\pgfsetroundjoin%
\definecolor{currentfill}{rgb}{1.000000,1.000000,1.000000}%
\pgfsetfillcolor{currentfill}%
\pgfsetlinewidth{0.803000pt}%
\definecolor{currentstroke}{rgb}{0.000000,0.000000,0.000000}%
\pgfsetstrokecolor{currentstroke}%
\pgfsetdash{}{0pt}%
\pgfsys@defobject{currentmarker}{\pgfqpoint{0.000000in}{0.000000in}}{\pgfqpoint{0.048611in}{0.000000in}}{%
\pgfpathmoveto{\pgfqpoint{0.000000in}{0.000000in}}%
\pgfpathlineto{\pgfqpoint{0.048611in}{0.000000in}}%
\pgfusepath{stroke,fill}%
}%
\begin{pgfscope}%
\pgfsys@transformshift{0.571024in}{1.201781in}%
\pgfsys@useobject{currentmarker}{}%
\end{pgfscope}%
\end{pgfscope}%
\begin{pgfscope}%
\definecolor{textcolor}{rgb}{0.000000,0.000000,0.000000}%
\pgfsetstrokecolor{textcolor}%
\pgfsetfillcolor{textcolor}%
\pgftext[x=0.272222in, y=1.143920in, left, base]{\color{textcolor}{\rmfamily\fontsize{12.000000}{14.400000}\selectfont\catcode`\^=\active\def^{\ifmmode\sp\else\^{}\fi}\catcode`\%=\active\def%{\%}$\mathdefault{0.2}$}}%
\end{pgfscope}%
\begin{pgfscope}%
\pgfsetbuttcap%
\pgfsetroundjoin%
\definecolor{currentfill}{rgb}{1.000000,1.000000,1.000000}%
\pgfsetfillcolor{currentfill}%
\pgfsetlinewidth{0.803000pt}%
\definecolor{currentstroke}{rgb}{0.000000,0.000000,0.000000}%
\pgfsetstrokecolor{currentstroke}%
\pgfsetdash{}{0pt}%
\pgfsys@defobject{currentmarker}{\pgfqpoint{0.000000in}{0.000000in}}{\pgfqpoint{0.048611in}{0.000000in}}{%
\pgfpathmoveto{\pgfqpoint{0.000000in}{0.000000in}}%
\pgfpathlineto{\pgfqpoint{0.048611in}{0.000000in}}%
\pgfusepath{stroke,fill}%
}%
\begin{pgfscope}%
\pgfsys@transformshift{0.571024in}{1.574221in}%
\pgfsys@useobject{currentmarker}{}%
\end{pgfscope}%
\end{pgfscope}%
\begin{pgfscope}%
\definecolor{textcolor}{rgb}{0.000000,0.000000,0.000000}%
\pgfsetstrokecolor{textcolor}%
\pgfsetfillcolor{textcolor}%
\pgftext[x=0.272222in, y=1.516360in, left, base]{\color{textcolor}{\rmfamily\fontsize{12.000000}{14.400000}\selectfont\catcode`\^=\active\def^{\ifmmode\sp\else\^{}\fi}\catcode`\%=\active\def%{\%}$\mathdefault{0.3}$}}%
\end{pgfscope}%
\begin{pgfscope}%
\pgfsetbuttcap%
\pgfsetroundjoin%
\definecolor{currentfill}{rgb}{1.000000,1.000000,1.000000}%
\pgfsetfillcolor{currentfill}%
\pgfsetlinewidth{0.803000pt}%
\definecolor{currentstroke}{rgb}{0.000000,0.000000,0.000000}%
\pgfsetstrokecolor{currentstroke}%
\pgfsetdash{}{0pt}%
\pgfsys@defobject{currentmarker}{\pgfqpoint{0.000000in}{0.000000in}}{\pgfqpoint{0.048611in}{0.000000in}}{%
\pgfpathmoveto{\pgfqpoint{0.000000in}{0.000000in}}%
\pgfpathlineto{\pgfqpoint{0.048611in}{0.000000in}}%
\pgfusepath{stroke,fill}%
}%
\begin{pgfscope}%
\pgfsys@transformshift{0.571024in}{1.946662in}%
\pgfsys@useobject{currentmarker}{}%
\end{pgfscope}%
\end{pgfscope}%
\begin{pgfscope}%
\definecolor{textcolor}{rgb}{0.000000,0.000000,0.000000}%
\pgfsetstrokecolor{textcolor}%
\pgfsetfillcolor{textcolor}%
\pgftext[x=0.272222in, y=1.888800in, left, base]{\color{textcolor}{\rmfamily\fontsize{12.000000}{14.400000}\selectfont\catcode`\^=\active\def^{\ifmmode\sp\else\^{}\fi}\catcode`\%=\active\def%{\%}$\mathdefault{0.4}$}}%
\end{pgfscope}%
\begin{pgfscope}%
\pgfsetbuttcap%
\pgfsetroundjoin%
\definecolor{currentfill}{rgb}{1.000000,1.000000,1.000000}%
\pgfsetfillcolor{currentfill}%
\pgfsetlinewidth{0.803000pt}%
\definecolor{currentstroke}{rgb}{0.000000,0.000000,0.000000}%
\pgfsetstrokecolor{currentstroke}%
\pgfsetdash{}{0pt}%
\pgfsys@defobject{currentmarker}{\pgfqpoint{0.000000in}{0.000000in}}{\pgfqpoint{0.048611in}{0.000000in}}{%
\pgfpathmoveto{\pgfqpoint{0.000000in}{0.000000in}}%
\pgfpathlineto{\pgfqpoint{0.048611in}{0.000000in}}%
\pgfusepath{stroke,fill}%
}%
\begin{pgfscope}%
\pgfsys@transformshift{0.571024in}{2.319102in}%
\pgfsys@useobject{currentmarker}{}%
\end{pgfscope}%
\end{pgfscope}%
\begin{pgfscope}%
\definecolor{textcolor}{rgb}{0.000000,0.000000,0.000000}%
\pgfsetstrokecolor{textcolor}%
\pgfsetfillcolor{textcolor}%
\pgftext[x=0.272222in, y=2.261240in, left, base]{\color{textcolor}{\rmfamily\fontsize{12.000000}{14.400000}\selectfont\catcode`\^=\active\def^{\ifmmode\sp\else\^{}\fi}\catcode`\%=\active\def%{\%}$\mathdefault{0.5}$}}%
\end{pgfscope}%
\begin{pgfscope}%
\pgfsetbuttcap%
\pgfsetroundjoin%
\definecolor{currentfill}{rgb}{1.000000,1.000000,1.000000}%
\pgfsetfillcolor{currentfill}%
\pgfsetlinewidth{0.803000pt}%
\definecolor{currentstroke}{rgb}{0.000000,0.000000,0.000000}%
\pgfsetstrokecolor{currentstroke}%
\pgfsetdash{}{0pt}%
\pgfsys@defobject{currentmarker}{\pgfqpoint{0.000000in}{0.000000in}}{\pgfqpoint{0.048611in}{0.000000in}}{%
\pgfpathmoveto{\pgfqpoint{0.000000in}{0.000000in}}%
\pgfpathlineto{\pgfqpoint{0.048611in}{0.000000in}}%
\pgfusepath{stroke,fill}%
}%
\begin{pgfscope}%
\pgfsys@transformshift{0.571024in}{2.691542in}%
\pgfsys@useobject{currentmarker}{}%
\end{pgfscope}%
\end{pgfscope}%
\begin{pgfscope}%
\definecolor{textcolor}{rgb}{0.000000,0.000000,0.000000}%
\pgfsetstrokecolor{textcolor}%
\pgfsetfillcolor{textcolor}%
\pgftext[x=0.272222in, y=2.633681in, left, base]{\color{textcolor}{\rmfamily\fontsize{12.000000}{14.400000}\selectfont\catcode`\^=\active\def^{\ifmmode\sp\else\^{}\fi}\catcode`\%=\active\def%{\%}$\mathdefault{0.6}$}}%
\end{pgfscope}%
\begin{pgfscope}%
\pgfsetbuttcap%
\pgfsetroundjoin%
\definecolor{currentfill}{rgb}{1.000000,1.000000,1.000000}%
\pgfsetfillcolor{currentfill}%
\pgfsetlinewidth{0.803000pt}%
\definecolor{currentstroke}{rgb}{0.000000,0.000000,0.000000}%
\pgfsetstrokecolor{currentstroke}%
\pgfsetdash{}{0pt}%
\pgfsys@defobject{currentmarker}{\pgfqpoint{0.000000in}{0.000000in}}{\pgfqpoint{0.048611in}{0.000000in}}{%
\pgfpathmoveto{\pgfqpoint{0.000000in}{0.000000in}}%
\pgfpathlineto{\pgfqpoint{0.048611in}{0.000000in}}%
\pgfusepath{stroke,fill}%
}%
\begin{pgfscope}%
\pgfsys@transformshift{0.571024in}{3.063982in}%
\pgfsys@useobject{currentmarker}{}%
\end{pgfscope}%
\end{pgfscope}%
\begin{pgfscope}%
\definecolor{textcolor}{rgb}{0.000000,0.000000,0.000000}%
\pgfsetstrokecolor{textcolor}%
\pgfsetfillcolor{textcolor}%
\pgftext[x=0.272222in, y=3.006121in, left, base]{\color{textcolor}{\rmfamily\fontsize{12.000000}{14.400000}\selectfont\catcode`\^=\active\def^{\ifmmode\sp\else\^{}\fi}\catcode`\%=\active\def%{\%}$\mathdefault{0.7}$}}%
\end{pgfscope}%
\begin{pgfscope}%
\definecolor{textcolor}{rgb}{0.000000,0.000000,0.000000}%
\pgfsetstrokecolor{textcolor}%
\pgfsetfillcolor{textcolor}%
\pgftext[x=0.216667in,y=1.865117in,,bottom,rotate=90.000000]{\color{textcolor}{\rmfamily\fontsize{12.000000}{14.400000}\selectfont\catcode`\^=\active\def^{\ifmmode\sp\else\^{}\fi}\catcode`\%=\active\def%{\%}$A_{\mathrm{вых}}(\omega)$}}%
\end{pgfscope}%
\begin{pgfscope}%
\pgfpathrectangle{\pgfqpoint{0.571024in}{0.456901in}}{\pgfqpoint{5.868976in}{2.816432in}}%
\pgfusepath{clip}%
\pgfsetbuttcap%
\pgfsetroundjoin%
\pgfsetlinewidth{2.007500pt}%
\definecolor{currentstroke}{rgb}{1.000000,0.000000,0.000000}%
\pgfsetstrokecolor{currentstroke}%
\pgfsetdash{{7.400000pt}{3.200000pt}}{0.000000pt}%
\pgfpathmoveto{\pgfqpoint{0.601339in}{0.456901in}}%
\pgfpathlineto{\pgfqpoint{0.779636in}{0.969466in}}%
\pgfpathlineto{\pgfqpoint{0.879928in}{1.250416in}}%
\pgfpathlineto{\pgfqpoint{0.957933in}{1.462476in}}%
\pgfpathlineto{\pgfqpoint{1.024794in}{1.638463in}}%
\pgfpathlineto{\pgfqpoint{1.091655in}{1.808139in}}%
\pgfpathlineto{\pgfqpoint{1.147373in}{1.944075in}}%
\pgfpathlineto{\pgfqpoint{1.203091in}{2.074529in}}%
\pgfpathlineto{\pgfqpoint{1.258809in}{2.199041in}}%
\pgfpathlineto{\pgfqpoint{1.314526in}{2.317178in}}%
\pgfpathlineto{\pgfqpoint{1.359100in}{2.406826in}}%
\pgfpathlineto{\pgfqpoint{1.403675in}{2.491944in}}%
\pgfpathlineto{\pgfqpoint{1.448249in}{2.572354in}}%
\pgfpathlineto{\pgfqpoint{1.492823in}{2.647891in}}%
\pgfpathlineto{\pgfqpoint{1.537397in}{2.718404in}}%
\pgfpathlineto{\pgfqpoint{1.581971in}{2.783760in}}%
\pgfpathlineto{\pgfqpoint{1.615402in}{2.829318in}}%
\pgfpathlineto{\pgfqpoint{1.648833in}{2.871861in}}%
\pgfpathlineto{\pgfqpoint{1.682263in}{2.911349in}}%
\pgfpathlineto{\pgfqpoint{1.715694in}{2.947747in}}%
\pgfpathlineto{\pgfqpoint{1.749125in}{2.981025in}}%
\pgfpathlineto{\pgfqpoint{1.782555in}{3.011157in}}%
\pgfpathlineto{\pgfqpoint{1.815986in}{3.038125in}}%
\pgfpathlineto{\pgfqpoint{1.849416in}{3.061912in}}%
\pgfpathlineto{\pgfqpoint{1.882847in}{3.082509in}}%
\pgfpathlineto{\pgfqpoint{1.916278in}{3.099912in}}%
\pgfpathlineto{\pgfqpoint{1.949708in}{3.114121in}}%
\pgfpathlineto{\pgfqpoint{1.983139in}{3.125143in}}%
\pgfpathlineto{\pgfqpoint{2.016570in}{3.132988in}}%
\pgfpathlineto{\pgfqpoint{2.050000in}{3.137673in}}%
\pgfpathlineto{\pgfqpoint{2.083431in}{3.139217in}}%
\pgfpathlineto{\pgfqpoint{2.116862in}{3.137649in}}%
\pgfpathlineto{\pgfqpoint{2.150292in}{3.132998in}}%
\pgfpathlineto{\pgfqpoint{2.183723in}{3.125301in}}%
\pgfpathlineto{\pgfqpoint{2.217153in}{3.114599in}}%
\pgfpathlineto{\pgfqpoint{2.250584in}{3.100938in}}%
\pgfpathlineto{\pgfqpoint{2.284015in}{3.084368in}}%
\pgfpathlineto{\pgfqpoint{2.317445in}{3.064944in}}%
\pgfpathlineto{\pgfqpoint{2.350876in}{3.042727in}}%
\pgfpathlineto{\pgfqpoint{2.384307in}{3.017780in}}%
\pgfpathlineto{\pgfqpoint{2.417737in}{2.990173in}}%
\pgfpathlineto{\pgfqpoint{2.451168in}{2.959978in}}%
\pgfpathlineto{\pgfqpoint{2.484599in}{2.927274in}}%
\pgfpathlineto{\pgfqpoint{2.529173in}{2.879905in}}%
\pgfpathlineto{\pgfqpoint{2.573747in}{2.828424in}}%
\pgfpathlineto{\pgfqpoint{2.618321in}{2.773048in}}%
\pgfpathlineto{\pgfqpoint{2.662895in}{2.714004in}}%
\pgfpathlineto{\pgfqpoint{2.707469in}{2.651534in}}%
\pgfpathlineto{\pgfqpoint{2.763187in}{2.569012in}}%
\pgfpathlineto{\pgfqpoint{2.818905in}{2.482041in}}%
\pgfpathlineto{\pgfqpoint{2.874623in}{2.391159in}}%
\pgfpathlineto{\pgfqpoint{2.941484in}{2.277727in}}%
\pgfpathlineto{\pgfqpoint{3.019489in}{2.140629in}}%
\pgfpathlineto{\pgfqpoint{3.119781in}{1.959385in}}%
\pgfpathlineto{\pgfqpoint{3.353795in}{1.534246in}}%
\pgfpathlineto{\pgfqpoint{3.431800in}{1.397981in}}%
\pgfpathlineto{\pgfqpoint{3.498661in}{1.285539in}}%
\pgfpathlineto{\pgfqpoint{3.554379in}{1.195673in}}%
\pgfpathlineto{\pgfqpoint{3.610097in}{1.109868in}}%
\pgfpathlineto{\pgfqpoint{3.665814in}{1.028615in}}%
\pgfpathlineto{\pgfqpoint{3.710389in}{0.967182in}}%
\pgfpathlineto{\pgfqpoint{3.754963in}{0.909133in}}%
\pgfpathlineto{\pgfqpoint{3.799537in}{0.854630in}}%
\pgfpathlineto{\pgfqpoint{3.844111in}{0.803802in}}%
\pgfpathlineto{\pgfqpoint{3.888685in}{0.756743in}}%
\pgfpathlineto{\pgfqpoint{3.933259in}{0.713512in}}%
\pgfpathlineto{\pgfqpoint{3.977834in}{0.674127in}}%
\pgfpathlineto{\pgfqpoint{4.022408in}{0.638570in}}%
\pgfpathlineto{\pgfqpoint{4.055838in}{0.614381in}}%
\pgfpathlineto{\pgfqpoint{4.089269in}{0.592276in}}%
\pgfpathlineto{\pgfqpoint{4.122700in}{0.572205in}}%
\pgfpathlineto{\pgfqpoint{4.156130in}{0.554110in}}%
\pgfpathlineto{\pgfqpoint{4.200705in}{0.532933in}}%
\pgfpathlineto{\pgfqpoint{4.245279in}{0.514952in}}%
\pgfpathlineto{\pgfqpoint{4.289853in}{0.499953in}}%
\pgfpathlineto{\pgfqpoint{4.334427in}{0.487699in}}%
\pgfpathlineto{\pgfqpoint{4.379001in}{0.477939in}}%
\pgfpathlineto{\pgfqpoint{4.423575in}{0.470413in}}%
\pgfpathlineto{\pgfqpoint{4.468150in}{0.464850in}}%
\pgfpathlineto{\pgfqpoint{4.523867in}{0.460247in}}%
\pgfpathlineto{\pgfqpoint{4.579585in}{0.457771in}}%
\pgfpathlineto{\pgfqpoint{4.646446in}{0.456901in}}%
\pgfpathlineto{\pgfqpoint{4.746738in}{0.458224in}}%
\pgfpathlineto{\pgfqpoint{4.958466in}{0.462103in}}%
\pgfpathlineto{\pgfqpoint{5.047614in}{0.461206in}}%
\pgfpathlineto{\pgfqpoint{5.125619in}{0.458312in}}%
\pgfpathlineto{\pgfqpoint{5.159049in}{0.457406in}}%
\pgfpathlineto{\pgfqpoint{5.237054in}{0.463538in}}%
\pgfpathlineto{\pgfqpoint{5.315059in}{0.472018in}}%
\pgfpathlineto{\pgfqpoint{5.393064in}{0.482776in}}%
\pgfpathlineto{\pgfqpoint{5.471069in}{0.495640in}}%
\pgfpathlineto{\pgfqpoint{5.560217in}{0.512598in}}%
\pgfpathlineto{\pgfqpoint{5.660509in}{0.533985in}}%
\pgfpathlineto{\pgfqpoint{5.794232in}{0.565003in}}%
\pgfpathlineto{\pgfqpoint{6.083964in}{0.633153in}}%
\pgfpathlineto{\pgfqpoint{6.161969in}{0.649807in}}%
\pgfpathlineto{\pgfqpoint{6.161969in}{0.649807in}}%
\pgfusepath{stroke}%
\end{pgfscope}%
\begin{pgfscope}%
\pgfpathrectangle{\pgfqpoint{0.571024in}{0.456901in}}{\pgfqpoint{5.868976in}{2.816432in}}%
\pgfusepath{clip}%
\pgfsetbuttcap%
\pgfsetroundjoin%
\pgfsetlinewidth{2.007500pt}%
\definecolor{currentstroke}{rgb}{0.000000,0.000000,1.000000}%
\pgfsetstrokecolor{currentstroke}%
\pgfsetdash{}{0pt}%
\pgfpathmoveto{\pgfqpoint{0.601339in}{0.456901in}}%
\pgfpathlineto{\pgfqpoint{0.601339in}{0.456901in}}%
\pgfusepath{stroke}%
\end{pgfscope}%
\begin{pgfscope}%
\pgfpathrectangle{\pgfqpoint{0.571024in}{0.456901in}}{\pgfqpoint{5.868976in}{2.816432in}}%
\pgfusepath{clip}%
\pgfsetbuttcap%
\pgfsetroundjoin%
\pgfsetlinewidth{2.007500pt}%
\definecolor{currentstroke}{rgb}{0.000000,0.000000,1.000000}%
\pgfsetstrokecolor{currentstroke}%
\pgfsetdash{}{0pt}%
\pgfpathmoveto{\pgfqpoint{1.612363in}{0.456901in}}%
\pgfpathlineto{\pgfqpoint{1.612363in}{2.825300in}}%
\pgfusepath{stroke}%
\end{pgfscope}%
\begin{pgfscope}%
\pgfpathrectangle{\pgfqpoint{0.571024in}{0.456901in}}{\pgfqpoint{5.868976in}{2.816432in}}%
\pgfusepath{clip}%
\pgfsetbuttcap%
\pgfsetroundjoin%
\pgfsetlinewidth{2.007500pt}%
\definecolor{currentstroke}{rgb}{0.000000,0.000000,1.000000}%
\pgfsetstrokecolor{currentstroke}%
\pgfsetdash{}{0pt}%
\pgfpathmoveto{\pgfqpoint{2.623386in}{0.456901in}}%
\pgfpathlineto{\pgfqpoint{2.623386in}{2.766518in}}%
\pgfusepath{stroke}%
\end{pgfscope}%
\begin{pgfscope}%
\pgfpathrectangle{\pgfqpoint{0.571024in}{0.456901in}}{\pgfqpoint{5.868976in}{2.816432in}}%
\pgfusepath{clip}%
\pgfsetbuttcap%
\pgfsetroundjoin%
\pgfsetlinewidth{2.007500pt}%
\definecolor{currentstroke}{rgb}{0.000000,0.000000,1.000000}%
\pgfsetstrokecolor{currentstroke}%
\pgfsetdash{}{0pt}%
\pgfpathmoveto{\pgfqpoint{3.634410in}{0.456901in}}%
\pgfpathlineto{\pgfqpoint{3.634410in}{1.073825in}}%
\pgfusepath{stroke}%
\end{pgfscope}%
\begin{pgfscope}%
\pgfpathrectangle{\pgfqpoint{0.571024in}{0.456901in}}{\pgfqpoint{5.868976in}{2.816432in}}%
\pgfusepath{clip}%
\pgfsetbuttcap%
\pgfsetroundjoin%
\pgfsetlinewidth{2.007500pt}%
\definecolor{currentstroke}{rgb}{0.000000,0.000000,1.000000}%
\pgfsetstrokecolor{currentstroke}%
\pgfsetdash{}{0pt}%
\pgfpathmoveto{\pgfqpoint{4.645433in}{0.456901in}}%
\pgfpathlineto{\pgfqpoint{4.645433in}{0.456901in}}%
\pgfusepath{stroke}%
\end{pgfscope}%
\begin{pgfscope}%
\pgfpathrectangle{\pgfqpoint{0.571024in}{0.456901in}}{\pgfqpoint{5.868976in}{2.816432in}}%
\pgfusepath{clip}%
\pgfsetbuttcap%
\pgfsetroundjoin%
\pgfsetlinewidth{2.007500pt}%
\definecolor{currentstroke}{rgb}{0.000000,0.000000,1.000000}%
\pgfsetstrokecolor{currentstroke}%
\pgfsetdash{}{0pt}%
\pgfpathmoveto{\pgfqpoint{5.656457in}{0.456901in}}%
\pgfpathlineto{\pgfqpoint{5.656457in}{0.533082in}}%
\pgfusepath{stroke}%
\end{pgfscope}%
\begin{pgfscope}%
\pgfpathrectangle{\pgfqpoint{0.571024in}{0.456901in}}{\pgfqpoint{5.868976in}{2.816432in}}%
\pgfusepath{clip}%
\pgfsetbuttcap%
\pgfsetroundjoin%
\definecolor{currentfill}{rgb}{1.000000,1.000000,1.000000}%
\pgfsetfillcolor{currentfill}%
\pgfsetlinewidth{1.505625pt}%
\definecolor{currentstroke}{rgb}{0.000000,0.000000,1.000000}%
\pgfsetstrokecolor{currentstroke}%
\pgfsetdash{}{0pt}%
\pgfsys@defobject{currentmarker}{\pgfqpoint{-0.027778in}{-0.027778in}}{\pgfqpoint{0.027778in}{0.027778in}}{%
\pgfpathmoveto{\pgfqpoint{0.000000in}{-0.027778in}}%
\pgfpathcurveto{\pgfqpoint{0.007367in}{-0.027778in}}{\pgfqpoint{0.014433in}{-0.024851in}}{\pgfqpoint{0.019642in}{-0.019642in}}%
\pgfpathcurveto{\pgfqpoint{0.024851in}{-0.014433in}}{\pgfqpoint{0.027778in}{-0.007367in}}{\pgfqpoint{0.027778in}{0.000000in}}%
\pgfpathcurveto{\pgfqpoint{0.027778in}{0.007367in}}{\pgfqpoint{0.024851in}{0.014433in}}{\pgfqpoint{0.019642in}{0.019642in}}%
\pgfpathcurveto{\pgfqpoint{0.014433in}{0.024851in}}{\pgfqpoint{0.007367in}{0.027778in}}{\pgfqpoint{0.000000in}{0.027778in}}%
\pgfpathcurveto{\pgfqpoint{-0.007367in}{0.027778in}}{\pgfqpoint{-0.014433in}{0.024851in}}{\pgfqpoint{-0.019642in}{0.019642in}}%
\pgfpathcurveto{\pgfqpoint{-0.024851in}{0.014433in}}{\pgfqpoint{-0.027778in}{0.007367in}}{\pgfqpoint{-0.027778in}{0.000000in}}%
\pgfpathcurveto{\pgfqpoint{-0.027778in}{-0.007367in}}{\pgfqpoint{-0.024851in}{-0.014433in}}{\pgfqpoint{-0.019642in}{-0.019642in}}%
\pgfpathcurveto{\pgfqpoint{-0.014433in}{-0.024851in}}{\pgfqpoint{-0.007367in}{-0.027778in}}{\pgfqpoint{0.000000in}{-0.027778in}}%
\pgfpathlineto{\pgfqpoint{0.000000in}{-0.027778in}}%
\pgfpathclose%
\pgfusepath{stroke,fill}%
}%
\begin{pgfscope}%
\pgfsys@transformshift{0.601339in}{0.456901in}%
\pgfsys@useobject{currentmarker}{}%
\end{pgfscope}%
\begin{pgfscope}%
\pgfsys@transformshift{1.612363in}{2.825300in}%
\pgfsys@useobject{currentmarker}{}%
\end{pgfscope}%
\begin{pgfscope}%
\pgfsys@transformshift{2.623386in}{2.766518in}%
\pgfsys@useobject{currentmarker}{}%
\end{pgfscope}%
\begin{pgfscope}%
\pgfsys@transformshift{3.634410in}{1.073825in}%
\pgfsys@useobject{currentmarker}{}%
\end{pgfscope}%
\begin{pgfscope}%
\pgfsys@transformshift{4.645433in}{0.456901in}%
\pgfsys@useobject{currentmarker}{}%
\end{pgfscope}%
\begin{pgfscope}%
\pgfsys@transformshift{5.656457in}{0.533082in}%
\pgfsys@useobject{currentmarker}{}%
\end{pgfscope}%
\end{pgfscope}%
\begin{pgfscope}%
\pgfpathrectangle{\pgfqpoint{0.571024in}{0.456901in}}{\pgfqpoint{5.868976in}{2.816432in}}%
\pgfusepath{clip}%
\pgfsetrectcap%
\pgfsetroundjoin%
\pgfsetlinewidth{2.007500pt}%
\definecolor{currentstroke}{rgb}{0.000000,0.000000,0.000000}%
\pgfsetstrokecolor{currentstroke}%
\pgfsetdash{}{0pt}%
\pgfpathmoveto{\pgfqpoint{0.601339in}{0.456901in}}%
\pgfpathlineto{\pgfqpoint{5.656457in}{0.456901in}}%
\pgfusepath{stroke}%
\end{pgfscope}%
\begin{pgfscope}%
\pgfsetrectcap%
\pgfsetmiterjoin%
\pgfsetlinewidth{0.602250pt}%
\definecolor{currentstroke}{rgb}{0.000000,0.000000,0.000000}%
\pgfsetstrokecolor{currentstroke}%
\pgfsetdash{}{0pt}%
\pgfpathmoveto{\pgfqpoint{0.571024in}{0.456901in}}%
\pgfpathlineto{\pgfqpoint{0.571024in}{3.273333in}}%
\pgfusepath{stroke}%
\end{pgfscope}%
\begin{pgfscope}%
\pgfsetrectcap%
\pgfsetmiterjoin%
\pgfsetlinewidth{0.602250pt}%
\definecolor{currentstroke}{rgb}{0.000000,0.000000,0.000000}%
\pgfsetstrokecolor{currentstroke}%
\pgfsetdash{}{0pt}%
\pgfpathmoveto{\pgfqpoint{0.571024in}{0.456901in}}%
\pgfpathlineto{\pgfqpoint{6.440000in}{0.456901in}}%
\pgfusepath{stroke}%
\end{pgfscope}%
\begin{pgfscope}%
\pgfsetbuttcap%
\pgfsetroundjoin%
\pgfsetlinewidth{2.007500pt}%
\definecolor{currentstroke}{rgb}{1.000000,0.000000,0.000000}%
\pgfsetstrokecolor{currentstroke}%
\pgfsetdash{{7.400000pt}{3.200000pt}}{0.000000pt}%
\pgfpathmoveto{\pgfqpoint{2.474688in}{3.556667in}}%
\pgfpathlineto{\pgfqpoint{2.724688in}{3.556667in}}%
\pgfpathlineto{\pgfqpoint{2.974688in}{3.556667in}}%
\pgfusepath{stroke}%
\end{pgfscope}%
\begin{pgfscope}%
\definecolor{textcolor}{rgb}{0.000000,0.000000,0.000000}%
\pgfsetstrokecolor{textcolor}%
\pgfsetfillcolor{textcolor}%
\pgftext[x=3.108021in,y=3.498333in,left,base]{\color{textcolor}{\rmfamily\fontsize{12.000000}{14.400000}\selectfont\catcode`\^=\active\def^{\ifmmode\sp\else\^{}\fi}\catcode`\%=\active\def%{\%}Огибающая $A_{\mathrm{вых}}(\omega)$}}%
\end{pgfscope}%
\end{pgfpicture}%
\makeatother%
\endgroup%

    \caption{Амплитудный дискретный спектр реакции на второй входной сигнал}
    \label{fig:plot_disc_spec_A_out_2}
\end{figure}

\begin{figure}[H]
    \centering
    %% Creator: Matplotlib, PGF backend
%%
%% To include the figure in your LaTeX document, write
%%   \input{<filename>.pgf}
%%
%% Make sure the required packages are loaded in your preamble
%%   \usepackage{pgf}
%%
%% Also ensure that all the required font packages are loaded; for instance,
%% the lmodern package is sometimes necessary when using math font.
%%   \usepackage{lmodern}
%%
%% Figures using additional raster images can only be included by \input if
%% they are in the same directory as the main LaTeX file. For loading figures
%% from other directories you can use the `import` package
%%   \usepackage{import}
%%
%% and then include the figures with
%%   \import{<path to file>}{<filename>.pgf}
%%
%% Matplotlib used the following preamble
%%   \def\mathdefault#1{#1}
%%   \everymath=\expandafter{\the\everymath\displaystyle}
%%   \IfFileExists{scrextend.sty}{
%%     \usepackage[fontsize=12.000000pt]{scrextend}
%%   }{
%%     \renewcommand{\normalsize}{\fontsize{12.000000}{14.400000}\selectfont}
%%     \normalsize
%%   }
%%   \usepackage{fontspec}\setmainfont{Times New Roman}\usepackage[english,russian]{babel}\usepackage{amsmath}
%%   \ifdefined\pdftexversion\else  % non-pdftex case.
%%     \usepackage{fontspec}
%%     \setmainfont{times.ttf}[Path=\detokenize{/usr/local/share/fonts/truetype/times-new-roman/}]
%%     \setsansfont{DejaVuSans.ttf}[Path=\detokenize{/usr/share/matplotlib/mpl-data/fonts/ttf/}]
%%     \setmonofont{DejaVuSansMono.ttf}[Path=\detokenize{/usr/share/matplotlib/mpl-data/fonts/ttf/}]
%%   \fi
%%   \makeatletter\@ifpackageloaded{underscore}{}{\usepackage[strings]{underscore}}\makeatother
%%
\begingroup%
\makeatletter%
\begin{pgfpicture}%
\pgfpathrectangle{\pgfpointorigin}{\pgfqpoint{6.490000in}{3.740000in}}%
\pgfusepath{use as bounding box, clip}%
\begin{pgfscope}%
\pgfsetbuttcap%
\pgfsetmiterjoin%
\definecolor{currentfill}{rgb}{1.000000,1.000000,1.000000}%
\pgfsetfillcolor{currentfill}%
\pgfsetlinewidth{0.000000pt}%
\definecolor{currentstroke}{rgb}{1.000000,1.000000,1.000000}%
\pgfsetstrokecolor{currentstroke}%
\pgfsetdash{}{0pt}%
\pgfpathmoveto{\pgfqpoint{0.000000in}{0.000000in}}%
\pgfpathlineto{\pgfqpoint{6.490000in}{0.000000in}}%
\pgfpathlineto{\pgfqpoint{6.490000in}{3.740000in}}%
\pgfpathlineto{\pgfqpoint{0.000000in}{3.740000in}}%
\pgfpathlineto{\pgfqpoint{0.000000in}{0.000000in}}%
\pgfpathclose%
\pgfusepath{fill}%
\end{pgfscope}%
\begin{pgfscope}%
\pgfsetbuttcap%
\pgfsetmiterjoin%
\definecolor{currentfill}{rgb}{1.000000,1.000000,1.000000}%
\pgfsetfillcolor{currentfill}%
\pgfsetlinewidth{0.000000pt}%
\definecolor{currentstroke}{rgb}{0.000000,0.000000,0.000000}%
\pgfsetstrokecolor{currentstroke}%
\pgfsetstrokeopacity{0.000000}%
\pgfsetdash{}{0pt}%
\pgfpathmoveto{\pgfqpoint{0.695253in}{0.456901in}}%
\pgfpathlineto{\pgfqpoint{6.440000in}{0.456901in}}%
\pgfpathlineto{\pgfqpoint{6.440000in}{3.273333in}}%
\pgfpathlineto{\pgfqpoint{0.695253in}{3.273333in}}%
\pgfpathlineto{\pgfqpoint{0.695253in}{0.456901in}}%
\pgfpathclose%
\pgfusepath{fill}%
\end{pgfscope}%
\begin{pgfscope}%
\pgfsetbuttcap%
\pgfsetroundjoin%
\definecolor{currentfill}{rgb}{1.000000,1.000000,1.000000}%
\pgfsetfillcolor{currentfill}%
\pgfsetlinewidth{0.803000pt}%
\definecolor{currentstroke}{rgb}{0.000000,0.000000,0.000000}%
\pgfsetstrokecolor{currentstroke}%
\pgfsetdash{}{0pt}%
\pgfsys@defobject{currentmarker}{\pgfqpoint{0.000000in}{0.000000in}}{\pgfqpoint{0.000000in}{0.048611in}}{%
\pgfpathmoveto{\pgfqpoint{0.000000in}{0.000000in}}%
\pgfpathlineto{\pgfqpoint{0.000000in}{0.048611in}}%
\pgfusepath{stroke,fill}%
}%
\begin{pgfscope}%
\pgfsys@transformshift{0.710128in}{0.456901in}%
\pgfsys@useobject{currentmarker}{}%
\end{pgfscope}%
\end{pgfscope}%
\begin{pgfscope}%
\definecolor{textcolor}{rgb}{0.000000,0.000000,0.000000}%
\pgfsetstrokecolor{textcolor}%
\pgfsetfillcolor{textcolor}%
\pgftext[x=0.710128in,y=0.408290in,,top]{\color{textcolor}{\rmfamily\fontsize{12.000000}{14.400000}\selectfont\catcode`\^=\active\def^{\ifmmode\sp\else\^{}\fi}\catcode`\%=\active\def%{\%}$0$}}%
\end{pgfscope}%
\begin{pgfscope}%
\pgfsetbuttcap%
\pgfsetroundjoin%
\definecolor{currentfill}{rgb}{1.000000,1.000000,1.000000}%
\pgfsetfillcolor{currentfill}%
\pgfsetlinewidth{0.803000pt}%
\definecolor{currentstroke}{rgb}{0.000000,0.000000,0.000000}%
\pgfsetstrokecolor{currentstroke}%
\pgfsetdash{}{0pt}%
\pgfsys@defobject{currentmarker}{\pgfqpoint{0.000000in}{0.000000in}}{\pgfqpoint{0.000000in}{0.048611in}}{%
\pgfpathmoveto{\pgfqpoint{0.000000in}{0.000000in}}%
\pgfpathlineto{\pgfqpoint{0.000000in}{0.048611in}}%
\pgfusepath{stroke,fill}%
}%
\begin{pgfscope}%
\pgfsys@transformshift{1.702313in}{0.456901in}%
\pgfsys@useobject{currentmarker}{}%
\end{pgfscope}%
\end{pgfscope}%
\begin{pgfscope}%
\definecolor{textcolor}{rgb}{0.000000,0.000000,0.000000}%
\pgfsetstrokecolor{textcolor}%
\pgfsetfillcolor{textcolor}%
\pgftext[x=1.702313in,y=0.408290in,,top]{\color{textcolor}{\rmfamily\fontsize{12.000000}{14.400000}\selectfont\catcode`\^=\active\def^{\ifmmode\sp\else\^{}\fi}\catcode`\%=\active\def%{\%}$1\omega_1$}}%
\end{pgfscope}%
\begin{pgfscope}%
\pgfsetbuttcap%
\pgfsetroundjoin%
\definecolor{currentfill}{rgb}{1.000000,1.000000,1.000000}%
\pgfsetfillcolor{currentfill}%
\pgfsetlinewidth{0.803000pt}%
\definecolor{currentstroke}{rgb}{0.000000,0.000000,0.000000}%
\pgfsetstrokecolor{currentstroke}%
\pgfsetdash{}{0pt}%
\pgfsys@defobject{currentmarker}{\pgfqpoint{0.000000in}{0.000000in}}{\pgfqpoint{0.000000in}{0.048611in}}{%
\pgfpathmoveto{\pgfqpoint{0.000000in}{0.000000in}}%
\pgfpathlineto{\pgfqpoint{0.000000in}{0.048611in}}%
\pgfusepath{stroke,fill}%
}%
\begin{pgfscope}%
\pgfsys@transformshift{2.694499in}{0.456901in}%
\pgfsys@useobject{currentmarker}{}%
\end{pgfscope}%
\end{pgfscope}%
\begin{pgfscope}%
\definecolor{textcolor}{rgb}{0.000000,0.000000,0.000000}%
\pgfsetstrokecolor{textcolor}%
\pgfsetfillcolor{textcolor}%
\pgftext[x=2.694499in,y=0.408290in,,top]{\color{textcolor}{\rmfamily\fontsize{12.000000}{14.400000}\selectfont\catcode`\^=\active\def^{\ifmmode\sp\else\^{}\fi}\catcode`\%=\active\def%{\%}$2\omega_1$}}%
\end{pgfscope}%
\begin{pgfscope}%
\pgfsetbuttcap%
\pgfsetroundjoin%
\definecolor{currentfill}{rgb}{1.000000,1.000000,1.000000}%
\pgfsetfillcolor{currentfill}%
\pgfsetlinewidth{0.803000pt}%
\definecolor{currentstroke}{rgb}{0.000000,0.000000,0.000000}%
\pgfsetstrokecolor{currentstroke}%
\pgfsetdash{}{0pt}%
\pgfsys@defobject{currentmarker}{\pgfqpoint{0.000000in}{0.000000in}}{\pgfqpoint{0.000000in}{0.048611in}}{%
\pgfpathmoveto{\pgfqpoint{0.000000in}{0.000000in}}%
\pgfpathlineto{\pgfqpoint{0.000000in}{0.048611in}}%
\pgfusepath{stroke,fill}%
}%
\begin{pgfscope}%
\pgfsys@transformshift{3.686685in}{0.456901in}%
\pgfsys@useobject{currentmarker}{}%
\end{pgfscope}%
\end{pgfscope}%
\begin{pgfscope}%
\definecolor{textcolor}{rgb}{0.000000,0.000000,0.000000}%
\pgfsetstrokecolor{textcolor}%
\pgfsetfillcolor{textcolor}%
\pgftext[x=3.686685in,y=0.408290in,,top]{\color{textcolor}{\rmfamily\fontsize{12.000000}{14.400000}\selectfont\catcode`\^=\active\def^{\ifmmode\sp\else\^{}\fi}\catcode`\%=\active\def%{\%}$3\omega_1$}}%
\end{pgfscope}%
\begin{pgfscope}%
\pgfsetbuttcap%
\pgfsetroundjoin%
\definecolor{currentfill}{rgb}{1.000000,1.000000,1.000000}%
\pgfsetfillcolor{currentfill}%
\pgfsetlinewidth{0.803000pt}%
\definecolor{currentstroke}{rgb}{0.000000,0.000000,0.000000}%
\pgfsetstrokecolor{currentstroke}%
\pgfsetdash{}{0pt}%
\pgfsys@defobject{currentmarker}{\pgfqpoint{0.000000in}{0.000000in}}{\pgfqpoint{0.000000in}{0.048611in}}{%
\pgfpathmoveto{\pgfqpoint{0.000000in}{0.000000in}}%
\pgfpathlineto{\pgfqpoint{0.000000in}{0.048611in}}%
\pgfusepath{stroke,fill}%
}%
\begin{pgfscope}%
\pgfsys@transformshift{4.678870in}{0.456901in}%
\pgfsys@useobject{currentmarker}{}%
\end{pgfscope}%
\end{pgfscope}%
\begin{pgfscope}%
\definecolor{textcolor}{rgb}{0.000000,0.000000,0.000000}%
\pgfsetstrokecolor{textcolor}%
\pgfsetfillcolor{textcolor}%
\pgftext[x=4.678870in,y=0.408290in,,top]{\color{textcolor}{\rmfamily\fontsize{12.000000}{14.400000}\selectfont\catcode`\^=\active\def^{\ifmmode\sp\else\^{}\fi}\catcode`\%=\active\def%{\%}$4\omega_1$}}%
\end{pgfscope}%
\begin{pgfscope}%
\pgfsetbuttcap%
\pgfsetroundjoin%
\definecolor{currentfill}{rgb}{1.000000,1.000000,1.000000}%
\pgfsetfillcolor{currentfill}%
\pgfsetlinewidth{0.803000pt}%
\definecolor{currentstroke}{rgb}{0.000000,0.000000,0.000000}%
\pgfsetstrokecolor{currentstroke}%
\pgfsetdash{}{0pt}%
\pgfsys@defobject{currentmarker}{\pgfqpoint{0.000000in}{0.000000in}}{\pgfqpoint{0.000000in}{0.048611in}}{%
\pgfpathmoveto{\pgfqpoint{0.000000in}{0.000000in}}%
\pgfpathlineto{\pgfqpoint{0.000000in}{0.048611in}}%
\pgfusepath{stroke,fill}%
}%
\begin{pgfscope}%
\pgfsys@transformshift{5.671056in}{0.456901in}%
\pgfsys@useobject{currentmarker}{}%
\end{pgfscope}%
\end{pgfscope}%
\begin{pgfscope}%
\definecolor{textcolor}{rgb}{0.000000,0.000000,0.000000}%
\pgfsetstrokecolor{textcolor}%
\pgfsetfillcolor{textcolor}%
\pgftext[x=5.671056in,y=0.408290in,,top]{\color{textcolor}{\rmfamily\fontsize{12.000000}{14.400000}\selectfont\catcode`\^=\active\def^{\ifmmode\sp\else\^{}\fi}\catcode`\%=\active\def%{\%}$5\omega_1$}}%
\end{pgfscope}%
\begin{pgfscope}%
\pgfpathrectangle{\pgfqpoint{0.695253in}{0.456901in}}{\pgfqpoint{5.744747in}{2.816432in}}%
\pgfusepath{clip}%
\pgfsetbuttcap%
\pgfsetroundjoin%
\pgfsetlinewidth{0.501875pt}%
\definecolor{currentstroke}{rgb}{0.501961,0.501961,0.501961}%
\pgfsetstrokecolor{currentstroke}%
\pgfsetstrokeopacity{0.400000}%
\pgfsetdash{{1.850000pt}{0.800000pt}}{0.000000pt}%
\pgfpathmoveto{\pgfqpoint{0.958174in}{0.456901in}}%
\pgfpathlineto{\pgfqpoint{0.958174in}{3.273333in}}%
\pgfusepath{stroke}%
\end{pgfscope}%
\begin{pgfscope}%
\pgfsetbuttcap%
\pgfsetroundjoin%
\definecolor{currentfill}{rgb}{1.000000,1.000000,1.000000}%
\pgfsetfillcolor{currentfill}%
\pgfsetlinewidth{0.602250pt}%
\definecolor{currentstroke}{rgb}{0.000000,0.000000,0.000000}%
\pgfsetstrokecolor{currentstroke}%
\pgfsetdash{}{0pt}%
\pgfsys@defobject{currentmarker}{\pgfqpoint{0.000000in}{0.000000in}}{\pgfqpoint{0.000000in}{0.027778in}}{%
\pgfpathmoveto{\pgfqpoint{0.000000in}{0.000000in}}%
\pgfpathlineto{\pgfqpoint{0.000000in}{0.027778in}}%
\pgfusepath{stroke,fill}%
}%
\begin{pgfscope}%
\pgfsys@transformshift{0.958174in}{0.456901in}%
\pgfsys@useobject{currentmarker}{}%
\end{pgfscope}%
\end{pgfscope}%
\begin{pgfscope}%
\pgfpathrectangle{\pgfqpoint{0.695253in}{0.456901in}}{\pgfqpoint{5.744747in}{2.816432in}}%
\pgfusepath{clip}%
\pgfsetbuttcap%
\pgfsetroundjoin%
\pgfsetlinewidth{0.501875pt}%
\definecolor{currentstroke}{rgb}{0.501961,0.501961,0.501961}%
\pgfsetstrokecolor{currentstroke}%
\pgfsetstrokeopacity{0.400000}%
\pgfsetdash{{1.850000pt}{0.800000pt}}{0.000000pt}%
\pgfpathmoveto{\pgfqpoint{1.206221in}{0.456901in}}%
\pgfpathlineto{\pgfqpoint{1.206221in}{3.273333in}}%
\pgfusepath{stroke}%
\end{pgfscope}%
\begin{pgfscope}%
\pgfsetbuttcap%
\pgfsetroundjoin%
\definecolor{currentfill}{rgb}{1.000000,1.000000,1.000000}%
\pgfsetfillcolor{currentfill}%
\pgfsetlinewidth{0.602250pt}%
\definecolor{currentstroke}{rgb}{0.000000,0.000000,0.000000}%
\pgfsetstrokecolor{currentstroke}%
\pgfsetdash{}{0pt}%
\pgfsys@defobject{currentmarker}{\pgfqpoint{0.000000in}{0.000000in}}{\pgfqpoint{0.000000in}{0.027778in}}{%
\pgfpathmoveto{\pgfqpoint{0.000000in}{0.000000in}}%
\pgfpathlineto{\pgfqpoint{0.000000in}{0.027778in}}%
\pgfusepath{stroke,fill}%
}%
\begin{pgfscope}%
\pgfsys@transformshift{1.206221in}{0.456901in}%
\pgfsys@useobject{currentmarker}{}%
\end{pgfscope}%
\end{pgfscope}%
\begin{pgfscope}%
\pgfpathrectangle{\pgfqpoint{0.695253in}{0.456901in}}{\pgfqpoint{5.744747in}{2.816432in}}%
\pgfusepath{clip}%
\pgfsetbuttcap%
\pgfsetroundjoin%
\pgfsetlinewidth{0.501875pt}%
\definecolor{currentstroke}{rgb}{0.501961,0.501961,0.501961}%
\pgfsetstrokecolor{currentstroke}%
\pgfsetstrokeopacity{0.400000}%
\pgfsetdash{{1.850000pt}{0.800000pt}}{0.000000pt}%
\pgfpathmoveto{\pgfqpoint{1.454267in}{0.456901in}}%
\pgfpathlineto{\pgfqpoint{1.454267in}{3.273333in}}%
\pgfusepath{stroke}%
\end{pgfscope}%
\begin{pgfscope}%
\pgfsetbuttcap%
\pgfsetroundjoin%
\definecolor{currentfill}{rgb}{1.000000,1.000000,1.000000}%
\pgfsetfillcolor{currentfill}%
\pgfsetlinewidth{0.602250pt}%
\definecolor{currentstroke}{rgb}{0.000000,0.000000,0.000000}%
\pgfsetstrokecolor{currentstroke}%
\pgfsetdash{}{0pt}%
\pgfsys@defobject{currentmarker}{\pgfqpoint{0.000000in}{0.000000in}}{\pgfqpoint{0.000000in}{0.027778in}}{%
\pgfpathmoveto{\pgfqpoint{0.000000in}{0.000000in}}%
\pgfpathlineto{\pgfqpoint{0.000000in}{0.027778in}}%
\pgfusepath{stroke,fill}%
}%
\begin{pgfscope}%
\pgfsys@transformshift{1.454267in}{0.456901in}%
\pgfsys@useobject{currentmarker}{}%
\end{pgfscope}%
\end{pgfscope}%
\begin{pgfscope}%
\pgfpathrectangle{\pgfqpoint{0.695253in}{0.456901in}}{\pgfqpoint{5.744747in}{2.816432in}}%
\pgfusepath{clip}%
\pgfsetbuttcap%
\pgfsetroundjoin%
\pgfsetlinewidth{0.501875pt}%
\definecolor{currentstroke}{rgb}{0.501961,0.501961,0.501961}%
\pgfsetstrokecolor{currentstroke}%
\pgfsetstrokeopacity{0.400000}%
\pgfsetdash{{1.850000pt}{0.800000pt}}{0.000000pt}%
\pgfpathmoveto{\pgfqpoint{1.950360in}{0.456901in}}%
\pgfpathlineto{\pgfqpoint{1.950360in}{3.273333in}}%
\pgfusepath{stroke}%
\end{pgfscope}%
\begin{pgfscope}%
\pgfsetbuttcap%
\pgfsetroundjoin%
\definecolor{currentfill}{rgb}{1.000000,1.000000,1.000000}%
\pgfsetfillcolor{currentfill}%
\pgfsetlinewidth{0.602250pt}%
\definecolor{currentstroke}{rgb}{0.000000,0.000000,0.000000}%
\pgfsetstrokecolor{currentstroke}%
\pgfsetdash{}{0pt}%
\pgfsys@defobject{currentmarker}{\pgfqpoint{0.000000in}{0.000000in}}{\pgfqpoint{0.000000in}{0.027778in}}{%
\pgfpathmoveto{\pgfqpoint{0.000000in}{0.000000in}}%
\pgfpathlineto{\pgfqpoint{0.000000in}{0.027778in}}%
\pgfusepath{stroke,fill}%
}%
\begin{pgfscope}%
\pgfsys@transformshift{1.950360in}{0.456901in}%
\pgfsys@useobject{currentmarker}{}%
\end{pgfscope}%
\end{pgfscope}%
\begin{pgfscope}%
\pgfpathrectangle{\pgfqpoint{0.695253in}{0.456901in}}{\pgfqpoint{5.744747in}{2.816432in}}%
\pgfusepath{clip}%
\pgfsetbuttcap%
\pgfsetroundjoin%
\pgfsetlinewidth{0.501875pt}%
\definecolor{currentstroke}{rgb}{0.501961,0.501961,0.501961}%
\pgfsetstrokecolor{currentstroke}%
\pgfsetstrokeopacity{0.400000}%
\pgfsetdash{{1.850000pt}{0.800000pt}}{0.000000pt}%
\pgfpathmoveto{\pgfqpoint{2.198406in}{0.456901in}}%
\pgfpathlineto{\pgfqpoint{2.198406in}{3.273333in}}%
\pgfusepath{stroke}%
\end{pgfscope}%
\begin{pgfscope}%
\pgfsetbuttcap%
\pgfsetroundjoin%
\definecolor{currentfill}{rgb}{1.000000,1.000000,1.000000}%
\pgfsetfillcolor{currentfill}%
\pgfsetlinewidth{0.602250pt}%
\definecolor{currentstroke}{rgb}{0.000000,0.000000,0.000000}%
\pgfsetstrokecolor{currentstroke}%
\pgfsetdash{}{0pt}%
\pgfsys@defobject{currentmarker}{\pgfqpoint{0.000000in}{0.000000in}}{\pgfqpoint{0.000000in}{0.027778in}}{%
\pgfpathmoveto{\pgfqpoint{0.000000in}{0.000000in}}%
\pgfpathlineto{\pgfqpoint{0.000000in}{0.027778in}}%
\pgfusepath{stroke,fill}%
}%
\begin{pgfscope}%
\pgfsys@transformshift{2.198406in}{0.456901in}%
\pgfsys@useobject{currentmarker}{}%
\end{pgfscope}%
\end{pgfscope}%
\begin{pgfscope}%
\pgfpathrectangle{\pgfqpoint{0.695253in}{0.456901in}}{\pgfqpoint{5.744747in}{2.816432in}}%
\pgfusepath{clip}%
\pgfsetbuttcap%
\pgfsetroundjoin%
\pgfsetlinewidth{0.501875pt}%
\definecolor{currentstroke}{rgb}{0.501961,0.501961,0.501961}%
\pgfsetstrokecolor{currentstroke}%
\pgfsetstrokeopacity{0.400000}%
\pgfsetdash{{1.850000pt}{0.800000pt}}{0.000000pt}%
\pgfpathmoveto{\pgfqpoint{2.446453in}{0.456901in}}%
\pgfpathlineto{\pgfqpoint{2.446453in}{3.273333in}}%
\pgfusepath{stroke}%
\end{pgfscope}%
\begin{pgfscope}%
\pgfsetbuttcap%
\pgfsetroundjoin%
\definecolor{currentfill}{rgb}{1.000000,1.000000,1.000000}%
\pgfsetfillcolor{currentfill}%
\pgfsetlinewidth{0.602250pt}%
\definecolor{currentstroke}{rgb}{0.000000,0.000000,0.000000}%
\pgfsetstrokecolor{currentstroke}%
\pgfsetdash{}{0pt}%
\pgfsys@defobject{currentmarker}{\pgfqpoint{0.000000in}{0.000000in}}{\pgfqpoint{0.000000in}{0.027778in}}{%
\pgfpathmoveto{\pgfqpoint{0.000000in}{0.000000in}}%
\pgfpathlineto{\pgfqpoint{0.000000in}{0.027778in}}%
\pgfusepath{stroke,fill}%
}%
\begin{pgfscope}%
\pgfsys@transformshift{2.446453in}{0.456901in}%
\pgfsys@useobject{currentmarker}{}%
\end{pgfscope}%
\end{pgfscope}%
\begin{pgfscope}%
\pgfpathrectangle{\pgfqpoint{0.695253in}{0.456901in}}{\pgfqpoint{5.744747in}{2.816432in}}%
\pgfusepath{clip}%
\pgfsetbuttcap%
\pgfsetroundjoin%
\pgfsetlinewidth{0.501875pt}%
\definecolor{currentstroke}{rgb}{0.501961,0.501961,0.501961}%
\pgfsetstrokecolor{currentstroke}%
\pgfsetstrokeopacity{0.400000}%
\pgfsetdash{{1.850000pt}{0.800000pt}}{0.000000pt}%
\pgfpathmoveto{\pgfqpoint{2.942546in}{0.456901in}}%
\pgfpathlineto{\pgfqpoint{2.942546in}{3.273333in}}%
\pgfusepath{stroke}%
\end{pgfscope}%
\begin{pgfscope}%
\pgfsetbuttcap%
\pgfsetroundjoin%
\definecolor{currentfill}{rgb}{1.000000,1.000000,1.000000}%
\pgfsetfillcolor{currentfill}%
\pgfsetlinewidth{0.602250pt}%
\definecolor{currentstroke}{rgb}{0.000000,0.000000,0.000000}%
\pgfsetstrokecolor{currentstroke}%
\pgfsetdash{}{0pt}%
\pgfsys@defobject{currentmarker}{\pgfqpoint{0.000000in}{0.000000in}}{\pgfqpoint{0.000000in}{0.027778in}}{%
\pgfpathmoveto{\pgfqpoint{0.000000in}{0.000000in}}%
\pgfpathlineto{\pgfqpoint{0.000000in}{0.027778in}}%
\pgfusepath{stroke,fill}%
}%
\begin{pgfscope}%
\pgfsys@transformshift{2.942546in}{0.456901in}%
\pgfsys@useobject{currentmarker}{}%
\end{pgfscope}%
\end{pgfscope}%
\begin{pgfscope}%
\pgfpathrectangle{\pgfqpoint{0.695253in}{0.456901in}}{\pgfqpoint{5.744747in}{2.816432in}}%
\pgfusepath{clip}%
\pgfsetbuttcap%
\pgfsetroundjoin%
\pgfsetlinewidth{0.501875pt}%
\definecolor{currentstroke}{rgb}{0.501961,0.501961,0.501961}%
\pgfsetstrokecolor{currentstroke}%
\pgfsetstrokeopacity{0.400000}%
\pgfsetdash{{1.850000pt}{0.800000pt}}{0.000000pt}%
\pgfpathmoveto{\pgfqpoint{3.190592in}{0.456901in}}%
\pgfpathlineto{\pgfqpoint{3.190592in}{3.273333in}}%
\pgfusepath{stroke}%
\end{pgfscope}%
\begin{pgfscope}%
\pgfsetbuttcap%
\pgfsetroundjoin%
\definecolor{currentfill}{rgb}{1.000000,1.000000,1.000000}%
\pgfsetfillcolor{currentfill}%
\pgfsetlinewidth{0.602250pt}%
\definecolor{currentstroke}{rgb}{0.000000,0.000000,0.000000}%
\pgfsetstrokecolor{currentstroke}%
\pgfsetdash{}{0pt}%
\pgfsys@defobject{currentmarker}{\pgfqpoint{0.000000in}{0.000000in}}{\pgfqpoint{0.000000in}{0.027778in}}{%
\pgfpathmoveto{\pgfqpoint{0.000000in}{0.000000in}}%
\pgfpathlineto{\pgfqpoint{0.000000in}{0.027778in}}%
\pgfusepath{stroke,fill}%
}%
\begin{pgfscope}%
\pgfsys@transformshift{3.190592in}{0.456901in}%
\pgfsys@useobject{currentmarker}{}%
\end{pgfscope}%
\end{pgfscope}%
\begin{pgfscope}%
\pgfpathrectangle{\pgfqpoint{0.695253in}{0.456901in}}{\pgfqpoint{5.744747in}{2.816432in}}%
\pgfusepath{clip}%
\pgfsetbuttcap%
\pgfsetroundjoin%
\pgfsetlinewidth{0.501875pt}%
\definecolor{currentstroke}{rgb}{0.501961,0.501961,0.501961}%
\pgfsetstrokecolor{currentstroke}%
\pgfsetstrokeopacity{0.400000}%
\pgfsetdash{{1.850000pt}{0.800000pt}}{0.000000pt}%
\pgfpathmoveto{\pgfqpoint{3.438638in}{0.456901in}}%
\pgfpathlineto{\pgfqpoint{3.438638in}{3.273333in}}%
\pgfusepath{stroke}%
\end{pgfscope}%
\begin{pgfscope}%
\pgfsetbuttcap%
\pgfsetroundjoin%
\definecolor{currentfill}{rgb}{1.000000,1.000000,1.000000}%
\pgfsetfillcolor{currentfill}%
\pgfsetlinewidth{0.602250pt}%
\definecolor{currentstroke}{rgb}{0.000000,0.000000,0.000000}%
\pgfsetstrokecolor{currentstroke}%
\pgfsetdash{}{0pt}%
\pgfsys@defobject{currentmarker}{\pgfqpoint{0.000000in}{0.000000in}}{\pgfqpoint{0.000000in}{0.027778in}}{%
\pgfpathmoveto{\pgfqpoint{0.000000in}{0.000000in}}%
\pgfpathlineto{\pgfqpoint{0.000000in}{0.027778in}}%
\pgfusepath{stroke,fill}%
}%
\begin{pgfscope}%
\pgfsys@transformshift{3.438638in}{0.456901in}%
\pgfsys@useobject{currentmarker}{}%
\end{pgfscope}%
\end{pgfscope}%
\begin{pgfscope}%
\pgfpathrectangle{\pgfqpoint{0.695253in}{0.456901in}}{\pgfqpoint{5.744747in}{2.816432in}}%
\pgfusepath{clip}%
\pgfsetbuttcap%
\pgfsetroundjoin%
\pgfsetlinewidth{0.501875pt}%
\definecolor{currentstroke}{rgb}{0.501961,0.501961,0.501961}%
\pgfsetstrokecolor{currentstroke}%
\pgfsetstrokeopacity{0.400000}%
\pgfsetdash{{1.850000pt}{0.800000pt}}{0.000000pt}%
\pgfpathmoveto{\pgfqpoint{3.934731in}{0.456901in}}%
\pgfpathlineto{\pgfqpoint{3.934731in}{3.273333in}}%
\pgfusepath{stroke}%
\end{pgfscope}%
\begin{pgfscope}%
\pgfsetbuttcap%
\pgfsetroundjoin%
\definecolor{currentfill}{rgb}{1.000000,1.000000,1.000000}%
\pgfsetfillcolor{currentfill}%
\pgfsetlinewidth{0.602250pt}%
\definecolor{currentstroke}{rgb}{0.000000,0.000000,0.000000}%
\pgfsetstrokecolor{currentstroke}%
\pgfsetdash{}{0pt}%
\pgfsys@defobject{currentmarker}{\pgfqpoint{0.000000in}{0.000000in}}{\pgfqpoint{0.000000in}{0.027778in}}{%
\pgfpathmoveto{\pgfqpoint{0.000000in}{0.000000in}}%
\pgfpathlineto{\pgfqpoint{0.000000in}{0.027778in}}%
\pgfusepath{stroke,fill}%
}%
\begin{pgfscope}%
\pgfsys@transformshift{3.934731in}{0.456901in}%
\pgfsys@useobject{currentmarker}{}%
\end{pgfscope}%
\end{pgfscope}%
\begin{pgfscope}%
\pgfpathrectangle{\pgfqpoint{0.695253in}{0.456901in}}{\pgfqpoint{5.744747in}{2.816432in}}%
\pgfusepath{clip}%
\pgfsetbuttcap%
\pgfsetroundjoin%
\pgfsetlinewidth{0.501875pt}%
\definecolor{currentstroke}{rgb}{0.501961,0.501961,0.501961}%
\pgfsetstrokecolor{currentstroke}%
\pgfsetstrokeopacity{0.400000}%
\pgfsetdash{{1.850000pt}{0.800000pt}}{0.000000pt}%
\pgfpathmoveto{\pgfqpoint{4.182778in}{0.456901in}}%
\pgfpathlineto{\pgfqpoint{4.182778in}{3.273333in}}%
\pgfusepath{stroke}%
\end{pgfscope}%
\begin{pgfscope}%
\pgfsetbuttcap%
\pgfsetroundjoin%
\definecolor{currentfill}{rgb}{1.000000,1.000000,1.000000}%
\pgfsetfillcolor{currentfill}%
\pgfsetlinewidth{0.602250pt}%
\definecolor{currentstroke}{rgb}{0.000000,0.000000,0.000000}%
\pgfsetstrokecolor{currentstroke}%
\pgfsetdash{}{0pt}%
\pgfsys@defobject{currentmarker}{\pgfqpoint{0.000000in}{0.000000in}}{\pgfqpoint{0.000000in}{0.027778in}}{%
\pgfpathmoveto{\pgfqpoint{0.000000in}{0.000000in}}%
\pgfpathlineto{\pgfqpoint{0.000000in}{0.027778in}}%
\pgfusepath{stroke,fill}%
}%
\begin{pgfscope}%
\pgfsys@transformshift{4.182778in}{0.456901in}%
\pgfsys@useobject{currentmarker}{}%
\end{pgfscope}%
\end{pgfscope}%
\begin{pgfscope}%
\pgfpathrectangle{\pgfqpoint{0.695253in}{0.456901in}}{\pgfqpoint{5.744747in}{2.816432in}}%
\pgfusepath{clip}%
\pgfsetbuttcap%
\pgfsetroundjoin%
\pgfsetlinewidth{0.501875pt}%
\definecolor{currentstroke}{rgb}{0.501961,0.501961,0.501961}%
\pgfsetstrokecolor{currentstroke}%
\pgfsetstrokeopacity{0.400000}%
\pgfsetdash{{1.850000pt}{0.800000pt}}{0.000000pt}%
\pgfpathmoveto{\pgfqpoint{4.430824in}{0.456901in}}%
\pgfpathlineto{\pgfqpoint{4.430824in}{3.273333in}}%
\pgfusepath{stroke}%
\end{pgfscope}%
\begin{pgfscope}%
\pgfsetbuttcap%
\pgfsetroundjoin%
\definecolor{currentfill}{rgb}{1.000000,1.000000,1.000000}%
\pgfsetfillcolor{currentfill}%
\pgfsetlinewidth{0.602250pt}%
\definecolor{currentstroke}{rgb}{0.000000,0.000000,0.000000}%
\pgfsetstrokecolor{currentstroke}%
\pgfsetdash{}{0pt}%
\pgfsys@defobject{currentmarker}{\pgfqpoint{0.000000in}{0.000000in}}{\pgfqpoint{0.000000in}{0.027778in}}{%
\pgfpathmoveto{\pgfqpoint{0.000000in}{0.000000in}}%
\pgfpathlineto{\pgfqpoint{0.000000in}{0.027778in}}%
\pgfusepath{stroke,fill}%
}%
\begin{pgfscope}%
\pgfsys@transformshift{4.430824in}{0.456901in}%
\pgfsys@useobject{currentmarker}{}%
\end{pgfscope}%
\end{pgfscope}%
\begin{pgfscope}%
\pgfpathrectangle{\pgfqpoint{0.695253in}{0.456901in}}{\pgfqpoint{5.744747in}{2.816432in}}%
\pgfusepath{clip}%
\pgfsetbuttcap%
\pgfsetroundjoin%
\pgfsetlinewidth{0.501875pt}%
\definecolor{currentstroke}{rgb}{0.501961,0.501961,0.501961}%
\pgfsetstrokecolor{currentstroke}%
\pgfsetstrokeopacity{0.400000}%
\pgfsetdash{{1.850000pt}{0.800000pt}}{0.000000pt}%
\pgfpathmoveto{\pgfqpoint{4.926917in}{0.456901in}}%
\pgfpathlineto{\pgfqpoint{4.926917in}{3.273333in}}%
\pgfusepath{stroke}%
\end{pgfscope}%
\begin{pgfscope}%
\pgfsetbuttcap%
\pgfsetroundjoin%
\definecolor{currentfill}{rgb}{1.000000,1.000000,1.000000}%
\pgfsetfillcolor{currentfill}%
\pgfsetlinewidth{0.602250pt}%
\definecolor{currentstroke}{rgb}{0.000000,0.000000,0.000000}%
\pgfsetstrokecolor{currentstroke}%
\pgfsetdash{}{0pt}%
\pgfsys@defobject{currentmarker}{\pgfqpoint{0.000000in}{0.000000in}}{\pgfqpoint{0.000000in}{0.027778in}}{%
\pgfpathmoveto{\pgfqpoint{0.000000in}{0.000000in}}%
\pgfpathlineto{\pgfqpoint{0.000000in}{0.027778in}}%
\pgfusepath{stroke,fill}%
}%
\begin{pgfscope}%
\pgfsys@transformshift{4.926917in}{0.456901in}%
\pgfsys@useobject{currentmarker}{}%
\end{pgfscope}%
\end{pgfscope}%
\begin{pgfscope}%
\pgfpathrectangle{\pgfqpoint{0.695253in}{0.456901in}}{\pgfqpoint{5.744747in}{2.816432in}}%
\pgfusepath{clip}%
\pgfsetbuttcap%
\pgfsetroundjoin%
\pgfsetlinewidth{0.501875pt}%
\definecolor{currentstroke}{rgb}{0.501961,0.501961,0.501961}%
\pgfsetstrokecolor{currentstroke}%
\pgfsetstrokeopacity{0.400000}%
\pgfsetdash{{1.850000pt}{0.800000pt}}{0.000000pt}%
\pgfpathmoveto{\pgfqpoint{5.174963in}{0.456901in}}%
\pgfpathlineto{\pgfqpoint{5.174963in}{3.273333in}}%
\pgfusepath{stroke}%
\end{pgfscope}%
\begin{pgfscope}%
\pgfsetbuttcap%
\pgfsetroundjoin%
\definecolor{currentfill}{rgb}{1.000000,1.000000,1.000000}%
\pgfsetfillcolor{currentfill}%
\pgfsetlinewidth{0.602250pt}%
\definecolor{currentstroke}{rgb}{0.000000,0.000000,0.000000}%
\pgfsetstrokecolor{currentstroke}%
\pgfsetdash{}{0pt}%
\pgfsys@defobject{currentmarker}{\pgfqpoint{0.000000in}{0.000000in}}{\pgfqpoint{0.000000in}{0.027778in}}{%
\pgfpathmoveto{\pgfqpoint{0.000000in}{0.000000in}}%
\pgfpathlineto{\pgfqpoint{0.000000in}{0.027778in}}%
\pgfusepath{stroke,fill}%
}%
\begin{pgfscope}%
\pgfsys@transformshift{5.174963in}{0.456901in}%
\pgfsys@useobject{currentmarker}{}%
\end{pgfscope}%
\end{pgfscope}%
\begin{pgfscope}%
\pgfpathrectangle{\pgfqpoint{0.695253in}{0.456901in}}{\pgfqpoint{5.744747in}{2.816432in}}%
\pgfusepath{clip}%
\pgfsetbuttcap%
\pgfsetroundjoin%
\pgfsetlinewidth{0.501875pt}%
\definecolor{currentstroke}{rgb}{0.501961,0.501961,0.501961}%
\pgfsetstrokecolor{currentstroke}%
\pgfsetstrokeopacity{0.400000}%
\pgfsetdash{{1.850000pt}{0.800000pt}}{0.000000pt}%
\pgfpathmoveto{\pgfqpoint{5.423010in}{0.456901in}}%
\pgfpathlineto{\pgfqpoint{5.423010in}{3.273333in}}%
\pgfusepath{stroke}%
\end{pgfscope}%
\begin{pgfscope}%
\pgfsetbuttcap%
\pgfsetroundjoin%
\definecolor{currentfill}{rgb}{1.000000,1.000000,1.000000}%
\pgfsetfillcolor{currentfill}%
\pgfsetlinewidth{0.602250pt}%
\definecolor{currentstroke}{rgb}{0.000000,0.000000,0.000000}%
\pgfsetstrokecolor{currentstroke}%
\pgfsetdash{}{0pt}%
\pgfsys@defobject{currentmarker}{\pgfqpoint{0.000000in}{0.000000in}}{\pgfqpoint{0.000000in}{0.027778in}}{%
\pgfpathmoveto{\pgfqpoint{0.000000in}{0.000000in}}%
\pgfpathlineto{\pgfqpoint{0.000000in}{0.027778in}}%
\pgfusepath{stroke,fill}%
}%
\begin{pgfscope}%
\pgfsys@transformshift{5.423010in}{0.456901in}%
\pgfsys@useobject{currentmarker}{}%
\end{pgfscope}%
\end{pgfscope}%
\begin{pgfscope}%
\pgfpathrectangle{\pgfqpoint{0.695253in}{0.456901in}}{\pgfqpoint{5.744747in}{2.816432in}}%
\pgfusepath{clip}%
\pgfsetbuttcap%
\pgfsetroundjoin%
\pgfsetlinewidth{0.501875pt}%
\definecolor{currentstroke}{rgb}{0.501961,0.501961,0.501961}%
\pgfsetstrokecolor{currentstroke}%
\pgfsetstrokeopacity{0.400000}%
\pgfsetdash{{1.850000pt}{0.800000pt}}{0.000000pt}%
\pgfpathmoveto{\pgfqpoint{5.919103in}{0.456901in}}%
\pgfpathlineto{\pgfqpoint{5.919103in}{3.273333in}}%
\pgfusepath{stroke}%
\end{pgfscope}%
\begin{pgfscope}%
\pgfsetbuttcap%
\pgfsetroundjoin%
\definecolor{currentfill}{rgb}{1.000000,1.000000,1.000000}%
\pgfsetfillcolor{currentfill}%
\pgfsetlinewidth{0.602250pt}%
\definecolor{currentstroke}{rgb}{0.000000,0.000000,0.000000}%
\pgfsetstrokecolor{currentstroke}%
\pgfsetdash{}{0pt}%
\pgfsys@defobject{currentmarker}{\pgfqpoint{0.000000in}{0.000000in}}{\pgfqpoint{0.000000in}{0.027778in}}{%
\pgfpathmoveto{\pgfqpoint{0.000000in}{0.000000in}}%
\pgfpathlineto{\pgfqpoint{0.000000in}{0.027778in}}%
\pgfusepath{stroke,fill}%
}%
\begin{pgfscope}%
\pgfsys@transformshift{5.919103in}{0.456901in}%
\pgfsys@useobject{currentmarker}{}%
\end{pgfscope}%
\end{pgfscope}%
\begin{pgfscope}%
\pgfpathrectangle{\pgfqpoint{0.695253in}{0.456901in}}{\pgfqpoint{5.744747in}{2.816432in}}%
\pgfusepath{clip}%
\pgfsetbuttcap%
\pgfsetroundjoin%
\pgfsetlinewidth{0.501875pt}%
\definecolor{currentstroke}{rgb}{0.501961,0.501961,0.501961}%
\pgfsetstrokecolor{currentstroke}%
\pgfsetstrokeopacity{0.400000}%
\pgfsetdash{{1.850000pt}{0.800000pt}}{0.000000pt}%
\pgfpathmoveto{\pgfqpoint{6.167149in}{0.456901in}}%
\pgfpathlineto{\pgfqpoint{6.167149in}{3.273333in}}%
\pgfusepath{stroke}%
\end{pgfscope}%
\begin{pgfscope}%
\pgfsetbuttcap%
\pgfsetroundjoin%
\definecolor{currentfill}{rgb}{1.000000,1.000000,1.000000}%
\pgfsetfillcolor{currentfill}%
\pgfsetlinewidth{0.602250pt}%
\definecolor{currentstroke}{rgb}{0.000000,0.000000,0.000000}%
\pgfsetstrokecolor{currentstroke}%
\pgfsetdash{}{0pt}%
\pgfsys@defobject{currentmarker}{\pgfqpoint{0.000000in}{0.000000in}}{\pgfqpoint{0.000000in}{0.027778in}}{%
\pgfpathmoveto{\pgfqpoint{0.000000in}{0.000000in}}%
\pgfpathlineto{\pgfqpoint{0.000000in}{0.027778in}}%
\pgfusepath{stroke,fill}%
}%
\begin{pgfscope}%
\pgfsys@transformshift{6.167149in}{0.456901in}%
\pgfsys@useobject{currentmarker}{}%
\end{pgfscope}%
\end{pgfscope}%
\begin{pgfscope}%
\pgfpathrectangle{\pgfqpoint{0.695253in}{0.456901in}}{\pgfqpoint{5.744747in}{2.816432in}}%
\pgfusepath{clip}%
\pgfsetbuttcap%
\pgfsetroundjoin%
\pgfsetlinewidth{0.501875pt}%
\definecolor{currentstroke}{rgb}{0.501961,0.501961,0.501961}%
\pgfsetstrokecolor{currentstroke}%
\pgfsetstrokeopacity{0.400000}%
\pgfsetdash{{1.850000pt}{0.800000pt}}{0.000000pt}%
\pgfpathmoveto{\pgfqpoint{6.415195in}{0.456901in}}%
\pgfpathlineto{\pgfqpoint{6.415195in}{3.273333in}}%
\pgfusepath{stroke}%
\end{pgfscope}%
\begin{pgfscope}%
\pgfsetbuttcap%
\pgfsetroundjoin%
\definecolor{currentfill}{rgb}{1.000000,1.000000,1.000000}%
\pgfsetfillcolor{currentfill}%
\pgfsetlinewidth{0.602250pt}%
\definecolor{currentstroke}{rgb}{0.000000,0.000000,0.000000}%
\pgfsetstrokecolor{currentstroke}%
\pgfsetdash{}{0pt}%
\pgfsys@defobject{currentmarker}{\pgfqpoint{0.000000in}{0.000000in}}{\pgfqpoint{0.000000in}{0.027778in}}{%
\pgfpathmoveto{\pgfqpoint{0.000000in}{0.000000in}}%
\pgfpathlineto{\pgfqpoint{0.000000in}{0.027778in}}%
\pgfusepath{stroke,fill}%
}%
\begin{pgfscope}%
\pgfsys@transformshift{6.415195in}{0.456901in}%
\pgfsys@useobject{currentmarker}{}%
\end{pgfscope}%
\end{pgfscope}%
\begin{pgfscope}%
\definecolor{textcolor}{rgb}{0.000000,0.000000,0.000000}%
\pgfsetstrokecolor{textcolor}%
\pgfsetfillcolor{textcolor}%
\pgftext[x=3.567626in,y=0.201367in,,top]{\color{textcolor}{\rmfamily\fontsize{12.000000}{14.400000}\selectfont\catcode`\^=\active\def^{\ifmmode\sp\else\^{}\fi}\catcode`\%=\active\def%{\%}$\omega$}}%
\end{pgfscope}%
\begin{pgfscope}%
\pgfsetbuttcap%
\pgfsetroundjoin%
\definecolor{currentfill}{rgb}{1.000000,1.000000,1.000000}%
\pgfsetfillcolor{currentfill}%
\pgfsetlinewidth{0.803000pt}%
\definecolor{currentstroke}{rgb}{0.000000,0.000000,0.000000}%
\pgfsetstrokecolor{currentstroke}%
\pgfsetdash{}{0pt}%
\pgfsys@defobject{currentmarker}{\pgfqpoint{0.000000in}{0.000000in}}{\pgfqpoint{0.048611in}{0.000000in}}{%
\pgfpathmoveto{\pgfqpoint{0.000000in}{0.000000in}}%
\pgfpathlineto{\pgfqpoint{0.048611in}{0.000000in}}%
\pgfusepath{stroke,fill}%
}%
\begin{pgfscope}%
\pgfsys@transformshift{0.695253in}{0.795924in}%
\pgfsys@useobject{currentmarker}{}%
\end{pgfscope}%
\end{pgfscope}%
\begin{pgfscope}%
\definecolor{textcolor}{rgb}{0.000000,0.000000,0.000000}%
\pgfsetstrokecolor{textcolor}%
\pgfsetfillcolor{textcolor}%
\pgftext[x=0.272222in, y=0.738063in, left, base]{\color{textcolor}{\rmfamily\fontsize{12.000000}{14.400000}\selectfont\catcode`\^=\active\def^{\ifmmode\sp\else\^{}\fi}\catcode`\%=\active\def%{\%}$\mathdefault{\ensuremath{-}150}$}}%
\end{pgfscope}%
\begin{pgfscope}%
\pgfsetbuttcap%
\pgfsetroundjoin%
\definecolor{currentfill}{rgb}{1.000000,1.000000,1.000000}%
\pgfsetfillcolor{currentfill}%
\pgfsetlinewidth{0.803000pt}%
\definecolor{currentstroke}{rgb}{0.000000,0.000000,0.000000}%
\pgfsetstrokecolor{currentstroke}%
\pgfsetdash{}{0pt}%
\pgfsys@defobject{currentmarker}{\pgfqpoint{0.000000in}{0.000000in}}{\pgfqpoint{0.048611in}{0.000000in}}{%
\pgfpathmoveto{\pgfqpoint{0.000000in}{0.000000in}}%
\pgfpathlineto{\pgfqpoint{0.048611in}{0.000000in}}%
\pgfusepath{stroke,fill}%
}%
\begin{pgfscope}%
\pgfsys@transformshift{0.695253in}{1.153272in}%
\pgfsys@useobject{currentmarker}{}%
\end{pgfscope}%
\end{pgfscope}%
\begin{pgfscope}%
\definecolor{textcolor}{rgb}{0.000000,0.000000,0.000000}%
\pgfsetstrokecolor{textcolor}%
\pgfsetfillcolor{textcolor}%
\pgftext[x=0.272222in, y=1.095411in, left, base]{\color{textcolor}{\rmfamily\fontsize{12.000000}{14.400000}\selectfont\catcode`\^=\active\def^{\ifmmode\sp\else\^{}\fi}\catcode`\%=\active\def%{\%}$\mathdefault{\ensuremath{-}100}$}}%
\end{pgfscope}%
\begin{pgfscope}%
\pgfsetbuttcap%
\pgfsetroundjoin%
\definecolor{currentfill}{rgb}{1.000000,1.000000,1.000000}%
\pgfsetfillcolor{currentfill}%
\pgfsetlinewidth{0.803000pt}%
\definecolor{currentstroke}{rgb}{0.000000,0.000000,0.000000}%
\pgfsetstrokecolor{currentstroke}%
\pgfsetdash{}{0pt}%
\pgfsys@defobject{currentmarker}{\pgfqpoint{0.000000in}{0.000000in}}{\pgfqpoint{0.048611in}{0.000000in}}{%
\pgfpathmoveto{\pgfqpoint{0.000000in}{0.000000in}}%
\pgfpathlineto{\pgfqpoint{0.048611in}{0.000000in}}%
\pgfusepath{stroke,fill}%
}%
\begin{pgfscope}%
\pgfsys@transformshift{0.695253in}{1.510619in}%
\pgfsys@useobject{currentmarker}{}%
\end{pgfscope}%
\end{pgfscope}%
\begin{pgfscope}%
\definecolor{textcolor}{rgb}{0.000000,0.000000,0.000000}%
\pgfsetstrokecolor{textcolor}%
\pgfsetfillcolor{textcolor}%
\pgftext[x=0.353819in, y=1.452758in, left, base]{\color{textcolor}{\rmfamily\fontsize{12.000000}{14.400000}\selectfont\catcode`\^=\active\def^{\ifmmode\sp\else\^{}\fi}\catcode`\%=\active\def%{\%}$\mathdefault{\ensuremath{-}50}$}}%
\end{pgfscope}%
\begin{pgfscope}%
\pgfsetbuttcap%
\pgfsetroundjoin%
\definecolor{currentfill}{rgb}{1.000000,1.000000,1.000000}%
\pgfsetfillcolor{currentfill}%
\pgfsetlinewidth{0.803000pt}%
\definecolor{currentstroke}{rgb}{0.000000,0.000000,0.000000}%
\pgfsetstrokecolor{currentstroke}%
\pgfsetdash{}{0pt}%
\pgfsys@defobject{currentmarker}{\pgfqpoint{0.000000in}{0.000000in}}{\pgfqpoint{0.048611in}{0.000000in}}{%
\pgfpathmoveto{\pgfqpoint{0.000000in}{0.000000in}}%
\pgfpathlineto{\pgfqpoint{0.048611in}{0.000000in}}%
\pgfusepath{stroke,fill}%
}%
\begin{pgfscope}%
\pgfsys@transformshift{0.695253in}{1.867967in}%
\pgfsys@useobject{currentmarker}{}%
\end{pgfscope}%
\end{pgfscope}%
\begin{pgfscope}%
\definecolor{textcolor}{rgb}{0.000000,0.000000,0.000000}%
\pgfsetstrokecolor{textcolor}%
\pgfsetfillcolor{textcolor}%
\pgftext[x=0.565045in, y=1.810106in, left, base]{\color{textcolor}{\rmfamily\fontsize{12.000000}{14.400000}\selectfont\catcode`\^=\active\def^{\ifmmode\sp\else\^{}\fi}\catcode`\%=\active\def%{\%}$\mathdefault{0}$}}%
\end{pgfscope}%
\begin{pgfscope}%
\pgfsetbuttcap%
\pgfsetroundjoin%
\definecolor{currentfill}{rgb}{1.000000,1.000000,1.000000}%
\pgfsetfillcolor{currentfill}%
\pgfsetlinewidth{0.803000pt}%
\definecolor{currentstroke}{rgb}{0.000000,0.000000,0.000000}%
\pgfsetstrokecolor{currentstroke}%
\pgfsetdash{}{0pt}%
\pgfsys@defobject{currentmarker}{\pgfqpoint{0.000000in}{0.000000in}}{\pgfqpoint{0.048611in}{0.000000in}}{%
\pgfpathmoveto{\pgfqpoint{0.000000in}{0.000000in}}%
\pgfpathlineto{\pgfqpoint{0.048611in}{0.000000in}}%
\pgfusepath{stroke,fill}%
}%
\begin{pgfscope}%
\pgfsys@transformshift{0.695253in}{2.225314in}%
\pgfsys@useobject{currentmarker}{}%
\end{pgfscope}%
\end{pgfscope}%
\begin{pgfscope}%
\definecolor{textcolor}{rgb}{0.000000,0.000000,0.000000}%
\pgfsetstrokecolor{textcolor}%
\pgfsetfillcolor{textcolor}%
\pgftext[x=0.483449in, y=2.167453in, left, base]{\color{textcolor}{\rmfamily\fontsize{12.000000}{14.400000}\selectfont\catcode`\^=\active\def^{\ifmmode\sp\else\^{}\fi}\catcode`\%=\active\def%{\%}$\mathdefault{50}$}}%
\end{pgfscope}%
\begin{pgfscope}%
\pgfsetbuttcap%
\pgfsetroundjoin%
\definecolor{currentfill}{rgb}{1.000000,1.000000,1.000000}%
\pgfsetfillcolor{currentfill}%
\pgfsetlinewidth{0.803000pt}%
\definecolor{currentstroke}{rgb}{0.000000,0.000000,0.000000}%
\pgfsetstrokecolor{currentstroke}%
\pgfsetdash{}{0pt}%
\pgfsys@defobject{currentmarker}{\pgfqpoint{0.000000in}{0.000000in}}{\pgfqpoint{0.048611in}{0.000000in}}{%
\pgfpathmoveto{\pgfqpoint{0.000000in}{0.000000in}}%
\pgfpathlineto{\pgfqpoint{0.048611in}{0.000000in}}%
\pgfusepath{stroke,fill}%
}%
\begin{pgfscope}%
\pgfsys@transformshift{0.695253in}{2.582662in}%
\pgfsys@useobject{currentmarker}{}%
\end{pgfscope}%
\end{pgfscope}%
\begin{pgfscope}%
\definecolor{textcolor}{rgb}{0.000000,0.000000,0.000000}%
\pgfsetstrokecolor{textcolor}%
\pgfsetfillcolor{textcolor}%
\pgftext[x=0.401852in, y=2.524801in, left, base]{\color{textcolor}{\rmfamily\fontsize{12.000000}{14.400000}\selectfont\catcode`\^=\active\def^{\ifmmode\sp\else\^{}\fi}\catcode`\%=\active\def%{\%}$\mathdefault{100}$}}%
\end{pgfscope}%
\begin{pgfscope}%
\pgfsetbuttcap%
\pgfsetroundjoin%
\definecolor{currentfill}{rgb}{1.000000,1.000000,1.000000}%
\pgfsetfillcolor{currentfill}%
\pgfsetlinewidth{0.803000pt}%
\definecolor{currentstroke}{rgb}{0.000000,0.000000,0.000000}%
\pgfsetstrokecolor{currentstroke}%
\pgfsetdash{}{0pt}%
\pgfsys@defobject{currentmarker}{\pgfqpoint{0.000000in}{0.000000in}}{\pgfqpoint{0.048611in}{0.000000in}}{%
\pgfpathmoveto{\pgfqpoint{0.000000in}{0.000000in}}%
\pgfpathlineto{\pgfqpoint{0.048611in}{0.000000in}}%
\pgfusepath{stroke,fill}%
}%
\begin{pgfscope}%
\pgfsys@transformshift{0.695253in}{2.940009in}%
\pgfsys@useobject{currentmarker}{}%
\end{pgfscope}%
\end{pgfscope}%
\begin{pgfscope}%
\definecolor{textcolor}{rgb}{0.000000,0.000000,0.000000}%
\pgfsetstrokecolor{textcolor}%
\pgfsetfillcolor{textcolor}%
\pgftext[x=0.401852in, y=2.882148in, left, base]{\color{textcolor}{\rmfamily\fontsize{12.000000}{14.400000}\selectfont\catcode`\^=\active\def^{\ifmmode\sp\else\^{}\fi}\catcode`\%=\active\def%{\%}$\mathdefault{150}$}}%
\end{pgfscope}%
\begin{pgfscope}%
\definecolor{textcolor}{rgb}{0.000000,0.000000,0.000000}%
\pgfsetstrokecolor{textcolor}%
\pgfsetfillcolor{textcolor}%
\pgftext[x=0.216667in,y=1.865117in,,bottom,rotate=90.000000]{\color{textcolor}{\rmfamily\fontsize{12.000000}{14.400000}\selectfont\catcode`\^=\active\def^{\ifmmode\sp\else\^{}\fi}\catcode`\%=\active\def%{\%}$\Phi_{\mathrm{вых}}(\omega), ^\circ$}}%
\end{pgfscope}%
\begin{pgfscope}%
\pgfpathrectangle{\pgfqpoint{0.695253in}{0.456901in}}{\pgfqpoint{5.744747in}{2.816432in}}%
\pgfusepath{clip}%
\pgfsetbuttcap%
\pgfsetroundjoin%
\pgfsetlinewidth{2.007500pt}%
\definecolor{currentstroke}{rgb}{1.000000,0.000000,0.000000}%
\pgfsetstrokecolor{currentstroke}%
\pgfsetdash{{7.400000pt}{3.200000pt}}{0.000000pt}%
\pgfpathmoveto{\pgfqpoint{0.710128in}{2.511192in}}%
\pgfpathlineto{\pgfqpoint{1.071013in}{1.952517in}}%
\pgfpathlineto{\pgfqpoint{1.256923in}{1.660761in}}%
\pgfpathlineto{\pgfqpoint{1.410026in}{1.416508in}}%
\pgfpathlineto{\pgfqpoint{1.541257in}{1.203009in}}%
\pgfpathlineto{\pgfqpoint{1.661552in}{1.002829in}}%
\pgfpathlineto{\pgfqpoint{1.770911in}{0.816157in}}%
\pgfpathlineto{\pgfqpoint{1.869335in}{0.643546in}}%
\pgfpathlineto{\pgfqpoint{1.902142in}{0.584920in}}%
\pgfpathlineto{\pgfqpoint{1.913078in}{3.138153in}}%
\pgfpathlineto{\pgfqpoint{2.000566in}{2.978400in}}%
\pgfpathlineto{\pgfqpoint{2.088053in}{2.814266in}}%
\pgfpathlineto{\pgfqpoint{2.186476in}{2.624514in}}%
\pgfpathlineto{\pgfqpoint{2.295835in}{2.408394in}}%
\pgfpathlineto{\pgfqpoint{2.645785in}{1.711842in}}%
\pgfpathlineto{\pgfqpoint{2.733272in}{1.544721in}}%
\pgfpathlineto{\pgfqpoint{2.820759in}{1.382246in}}%
\pgfpathlineto{\pgfqpoint{2.908247in}{1.224497in}}%
\pgfpathlineto{\pgfqpoint{2.941054in}{1.166518in}}%
\pgfpathlineto{\pgfqpoint{2.951990in}{2.433781in}}%
\pgfpathlineto{\pgfqpoint{3.039477in}{2.282631in}}%
\pgfpathlineto{\pgfqpoint{3.137901in}{2.117191in}}%
\pgfpathlineto{\pgfqpoint{3.236324in}{1.955989in}}%
\pgfpathlineto{\pgfqpoint{3.345683in}{1.781093in}}%
\pgfpathlineto{\pgfqpoint{3.465978in}{1.592975in}}%
\pgfpathlineto{\pgfqpoint{3.597209in}{1.391929in}}%
\pgfpathlineto{\pgfqpoint{3.739376in}{1.178095in}}%
\pgfpathlineto{\pgfqpoint{3.903415in}{0.935452in}}%
\pgfpathlineto{\pgfqpoint{4.089325in}{0.664630in}}%
\pgfpathlineto{\pgfqpoint{4.144005in}{0.585677in}}%
\pgfpathlineto{\pgfqpoint{4.154941in}{3.142823in}}%
\pgfpathlineto{\pgfqpoint{4.373659in}{2.829832in}}%
\pgfpathlineto{\pgfqpoint{4.614249in}{2.489490in}}%
\pgfpathlineto{\pgfqpoint{4.887647in}{2.106634in}}%
\pgfpathlineto{\pgfqpoint{5.193853in}{1.681672in}}%
\pgfpathlineto{\pgfqpoint{5.543802in}{1.199883in}}%
\pgfpathlineto{\pgfqpoint{5.937495in}{0.661760in}}%
\pgfpathlineto{\pgfqpoint{5.992174in}{0.587297in}}%
\pgfpathlineto{\pgfqpoint{6.003110in}{3.145314in}}%
\pgfpathlineto{\pgfqpoint{6.167149in}{2.922317in}}%
\pgfpathlineto{\pgfqpoint{6.167149in}{2.922317in}}%
\pgfusepath{stroke}%
\end{pgfscope}%
\begin{pgfscope}%
\pgfpathrectangle{\pgfqpoint{0.695253in}{0.456901in}}{\pgfqpoint{5.744747in}{2.816432in}}%
\pgfusepath{clip}%
\pgfsetbuttcap%
\pgfsetroundjoin%
\pgfsetlinewidth{2.007500pt}%
\definecolor{currentstroke}{rgb}{0.000000,0.000000,1.000000}%
\pgfsetstrokecolor{currentstroke}%
\pgfsetdash{}{0pt}%
\pgfpathmoveto{\pgfqpoint{0.710128in}{1.867967in}}%
\pgfpathlineto{\pgfqpoint{0.710128in}{2.511192in}}%
\pgfusepath{stroke}%
\end{pgfscope}%
\begin{pgfscope}%
\pgfpathrectangle{\pgfqpoint{0.695253in}{0.456901in}}{\pgfqpoint{5.744747in}{2.816432in}}%
\pgfusepath{clip}%
\pgfsetbuttcap%
\pgfsetroundjoin%
\pgfsetlinewidth{2.007500pt}%
\definecolor{currentstroke}{rgb}{0.000000,0.000000,1.000000}%
\pgfsetstrokecolor{currentstroke}%
\pgfsetdash{}{0pt}%
\pgfpathmoveto{\pgfqpoint{1.702313in}{1.867967in}}%
\pgfpathlineto{\pgfqpoint{1.702313in}{0.933826in}}%
\pgfusepath{stroke}%
\end{pgfscope}%
\begin{pgfscope}%
\pgfpathrectangle{\pgfqpoint{0.695253in}{0.456901in}}{\pgfqpoint{5.744747in}{2.816432in}}%
\pgfusepath{clip}%
\pgfsetbuttcap%
\pgfsetroundjoin%
\pgfsetlinewidth{2.007500pt}%
\definecolor{currentstroke}{rgb}{0.000000,0.000000,1.000000}%
\pgfsetstrokecolor{currentstroke}%
\pgfsetdash{}{0pt}%
\pgfpathmoveto{\pgfqpoint{2.694499in}{1.867967in}}%
\pgfpathlineto{\pgfqpoint{2.694499in}{1.618233in}}%
\pgfusepath{stroke}%
\end{pgfscope}%
\begin{pgfscope}%
\pgfpathrectangle{\pgfqpoint{0.695253in}{0.456901in}}{\pgfqpoint{5.744747in}{2.816432in}}%
\pgfusepath{clip}%
\pgfsetbuttcap%
\pgfsetroundjoin%
\pgfsetlinewidth{2.007500pt}%
\definecolor{currentstroke}{rgb}{0.000000,0.000000,1.000000}%
\pgfsetstrokecolor{currentstroke}%
\pgfsetdash{}{0pt}%
\pgfpathmoveto{\pgfqpoint{3.686685in}{1.867967in}}%
\pgfpathlineto{\pgfqpoint{3.686685in}{1.256921in}}%
\pgfusepath{stroke}%
\end{pgfscope}%
\begin{pgfscope}%
\pgfpathrectangle{\pgfqpoint{0.695253in}{0.456901in}}{\pgfqpoint{5.744747in}{2.816432in}}%
\pgfusepath{clip}%
\pgfsetbuttcap%
\pgfsetroundjoin%
\pgfsetlinewidth{2.007500pt}%
\definecolor{currentstroke}{rgb}{0.000000,0.000000,1.000000}%
\pgfsetstrokecolor{currentstroke}%
\pgfsetdash{}{0pt}%
\pgfpathmoveto{\pgfqpoint{4.678870in}{1.867967in}}%
\pgfpathlineto{\pgfqpoint{4.678870in}{2.398659in}}%
\pgfusepath{stroke}%
\end{pgfscope}%
\begin{pgfscope}%
\pgfpathrectangle{\pgfqpoint{0.695253in}{0.456901in}}{\pgfqpoint{5.744747in}{2.816432in}}%
\pgfusepath{clip}%
\pgfsetbuttcap%
\pgfsetroundjoin%
\pgfsetlinewidth{2.007500pt}%
\definecolor{currentstroke}{rgb}{0.000000,0.000000,1.000000}%
\pgfsetstrokecolor{currentstroke}%
\pgfsetdash{}{0pt}%
\pgfpathmoveto{\pgfqpoint{5.671056in}{1.867967in}}%
\pgfpathlineto{\pgfqpoint{5.671056in}{1.025539in}}%
\pgfusepath{stroke}%
\end{pgfscope}%
\begin{pgfscope}%
\pgfpathrectangle{\pgfqpoint{0.695253in}{0.456901in}}{\pgfqpoint{5.744747in}{2.816432in}}%
\pgfusepath{clip}%
\pgfsetbuttcap%
\pgfsetroundjoin%
\definecolor{currentfill}{rgb}{1.000000,1.000000,1.000000}%
\pgfsetfillcolor{currentfill}%
\pgfsetlinewidth{1.505625pt}%
\definecolor{currentstroke}{rgb}{0.000000,0.000000,1.000000}%
\pgfsetstrokecolor{currentstroke}%
\pgfsetdash{}{0pt}%
\pgfsys@defobject{currentmarker}{\pgfqpoint{-0.027778in}{-0.027778in}}{\pgfqpoint{0.027778in}{0.027778in}}{%
\pgfpathmoveto{\pgfqpoint{0.000000in}{-0.027778in}}%
\pgfpathcurveto{\pgfqpoint{0.007367in}{-0.027778in}}{\pgfqpoint{0.014433in}{-0.024851in}}{\pgfqpoint{0.019642in}{-0.019642in}}%
\pgfpathcurveto{\pgfqpoint{0.024851in}{-0.014433in}}{\pgfqpoint{0.027778in}{-0.007367in}}{\pgfqpoint{0.027778in}{0.000000in}}%
\pgfpathcurveto{\pgfqpoint{0.027778in}{0.007367in}}{\pgfqpoint{0.024851in}{0.014433in}}{\pgfqpoint{0.019642in}{0.019642in}}%
\pgfpathcurveto{\pgfqpoint{0.014433in}{0.024851in}}{\pgfqpoint{0.007367in}{0.027778in}}{\pgfqpoint{0.000000in}{0.027778in}}%
\pgfpathcurveto{\pgfqpoint{-0.007367in}{0.027778in}}{\pgfqpoint{-0.014433in}{0.024851in}}{\pgfqpoint{-0.019642in}{0.019642in}}%
\pgfpathcurveto{\pgfqpoint{-0.024851in}{0.014433in}}{\pgfqpoint{-0.027778in}{0.007367in}}{\pgfqpoint{-0.027778in}{0.000000in}}%
\pgfpathcurveto{\pgfqpoint{-0.027778in}{-0.007367in}}{\pgfqpoint{-0.024851in}{-0.014433in}}{\pgfqpoint{-0.019642in}{-0.019642in}}%
\pgfpathcurveto{\pgfqpoint{-0.014433in}{-0.024851in}}{\pgfqpoint{-0.007367in}{-0.027778in}}{\pgfqpoint{0.000000in}{-0.027778in}}%
\pgfpathlineto{\pgfqpoint{0.000000in}{-0.027778in}}%
\pgfpathclose%
\pgfusepath{stroke,fill}%
}%
\begin{pgfscope}%
\pgfsys@transformshift{0.710128in}{2.511192in}%
\pgfsys@useobject{currentmarker}{}%
\end{pgfscope}%
\begin{pgfscope}%
\pgfsys@transformshift{1.702313in}{0.933826in}%
\pgfsys@useobject{currentmarker}{}%
\end{pgfscope}%
\begin{pgfscope}%
\pgfsys@transformshift{2.694499in}{1.618233in}%
\pgfsys@useobject{currentmarker}{}%
\end{pgfscope}%
\begin{pgfscope}%
\pgfsys@transformshift{3.686685in}{1.256921in}%
\pgfsys@useobject{currentmarker}{}%
\end{pgfscope}%
\begin{pgfscope}%
\pgfsys@transformshift{4.678870in}{2.398659in}%
\pgfsys@useobject{currentmarker}{}%
\end{pgfscope}%
\begin{pgfscope}%
\pgfsys@transformshift{5.671056in}{1.025539in}%
\pgfsys@useobject{currentmarker}{}%
\end{pgfscope}%
\end{pgfscope}%
\begin{pgfscope}%
\pgfpathrectangle{\pgfqpoint{0.695253in}{0.456901in}}{\pgfqpoint{5.744747in}{2.816432in}}%
\pgfusepath{clip}%
\pgfsetrectcap%
\pgfsetroundjoin%
\pgfsetlinewidth{2.007500pt}%
\definecolor{currentstroke}{rgb}{0.000000,0.000000,0.000000}%
\pgfsetstrokecolor{currentstroke}%
\pgfsetdash{}{0pt}%
\pgfpathmoveto{\pgfqpoint{0.710128in}{1.867967in}}%
\pgfpathlineto{\pgfqpoint{5.671056in}{1.867967in}}%
\pgfusepath{stroke}%
\end{pgfscope}%
\begin{pgfscope}%
\pgfsetrectcap%
\pgfsetmiterjoin%
\pgfsetlinewidth{0.602250pt}%
\definecolor{currentstroke}{rgb}{0.000000,0.000000,0.000000}%
\pgfsetstrokecolor{currentstroke}%
\pgfsetdash{}{0pt}%
\pgfpathmoveto{\pgfqpoint{0.695253in}{0.456901in}}%
\pgfpathlineto{\pgfqpoint{0.695253in}{3.273333in}}%
\pgfusepath{stroke}%
\end{pgfscope}%
\begin{pgfscope}%
\pgfsetrectcap%
\pgfsetmiterjoin%
\pgfsetlinewidth{0.602250pt}%
\definecolor{currentstroke}{rgb}{0.000000,0.000000,0.000000}%
\pgfsetstrokecolor{currentstroke}%
\pgfsetdash{}{0pt}%
\pgfpathmoveto{\pgfqpoint{0.695253in}{0.456901in}}%
\pgfpathlineto{\pgfqpoint{6.440000in}{0.456901in}}%
\pgfusepath{stroke}%
\end{pgfscope}%
\begin{pgfscope}%
\pgfsetbuttcap%
\pgfsetroundjoin%
\pgfsetlinewidth{2.007500pt}%
\definecolor{currentstroke}{rgb}{1.000000,0.000000,0.000000}%
\pgfsetstrokecolor{currentstroke}%
\pgfsetdash{{7.400000pt}{3.200000pt}}{0.000000pt}%
\pgfpathmoveto{\pgfqpoint{2.589040in}{3.556667in}}%
\pgfpathlineto{\pgfqpoint{2.989040in}{3.556667in}}%
\pgfusepath{stroke}%
\end{pgfscope}%
\begin{pgfscope}%
\definecolor{textcolor}{rgb}{0.000000,0.000000,0.000000}%
\pgfsetstrokecolor{textcolor}%
\pgfsetfillcolor{textcolor}%
\pgftext[x=3.172373in,y=3.498333in,left,base]{\color{textcolor}{\rmfamily\fontsize{12.000000}{14.400000}\selectfont\catcode`\^=\active\def^{\ifmmode\sp\else\^{}\fi}\catcode`\%=\active\def%{\%}Огибающая $\Phi_{\mathrm{вых}}(\omega)$}}%
\end{pgfscope}%
\end{pgfpicture}%
\makeatother%
\endgroup%

    \caption{Фазовый дискретный спектр реакции на второй входной сигнал}
    \label{fig:plot_disc_spec_Phi_out_2}
\end{figure}

Результаты расчетов для двух импульсов представим в виде таблицы:

\begin{longtblr}[
    caption = {Результаты расчета дискретных спектров реакции},
    label = {tab:spec_out_results},
    expand = \expanded,
]{
    colspec = { Q[c,m] X[c,m] X[c,m] X[c,m] X[c,m] X[c,m] X[c,m] },
    hlines, vlines,
    rowhead = 1,
}
    $k$ & 0 & 1 & 2 & 3 & 4 & 5 \\
    \midrule
    \SetCell[c=7]{c} Первый импульс & & & & & & \\

    $\omega = \omega_1 k$ & \pv*[.3]{ex11['tab_wk_1'][0]} & \pv*[.3]{ex11['tab_wk_1'][1]} & \pv*[.3]{ex11['tab_wk_1'][2]} & \pv*[.3]{ex11['tab_wk_1'][3]} & \pv*[.3]{ex11['tab_wk_1'][4]} & \pv*[.3]{ex11['tab_wk_1'][5]} \\
    
    $A_{\text{вых1}}(k)$ & \pv*[.3]{ex11['tab_Ak_1'][0]} & \pv*[.3]{ex11['tab_Ak_1'][1]} & \pv*[.3]{ex11['tab_Ak_1'][2]} & \pv*[.3]{ex11['tab_Ak_1'][3]} & \pv*[.3]{ex11['tab_Ak_1'][4]} & \pv*[.3]{ex11['tab_Ak_1'][5]} \\
    
    $\Phi_{\text{вых1}}(k),\,{}^\circ$ & \pv*[.1]{ex11['tab_Phik_1'][0]} & \pv*[.1]{ex11['tab_Phik_1'][1]} & \pv*[.1]{ex11['tab_Phik_1'][2]} & \pv*[.1]{ex11['tab_Phik_1'][3]} & \pv*[.1]{ex11['tab_Phik_1'][4]} & \pv*[.1]{ex11['tab_Phik_1'][5]} \\

    \midrule
    \SetCell[c=7]{c} Второй импульс & & & & & & \\
    
    $\omega = \omega_1 k$ & \pv*[.3]{ex11['tab_wk_2'][0]} & \pv*[.3]{ex11['tab_wk_2'][1]} & \pv*[.3]{ex11['tab_wk_2'][2]} & \pv*[.3]{ex11['tab_wk_2'][3]} & \pv*[.3]{ex11['tab_wk_2'][4]} & \pv*[.3]{ex11['tab_wk_2'][5]} \\
    
    $A_{\text{вых2}}(k)$ & \pv*[.3]{ex11['tab_Ak_2'][0]} & \pv*[.3]{ex11['tab_Ak_2'][1]} & \pv*[.3]{ex11['tab_Ak_2'][2]} & \pv*[.3]{ex11['tab_Ak_2'][3]} & \pv*[.3]{ex11['tab_Ak_2'][4]} & \pv*[.3]{ex11['tab_Ak_2'][5]} \\
    
    $\Phi_{\text{вых2}}(k),\,{}^\circ$ & \pv*[.1]{ex11['tab_Phik_2'][0]} & \pv*[.1]{ex11['tab_Phik_2'][1]} & \pv*[.1]{ex11['tab_Phik_2'][2]} & \pv*[.1]{ex11['tab_Phik_2'][3]} & \pv*[.1]{ex11['tab_Phik_2'][4]} & \pv*[.1]{ex11['tab_Phik_2'][5]} \\
\end{longtblr}

Таким образом, выходные сигналы в виде отрезка ряда Фурье имеют следующий вид:
{\sloppy
    \begin{enumerate}
        \item Первый: $i_{\text{н1}}(t) = \pv{ex11['fs_1']}$
        \item Второй: $i_{\text{н2}}(t) = \pv{ex11['fs_2']}$
    \end{enumerate}
\par}

На \cref{fig:plot_fourier_approx_out_1,fig:plot_fourier_approx_out_2} приведены графики ряда Фурье реакции, а также отдельные гармоники.

\begin{figure}[H]
    \centering
    %% Creator: Matplotlib, PGF backend
%%
%% To include the figure in your LaTeX document, write
%%   \input{<filename>.pgf}
%%
%% Make sure the required packages are loaded in your preamble
%%   \usepackage{pgf}
%%
%% Also ensure that all the required font packages are loaded; for instance,
%% the lmodern package is sometimes necessary when using math font.
%%   \usepackage{lmodern}
%%
%% Figures using additional raster images can only be included by \input if
%% they are in the same directory as the main LaTeX file. For loading figures
%% from other directories you can use the `import` package
%%   \usepackage{import}
%%
%% and then include the figures with
%%   \import{<path to file>}{<filename>.pgf}
%%
%% Matplotlib used the following preamble
%%   \def\mathdefault#1{#1}
%%   \everymath=\expandafter{\the\everymath\displaystyle}
%%   \IfFileExists{scrextend.sty}{
%%     \usepackage[fontsize=12.000000pt]{scrextend}
%%   }{
%%     \renewcommand{\normalsize}{\fontsize{12.000000}{14.400000}\selectfont}
%%     \normalsize
%%   }
%%   \usepackage{fontspec}\setmainfont{Times New Roman}\usepackage[english,russian]{babel}\usepackage{amsmath}
%%   \ifdefined\pdftexversion\else  % non-pdftex case.
%%     \usepackage{fontspec}
%%     \setmainfont{times.ttf}[Path=\detokenize{/usr/local/share/fonts/truetype/times-new-roman/}]
%%     \setsansfont{DejaVuSans.ttf}[Path=\detokenize{/usr/share/matplotlib/mpl-data/fonts/ttf/}]
%%     \setmonofont{DejaVuSansMono.ttf}[Path=\detokenize{/usr/share/matplotlib/mpl-data/fonts/ttf/}]
%%   \fi
%%   \makeatletter\@ifpackageloaded{underscore}{}{\usepackage[strings]{underscore}}\makeatother
%%
\begingroup%
\makeatletter%
\begin{pgfpicture}%
\pgfpathrectangle{\pgfpointorigin}{\pgfqpoint{6.316267in}{3.740000in}}%
\pgfusepath{use as bounding box, clip}%
\begin{pgfscope}%
\pgfsetbuttcap%
\pgfsetmiterjoin%
\definecolor{currentfill}{rgb}{1.000000,1.000000,1.000000}%
\pgfsetfillcolor{currentfill}%
\pgfsetlinewidth{0.000000pt}%
\definecolor{currentstroke}{rgb}{1.000000,1.000000,1.000000}%
\pgfsetstrokecolor{currentstroke}%
\pgfsetdash{}{0pt}%
\pgfpathmoveto{\pgfqpoint{0.000000in}{0.000000in}}%
\pgfpathlineto{\pgfqpoint{6.316267in}{0.000000in}}%
\pgfpathlineto{\pgfqpoint{6.316267in}{3.740000in}}%
\pgfpathlineto{\pgfqpoint{0.000000in}{3.740000in}}%
\pgfpathlineto{\pgfqpoint{0.000000in}{0.000000in}}%
\pgfpathclose%
\pgfusepath{fill}%
\end{pgfscope}%
\begin{pgfscope}%
\pgfsetbuttcap%
\pgfsetmiterjoin%
\definecolor{currentfill}{rgb}{1.000000,1.000000,1.000000}%
\pgfsetfillcolor{currentfill}%
\pgfsetlinewidth{0.000000pt}%
\definecolor{currentstroke}{rgb}{0.000000,0.000000,0.000000}%
\pgfsetstrokecolor{currentstroke}%
\pgfsetstrokeopacity{0.000000}%
\pgfsetdash{}{0pt}%
\pgfpathmoveto{\pgfqpoint{0.573726in}{0.456901in}}%
\pgfpathlineto{\pgfqpoint{6.152443in}{0.456901in}}%
\pgfpathlineto{\pgfqpoint{6.152443in}{3.273333in}}%
\pgfpathlineto{\pgfqpoint{0.573726in}{3.273333in}}%
\pgfpathlineto{\pgfqpoint{0.573726in}{0.456901in}}%
\pgfpathclose%
\pgfusepath{fill}%
\end{pgfscope}%
\begin{pgfscope}%
\pgfsetbuttcap%
\pgfsetroundjoin%
\definecolor{currentfill}{rgb}{1.000000,1.000000,1.000000}%
\pgfsetfillcolor{currentfill}%
\pgfsetlinewidth{0.803000pt}%
\definecolor{currentstroke}{rgb}{0.000000,0.000000,0.000000}%
\pgfsetstrokecolor{currentstroke}%
\pgfsetdash{}{0pt}%
\pgfsys@defobject{currentmarker}{\pgfqpoint{0.000000in}{0.000000in}}{\pgfqpoint{0.000000in}{0.048611in}}{%
\pgfpathmoveto{\pgfqpoint{0.000000in}{0.000000in}}%
\pgfpathlineto{\pgfqpoint{0.000000in}{0.048611in}}%
\pgfusepath{stroke,fill}%
}%
\begin{pgfscope}%
\pgfsys@transformshift{0.785041in}{0.456901in}%
\pgfsys@useobject{currentmarker}{}%
\end{pgfscope}%
\end{pgfscope}%
\begin{pgfscope}%
\definecolor{textcolor}{rgb}{0.000000,0.000000,0.000000}%
\pgfsetstrokecolor{textcolor}%
\pgfsetfillcolor{textcolor}%
\pgftext[x=0.785041in,y=0.408290in,,top]{\color{textcolor}{\rmfamily\fontsize{12.000000}{14.400000}\selectfont\catcode`\^=\active\def^{\ifmmode\sp\else\^{}\fi}\catcode`\%=\active\def%{\%}$\mathdefault{0}$}}%
\end{pgfscope}%
\begin{pgfscope}%
\pgfsetbuttcap%
\pgfsetroundjoin%
\definecolor{currentfill}{rgb}{1.000000,1.000000,1.000000}%
\pgfsetfillcolor{currentfill}%
\pgfsetlinewidth{0.803000pt}%
\definecolor{currentstroke}{rgb}{0.000000,0.000000,0.000000}%
\pgfsetstrokecolor{currentstroke}%
\pgfsetdash{}{0pt}%
\pgfsys@defobject{currentmarker}{\pgfqpoint{0.000000in}{0.000000in}}{\pgfqpoint{0.000000in}{0.048611in}}{%
\pgfpathmoveto{\pgfqpoint{0.000000in}{0.000000in}}%
\pgfpathlineto{\pgfqpoint{0.000000in}{0.048611in}}%
\pgfusepath{stroke,fill}%
}%
\begin{pgfscope}%
\pgfsys@transformshift{1.682345in}{0.456901in}%
\pgfsys@useobject{currentmarker}{}%
\end{pgfscope}%
\end{pgfscope}%
\begin{pgfscope}%
\definecolor{textcolor}{rgb}{0.000000,0.000000,0.000000}%
\pgfsetstrokecolor{textcolor}%
\pgfsetfillcolor{textcolor}%
\pgftext[x=1.682345in,y=0.408290in,,top]{\color{textcolor}{\rmfamily\fontsize{12.000000}{14.400000}\selectfont\catcode`\^=\active\def^{\ifmmode\sp\else\^{}\fi}\catcode`\%=\active\def%{\%}$\mathdefault{20}$}}%
\end{pgfscope}%
\begin{pgfscope}%
\pgfsetbuttcap%
\pgfsetroundjoin%
\definecolor{currentfill}{rgb}{1.000000,1.000000,1.000000}%
\pgfsetfillcolor{currentfill}%
\pgfsetlinewidth{0.803000pt}%
\definecolor{currentstroke}{rgb}{0.000000,0.000000,0.000000}%
\pgfsetstrokecolor{currentstroke}%
\pgfsetdash{}{0pt}%
\pgfsys@defobject{currentmarker}{\pgfqpoint{0.000000in}{0.000000in}}{\pgfqpoint{0.000000in}{0.048611in}}{%
\pgfpathmoveto{\pgfqpoint{0.000000in}{0.000000in}}%
\pgfpathlineto{\pgfqpoint{0.000000in}{0.048611in}}%
\pgfusepath{stroke,fill}%
}%
\begin{pgfscope}%
\pgfsys@transformshift{2.579649in}{0.456901in}%
\pgfsys@useobject{currentmarker}{}%
\end{pgfscope}%
\end{pgfscope}%
\begin{pgfscope}%
\definecolor{textcolor}{rgb}{0.000000,0.000000,0.000000}%
\pgfsetstrokecolor{textcolor}%
\pgfsetfillcolor{textcolor}%
\pgftext[x=2.579649in,y=0.408290in,,top]{\color{textcolor}{\rmfamily\fontsize{12.000000}{14.400000}\selectfont\catcode`\^=\active\def^{\ifmmode\sp\else\^{}\fi}\catcode`\%=\active\def%{\%}$\mathdefault{40}$}}%
\end{pgfscope}%
\begin{pgfscope}%
\pgfsetbuttcap%
\pgfsetroundjoin%
\definecolor{currentfill}{rgb}{1.000000,1.000000,1.000000}%
\pgfsetfillcolor{currentfill}%
\pgfsetlinewidth{0.803000pt}%
\definecolor{currentstroke}{rgb}{0.000000,0.000000,0.000000}%
\pgfsetstrokecolor{currentstroke}%
\pgfsetdash{}{0pt}%
\pgfsys@defobject{currentmarker}{\pgfqpoint{0.000000in}{0.000000in}}{\pgfqpoint{0.000000in}{0.048611in}}{%
\pgfpathmoveto{\pgfqpoint{0.000000in}{0.000000in}}%
\pgfpathlineto{\pgfqpoint{0.000000in}{0.048611in}}%
\pgfusepath{stroke,fill}%
}%
\begin{pgfscope}%
\pgfsys@transformshift{3.476953in}{0.456901in}%
\pgfsys@useobject{currentmarker}{}%
\end{pgfscope}%
\end{pgfscope}%
\begin{pgfscope}%
\definecolor{textcolor}{rgb}{0.000000,0.000000,0.000000}%
\pgfsetstrokecolor{textcolor}%
\pgfsetfillcolor{textcolor}%
\pgftext[x=3.476953in,y=0.408290in,,top]{\color{textcolor}{\rmfamily\fontsize{12.000000}{14.400000}\selectfont\catcode`\^=\active\def^{\ifmmode\sp\else\^{}\fi}\catcode`\%=\active\def%{\%}$\mathdefault{60}$}}%
\end{pgfscope}%
\begin{pgfscope}%
\pgfsetbuttcap%
\pgfsetroundjoin%
\definecolor{currentfill}{rgb}{1.000000,1.000000,1.000000}%
\pgfsetfillcolor{currentfill}%
\pgfsetlinewidth{0.803000pt}%
\definecolor{currentstroke}{rgb}{0.000000,0.000000,0.000000}%
\pgfsetstrokecolor{currentstroke}%
\pgfsetdash{}{0pt}%
\pgfsys@defobject{currentmarker}{\pgfqpoint{0.000000in}{0.000000in}}{\pgfqpoint{0.000000in}{0.048611in}}{%
\pgfpathmoveto{\pgfqpoint{0.000000in}{0.000000in}}%
\pgfpathlineto{\pgfqpoint{0.000000in}{0.048611in}}%
\pgfusepath{stroke,fill}%
}%
\begin{pgfscope}%
\pgfsys@transformshift{4.374256in}{0.456901in}%
\pgfsys@useobject{currentmarker}{}%
\end{pgfscope}%
\end{pgfscope}%
\begin{pgfscope}%
\definecolor{textcolor}{rgb}{0.000000,0.000000,0.000000}%
\pgfsetstrokecolor{textcolor}%
\pgfsetfillcolor{textcolor}%
\pgftext[x=4.374256in,y=0.408290in,,top]{\color{textcolor}{\rmfamily\fontsize{12.000000}{14.400000}\selectfont\catcode`\^=\active\def^{\ifmmode\sp\else\^{}\fi}\catcode`\%=\active\def%{\%}$\mathdefault{80}$}}%
\end{pgfscope}%
\begin{pgfscope}%
\pgfsetbuttcap%
\pgfsetroundjoin%
\definecolor{currentfill}{rgb}{1.000000,1.000000,1.000000}%
\pgfsetfillcolor{currentfill}%
\pgfsetlinewidth{0.803000pt}%
\definecolor{currentstroke}{rgb}{0.000000,0.000000,0.000000}%
\pgfsetstrokecolor{currentstroke}%
\pgfsetdash{}{0pt}%
\pgfsys@defobject{currentmarker}{\pgfqpoint{0.000000in}{0.000000in}}{\pgfqpoint{0.000000in}{0.048611in}}{%
\pgfpathmoveto{\pgfqpoint{0.000000in}{0.000000in}}%
\pgfpathlineto{\pgfqpoint{0.000000in}{0.048611in}}%
\pgfusepath{stroke,fill}%
}%
\begin{pgfscope}%
\pgfsys@transformshift{5.271560in}{0.456901in}%
\pgfsys@useobject{currentmarker}{}%
\end{pgfscope}%
\end{pgfscope}%
\begin{pgfscope}%
\definecolor{textcolor}{rgb}{0.000000,0.000000,0.000000}%
\pgfsetstrokecolor{textcolor}%
\pgfsetfillcolor{textcolor}%
\pgftext[x=5.271560in,y=0.408290in,,top]{\color{textcolor}{\rmfamily\fontsize{12.000000}{14.400000}\selectfont\catcode`\^=\active\def^{\ifmmode\sp\else\^{}\fi}\catcode`\%=\active\def%{\%}$\mathdefault{100}$}}%
\end{pgfscope}%
\begin{pgfscope}%
\pgfpathrectangle{\pgfqpoint{0.573726in}{0.456901in}}{\pgfqpoint{5.578717in}{2.816433in}}%
\pgfusepath{clip}%
\pgfsetbuttcap%
\pgfsetroundjoin%
\pgfsetlinewidth{0.501875pt}%
\definecolor{currentstroke}{rgb}{0.501961,0.501961,0.501961}%
\pgfsetstrokecolor{currentstroke}%
\pgfsetstrokeopacity{0.400000}%
\pgfsetdash{{1.850000pt}{0.800000pt}}{0.000000pt}%
\pgfpathmoveto{\pgfqpoint{1.009367in}{0.456901in}}%
\pgfpathlineto{\pgfqpoint{1.009367in}{3.273333in}}%
\pgfusepath{stroke}%
\end{pgfscope}%
\begin{pgfscope}%
\pgfsetbuttcap%
\pgfsetroundjoin%
\definecolor{currentfill}{rgb}{1.000000,1.000000,1.000000}%
\pgfsetfillcolor{currentfill}%
\pgfsetlinewidth{0.602250pt}%
\definecolor{currentstroke}{rgb}{0.000000,0.000000,0.000000}%
\pgfsetstrokecolor{currentstroke}%
\pgfsetdash{}{0pt}%
\pgfsys@defobject{currentmarker}{\pgfqpoint{0.000000in}{0.000000in}}{\pgfqpoint{0.000000in}{0.027778in}}{%
\pgfpathmoveto{\pgfqpoint{0.000000in}{0.000000in}}%
\pgfpathlineto{\pgfqpoint{0.000000in}{0.027778in}}%
\pgfusepath{stroke,fill}%
}%
\begin{pgfscope}%
\pgfsys@transformshift{1.009367in}{0.456901in}%
\pgfsys@useobject{currentmarker}{}%
\end{pgfscope}%
\end{pgfscope}%
\begin{pgfscope}%
\pgfpathrectangle{\pgfqpoint{0.573726in}{0.456901in}}{\pgfqpoint{5.578717in}{2.816433in}}%
\pgfusepath{clip}%
\pgfsetbuttcap%
\pgfsetroundjoin%
\pgfsetlinewidth{0.501875pt}%
\definecolor{currentstroke}{rgb}{0.501961,0.501961,0.501961}%
\pgfsetstrokecolor{currentstroke}%
\pgfsetstrokeopacity{0.400000}%
\pgfsetdash{{1.850000pt}{0.800000pt}}{0.000000pt}%
\pgfpathmoveto{\pgfqpoint{1.233693in}{0.456901in}}%
\pgfpathlineto{\pgfqpoint{1.233693in}{3.273333in}}%
\pgfusepath{stroke}%
\end{pgfscope}%
\begin{pgfscope}%
\pgfsetbuttcap%
\pgfsetroundjoin%
\definecolor{currentfill}{rgb}{1.000000,1.000000,1.000000}%
\pgfsetfillcolor{currentfill}%
\pgfsetlinewidth{0.602250pt}%
\definecolor{currentstroke}{rgb}{0.000000,0.000000,0.000000}%
\pgfsetstrokecolor{currentstroke}%
\pgfsetdash{}{0pt}%
\pgfsys@defobject{currentmarker}{\pgfqpoint{0.000000in}{0.000000in}}{\pgfqpoint{0.000000in}{0.027778in}}{%
\pgfpathmoveto{\pgfqpoint{0.000000in}{0.000000in}}%
\pgfpathlineto{\pgfqpoint{0.000000in}{0.027778in}}%
\pgfusepath{stroke,fill}%
}%
\begin{pgfscope}%
\pgfsys@transformshift{1.233693in}{0.456901in}%
\pgfsys@useobject{currentmarker}{}%
\end{pgfscope}%
\end{pgfscope}%
\begin{pgfscope}%
\pgfpathrectangle{\pgfqpoint{0.573726in}{0.456901in}}{\pgfqpoint{5.578717in}{2.816433in}}%
\pgfusepath{clip}%
\pgfsetbuttcap%
\pgfsetroundjoin%
\pgfsetlinewidth{0.501875pt}%
\definecolor{currentstroke}{rgb}{0.501961,0.501961,0.501961}%
\pgfsetstrokecolor{currentstroke}%
\pgfsetstrokeopacity{0.400000}%
\pgfsetdash{{1.850000pt}{0.800000pt}}{0.000000pt}%
\pgfpathmoveto{\pgfqpoint{1.458019in}{0.456901in}}%
\pgfpathlineto{\pgfqpoint{1.458019in}{3.273333in}}%
\pgfusepath{stroke}%
\end{pgfscope}%
\begin{pgfscope}%
\pgfsetbuttcap%
\pgfsetroundjoin%
\definecolor{currentfill}{rgb}{1.000000,1.000000,1.000000}%
\pgfsetfillcolor{currentfill}%
\pgfsetlinewidth{0.602250pt}%
\definecolor{currentstroke}{rgb}{0.000000,0.000000,0.000000}%
\pgfsetstrokecolor{currentstroke}%
\pgfsetdash{}{0pt}%
\pgfsys@defobject{currentmarker}{\pgfqpoint{0.000000in}{0.000000in}}{\pgfqpoint{0.000000in}{0.027778in}}{%
\pgfpathmoveto{\pgfqpoint{0.000000in}{0.000000in}}%
\pgfpathlineto{\pgfqpoint{0.000000in}{0.027778in}}%
\pgfusepath{stroke,fill}%
}%
\begin{pgfscope}%
\pgfsys@transformshift{1.458019in}{0.456901in}%
\pgfsys@useobject{currentmarker}{}%
\end{pgfscope}%
\end{pgfscope}%
\begin{pgfscope}%
\pgfpathrectangle{\pgfqpoint{0.573726in}{0.456901in}}{\pgfqpoint{5.578717in}{2.816433in}}%
\pgfusepath{clip}%
\pgfsetbuttcap%
\pgfsetroundjoin%
\pgfsetlinewidth{0.501875pt}%
\definecolor{currentstroke}{rgb}{0.501961,0.501961,0.501961}%
\pgfsetstrokecolor{currentstroke}%
\pgfsetstrokeopacity{0.400000}%
\pgfsetdash{{1.850000pt}{0.800000pt}}{0.000000pt}%
\pgfpathmoveto{\pgfqpoint{1.906671in}{0.456901in}}%
\pgfpathlineto{\pgfqpoint{1.906671in}{3.273333in}}%
\pgfusepath{stroke}%
\end{pgfscope}%
\begin{pgfscope}%
\pgfsetbuttcap%
\pgfsetroundjoin%
\definecolor{currentfill}{rgb}{1.000000,1.000000,1.000000}%
\pgfsetfillcolor{currentfill}%
\pgfsetlinewidth{0.602250pt}%
\definecolor{currentstroke}{rgb}{0.000000,0.000000,0.000000}%
\pgfsetstrokecolor{currentstroke}%
\pgfsetdash{}{0pt}%
\pgfsys@defobject{currentmarker}{\pgfqpoint{0.000000in}{0.000000in}}{\pgfqpoint{0.000000in}{0.027778in}}{%
\pgfpathmoveto{\pgfqpoint{0.000000in}{0.000000in}}%
\pgfpathlineto{\pgfqpoint{0.000000in}{0.027778in}}%
\pgfusepath{stroke,fill}%
}%
\begin{pgfscope}%
\pgfsys@transformshift{1.906671in}{0.456901in}%
\pgfsys@useobject{currentmarker}{}%
\end{pgfscope}%
\end{pgfscope}%
\begin{pgfscope}%
\pgfpathrectangle{\pgfqpoint{0.573726in}{0.456901in}}{\pgfqpoint{5.578717in}{2.816433in}}%
\pgfusepath{clip}%
\pgfsetbuttcap%
\pgfsetroundjoin%
\pgfsetlinewidth{0.501875pt}%
\definecolor{currentstroke}{rgb}{0.501961,0.501961,0.501961}%
\pgfsetstrokecolor{currentstroke}%
\pgfsetstrokeopacity{0.400000}%
\pgfsetdash{{1.850000pt}{0.800000pt}}{0.000000pt}%
\pgfpathmoveto{\pgfqpoint{2.130997in}{0.456901in}}%
\pgfpathlineto{\pgfqpoint{2.130997in}{3.273333in}}%
\pgfusepath{stroke}%
\end{pgfscope}%
\begin{pgfscope}%
\pgfsetbuttcap%
\pgfsetroundjoin%
\definecolor{currentfill}{rgb}{1.000000,1.000000,1.000000}%
\pgfsetfillcolor{currentfill}%
\pgfsetlinewidth{0.602250pt}%
\definecolor{currentstroke}{rgb}{0.000000,0.000000,0.000000}%
\pgfsetstrokecolor{currentstroke}%
\pgfsetdash{}{0pt}%
\pgfsys@defobject{currentmarker}{\pgfqpoint{0.000000in}{0.000000in}}{\pgfqpoint{0.000000in}{0.027778in}}{%
\pgfpathmoveto{\pgfqpoint{0.000000in}{0.000000in}}%
\pgfpathlineto{\pgfqpoint{0.000000in}{0.027778in}}%
\pgfusepath{stroke,fill}%
}%
\begin{pgfscope}%
\pgfsys@transformshift{2.130997in}{0.456901in}%
\pgfsys@useobject{currentmarker}{}%
\end{pgfscope}%
\end{pgfscope}%
\begin{pgfscope}%
\pgfpathrectangle{\pgfqpoint{0.573726in}{0.456901in}}{\pgfqpoint{5.578717in}{2.816433in}}%
\pgfusepath{clip}%
\pgfsetbuttcap%
\pgfsetroundjoin%
\pgfsetlinewidth{0.501875pt}%
\definecolor{currentstroke}{rgb}{0.501961,0.501961,0.501961}%
\pgfsetstrokecolor{currentstroke}%
\pgfsetstrokeopacity{0.400000}%
\pgfsetdash{{1.850000pt}{0.800000pt}}{0.000000pt}%
\pgfpathmoveto{\pgfqpoint{2.355323in}{0.456901in}}%
\pgfpathlineto{\pgfqpoint{2.355323in}{3.273333in}}%
\pgfusepath{stroke}%
\end{pgfscope}%
\begin{pgfscope}%
\pgfsetbuttcap%
\pgfsetroundjoin%
\definecolor{currentfill}{rgb}{1.000000,1.000000,1.000000}%
\pgfsetfillcolor{currentfill}%
\pgfsetlinewidth{0.602250pt}%
\definecolor{currentstroke}{rgb}{0.000000,0.000000,0.000000}%
\pgfsetstrokecolor{currentstroke}%
\pgfsetdash{}{0pt}%
\pgfsys@defobject{currentmarker}{\pgfqpoint{0.000000in}{0.000000in}}{\pgfqpoint{0.000000in}{0.027778in}}{%
\pgfpathmoveto{\pgfqpoint{0.000000in}{0.000000in}}%
\pgfpathlineto{\pgfqpoint{0.000000in}{0.027778in}}%
\pgfusepath{stroke,fill}%
}%
\begin{pgfscope}%
\pgfsys@transformshift{2.355323in}{0.456901in}%
\pgfsys@useobject{currentmarker}{}%
\end{pgfscope}%
\end{pgfscope}%
\begin{pgfscope}%
\pgfpathrectangle{\pgfqpoint{0.573726in}{0.456901in}}{\pgfqpoint{5.578717in}{2.816433in}}%
\pgfusepath{clip}%
\pgfsetbuttcap%
\pgfsetroundjoin%
\pgfsetlinewidth{0.501875pt}%
\definecolor{currentstroke}{rgb}{0.501961,0.501961,0.501961}%
\pgfsetstrokecolor{currentstroke}%
\pgfsetstrokeopacity{0.400000}%
\pgfsetdash{{1.850000pt}{0.800000pt}}{0.000000pt}%
\pgfpathmoveto{\pgfqpoint{2.803975in}{0.456901in}}%
\pgfpathlineto{\pgfqpoint{2.803975in}{3.273333in}}%
\pgfusepath{stroke}%
\end{pgfscope}%
\begin{pgfscope}%
\pgfsetbuttcap%
\pgfsetroundjoin%
\definecolor{currentfill}{rgb}{1.000000,1.000000,1.000000}%
\pgfsetfillcolor{currentfill}%
\pgfsetlinewidth{0.602250pt}%
\definecolor{currentstroke}{rgb}{0.000000,0.000000,0.000000}%
\pgfsetstrokecolor{currentstroke}%
\pgfsetdash{}{0pt}%
\pgfsys@defobject{currentmarker}{\pgfqpoint{0.000000in}{0.000000in}}{\pgfqpoint{0.000000in}{0.027778in}}{%
\pgfpathmoveto{\pgfqpoint{0.000000in}{0.000000in}}%
\pgfpathlineto{\pgfqpoint{0.000000in}{0.027778in}}%
\pgfusepath{stroke,fill}%
}%
\begin{pgfscope}%
\pgfsys@transformshift{2.803975in}{0.456901in}%
\pgfsys@useobject{currentmarker}{}%
\end{pgfscope}%
\end{pgfscope}%
\begin{pgfscope}%
\pgfpathrectangle{\pgfqpoint{0.573726in}{0.456901in}}{\pgfqpoint{5.578717in}{2.816433in}}%
\pgfusepath{clip}%
\pgfsetbuttcap%
\pgfsetroundjoin%
\pgfsetlinewidth{0.501875pt}%
\definecolor{currentstroke}{rgb}{0.501961,0.501961,0.501961}%
\pgfsetstrokecolor{currentstroke}%
\pgfsetstrokeopacity{0.400000}%
\pgfsetdash{{1.850000pt}{0.800000pt}}{0.000000pt}%
\pgfpathmoveto{\pgfqpoint{3.028301in}{0.456901in}}%
\pgfpathlineto{\pgfqpoint{3.028301in}{3.273333in}}%
\pgfusepath{stroke}%
\end{pgfscope}%
\begin{pgfscope}%
\pgfsetbuttcap%
\pgfsetroundjoin%
\definecolor{currentfill}{rgb}{1.000000,1.000000,1.000000}%
\pgfsetfillcolor{currentfill}%
\pgfsetlinewidth{0.602250pt}%
\definecolor{currentstroke}{rgb}{0.000000,0.000000,0.000000}%
\pgfsetstrokecolor{currentstroke}%
\pgfsetdash{}{0pt}%
\pgfsys@defobject{currentmarker}{\pgfqpoint{0.000000in}{0.000000in}}{\pgfqpoint{0.000000in}{0.027778in}}{%
\pgfpathmoveto{\pgfqpoint{0.000000in}{0.000000in}}%
\pgfpathlineto{\pgfqpoint{0.000000in}{0.027778in}}%
\pgfusepath{stroke,fill}%
}%
\begin{pgfscope}%
\pgfsys@transformshift{3.028301in}{0.456901in}%
\pgfsys@useobject{currentmarker}{}%
\end{pgfscope}%
\end{pgfscope}%
\begin{pgfscope}%
\pgfpathrectangle{\pgfqpoint{0.573726in}{0.456901in}}{\pgfqpoint{5.578717in}{2.816433in}}%
\pgfusepath{clip}%
\pgfsetbuttcap%
\pgfsetroundjoin%
\pgfsetlinewidth{0.501875pt}%
\definecolor{currentstroke}{rgb}{0.501961,0.501961,0.501961}%
\pgfsetstrokecolor{currentstroke}%
\pgfsetstrokeopacity{0.400000}%
\pgfsetdash{{1.850000pt}{0.800000pt}}{0.000000pt}%
\pgfpathmoveto{\pgfqpoint{3.252627in}{0.456901in}}%
\pgfpathlineto{\pgfqpoint{3.252627in}{3.273333in}}%
\pgfusepath{stroke}%
\end{pgfscope}%
\begin{pgfscope}%
\pgfsetbuttcap%
\pgfsetroundjoin%
\definecolor{currentfill}{rgb}{1.000000,1.000000,1.000000}%
\pgfsetfillcolor{currentfill}%
\pgfsetlinewidth{0.602250pt}%
\definecolor{currentstroke}{rgb}{0.000000,0.000000,0.000000}%
\pgfsetstrokecolor{currentstroke}%
\pgfsetdash{}{0pt}%
\pgfsys@defobject{currentmarker}{\pgfqpoint{0.000000in}{0.000000in}}{\pgfqpoint{0.000000in}{0.027778in}}{%
\pgfpathmoveto{\pgfqpoint{0.000000in}{0.000000in}}%
\pgfpathlineto{\pgfqpoint{0.000000in}{0.027778in}}%
\pgfusepath{stroke,fill}%
}%
\begin{pgfscope}%
\pgfsys@transformshift{3.252627in}{0.456901in}%
\pgfsys@useobject{currentmarker}{}%
\end{pgfscope}%
\end{pgfscope}%
\begin{pgfscope}%
\pgfpathrectangle{\pgfqpoint{0.573726in}{0.456901in}}{\pgfqpoint{5.578717in}{2.816433in}}%
\pgfusepath{clip}%
\pgfsetbuttcap%
\pgfsetroundjoin%
\pgfsetlinewidth{0.501875pt}%
\definecolor{currentstroke}{rgb}{0.501961,0.501961,0.501961}%
\pgfsetstrokecolor{currentstroke}%
\pgfsetstrokeopacity{0.400000}%
\pgfsetdash{{1.850000pt}{0.800000pt}}{0.000000pt}%
\pgfpathmoveto{\pgfqpoint{3.701278in}{0.456901in}}%
\pgfpathlineto{\pgfqpoint{3.701278in}{3.273333in}}%
\pgfusepath{stroke}%
\end{pgfscope}%
\begin{pgfscope}%
\pgfsetbuttcap%
\pgfsetroundjoin%
\definecolor{currentfill}{rgb}{1.000000,1.000000,1.000000}%
\pgfsetfillcolor{currentfill}%
\pgfsetlinewidth{0.602250pt}%
\definecolor{currentstroke}{rgb}{0.000000,0.000000,0.000000}%
\pgfsetstrokecolor{currentstroke}%
\pgfsetdash{}{0pt}%
\pgfsys@defobject{currentmarker}{\pgfqpoint{0.000000in}{0.000000in}}{\pgfqpoint{0.000000in}{0.027778in}}{%
\pgfpathmoveto{\pgfqpoint{0.000000in}{0.000000in}}%
\pgfpathlineto{\pgfqpoint{0.000000in}{0.027778in}}%
\pgfusepath{stroke,fill}%
}%
\begin{pgfscope}%
\pgfsys@transformshift{3.701278in}{0.456901in}%
\pgfsys@useobject{currentmarker}{}%
\end{pgfscope}%
\end{pgfscope}%
\begin{pgfscope}%
\pgfpathrectangle{\pgfqpoint{0.573726in}{0.456901in}}{\pgfqpoint{5.578717in}{2.816433in}}%
\pgfusepath{clip}%
\pgfsetbuttcap%
\pgfsetroundjoin%
\pgfsetlinewidth{0.501875pt}%
\definecolor{currentstroke}{rgb}{0.501961,0.501961,0.501961}%
\pgfsetstrokecolor{currentstroke}%
\pgfsetstrokeopacity{0.400000}%
\pgfsetdash{{1.850000pt}{0.800000pt}}{0.000000pt}%
\pgfpathmoveto{\pgfqpoint{3.925604in}{0.456901in}}%
\pgfpathlineto{\pgfqpoint{3.925604in}{3.273333in}}%
\pgfusepath{stroke}%
\end{pgfscope}%
\begin{pgfscope}%
\pgfsetbuttcap%
\pgfsetroundjoin%
\definecolor{currentfill}{rgb}{1.000000,1.000000,1.000000}%
\pgfsetfillcolor{currentfill}%
\pgfsetlinewidth{0.602250pt}%
\definecolor{currentstroke}{rgb}{0.000000,0.000000,0.000000}%
\pgfsetstrokecolor{currentstroke}%
\pgfsetdash{}{0pt}%
\pgfsys@defobject{currentmarker}{\pgfqpoint{0.000000in}{0.000000in}}{\pgfqpoint{0.000000in}{0.027778in}}{%
\pgfpathmoveto{\pgfqpoint{0.000000in}{0.000000in}}%
\pgfpathlineto{\pgfqpoint{0.000000in}{0.027778in}}%
\pgfusepath{stroke,fill}%
}%
\begin{pgfscope}%
\pgfsys@transformshift{3.925604in}{0.456901in}%
\pgfsys@useobject{currentmarker}{}%
\end{pgfscope}%
\end{pgfscope}%
\begin{pgfscope}%
\pgfpathrectangle{\pgfqpoint{0.573726in}{0.456901in}}{\pgfqpoint{5.578717in}{2.816433in}}%
\pgfusepath{clip}%
\pgfsetbuttcap%
\pgfsetroundjoin%
\pgfsetlinewidth{0.501875pt}%
\definecolor{currentstroke}{rgb}{0.501961,0.501961,0.501961}%
\pgfsetstrokecolor{currentstroke}%
\pgfsetstrokeopacity{0.400000}%
\pgfsetdash{{1.850000pt}{0.800000pt}}{0.000000pt}%
\pgfpathmoveto{\pgfqpoint{4.149930in}{0.456901in}}%
\pgfpathlineto{\pgfqpoint{4.149930in}{3.273333in}}%
\pgfusepath{stroke}%
\end{pgfscope}%
\begin{pgfscope}%
\pgfsetbuttcap%
\pgfsetroundjoin%
\definecolor{currentfill}{rgb}{1.000000,1.000000,1.000000}%
\pgfsetfillcolor{currentfill}%
\pgfsetlinewidth{0.602250pt}%
\definecolor{currentstroke}{rgb}{0.000000,0.000000,0.000000}%
\pgfsetstrokecolor{currentstroke}%
\pgfsetdash{}{0pt}%
\pgfsys@defobject{currentmarker}{\pgfqpoint{0.000000in}{0.000000in}}{\pgfqpoint{0.000000in}{0.027778in}}{%
\pgfpathmoveto{\pgfqpoint{0.000000in}{0.000000in}}%
\pgfpathlineto{\pgfqpoint{0.000000in}{0.027778in}}%
\pgfusepath{stroke,fill}%
}%
\begin{pgfscope}%
\pgfsys@transformshift{4.149930in}{0.456901in}%
\pgfsys@useobject{currentmarker}{}%
\end{pgfscope}%
\end{pgfscope}%
\begin{pgfscope}%
\pgfpathrectangle{\pgfqpoint{0.573726in}{0.456901in}}{\pgfqpoint{5.578717in}{2.816433in}}%
\pgfusepath{clip}%
\pgfsetbuttcap%
\pgfsetroundjoin%
\pgfsetlinewidth{0.501875pt}%
\definecolor{currentstroke}{rgb}{0.501961,0.501961,0.501961}%
\pgfsetstrokecolor{currentstroke}%
\pgfsetstrokeopacity{0.400000}%
\pgfsetdash{{1.850000pt}{0.800000pt}}{0.000000pt}%
\pgfpathmoveto{\pgfqpoint{4.598582in}{0.456901in}}%
\pgfpathlineto{\pgfqpoint{4.598582in}{3.273333in}}%
\pgfusepath{stroke}%
\end{pgfscope}%
\begin{pgfscope}%
\pgfsetbuttcap%
\pgfsetroundjoin%
\definecolor{currentfill}{rgb}{1.000000,1.000000,1.000000}%
\pgfsetfillcolor{currentfill}%
\pgfsetlinewidth{0.602250pt}%
\definecolor{currentstroke}{rgb}{0.000000,0.000000,0.000000}%
\pgfsetstrokecolor{currentstroke}%
\pgfsetdash{}{0pt}%
\pgfsys@defobject{currentmarker}{\pgfqpoint{0.000000in}{0.000000in}}{\pgfqpoint{0.000000in}{0.027778in}}{%
\pgfpathmoveto{\pgfqpoint{0.000000in}{0.000000in}}%
\pgfpathlineto{\pgfqpoint{0.000000in}{0.027778in}}%
\pgfusepath{stroke,fill}%
}%
\begin{pgfscope}%
\pgfsys@transformshift{4.598582in}{0.456901in}%
\pgfsys@useobject{currentmarker}{}%
\end{pgfscope}%
\end{pgfscope}%
\begin{pgfscope}%
\pgfpathrectangle{\pgfqpoint{0.573726in}{0.456901in}}{\pgfqpoint{5.578717in}{2.816433in}}%
\pgfusepath{clip}%
\pgfsetbuttcap%
\pgfsetroundjoin%
\pgfsetlinewidth{0.501875pt}%
\definecolor{currentstroke}{rgb}{0.501961,0.501961,0.501961}%
\pgfsetstrokecolor{currentstroke}%
\pgfsetstrokeopacity{0.400000}%
\pgfsetdash{{1.850000pt}{0.800000pt}}{0.000000pt}%
\pgfpathmoveto{\pgfqpoint{4.822908in}{0.456901in}}%
\pgfpathlineto{\pgfqpoint{4.822908in}{3.273333in}}%
\pgfusepath{stroke}%
\end{pgfscope}%
\begin{pgfscope}%
\pgfsetbuttcap%
\pgfsetroundjoin%
\definecolor{currentfill}{rgb}{1.000000,1.000000,1.000000}%
\pgfsetfillcolor{currentfill}%
\pgfsetlinewidth{0.602250pt}%
\definecolor{currentstroke}{rgb}{0.000000,0.000000,0.000000}%
\pgfsetstrokecolor{currentstroke}%
\pgfsetdash{}{0pt}%
\pgfsys@defobject{currentmarker}{\pgfqpoint{0.000000in}{0.000000in}}{\pgfqpoint{0.000000in}{0.027778in}}{%
\pgfpathmoveto{\pgfqpoint{0.000000in}{0.000000in}}%
\pgfpathlineto{\pgfqpoint{0.000000in}{0.027778in}}%
\pgfusepath{stroke,fill}%
}%
\begin{pgfscope}%
\pgfsys@transformshift{4.822908in}{0.456901in}%
\pgfsys@useobject{currentmarker}{}%
\end{pgfscope}%
\end{pgfscope}%
\begin{pgfscope}%
\pgfpathrectangle{\pgfqpoint{0.573726in}{0.456901in}}{\pgfqpoint{5.578717in}{2.816433in}}%
\pgfusepath{clip}%
\pgfsetbuttcap%
\pgfsetroundjoin%
\pgfsetlinewidth{0.501875pt}%
\definecolor{currentstroke}{rgb}{0.501961,0.501961,0.501961}%
\pgfsetstrokecolor{currentstroke}%
\pgfsetstrokeopacity{0.400000}%
\pgfsetdash{{1.850000pt}{0.800000pt}}{0.000000pt}%
\pgfpathmoveto{\pgfqpoint{5.047234in}{0.456901in}}%
\pgfpathlineto{\pgfqpoint{5.047234in}{3.273333in}}%
\pgfusepath{stroke}%
\end{pgfscope}%
\begin{pgfscope}%
\pgfsetbuttcap%
\pgfsetroundjoin%
\definecolor{currentfill}{rgb}{1.000000,1.000000,1.000000}%
\pgfsetfillcolor{currentfill}%
\pgfsetlinewidth{0.602250pt}%
\definecolor{currentstroke}{rgb}{0.000000,0.000000,0.000000}%
\pgfsetstrokecolor{currentstroke}%
\pgfsetdash{}{0pt}%
\pgfsys@defobject{currentmarker}{\pgfqpoint{0.000000in}{0.000000in}}{\pgfqpoint{0.000000in}{0.027778in}}{%
\pgfpathmoveto{\pgfqpoint{0.000000in}{0.000000in}}%
\pgfpathlineto{\pgfqpoint{0.000000in}{0.027778in}}%
\pgfusepath{stroke,fill}%
}%
\begin{pgfscope}%
\pgfsys@transformshift{5.047234in}{0.456901in}%
\pgfsys@useobject{currentmarker}{}%
\end{pgfscope}%
\end{pgfscope}%
\begin{pgfscope}%
\pgfpathrectangle{\pgfqpoint{0.573726in}{0.456901in}}{\pgfqpoint{5.578717in}{2.816433in}}%
\pgfusepath{clip}%
\pgfsetbuttcap%
\pgfsetroundjoin%
\pgfsetlinewidth{0.501875pt}%
\definecolor{currentstroke}{rgb}{0.501961,0.501961,0.501961}%
\pgfsetstrokecolor{currentstroke}%
\pgfsetstrokeopacity{0.400000}%
\pgfsetdash{{1.850000pt}{0.800000pt}}{0.000000pt}%
\pgfpathmoveto{\pgfqpoint{5.495886in}{0.456901in}}%
\pgfpathlineto{\pgfqpoint{5.495886in}{3.273333in}}%
\pgfusepath{stroke}%
\end{pgfscope}%
\begin{pgfscope}%
\pgfsetbuttcap%
\pgfsetroundjoin%
\definecolor{currentfill}{rgb}{1.000000,1.000000,1.000000}%
\pgfsetfillcolor{currentfill}%
\pgfsetlinewidth{0.602250pt}%
\definecolor{currentstroke}{rgb}{0.000000,0.000000,0.000000}%
\pgfsetstrokecolor{currentstroke}%
\pgfsetdash{}{0pt}%
\pgfsys@defobject{currentmarker}{\pgfqpoint{0.000000in}{0.000000in}}{\pgfqpoint{0.000000in}{0.027778in}}{%
\pgfpathmoveto{\pgfqpoint{0.000000in}{0.000000in}}%
\pgfpathlineto{\pgfqpoint{0.000000in}{0.027778in}}%
\pgfusepath{stroke,fill}%
}%
\begin{pgfscope}%
\pgfsys@transformshift{5.495886in}{0.456901in}%
\pgfsys@useobject{currentmarker}{}%
\end{pgfscope}%
\end{pgfscope}%
\begin{pgfscope}%
\pgfpathrectangle{\pgfqpoint{0.573726in}{0.456901in}}{\pgfqpoint{5.578717in}{2.816433in}}%
\pgfusepath{clip}%
\pgfsetbuttcap%
\pgfsetroundjoin%
\pgfsetlinewidth{0.501875pt}%
\definecolor{currentstroke}{rgb}{0.501961,0.501961,0.501961}%
\pgfsetstrokecolor{currentstroke}%
\pgfsetstrokeopacity{0.400000}%
\pgfsetdash{{1.850000pt}{0.800000pt}}{0.000000pt}%
\pgfpathmoveto{\pgfqpoint{5.720212in}{0.456901in}}%
\pgfpathlineto{\pgfqpoint{5.720212in}{3.273333in}}%
\pgfusepath{stroke}%
\end{pgfscope}%
\begin{pgfscope}%
\pgfsetbuttcap%
\pgfsetroundjoin%
\definecolor{currentfill}{rgb}{1.000000,1.000000,1.000000}%
\pgfsetfillcolor{currentfill}%
\pgfsetlinewidth{0.602250pt}%
\definecolor{currentstroke}{rgb}{0.000000,0.000000,0.000000}%
\pgfsetstrokecolor{currentstroke}%
\pgfsetdash{}{0pt}%
\pgfsys@defobject{currentmarker}{\pgfqpoint{0.000000in}{0.000000in}}{\pgfqpoint{0.000000in}{0.027778in}}{%
\pgfpathmoveto{\pgfqpoint{0.000000in}{0.000000in}}%
\pgfpathlineto{\pgfqpoint{0.000000in}{0.027778in}}%
\pgfusepath{stroke,fill}%
}%
\begin{pgfscope}%
\pgfsys@transformshift{5.720212in}{0.456901in}%
\pgfsys@useobject{currentmarker}{}%
\end{pgfscope}%
\end{pgfscope}%
\begin{pgfscope}%
\pgfpathrectangle{\pgfqpoint{0.573726in}{0.456901in}}{\pgfqpoint{5.578717in}{2.816433in}}%
\pgfusepath{clip}%
\pgfsetbuttcap%
\pgfsetroundjoin%
\pgfsetlinewidth{0.501875pt}%
\definecolor{currentstroke}{rgb}{0.501961,0.501961,0.501961}%
\pgfsetstrokecolor{currentstroke}%
\pgfsetstrokeopacity{0.400000}%
\pgfsetdash{{1.850000pt}{0.800000pt}}{0.000000pt}%
\pgfpathmoveto{\pgfqpoint{5.944538in}{0.456901in}}%
\pgfpathlineto{\pgfqpoint{5.944538in}{3.273333in}}%
\pgfusepath{stroke}%
\end{pgfscope}%
\begin{pgfscope}%
\pgfsetbuttcap%
\pgfsetroundjoin%
\definecolor{currentfill}{rgb}{1.000000,1.000000,1.000000}%
\pgfsetfillcolor{currentfill}%
\pgfsetlinewidth{0.602250pt}%
\definecolor{currentstroke}{rgb}{0.000000,0.000000,0.000000}%
\pgfsetstrokecolor{currentstroke}%
\pgfsetdash{}{0pt}%
\pgfsys@defobject{currentmarker}{\pgfqpoint{0.000000in}{0.000000in}}{\pgfqpoint{0.000000in}{0.027778in}}{%
\pgfpathmoveto{\pgfqpoint{0.000000in}{0.000000in}}%
\pgfpathlineto{\pgfqpoint{0.000000in}{0.027778in}}%
\pgfusepath{stroke,fill}%
}%
\begin{pgfscope}%
\pgfsys@transformshift{5.944538in}{0.456901in}%
\pgfsys@useobject{currentmarker}{}%
\end{pgfscope}%
\end{pgfscope}%
\begin{pgfscope}%
\definecolor{textcolor}{rgb}{0.000000,0.000000,0.000000}%
\pgfsetstrokecolor{textcolor}%
\pgfsetfillcolor{textcolor}%
\pgftext[x=3.363085in,y=0.201367in,,top]{\color{textcolor}{\rmfamily\fontsize{12.000000}{14.400000}\selectfont\catcode`\^=\active\def^{\ifmmode\sp\else\^{}\fi}\catcode`\%=\active\def%{\%}$t$}}%
\end{pgfscope}%
\begin{pgfscope}%
\pgfsetbuttcap%
\pgfsetroundjoin%
\definecolor{currentfill}{rgb}{1.000000,1.000000,1.000000}%
\pgfsetfillcolor{currentfill}%
\pgfsetlinewidth{0.803000pt}%
\definecolor{currentstroke}{rgb}{0.000000,0.000000,0.000000}%
\pgfsetstrokecolor{currentstroke}%
\pgfsetdash{}{0pt}%
\pgfsys@defobject{currentmarker}{\pgfqpoint{0.000000in}{0.000000in}}{\pgfqpoint{0.048611in}{0.000000in}}{%
\pgfpathmoveto{\pgfqpoint{0.000000in}{0.000000in}}%
\pgfpathlineto{\pgfqpoint{0.048611in}{0.000000in}}%
\pgfusepath{stroke,fill}%
}%
\begin{pgfscope}%
\pgfsys@transformshift{0.573726in}{0.787834in}%
\pgfsys@useobject{currentmarker}{}%
\end{pgfscope}%
\end{pgfscope}%
\begin{pgfscope}%
\definecolor{textcolor}{rgb}{0.000000,0.000000,0.000000}%
\pgfsetstrokecolor{textcolor}%
\pgfsetfillcolor{textcolor}%
\pgftext[x=0.272222in, y=0.729973in, left, base]{\color{textcolor}{\rmfamily\fontsize{12.000000}{14.400000}\selectfont\catcode`\^=\active\def^{\ifmmode\sp\else\^{}\fi}\catcode`\%=\active\def%{\%}$\mathdefault{\ensuremath{-}2}$}}%
\end{pgfscope}%
\begin{pgfscope}%
\pgfsetbuttcap%
\pgfsetroundjoin%
\definecolor{currentfill}{rgb}{1.000000,1.000000,1.000000}%
\pgfsetfillcolor{currentfill}%
\pgfsetlinewidth{0.803000pt}%
\definecolor{currentstroke}{rgb}{0.000000,0.000000,0.000000}%
\pgfsetstrokecolor{currentstroke}%
\pgfsetdash{}{0pt}%
\pgfsys@defobject{currentmarker}{\pgfqpoint{0.000000in}{0.000000in}}{\pgfqpoint{0.048611in}{0.000000in}}{%
\pgfpathmoveto{\pgfqpoint{0.000000in}{0.000000in}}%
\pgfpathlineto{\pgfqpoint{0.048611in}{0.000000in}}%
\pgfusepath{stroke,fill}%
}%
\begin{pgfscope}%
\pgfsys@transformshift{0.573726in}{1.326475in}%
\pgfsys@useobject{currentmarker}{}%
\end{pgfscope}%
\end{pgfscope}%
\begin{pgfscope}%
\definecolor{textcolor}{rgb}{0.000000,0.000000,0.000000}%
\pgfsetstrokecolor{textcolor}%
\pgfsetfillcolor{textcolor}%
\pgftext[x=0.272222in, y=1.268614in, left, base]{\color{textcolor}{\rmfamily\fontsize{12.000000}{14.400000}\selectfont\catcode`\^=\active\def^{\ifmmode\sp\else\^{}\fi}\catcode`\%=\active\def%{\%}$\mathdefault{\ensuremath{-}1}$}}%
\end{pgfscope}%
\begin{pgfscope}%
\pgfsetbuttcap%
\pgfsetroundjoin%
\definecolor{currentfill}{rgb}{1.000000,1.000000,1.000000}%
\pgfsetfillcolor{currentfill}%
\pgfsetlinewidth{0.803000pt}%
\definecolor{currentstroke}{rgb}{0.000000,0.000000,0.000000}%
\pgfsetstrokecolor{currentstroke}%
\pgfsetdash{}{0pt}%
\pgfsys@defobject{currentmarker}{\pgfqpoint{0.000000in}{0.000000in}}{\pgfqpoint{0.048611in}{0.000000in}}{%
\pgfpathmoveto{\pgfqpoint{0.000000in}{0.000000in}}%
\pgfpathlineto{\pgfqpoint{0.048611in}{0.000000in}}%
\pgfusepath{stroke,fill}%
}%
\begin{pgfscope}%
\pgfsys@transformshift{0.573726in}{1.865117in}%
\pgfsys@useobject{currentmarker}{}%
\end{pgfscope}%
\end{pgfscope}%
\begin{pgfscope}%
\definecolor{textcolor}{rgb}{0.000000,0.000000,0.000000}%
\pgfsetstrokecolor{textcolor}%
\pgfsetfillcolor{textcolor}%
\pgftext[x=0.401852in, y=1.807256in, left, base]{\color{textcolor}{\rmfamily\fontsize{12.000000}{14.400000}\selectfont\catcode`\^=\active\def^{\ifmmode\sp\else\^{}\fi}\catcode`\%=\active\def%{\%}$\mathdefault{0}$}}%
\end{pgfscope}%
\begin{pgfscope}%
\pgfsetbuttcap%
\pgfsetroundjoin%
\definecolor{currentfill}{rgb}{1.000000,1.000000,1.000000}%
\pgfsetfillcolor{currentfill}%
\pgfsetlinewidth{0.803000pt}%
\definecolor{currentstroke}{rgb}{0.000000,0.000000,0.000000}%
\pgfsetstrokecolor{currentstroke}%
\pgfsetdash{}{0pt}%
\pgfsys@defobject{currentmarker}{\pgfqpoint{0.000000in}{0.000000in}}{\pgfqpoint{0.048611in}{0.000000in}}{%
\pgfpathmoveto{\pgfqpoint{0.000000in}{0.000000in}}%
\pgfpathlineto{\pgfqpoint{0.048611in}{0.000000in}}%
\pgfusepath{stroke,fill}%
}%
\begin{pgfscope}%
\pgfsys@transformshift{0.573726in}{2.403759in}%
\pgfsys@useobject{currentmarker}{}%
\end{pgfscope}%
\end{pgfscope}%
\begin{pgfscope}%
\definecolor{textcolor}{rgb}{0.000000,0.000000,0.000000}%
\pgfsetstrokecolor{textcolor}%
\pgfsetfillcolor{textcolor}%
\pgftext[x=0.401852in, y=2.345897in, left, base]{\color{textcolor}{\rmfamily\fontsize{12.000000}{14.400000}\selectfont\catcode`\^=\active\def^{\ifmmode\sp\else\^{}\fi}\catcode`\%=\active\def%{\%}$\mathdefault{1}$}}%
\end{pgfscope}%
\begin{pgfscope}%
\pgfsetbuttcap%
\pgfsetroundjoin%
\definecolor{currentfill}{rgb}{1.000000,1.000000,1.000000}%
\pgfsetfillcolor{currentfill}%
\pgfsetlinewidth{0.803000pt}%
\definecolor{currentstroke}{rgb}{0.000000,0.000000,0.000000}%
\pgfsetstrokecolor{currentstroke}%
\pgfsetdash{}{0pt}%
\pgfsys@defobject{currentmarker}{\pgfqpoint{0.000000in}{0.000000in}}{\pgfqpoint{0.048611in}{0.000000in}}{%
\pgfpathmoveto{\pgfqpoint{0.000000in}{0.000000in}}%
\pgfpathlineto{\pgfqpoint{0.048611in}{0.000000in}}%
\pgfusepath{stroke,fill}%
}%
\begin{pgfscope}%
\pgfsys@transformshift{0.573726in}{2.942400in}%
\pgfsys@useobject{currentmarker}{}%
\end{pgfscope}%
\end{pgfscope}%
\begin{pgfscope}%
\definecolor{textcolor}{rgb}{0.000000,0.000000,0.000000}%
\pgfsetstrokecolor{textcolor}%
\pgfsetfillcolor{textcolor}%
\pgftext[x=0.401852in, y=2.884539in, left, base]{\color{textcolor}{\rmfamily\fontsize{12.000000}{14.400000}\selectfont\catcode`\^=\active\def^{\ifmmode\sp\else\^{}\fi}\catcode`\%=\active\def%{\%}$\mathdefault{2}$}}%
\end{pgfscope}%
\begin{pgfscope}%
\definecolor{textcolor}{rgb}{0.000000,0.000000,0.000000}%
\pgfsetstrokecolor{textcolor}%
\pgfsetfillcolor{textcolor}%
\pgftext[x=0.216667in,y=1.865117in,,bottom,rotate=90.000000]{\color{textcolor}{\rmfamily\fontsize{12.000000}{14.400000}\selectfont\catcode`\^=\active\def^{\ifmmode\sp\else\^{}\fi}\catcode`\%=\active\def%{\%}$i_{\mathrm{н}}(t)$}}%
\end{pgfscope}%
\begin{pgfscope}%
\pgfpathrectangle{\pgfqpoint{0.573726in}{0.456901in}}{\pgfqpoint{5.578717in}{2.816433in}}%
\pgfusepath{clip}%
\pgfsetbuttcap%
\pgfsetroundjoin%
\pgfsetlinewidth{2.007500pt}%
\definecolor{currentstroke}{rgb}{0.000000,0.000000,0.000000}%
\pgfsetstrokecolor{currentstroke}%
\pgfsetdash{{7.400000pt}{3.200000pt}}{0.000000pt}%
\pgfpathmoveto{\pgfqpoint{0.563726in}{0.787834in}}%
\pgfpathlineto{\pgfqpoint{0.780811in}{0.787834in}}%
\pgfpathlineto{\pgfqpoint{0.786734in}{2.942400in}}%
\pgfpathlineto{\pgfqpoint{2.895230in}{2.942400in}}%
\pgfpathlineto{\pgfqpoint{2.901153in}{0.787834in}}%
\pgfpathlineto{\pgfqpoint{5.009650in}{0.787834in}}%
\pgfpathlineto{\pgfqpoint{5.015572in}{2.942400in}}%
\pgfpathlineto{\pgfqpoint{5.856602in}{2.942400in}}%
\pgfpathlineto{\pgfqpoint{5.856602in}{2.942400in}}%
\pgfusepath{stroke}%
\end{pgfscope}%
\begin{pgfscope}%
\pgfpathrectangle{\pgfqpoint{0.573726in}{0.456901in}}{\pgfqpoint{5.578717in}{2.816433in}}%
\pgfusepath{clip}%
\pgfsetrectcap%
\pgfsetroundjoin%
\pgfsetlinewidth{2.007500pt}%
\definecolor{currentstroke}{rgb}{0.835294,0.368627,0.000000}%
\pgfsetstrokecolor{currentstroke}%
\pgfsetdash{}{0pt}%
\pgfpathmoveto{\pgfqpoint{0.563726in}{0.766927in}}%
\pgfpathlineto{\pgfqpoint{0.579438in}{0.815339in}}%
\pgfpathlineto{\pgfqpoint{0.597206in}{0.877106in}}%
\pgfpathlineto{\pgfqpoint{0.614974in}{0.945975in}}%
\pgfpathlineto{\pgfqpoint{0.638665in}{1.048260in}}%
\pgfpathlineto{\pgfqpoint{0.662356in}{1.161526in}}%
\pgfpathlineto{\pgfqpoint{0.691970in}{1.316533in}}%
\pgfpathlineto{\pgfqpoint{0.721583in}{1.483565in}}%
\pgfpathlineto{\pgfqpoint{0.763043in}{1.730817in}}%
\pgfpathlineto{\pgfqpoint{0.851884in}{2.266258in}}%
\pgfpathlineto{\pgfqpoint{0.887420in}{2.463868in}}%
\pgfpathlineto{\pgfqpoint{0.917034in}{2.614437in}}%
\pgfpathlineto{\pgfqpoint{0.940725in}{2.723531in}}%
\pgfpathlineto{\pgfqpoint{0.964416in}{2.821201in}}%
\pgfpathlineto{\pgfqpoint{0.982184in}{2.886385in}}%
\pgfpathlineto{\pgfqpoint{0.999952in}{2.944330in}}%
\pgfpathlineto{\pgfqpoint{1.017721in}{2.994850in}}%
\pgfpathlineto{\pgfqpoint{1.035489in}{3.037854in}}%
\pgfpathlineto{\pgfqpoint{1.053257in}{3.073353in}}%
\pgfpathlineto{\pgfqpoint{1.065103in}{3.092897in}}%
\pgfpathlineto{\pgfqpoint{1.076948in}{3.109205in}}%
\pgfpathlineto{\pgfqpoint{1.088794in}{3.122349in}}%
\pgfpathlineto{\pgfqpoint{1.100639in}{3.132418in}}%
\pgfpathlineto{\pgfqpoint{1.112485in}{3.139517in}}%
\pgfpathlineto{\pgfqpoint{1.124330in}{3.143767in}}%
\pgfpathlineto{\pgfqpoint{1.136175in}{3.145304in}}%
\pgfpathlineto{\pgfqpoint{1.148021in}{3.144276in}}%
\pgfpathlineto{\pgfqpoint{1.159866in}{3.140842in}}%
\pgfpathlineto{\pgfqpoint{1.171712in}{3.135175in}}%
\pgfpathlineto{\pgfqpoint{1.183557in}{3.127455in}}%
\pgfpathlineto{\pgfqpoint{1.201326in}{3.112437in}}%
\pgfpathlineto{\pgfqpoint{1.219094in}{3.093886in}}%
\pgfpathlineto{\pgfqpoint{1.242785in}{3.064835in}}%
\pgfpathlineto{\pgfqpoint{1.272399in}{3.023896in}}%
\pgfpathlineto{\pgfqpoint{1.349394in}{2.914510in}}%
\pgfpathlineto{\pgfqpoint{1.373085in}{2.885462in}}%
\pgfpathlineto{\pgfqpoint{1.390853in}{2.866358in}}%
\pgfpathlineto{\pgfqpoint{1.408622in}{2.849958in}}%
\pgfpathlineto{\pgfqpoint{1.426390in}{2.836545in}}%
\pgfpathlineto{\pgfqpoint{1.444158in}{2.826326in}}%
\pgfpathlineto{\pgfqpoint{1.461926in}{2.819433in}}%
\pgfpathlineto{\pgfqpoint{1.479695in}{2.815917in}}%
\pgfpathlineto{\pgfqpoint{1.497463in}{2.815754in}}%
\pgfpathlineto{\pgfqpoint{1.515231in}{2.818849in}}%
\pgfpathlineto{\pgfqpoint{1.532999in}{2.825034in}}%
\pgfpathlineto{\pgfqpoint{1.550768in}{2.834077in}}%
\pgfpathlineto{\pgfqpoint{1.568536in}{2.845691in}}%
\pgfpathlineto{\pgfqpoint{1.592227in}{2.864580in}}%
\pgfpathlineto{\pgfqpoint{1.615918in}{2.886510in}}%
\pgfpathlineto{\pgfqpoint{1.651454in}{2.922870in}}%
\pgfpathlineto{\pgfqpoint{1.704759in}{2.978012in}}%
\pgfpathlineto{\pgfqpoint{1.734373in}{3.005137in}}%
\pgfpathlineto{\pgfqpoint{1.758064in}{3.023452in}}%
\pgfpathlineto{\pgfqpoint{1.775832in}{3.034705in}}%
\pgfpathlineto{\pgfqpoint{1.793600in}{3.043533in}}%
\pgfpathlineto{\pgfqpoint{1.811368in}{3.049734in}}%
\pgfpathlineto{\pgfqpoint{1.829136in}{3.053164in}}%
\pgfpathlineto{\pgfqpoint{1.846905in}{3.053744in}}%
\pgfpathlineto{\pgfqpoint{1.864673in}{3.051462in}}%
\pgfpathlineto{\pgfqpoint{1.882441in}{3.046369in}}%
\pgfpathlineto{\pgfqpoint{1.900209in}{3.038583in}}%
\pgfpathlineto{\pgfqpoint{1.917978in}{3.028283in}}%
\pgfpathlineto{\pgfqpoint{1.935746in}{3.015708in}}%
\pgfpathlineto{\pgfqpoint{1.959437in}{2.995915in}}%
\pgfpathlineto{\pgfqpoint{1.989051in}{2.967488in}}%
\pgfpathlineto{\pgfqpoint{2.089737in}{2.865686in}}%
\pgfpathlineto{\pgfqpoint{2.113428in}{2.846610in}}%
\pgfpathlineto{\pgfqpoint{2.131196in}{2.834825in}}%
\pgfpathlineto{\pgfqpoint{2.148965in}{2.825588in}}%
\pgfpathlineto{\pgfqpoint{2.166733in}{2.819191in}}%
\pgfpathlineto{\pgfqpoint{2.184501in}{2.815869in}}%
\pgfpathlineto{\pgfqpoint{2.202269in}{2.815794in}}%
\pgfpathlineto{\pgfqpoint{2.214115in}{2.817602in}}%
\pgfpathlineto{\pgfqpoint{2.231883in}{2.823131in}}%
\pgfpathlineto{\pgfqpoint{2.249651in}{2.832015in}}%
\pgfpathlineto{\pgfqpoint{2.267419in}{2.844156in}}%
\pgfpathlineto{\pgfqpoint{2.285188in}{2.859377in}}%
\pgfpathlineto{\pgfqpoint{2.302956in}{2.877425in}}%
\pgfpathlineto{\pgfqpoint{2.326647in}{2.905310in}}%
\pgfpathlineto{\pgfqpoint{2.350338in}{2.936658in}}%
\pgfpathlineto{\pgfqpoint{2.391797in}{2.996226in}}%
\pgfpathlineto{\pgfqpoint{2.433256in}{3.055263in}}%
\pgfpathlineto{\pgfqpoint{2.456947in}{3.085527in}}%
\pgfpathlineto{\pgfqpoint{2.474716in}{3.105316in}}%
\pgfpathlineto{\pgfqpoint{2.492484in}{3.121845in}}%
\pgfpathlineto{\pgfqpoint{2.510252in}{3.134443in}}%
\pgfpathlineto{\pgfqpoint{2.522097in}{3.140339in}}%
\pgfpathlineto{\pgfqpoint{2.533943in}{3.144020in}}%
\pgfpathlineto{\pgfqpoint{2.545788in}{3.145314in}}%
\pgfpathlineto{\pgfqpoint{2.557634in}{3.144059in}}%
\pgfpathlineto{\pgfqpoint{2.569479in}{3.140106in}}%
\pgfpathlineto{\pgfqpoint{2.581325in}{3.133318in}}%
\pgfpathlineto{\pgfqpoint{2.593170in}{3.123573in}}%
\pgfpathlineto{\pgfqpoint{2.605016in}{3.110764in}}%
\pgfpathlineto{\pgfqpoint{2.616861in}{3.094798in}}%
\pgfpathlineto{\pgfqpoint{2.628707in}{3.075604in}}%
\pgfpathlineto{\pgfqpoint{2.640552in}{3.053123in}}%
\pgfpathlineto{\pgfqpoint{2.658321in}{3.013163in}}%
\pgfpathlineto{\pgfqpoint{2.676089in}{2.965681in}}%
\pgfpathlineto{\pgfqpoint{2.693857in}{2.910724in}}%
\pgfpathlineto{\pgfqpoint{2.711625in}{2.848441in}}%
\pgfpathlineto{\pgfqpoint{2.729393in}{2.779076in}}%
\pgfpathlineto{\pgfqpoint{2.753084in}{2.676174in}}%
\pgfpathlineto{\pgfqpoint{2.776775in}{2.562354in}}%
\pgfpathlineto{\pgfqpoint{2.806389in}{2.406765in}}%
\pgfpathlineto{\pgfqpoint{2.841926in}{2.204680in}}%
\pgfpathlineto{\pgfqpoint{2.889308in}{1.919448in}}%
\pgfpathlineto{\pgfqpoint{2.960380in}{1.490940in}}%
\pgfpathlineto{\pgfqpoint{2.995917in}{1.291294in}}%
\pgfpathlineto{\pgfqpoint{3.025531in}{1.138469in}}%
\pgfpathlineto{\pgfqpoint{3.049222in}{1.027265in}}%
\pgfpathlineto{\pgfqpoint{3.072913in}{0.927269in}}%
\pgfpathlineto{\pgfqpoint{3.090681in}{0.860231in}}%
\pgfpathlineto{\pgfqpoint{3.108449in}{0.800368in}}%
\pgfpathlineto{\pgfqpoint{3.126217in}{0.747891in}}%
\pgfpathlineto{\pgfqpoint{3.143986in}{0.702916in}}%
\pgfpathlineto{\pgfqpoint{3.161754in}{0.665457in}}%
\pgfpathlineto{\pgfqpoint{3.173599in}{0.644626in}}%
\pgfpathlineto{\pgfqpoint{3.185445in}{0.627054in}}%
\pgfpathlineto{\pgfqpoint{3.197290in}{0.612677in}}%
\pgfpathlineto{\pgfqpoint{3.209136in}{0.601412in}}%
\pgfpathlineto{\pgfqpoint{3.220981in}{0.593160in}}%
\pgfpathlineto{\pgfqpoint{3.232827in}{0.587806in}}%
\pgfpathlineto{\pgfqpoint{3.244672in}{0.585220in}}%
\pgfpathlineto{\pgfqpoint{3.256518in}{0.585259in}}%
\pgfpathlineto{\pgfqpoint{3.268363in}{0.587767in}}%
\pgfpathlineto{\pgfqpoint{3.280209in}{0.592577in}}%
\pgfpathlineto{\pgfqpoint{3.292054in}{0.599513in}}%
\pgfpathlineto{\pgfqpoint{3.309822in}{0.613495in}}%
\pgfpathlineto{\pgfqpoint{3.327591in}{0.631186in}}%
\pgfpathlineto{\pgfqpoint{3.345359in}{0.651909in}}%
\pgfpathlineto{\pgfqpoint{3.369050in}{0.683059in}}%
\pgfpathlineto{\pgfqpoint{3.410509in}{0.742673in}}%
\pgfpathlineto{\pgfqpoint{3.451968in}{0.801677in}}%
\pgfpathlineto{\pgfqpoint{3.475659in}{0.832262in}}%
\pgfpathlineto{\pgfqpoint{3.499350in}{0.859121in}}%
\pgfpathlineto{\pgfqpoint{3.517118in}{0.876259in}}%
\pgfpathlineto{\pgfqpoint{3.534887in}{0.890478in}}%
\pgfpathlineto{\pgfqpoint{3.552655in}{0.901549in}}%
\pgfpathlineto{\pgfqpoint{3.570423in}{0.909322in}}%
\pgfpathlineto{\pgfqpoint{3.588191in}{0.913723in}}%
\pgfpathlineto{\pgfqpoint{3.605960in}{0.914757in}}%
\pgfpathlineto{\pgfqpoint{3.623728in}{0.912502in}}%
\pgfpathlineto{\pgfqpoint{3.641496in}{0.907107in}}%
\pgfpathlineto{\pgfqpoint{3.659264in}{0.898786in}}%
\pgfpathlineto{\pgfqpoint{3.677032in}{0.887815in}}%
\pgfpathlineto{\pgfqpoint{3.694801in}{0.874520in}}%
\pgfpathlineto{\pgfqpoint{3.718492in}{0.853827in}}%
\pgfpathlineto{\pgfqpoint{3.748105in}{0.824583in}}%
\pgfpathlineto{\pgfqpoint{3.831024in}{0.739724in}}%
\pgfpathlineto{\pgfqpoint{3.854715in}{0.719174in}}%
\pgfpathlineto{\pgfqpoint{3.878406in}{0.701951in}}%
\pgfpathlineto{\pgfqpoint{3.896174in}{0.691651in}}%
\pgfpathlineto{\pgfqpoint{3.913942in}{0.683865in}}%
\pgfpathlineto{\pgfqpoint{3.931710in}{0.678772in}}%
\pgfpathlineto{\pgfqpoint{3.949479in}{0.676490in}}%
\pgfpathlineto{\pgfqpoint{3.967247in}{0.677070in}}%
\pgfpathlineto{\pgfqpoint{3.985015in}{0.680500in}}%
\pgfpathlineto{\pgfqpoint{4.002783in}{0.686701in}}%
\pgfpathlineto{\pgfqpoint{4.020552in}{0.695530in}}%
\pgfpathlineto{\pgfqpoint{4.038320in}{0.706782in}}%
\pgfpathlineto{\pgfqpoint{4.062011in}{0.725097in}}%
\pgfpathlineto{\pgfqpoint{4.085702in}{0.746494in}}%
\pgfpathlineto{\pgfqpoint{4.121238in}{0.782384in}}%
\pgfpathlineto{\pgfqpoint{4.186388in}{0.849438in}}%
\pgfpathlineto{\pgfqpoint{4.210079in}{0.870700in}}%
\pgfpathlineto{\pgfqpoint{4.233770in}{0.888680in}}%
\pgfpathlineto{\pgfqpoint{4.251539in}{0.899472in}}%
\pgfpathlineto{\pgfqpoint{4.269307in}{0.907592in}}%
\pgfpathlineto{\pgfqpoint{4.287075in}{0.912770in}}%
\pgfpathlineto{\pgfqpoint{4.304843in}{0.914794in}}%
\pgfpathlineto{\pgfqpoint{4.316689in}{0.914318in}}%
\pgfpathlineto{\pgfqpoint{4.328534in}{0.912349in}}%
\pgfpathlineto{\pgfqpoint{4.346302in}{0.906579in}}%
\pgfpathlineto{\pgfqpoint{4.364071in}{0.897459in}}%
\pgfpathlineto{\pgfqpoint{4.381839in}{0.885091in}}%
\pgfpathlineto{\pgfqpoint{4.399607in}{0.869658in}}%
\pgfpathlineto{\pgfqpoint{4.417375in}{0.851420in}}%
\pgfpathlineto{\pgfqpoint{4.441066in}{0.823323in}}%
\pgfpathlineto{\pgfqpoint{4.470680in}{0.783548in}}%
\pgfpathlineto{\pgfqpoint{4.565444in}{0.650343in}}%
\pgfpathlineto{\pgfqpoint{4.583212in}{0.629814in}}%
\pgfpathlineto{\pgfqpoint{4.600980in}{0.612365in}}%
\pgfpathlineto{\pgfqpoint{4.618749in}{0.598674in}}%
\pgfpathlineto{\pgfqpoint{4.630594in}{0.591957in}}%
\pgfpathlineto{\pgfqpoint{4.642440in}{0.587385in}}%
\pgfpathlineto{\pgfqpoint{4.654285in}{0.585134in}}%
\pgfpathlineto{\pgfqpoint{4.666131in}{0.585368in}}%
\pgfpathlineto{\pgfqpoint{4.677976in}{0.588244in}}%
\pgfpathlineto{\pgfqpoint{4.689822in}{0.593903in}}%
\pgfpathlineto{\pgfqpoint{4.701667in}{0.602473in}}%
\pgfpathlineto{\pgfqpoint{4.713513in}{0.614067in}}%
\pgfpathlineto{\pgfqpoint{4.725358in}{0.628782in}}%
\pgfpathlineto{\pgfqpoint{4.737204in}{0.646701in}}%
\pgfpathlineto{\pgfqpoint{4.749049in}{0.667885in}}%
\pgfpathlineto{\pgfqpoint{4.760894in}{0.692380in}}%
\pgfpathlineto{\pgfqpoint{4.778663in}{0.735385in}}%
\pgfpathlineto{\pgfqpoint{4.796431in}{0.785904in}}%
\pgfpathlineto{\pgfqpoint{4.814199in}{0.843849in}}%
\pgfpathlineto{\pgfqpoint{4.831967in}{0.909033in}}%
\pgfpathlineto{\pgfqpoint{4.855658in}{1.006703in}}%
\pgfpathlineto{\pgfqpoint{4.879349in}{1.115797in}}%
\pgfpathlineto{\pgfqpoint{4.903040in}{1.235108in}}%
\pgfpathlineto{\pgfqpoint{4.932654in}{1.396370in}}%
\pgfpathlineto{\pgfqpoint{4.968191in}{1.603238in}}%
\pgfpathlineto{\pgfqpoint{5.045186in}{2.071146in}}%
\pgfpathlineto{\pgfqpoint{5.086645in}{2.314707in}}%
\pgfpathlineto{\pgfqpoint{5.116259in}{2.477330in}}%
\pgfpathlineto{\pgfqpoint{5.145873in}{2.626638in}}%
\pgfpathlineto{\pgfqpoint{5.169564in}{2.734561in}}%
\pgfpathlineto{\pgfqpoint{5.193255in}{2.830946in}}%
\pgfpathlineto{\pgfqpoint{5.211023in}{2.895112in}}%
\pgfpathlineto{\pgfqpoint{5.228791in}{2.952005in}}%
\pgfpathlineto{\pgfqpoint{5.246559in}{3.001454in}}%
\pgfpathlineto{\pgfqpoint{5.264328in}{3.043383in}}%
\pgfpathlineto{\pgfqpoint{5.282096in}{3.077817in}}%
\pgfpathlineto{\pgfqpoint{5.293941in}{3.096662in}}%
\pgfpathlineto{\pgfqpoint{5.305787in}{3.112285in}}%
\pgfpathlineto{\pgfqpoint{5.317632in}{3.124762in}}%
\pgfpathlineto{\pgfqpoint{5.329478in}{3.134185in}}%
\pgfpathlineto{\pgfqpoint{5.341323in}{3.140663in}}%
\pgfpathlineto{\pgfqpoint{5.353169in}{3.144320in}}%
\pgfpathlineto{\pgfqpoint{5.365014in}{3.145294in}}%
\pgfpathlineto{\pgfqpoint{5.376860in}{3.143737in}}%
\pgfpathlineto{\pgfqpoint{5.388705in}{3.139810in}}%
\pgfpathlineto{\pgfqpoint{5.400551in}{3.133687in}}%
\pgfpathlineto{\pgfqpoint{5.412396in}{3.125551in}}%
\pgfpathlineto{\pgfqpoint{5.430164in}{3.109988in}}%
\pgfpathlineto{\pgfqpoint{5.447933in}{3.090987in}}%
\pgfpathlineto{\pgfqpoint{5.471624in}{3.061490in}}%
\pgfpathlineto{\pgfqpoint{5.501237in}{3.020241in}}%
\pgfpathlineto{\pgfqpoint{5.572310in}{2.918942in}}%
\pgfpathlineto{\pgfqpoint{5.596001in}{2.889379in}}%
\pgfpathlineto{\pgfqpoint{5.619692in}{2.863843in}}%
\pgfpathlineto{\pgfqpoint{5.637461in}{2.847853in}}%
\pgfpathlineto{\pgfqpoint{5.655229in}{2.834885in}}%
\pgfpathlineto{\pgfqpoint{5.672997in}{2.825136in}}%
\pgfpathlineto{\pgfqpoint{5.690765in}{2.818723in}}%
\pgfpathlineto{\pgfqpoint{5.708533in}{2.815690in}}%
\pgfpathlineto{\pgfqpoint{5.726302in}{2.816000in}}%
\pgfpathlineto{\pgfqpoint{5.744070in}{2.819548in}}%
\pgfpathlineto{\pgfqpoint{5.761838in}{2.826157in}}%
\pgfpathlineto{\pgfqpoint{5.779606in}{2.835587in}}%
\pgfpathlineto{\pgfqpoint{5.797375in}{2.847540in}}%
\pgfpathlineto{\pgfqpoint{5.821066in}{2.866801in}}%
\pgfpathlineto{\pgfqpoint{5.844757in}{2.888998in}}%
\pgfpathlineto{\pgfqpoint{5.856602in}{2.900881in}}%
\pgfpathlineto{\pgfqpoint{5.856602in}{2.900881in}}%
\pgfusepath{stroke}%
\end{pgfscope}%
\begin{pgfscope}%
\pgfpathrectangle{\pgfqpoint{0.573726in}{0.456901in}}{\pgfqpoint{5.578717in}{2.816433in}}%
\pgfusepath{clip}%
\pgfsetbuttcap%
\pgfsetroundjoin%
\pgfsetlinewidth{2.007500pt}%
\definecolor{currentstroke}{rgb}{0.000000,0.619608,0.450980}%
\pgfsetstrokecolor{currentstroke}%
\pgfsetdash{{12.800000pt}{3.200000pt}{2.000000pt}{3.200000pt}}{0.000000pt}%
\pgfpathmoveto{\pgfqpoint{0.563726in}{1.267873in}}%
\pgfpathlineto{\pgfqpoint{0.579438in}{1.284369in}}%
\pgfpathlineto{\pgfqpoint{0.597206in}{1.306612in}}%
\pgfpathlineto{\pgfqpoint{0.614974in}{1.332601in}}%
\pgfpathlineto{\pgfqpoint{0.632742in}{1.362283in}}%
\pgfpathlineto{\pgfqpoint{0.650511in}{1.395555in}}%
\pgfpathlineto{\pgfqpoint{0.674202in}{1.445236in}}%
\pgfpathlineto{\pgfqpoint{0.697892in}{1.500542in}}%
\pgfpathlineto{\pgfqpoint{0.727506in}{1.576628in}}%
\pgfpathlineto{\pgfqpoint{0.757120in}{1.659047in}}%
\pgfpathlineto{\pgfqpoint{0.798579in}{1.781754in}}%
\pgfpathlineto{\pgfqpoint{0.899266in}{2.084319in}}%
\pgfpathlineto{\pgfqpoint{0.928879in}{2.166094in}}%
\pgfpathlineto{\pgfqpoint{0.958493in}{2.241349in}}%
\pgfpathlineto{\pgfqpoint{0.982184in}{2.295874in}}%
\pgfpathlineto{\pgfqpoint{1.005875in}{2.344688in}}%
\pgfpathlineto{\pgfqpoint{1.029566in}{2.387323in}}%
\pgfpathlineto{\pgfqpoint{1.047334in}{2.415049in}}%
\pgfpathlineto{\pgfqpoint{1.065103in}{2.439041in}}%
\pgfpathlineto{\pgfqpoint{1.082871in}{2.459269in}}%
\pgfpathlineto{\pgfqpoint{1.100639in}{2.475756in}}%
\pgfpathlineto{\pgfqpoint{1.118407in}{2.488567in}}%
\pgfpathlineto{\pgfqpoint{1.136175in}{2.497818in}}%
\pgfpathlineto{\pgfqpoint{1.153944in}{2.503665in}}%
\pgfpathlineto{\pgfqpoint{1.171712in}{2.506305in}}%
\pgfpathlineto{\pgfqpoint{1.189480in}{2.505971in}}%
\pgfpathlineto{\pgfqpoint{1.207248in}{2.502929in}}%
\pgfpathlineto{\pgfqpoint{1.225017in}{2.497469in}}%
\pgfpathlineto{\pgfqpoint{1.242785in}{2.489908in}}%
\pgfpathlineto{\pgfqpoint{1.266476in}{2.477129in}}%
\pgfpathlineto{\pgfqpoint{1.296090in}{2.457964in}}%
\pgfpathlineto{\pgfqpoint{1.343472in}{2.423683in}}%
\pgfpathlineto{\pgfqpoint{1.390853in}{2.390333in}}%
\pgfpathlineto{\pgfqpoint{1.420467in}{2.372337in}}%
\pgfpathlineto{\pgfqpoint{1.444158in}{2.360364in}}%
\pgfpathlineto{\pgfqpoint{1.467849in}{2.350966in}}%
\pgfpathlineto{\pgfqpoint{1.491540in}{2.344406in}}%
\pgfpathlineto{\pgfqpoint{1.515231in}{2.340826in}}%
\pgfpathlineto{\pgfqpoint{1.532999in}{2.340112in}}%
\pgfpathlineto{\pgfqpoint{1.556690in}{2.341722in}}%
\pgfpathlineto{\pgfqpoint{1.580381in}{2.346085in}}%
\pgfpathlineto{\pgfqpoint{1.604072in}{2.352919in}}%
\pgfpathlineto{\pgfqpoint{1.627763in}{2.361852in}}%
\pgfpathlineto{\pgfqpoint{1.657377in}{2.375290in}}%
\pgfpathlineto{\pgfqpoint{1.704759in}{2.399697in}}%
\pgfpathlineto{\pgfqpoint{1.758064in}{2.426723in}}%
\pgfpathlineto{\pgfqpoint{1.787677in}{2.439483in}}%
\pgfpathlineto{\pgfqpoint{1.811368in}{2.447787in}}%
\pgfpathlineto{\pgfqpoint{1.835059in}{2.454023in}}%
\pgfpathlineto{\pgfqpoint{1.858750in}{2.457936in}}%
\pgfpathlineto{\pgfqpoint{1.882441in}{2.459368in}}%
\pgfpathlineto{\pgfqpoint{1.906132in}{2.458260in}}%
\pgfpathlineto{\pgfqpoint{1.929823in}{2.454659in}}%
\pgfpathlineto{\pgfqpoint{1.953514in}{2.448714in}}%
\pgfpathlineto{\pgfqpoint{1.977205in}{2.440670in}}%
\pgfpathlineto{\pgfqpoint{2.006819in}{2.428177in}}%
\pgfpathlineto{\pgfqpoint{2.048278in}{2.407633in}}%
\pgfpathlineto{\pgfqpoint{2.125274in}{2.368633in}}%
\pgfpathlineto{\pgfqpoint{2.154887in}{2.356449in}}%
\pgfpathlineto{\pgfqpoint{2.178578in}{2.348841in}}%
\pgfpathlineto{\pgfqpoint{2.202269in}{2.343567in}}%
\pgfpathlineto{\pgfqpoint{2.225960in}{2.340941in}}%
\pgfpathlineto{\pgfqpoint{2.249651in}{2.341180in}}%
\pgfpathlineto{\pgfqpoint{2.267419in}{2.343308in}}%
\pgfpathlineto{\pgfqpoint{2.291110in}{2.348750in}}%
\pgfpathlineto{\pgfqpoint{2.314801in}{2.357069in}}%
\pgfpathlineto{\pgfqpoint{2.338492in}{2.368044in}}%
\pgfpathlineto{\pgfqpoint{2.362183in}{2.381331in}}%
\pgfpathlineto{\pgfqpoint{2.391797in}{2.400488in}}%
\pgfpathlineto{\pgfqpoint{2.445102in}{2.438570in}}%
\pgfpathlineto{\pgfqpoint{2.486561in}{2.466993in}}%
\pgfpathlineto{\pgfqpoint{2.510252in}{2.480947in}}%
\pgfpathlineto{\pgfqpoint{2.533943in}{2.492216in}}%
\pgfpathlineto{\pgfqpoint{2.551711in}{2.498425in}}%
\pgfpathlineto{\pgfqpoint{2.569479in}{2.502353in}}%
\pgfpathlineto{\pgfqpoint{2.587248in}{2.503699in}}%
\pgfpathlineto{\pgfqpoint{2.605016in}{2.502187in}}%
\pgfpathlineto{\pgfqpoint{2.622784in}{2.497570in}}%
\pgfpathlineto{\pgfqpoint{2.640552in}{2.489636in}}%
\pgfpathlineto{\pgfqpoint{2.658321in}{2.478210in}}%
\pgfpathlineto{\pgfqpoint{2.676089in}{2.463161in}}%
\pgfpathlineto{\pgfqpoint{2.693857in}{2.444401in}}%
\pgfpathlineto{\pgfqpoint{2.711625in}{2.421889in}}%
\pgfpathlineto{\pgfqpoint{2.729393in}{2.395634in}}%
\pgfpathlineto{\pgfqpoint{2.747162in}{2.365692in}}%
\pgfpathlineto{\pgfqpoint{2.764930in}{2.332169in}}%
\pgfpathlineto{\pgfqpoint{2.788621in}{2.282173in}}%
\pgfpathlineto{\pgfqpoint{2.812312in}{2.226582in}}%
\pgfpathlineto{\pgfqpoint{2.841926in}{2.150192in}}%
\pgfpathlineto{\pgfqpoint{2.871539in}{2.067539in}}%
\pgfpathlineto{\pgfqpoint{2.912999in}{1.944637in}}%
\pgfpathlineto{\pgfqpoint{3.013685in}{1.642299in}}%
\pgfpathlineto{\pgfqpoint{3.043299in}{1.560769in}}%
\pgfpathlineto{\pgfqpoint{3.072913in}{1.485828in}}%
\pgfpathlineto{\pgfqpoint{3.096604in}{1.431595in}}%
\pgfpathlineto{\pgfqpoint{3.120295in}{1.383102in}}%
\pgfpathlineto{\pgfqpoint{3.138063in}{1.350781in}}%
\pgfpathlineto{\pgfqpoint{3.155831in}{1.322088in}}%
\pgfpathlineto{\pgfqpoint{3.173599in}{1.297113in}}%
\pgfpathlineto{\pgfqpoint{3.191367in}{1.275899in}}%
\pgfpathlineto{\pgfqpoint{3.209136in}{1.258438in}}%
\pgfpathlineto{\pgfqpoint{3.226904in}{1.244673in}}%
\pgfpathlineto{\pgfqpoint{3.244672in}{1.234503in}}%
\pgfpathlineto{\pgfqpoint{3.262440in}{1.227782in}}%
\pgfpathlineto{\pgfqpoint{3.280209in}{1.224323in}}%
\pgfpathlineto{\pgfqpoint{3.297977in}{1.223903in}}%
\pgfpathlineto{\pgfqpoint{3.315745in}{1.226263in}}%
\pgfpathlineto{\pgfqpoint{3.333513in}{1.231119in}}%
\pgfpathlineto{\pgfqpoint{3.351282in}{1.238161in}}%
\pgfpathlineto{\pgfqpoint{3.374973in}{1.250385in}}%
\pgfpathlineto{\pgfqpoint{3.398663in}{1.265108in}}%
\pgfpathlineto{\pgfqpoint{3.434200in}{1.290081in}}%
\pgfpathlineto{\pgfqpoint{3.499350in}{1.336824in}}%
\pgfpathlineto{\pgfqpoint{3.528964in}{1.355273in}}%
\pgfpathlineto{\pgfqpoint{3.552655in}{1.367711in}}%
\pgfpathlineto{\pgfqpoint{3.576346in}{1.377640in}}%
\pgfpathlineto{\pgfqpoint{3.600037in}{1.384772in}}%
\pgfpathlineto{\pgfqpoint{3.623728in}{1.388943in}}%
\pgfpathlineto{\pgfqpoint{3.641496in}{1.390097in}}%
\pgfpathlineto{\pgfqpoint{3.665187in}{1.389053in}}%
\pgfpathlineto{\pgfqpoint{3.688878in}{1.385211in}}%
\pgfpathlineto{\pgfqpoint{3.712569in}{1.378836in}}%
\pgfpathlineto{\pgfqpoint{3.736260in}{1.370282in}}%
\pgfpathlineto{\pgfqpoint{3.765874in}{1.357188in}}%
\pgfpathlineto{\pgfqpoint{3.807333in}{1.336122in}}%
\pgfpathlineto{\pgfqpoint{3.872483in}{1.302920in}}%
\pgfpathlineto{\pgfqpoint{3.902097in}{1.290259in}}%
\pgfpathlineto{\pgfqpoint{3.925788in}{1.282058in}}%
\pgfpathlineto{\pgfqpoint{3.949479in}{1.275941in}}%
\pgfpathlineto{\pgfqpoint{3.973170in}{1.272157in}}%
\pgfpathlineto{\pgfqpoint{3.996861in}{1.270861in}}%
\pgfpathlineto{\pgfqpoint{4.020552in}{1.272105in}}%
\pgfpathlineto{\pgfqpoint{4.044243in}{1.275836in}}%
\pgfpathlineto{\pgfqpoint{4.067933in}{1.281900in}}%
\pgfpathlineto{\pgfqpoint{4.091624in}{1.290048in}}%
\pgfpathlineto{\pgfqpoint{4.121238in}{1.302642in}}%
\pgfpathlineto{\pgfqpoint{4.162697in}{1.323262in}}%
\pgfpathlineto{\pgfqpoint{4.233770in}{1.359470in}}%
\pgfpathlineto{\pgfqpoint{4.263384in}{1.372049in}}%
\pgfpathlineto{\pgfqpoint{4.287075in}{1.380072in}}%
\pgfpathlineto{\pgfqpoint{4.310766in}{1.385832in}}%
\pgfpathlineto{\pgfqpoint{4.334457in}{1.388998in}}%
\pgfpathlineto{\pgfqpoint{4.358148in}{1.389334in}}%
\pgfpathlineto{\pgfqpoint{4.375916in}{1.387646in}}%
\pgfpathlineto{\pgfqpoint{4.399607in}{1.382785in}}%
\pgfpathlineto{\pgfqpoint{4.423298in}{1.375017in}}%
\pgfpathlineto{\pgfqpoint{4.446989in}{1.364541in}}%
\pgfpathlineto{\pgfqpoint{4.470680in}{1.351675in}}%
\pgfpathlineto{\pgfqpoint{4.500294in}{1.332905in}}%
\pgfpathlineto{\pgfqpoint{4.541753in}{1.303602in}}%
\pgfpathlineto{\pgfqpoint{4.595058in}{1.266228in}}%
\pgfpathlineto{\pgfqpoint{4.618749in}{1.251842in}}%
\pgfpathlineto{\pgfqpoint{4.642440in}{1.239982in}}%
\pgfpathlineto{\pgfqpoint{4.660208in}{1.233228in}}%
\pgfpathlineto{\pgfqpoint{4.677976in}{1.228672in}}%
\pgfpathlineto{\pgfqpoint{4.695744in}{1.226622in}}%
\pgfpathlineto{\pgfqpoint{4.713513in}{1.227360in}}%
\pgfpathlineto{\pgfqpoint{4.731281in}{1.231142in}}%
\pgfpathlineto{\pgfqpoint{4.749049in}{1.238189in}}%
\pgfpathlineto{\pgfqpoint{4.766817in}{1.248686in}}%
\pgfpathlineto{\pgfqpoint{4.784585in}{1.262776in}}%
\pgfpathlineto{\pgfqpoint{4.802354in}{1.280558in}}%
\pgfpathlineto{\pgfqpoint{4.820122in}{1.302086in}}%
\pgfpathlineto{\pgfqpoint{4.837890in}{1.327364in}}%
\pgfpathlineto{\pgfqpoint{4.855658in}{1.356349in}}%
\pgfpathlineto{\pgfqpoint{4.873427in}{1.388947in}}%
\pgfpathlineto{\pgfqpoint{4.897118in}{1.437781in}}%
\pgfpathlineto{\pgfqpoint{4.920809in}{1.492319in}}%
\pgfpathlineto{\pgfqpoint{4.944500in}{1.551958in}}%
\pgfpathlineto{\pgfqpoint{4.974113in}{1.632561in}}%
\pgfpathlineto{\pgfqpoint{5.009650in}{1.735954in}}%
\pgfpathlineto{\pgfqpoint{5.145873in}{2.141020in}}%
\pgfpathlineto{\pgfqpoint{5.175487in}{2.218494in}}%
\pgfpathlineto{\pgfqpoint{5.199177in}{2.275104in}}%
\pgfpathlineto{\pgfqpoint{5.222868in}{2.326225in}}%
\pgfpathlineto{\pgfqpoint{5.246559in}{2.371334in}}%
\pgfpathlineto{\pgfqpoint{5.264328in}{2.400988in}}%
\pgfpathlineto{\pgfqpoint{5.282096in}{2.426942in}}%
\pgfpathlineto{\pgfqpoint{5.299864in}{2.449143in}}%
\pgfpathlineto{\pgfqpoint{5.317632in}{2.467584in}}%
\pgfpathlineto{\pgfqpoint{5.335401in}{2.482309in}}%
\pgfpathlineto{\pgfqpoint{5.353169in}{2.493408in}}%
\pgfpathlineto{\pgfqpoint{5.370937in}{2.501015in}}%
\pgfpathlineto{\pgfqpoint{5.388705in}{2.505308in}}%
\pgfpathlineto{\pgfqpoint{5.406474in}{2.506501in}}%
\pgfpathlineto{\pgfqpoint{5.424242in}{2.504843in}}%
\pgfpathlineto{\pgfqpoint{5.442010in}{2.500612in}}%
\pgfpathlineto{\pgfqpoint{5.459778in}{2.494111in}}%
\pgfpathlineto{\pgfqpoint{5.483469in}{2.482473in}}%
\pgfpathlineto{\pgfqpoint{5.507160in}{2.468177in}}%
\pgfpathlineto{\pgfqpoint{5.542697in}{2.443540in}}%
\pgfpathlineto{\pgfqpoint{5.619692in}{2.388679in}}%
\pgfpathlineto{\pgfqpoint{5.649306in}{2.370942in}}%
\pgfpathlineto{\pgfqpoint{5.672997in}{2.359228in}}%
\pgfpathlineto{\pgfqpoint{5.696688in}{2.350124in}}%
\pgfpathlineto{\pgfqpoint{5.720379in}{2.343879in}}%
\pgfpathlineto{\pgfqpoint{5.744070in}{2.340621in}}%
\pgfpathlineto{\pgfqpoint{5.761838in}{2.340147in}}%
\pgfpathlineto{\pgfqpoint{5.785529in}{2.342062in}}%
\pgfpathlineto{\pgfqpoint{5.809220in}{2.346705in}}%
\pgfpathlineto{\pgfqpoint{5.832911in}{2.353783in}}%
\pgfpathlineto{\pgfqpoint{5.856602in}{2.362916in}}%
\pgfpathlineto{\pgfqpoint{5.856602in}{2.362916in}}%
\pgfusepath{stroke}%
\end{pgfscope}%
\begin{pgfscope}%
\pgfpathrectangle{\pgfqpoint{0.573726in}{0.456901in}}{\pgfqpoint{5.578717in}{2.816433in}}%
\pgfusepath{clip}%
\pgfsetrectcap%
\pgfsetroundjoin%
\pgfsetlinewidth{1.003750pt}%
\definecolor{currentstroke}{rgb}{0.000000,0.750000,0.750000}%
\pgfsetstrokecolor{currentstroke}%
\pgfsetdash{}{0pt}%
\pgfpathmoveto{\pgfqpoint{0.563726in}{1.605153in}}%
\pgfpathlineto{\pgfqpoint{0.638665in}{1.677321in}}%
\pgfpathlineto{\pgfqpoint{0.733429in}{1.771803in}}%
\pgfpathlineto{\pgfqpoint{1.041412in}{2.081756in}}%
\pgfpathlineto{\pgfqpoint{1.118407in}{2.154661in}}%
\pgfpathlineto{\pgfqpoint{1.183557in}{2.213426in}}%
\pgfpathlineto{\pgfqpoint{1.242785in}{2.264030in}}%
\pgfpathlineto{\pgfqpoint{1.296090in}{2.306941in}}%
\pgfpathlineto{\pgfqpoint{1.349394in}{2.347079in}}%
\pgfpathlineto{\pgfqpoint{1.396776in}{2.380225in}}%
\pgfpathlineto{\pgfqpoint{1.444158in}{2.410816in}}%
\pgfpathlineto{\pgfqpoint{1.491540in}{2.438700in}}%
\pgfpathlineto{\pgfqpoint{1.538922in}{2.463740in}}%
\pgfpathlineto{\pgfqpoint{1.580381in}{2.483217in}}%
\pgfpathlineto{\pgfqpoint{1.621840in}{2.500348in}}%
\pgfpathlineto{\pgfqpoint{1.663300in}{2.515065in}}%
\pgfpathlineto{\pgfqpoint{1.704759in}{2.527314in}}%
\pgfpathlineto{\pgfqpoint{1.746218in}{2.537049in}}%
\pgfpathlineto{\pgfqpoint{1.787677in}{2.544231in}}%
\pgfpathlineto{\pgfqpoint{1.829136in}{2.548834in}}%
\pgfpathlineto{\pgfqpoint{1.864673in}{2.550713in}}%
\pgfpathlineto{\pgfqpoint{1.900209in}{2.550679in}}%
\pgfpathlineto{\pgfqpoint{1.935746in}{2.548732in}}%
\pgfpathlineto{\pgfqpoint{1.977205in}{2.544050in}}%
\pgfpathlineto{\pgfqpoint{2.018664in}{2.536790in}}%
\pgfpathlineto{\pgfqpoint{2.060123in}{2.526978in}}%
\pgfpathlineto{\pgfqpoint{2.101583in}{2.514653in}}%
\pgfpathlineto{\pgfqpoint{2.143042in}{2.499861in}}%
\pgfpathlineto{\pgfqpoint{2.184501in}{2.482658in}}%
\pgfpathlineto{\pgfqpoint{2.225960in}{2.463110in}}%
\pgfpathlineto{\pgfqpoint{2.267419in}{2.441291in}}%
\pgfpathlineto{\pgfqpoint{2.314801in}{2.413681in}}%
\pgfpathlineto{\pgfqpoint{2.362183in}{2.383349in}}%
\pgfpathlineto{\pgfqpoint{2.409565in}{2.350447in}}%
\pgfpathlineto{\pgfqpoint{2.462870in}{2.310564in}}%
\pgfpathlineto{\pgfqpoint{2.516175in}{2.267885in}}%
\pgfpathlineto{\pgfqpoint{2.575402in}{2.217510in}}%
\pgfpathlineto{\pgfqpoint{2.640552in}{2.158960in}}%
\pgfpathlineto{\pgfqpoint{2.711625in}{2.091966in}}%
\pgfpathlineto{\pgfqpoint{2.794544in}{2.010662in}}%
\pgfpathlineto{\pgfqpoint{2.912999in}{1.890995in}}%
\pgfpathlineto{\pgfqpoint{3.102526in}{1.699415in}}%
\pgfpathlineto{\pgfqpoint{3.185445in}{1.618841in}}%
\pgfpathlineto{\pgfqpoint{3.256518in}{1.552708in}}%
\pgfpathlineto{\pgfqpoint{3.321668in}{1.495130in}}%
\pgfpathlineto{\pgfqpoint{3.380895in}{1.445782in}}%
\pgfpathlineto{\pgfqpoint{3.434200in}{1.404137in}}%
\pgfpathlineto{\pgfqpoint{3.487505in}{1.365386in}}%
\pgfpathlineto{\pgfqpoint{3.534887in}{1.333566in}}%
\pgfpathlineto{\pgfqpoint{3.582269in}{1.304383in}}%
\pgfpathlineto{\pgfqpoint{3.629650in}{1.277981in}}%
\pgfpathlineto{\pgfqpoint{3.671110in}{1.257264in}}%
\pgfpathlineto{\pgfqpoint{3.712569in}{1.238855in}}%
\pgfpathlineto{\pgfqpoint{3.754028in}{1.222825in}}%
\pgfpathlineto{\pgfqpoint{3.795487in}{1.209235in}}%
\pgfpathlineto{\pgfqpoint{3.836947in}{1.198135in}}%
\pgfpathlineto{\pgfqpoint{3.878406in}{1.189569in}}%
\pgfpathlineto{\pgfqpoint{3.919865in}{1.183568in}}%
\pgfpathlineto{\pgfqpoint{3.961324in}{1.180156in}}%
\pgfpathlineto{\pgfqpoint{3.996861in}{1.179301in}}%
\pgfpathlineto{\pgfqpoint{4.032397in}{1.180360in}}%
\pgfpathlineto{\pgfqpoint{4.073856in}{1.184011in}}%
\pgfpathlineto{\pgfqpoint{4.115315in}{1.190248in}}%
\pgfpathlineto{\pgfqpoint{4.156775in}{1.199048in}}%
\pgfpathlineto{\pgfqpoint{4.198234in}{1.210378in}}%
\pgfpathlineto{\pgfqpoint{4.239693in}{1.224194in}}%
\pgfpathlineto{\pgfqpoint{4.281152in}{1.240445in}}%
\pgfpathlineto{\pgfqpoint{4.322611in}{1.259068in}}%
\pgfpathlineto{\pgfqpoint{4.364071in}{1.279993in}}%
\pgfpathlineto{\pgfqpoint{4.411453in}{1.306623in}}%
\pgfpathlineto{\pgfqpoint{4.458835in}{1.336024in}}%
\pgfpathlineto{\pgfqpoint{4.506217in}{1.368048in}}%
\pgfpathlineto{\pgfqpoint{4.559521in}{1.407015in}}%
\pgfpathlineto{\pgfqpoint{4.612826in}{1.448856in}}%
\pgfpathlineto{\pgfqpoint{4.672053in}{1.498400in}}%
\pgfpathlineto{\pgfqpoint{4.731281in}{1.550784in}}%
\pgfpathlineto{\pgfqpoint{4.802354in}{1.616824in}}%
\pgfpathlineto{\pgfqpoint{4.879349in}{1.691468in}}%
\pgfpathlineto{\pgfqpoint{4.980036in}{1.792354in}}%
\pgfpathlineto{\pgfqpoint{5.246559in}{2.061189in}}%
\pgfpathlineto{\pgfqpoint{5.323555in}{2.134970in}}%
\pgfpathlineto{\pgfqpoint{5.394628in}{2.199961in}}%
\pgfpathlineto{\pgfqpoint{5.453855in}{2.251297in}}%
\pgfpathlineto{\pgfqpoint{5.513083in}{2.299641in}}%
\pgfpathlineto{\pgfqpoint{5.566388in}{2.340282in}}%
\pgfpathlineto{\pgfqpoint{5.619692in}{2.377940in}}%
\pgfpathlineto{\pgfqpoint{5.667074in}{2.408719in}}%
\pgfpathlineto{\pgfqpoint{5.714456in}{2.436802in}}%
\pgfpathlineto{\pgfqpoint{5.761838in}{2.462048in}}%
\pgfpathlineto{\pgfqpoint{5.803297in}{2.481714in}}%
\pgfpathlineto{\pgfqpoint{5.844757in}{2.499039in}}%
\pgfpathlineto{\pgfqpoint{5.856602in}{2.503549in}}%
\pgfpathlineto{\pgfqpoint{5.856602in}{2.503549in}}%
\pgfusepath{stroke}%
\end{pgfscope}%
\begin{pgfscope}%
\pgfpathrectangle{\pgfqpoint{0.573726in}{0.456901in}}{\pgfqpoint{5.578717in}{2.816433in}}%
\pgfusepath{clip}%
\pgfsetrectcap%
\pgfsetroundjoin%
\pgfsetlinewidth{1.003750pt}%
\definecolor{currentstroke}{rgb}{0.750000,0.000000,0.750000}%
\pgfsetstrokecolor{currentstroke}%
\pgfsetdash{}{0pt}%
\pgfpathmoveto{\pgfqpoint{0.563726in}{1.654794in}}%
\pgfpathlineto{\pgfqpoint{0.585360in}{1.664374in}}%
\pgfpathlineto{\pgfqpoint{0.609051in}{1.677023in}}%
\pgfpathlineto{\pgfqpoint{0.638665in}{1.695765in}}%
\pgfpathlineto{\pgfqpoint{0.668279in}{1.717457in}}%
\pgfpathlineto{\pgfqpoint{0.697892in}{1.741722in}}%
\pgfpathlineto{\pgfqpoint{0.733429in}{1.773637in}}%
\pgfpathlineto{\pgfqpoint{0.780811in}{1.819610in}}%
\pgfpathlineto{\pgfqpoint{0.928879in}{1.966586in}}%
\pgfpathlineto{\pgfqpoint{0.964416in}{1.997641in}}%
\pgfpathlineto{\pgfqpoint{0.994030in}{2.021014in}}%
\pgfpathlineto{\pgfqpoint{1.023643in}{2.041670in}}%
\pgfpathlineto{\pgfqpoint{1.053257in}{2.059252in}}%
\pgfpathlineto{\pgfqpoint{1.076948in}{2.070895in}}%
\pgfpathlineto{\pgfqpoint{1.100639in}{2.080242in}}%
\pgfpathlineto{\pgfqpoint{1.124330in}{2.087190in}}%
\pgfpathlineto{\pgfqpoint{1.148021in}{2.091661in}}%
\pgfpathlineto{\pgfqpoint{1.171712in}{2.093605in}}%
\pgfpathlineto{\pgfqpoint{1.195403in}{2.093000in}}%
\pgfpathlineto{\pgfqpoint{1.219094in}{2.089854in}}%
\pgfpathlineto{\pgfqpoint{1.242785in}{2.084201in}}%
\pgfpathlineto{\pgfqpoint{1.266476in}{2.076104in}}%
\pgfpathlineto{\pgfqpoint{1.290167in}{2.065653in}}%
\pgfpathlineto{\pgfqpoint{1.313858in}{2.052966in}}%
\pgfpathlineto{\pgfqpoint{1.343472in}{2.034179in}}%
\pgfpathlineto{\pgfqpoint{1.373085in}{2.012447in}}%
\pgfpathlineto{\pgfqpoint{1.402699in}{1.988149in}}%
\pgfpathlineto{\pgfqpoint{1.438235in}{1.956202in}}%
\pgfpathlineto{\pgfqpoint{1.485617in}{1.910202in}}%
\pgfpathlineto{\pgfqpoint{1.633686in}{1.763262in}}%
\pgfpathlineto{\pgfqpoint{1.669222in}{1.732242in}}%
\pgfpathlineto{\pgfqpoint{1.698836in}{1.708905in}}%
\pgfpathlineto{\pgfqpoint{1.728450in}{1.688290in}}%
\pgfpathlineto{\pgfqpoint{1.758064in}{1.670755in}}%
\pgfpathlineto{\pgfqpoint{1.781755in}{1.659152in}}%
\pgfpathlineto{\pgfqpoint{1.805446in}{1.649847in}}%
\pgfpathlineto{\pgfqpoint{1.829136in}{1.642942in}}%
\pgfpathlineto{\pgfqpoint{1.852827in}{1.638516in}}%
\pgfpathlineto{\pgfqpoint{1.876518in}{1.636618in}}%
\pgfpathlineto{\pgfqpoint{1.900209in}{1.637268in}}%
\pgfpathlineto{\pgfqpoint{1.923900in}{1.640459in}}%
\pgfpathlineto{\pgfqpoint{1.947591in}{1.646157in}}%
\pgfpathlineto{\pgfqpoint{1.971282in}{1.654297in}}%
\pgfpathlineto{\pgfqpoint{1.994973in}{1.664788in}}%
\pgfpathlineto{\pgfqpoint{2.018664in}{1.677514in}}%
\pgfpathlineto{\pgfqpoint{2.048278in}{1.696346in}}%
\pgfpathlineto{\pgfqpoint{2.077892in}{1.718117in}}%
\pgfpathlineto{\pgfqpoint{2.107505in}{1.742449in}}%
\pgfpathlineto{\pgfqpoint{2.143042in}{1.774428in}}%
\pgfpathlineto{\pgfqpoint{2.190424in}{1.820455in}}%
\pgfpathlineto{\pgfqpoint{2.338492in}{1.967358in}}%
\pgfpathlineto{\pgfqpoint{2.374029in}{1.998343in}}%
\pgfpathlineto{\pgfqpoint{2.403643in}{2.021643in}}%
\pgfpathlineto{\pgfqpoint{2.433256in}{2.042217in}}%
\pgfpathlineto{\pgfqpoint{2.462870in}{2.059706in}}%
\pgfpathlineto{\pgfqpoint{2.486561in}{2.071269in}}%
\pgfpathlineto{\pgfqpoint{2.510252in}{2.080532in}}%
\pgfpathlineto{\pgfqpoint{2.533943in}{2.087393in}}%
\pgfpathlineto{\pgfqpoint{2.557634in}{2.091774in}}%
\pgfpathlineto{\pgfqpoint{2.581325in}{2.093627in}}%
\pgfpathlineto{\pgfqpoint{2.605016in}{2.092932in}}%
\pgfpathlineto{\pgfqpoint{2.628707in}{2.089695in}}%
\pgfpathlineto{\pgfqpoint{2.652398in}{2.083953in}}%
\pgfpathlineto{\pgfqpoint{2.676089in}{2.075770in}}%
\pgfpathlineto{\pgfqpoint{2.699780in}{2.065238in}}%
\pgfpathlineto{\pgfqpoint{2.723471in}{2.052473in}}%
\pgfpathlineto{\pgfqpoint{2.753084in}{2.033597in}}%
\pgfpathlineto{\pgfqpoint{2.782698in}{2.011787in}}%
\pgfpathlineto{\pgfqpoint{2.812312in}{1.987421in}}%
\pgfpathlineto{\pgfqpoint{2.847848in}{1.955410in}}%
\pgfpathlineto{\pgfqpoint{2.895230in}{1.909356in}}%
\pgfpathlineto{\pgfqpoint{3.037376in}{1.767921in}}%
\pgfpathlineto{\pgfqpoint{3.072913in}{1.736487in}}%
\pgfpathlineto{\pgfqpoint{3.102526in}{1.712723in}}%
\pgfpathlineto{\pgfqpoint{3.132140in}{1.691615in}}%
\pgfpathlineto{\pgfqpoint{3.161754in}{1.673529in}}%
\pgfpathlineto{\pgfqpoint{3.185445in}{1.661449in}}%
\pgfpathlineto{\pgfqpoint{3.209136in}{1.651642in}}%
\pgfpathlineto{\pgfqpoint{3.232827in}{1.644215in}}%
\pgfpathlineto{\pgfqpoint{3.256518in}{1.639253in}}%
\pgfpathlineto{\pgfqpoint{3.280209in}{1.636810in}}%
\pgfpathlineto{\pgfqpoint{3.303900in}{1.636914in}}%
\pgfpathlineto{\pgfqpoint{3.327591in}{1.639563in}}%
\pgfpathlineto{\pgfqpoint{3.351282in}{1.644728in}}%
\pgfpathlineto{\pgfqpoint{3.374973in}{1.652351in}}%
\pgfpathlineto{\pgfqpoint{3.398663in}{1.662348in}}%
\pgfpathlineto{\pgfqpoint{3.422354in}{1.674606in}}%
\pgfpathlineto{\pgfqpoint{3.451968in}{1.692899in}}%
\pgfpathlineto{\pgfqpoint{3.481582in}{1.714192in}}%
\pgfpathlineto{\pgfqpoint{3.511196in}{1.738114in}}%
\pgfpathlineto{\pgfqpoint{3.546732in}{1.769701in}}%
\pgfpathlineto{\pgfqpoint{3.588191in}{1.809516in}}%
\pgfpathlineto{\pgfqpoint{3.665187in}{1.887412in}}%
\pgfpathlineto{\pgfqpoint{3.724414in}{1.946034in}}%
\pgfpathlineto{\pgfqpoint{3.765874in}{1.983959in}}%
\pgfpathlineto{\pgfqpoint{3.801410in}{2.013287in}}%
\pgfpathlineto{\pgfqpoint{3.831024in}{2.034917in}}%
\pgfpathlineto{\pgfqpoint{3.860637in}{2.053590in}}%
\pgfpathlineto{\pgfqpoint{3.884328in}{2.066179in}}%
\pgfpathlineto{\pgfqpoint{3.908019in}{2.076524in}}%
\pgfpathlineto{\pgfqpoint{3.931710in}{2.084512in}}%
\pgfpathlineto{\pgfqpoint{3.955401in}{2.090052in}}%
\pgfpathlineto{\pgfqpoint{3.979092in}{2.093083in}}%
\pgfpathlineto{\pgfqpoint{4.002783in}{2.093572in}}%
\pgfpathlineto{\pgfqpoint{4.026474in}{2.091512in}}%
\pgfpathlineto{\pgfqpoint{4.050165in}{2.086927in}}%
\pgfpathlineto{\pgfqpoint{4.073856in}{2.079868in}}%
\pgfpathlineto{\pgfqpoint{4.097547in}{2.070413in}}%
\pgfpathlineto{\pgfqpoint{4.121238in}{2.058669in}}%
\pgfpathlineto{\pgfqpoint{4.144929in}{2.044765in}}%
\pgfpathlineto{\pgfqpoint{4.174543in}{2.024588in}}%
\pgfpathlineto{\pgfqpoint{4.204157in}{2.001633in}}%
\pgfpathlineto{\pgfqpoint{4.239693in}{1.970987in}}%
\pgfpathlineto{\pgfqpoint{4.281152in}{1.931936in}}%
\pgfpathlineto{\pgfqpoint{4.340380in}{1.872545in}}%
\pgfpathlineto{\pgfqpoint{4.423298in}{1.789467in}}%
\pgfpathlineto{\pgfqpoint{4.464757in}{1.751100in}}%
\pgfpathlineto{\pgfqpoint{4.500294in}{1.721261in}}%
\pgfpathlineto{\pgfqpoint{4.529907in}{1.699121in}}%
\pgfpathlineto{\pgfqpoint{4.559521in}{1.679872in}}%
\pgfpathlineto{\pgfqpoint{4.583212in}{1.666783in}}%
\pgfpathlineto{\pgfqpoint{4.606903in}{1.655905in}}%
\pgfpathlineto{\pgfqpoint{4.630594in}{1.647361in}}%
\pgfpathlineto{\pgfqpoint{4.654285in}{1.641246in}}%
\pgfpathlineto{\pgfqpoint{4.677976in}{1.637629in}}%
\pgfpathlineto{\pgfqpoint{4.701667in}{1.636548in}}%
\pgfpathlineto{\pgfqpoint{4.725358in}{1.638018in}}%
\pgfpathlineto{\pgfqpoint{4.749049in}{1.642020in}}%
\pgfpathlineto{\pgfqpoint{4.772740in}{1.648511in}}%
\pgfpathlineto{\pgfqpoint{4.796431in}{1.657418in}}%
\pgfpathlineto{\pgfqpoint{4.820122in}{1.668641in}}%
\pgfpathlineto{\pgfqpoint{4.843813in}{1.682057in}}%
\pgfpathlineto{\pgfqpoint{4.873427in}{1.701678in}}%
\pgfpathlineto{\pgfqpoint{4.903040in}{1.724145in}}%
\pgfpathlineto{\pgfqpoint{4.938577in}{1.754311in}}%
\pgfpathlineto{\pgfqpoint{4.980036in}{1.792957in}}%
\pgfpathlineto{\pgfqpoint{5.033341in}{1.846066in}}%
\pgfpathlineto{\pgfqpoint{5.128105in}{1.941174in}}%
\pgfpathlineto{\pgfqpoint{5.169564in}{1.979508in}}%
\pgfpathlineto{\pgfqpoint{5.205100in}{2.009308in}}%
\pgfpathlineto{\pgfqpoint{5.234714in}{2.031410in}}%
\pgfpathlineto{\pgfqpoint{5.264328in}{2.050614in}}%
\pgfpathlineto{\pgfqpoint{5.288019in}{2.063666in}}%
\pgfpathlineto{\pgfqpoint{5.311710in}{2.074502in}}%
\pgfpathlineto{\pgfqpoint{5.335401in}{2.083004in}}%
\pgfpathlineto{\pgfqpoint{5.359092in}{2.089074in}}%
\pgfpathlineto{\pgfqpoint{5.382783in}{2.092647in}}%
\pgfpathlineto{\pgfqpoint{5.406474in}{2.093682in}}%
\pgfpathlineto{\pgfqpoint{5.430164in}{2.092167in}}%
\pgfpathlineto{\pgfqpoint{5.453855in}{2.088120in}}%
\pgfpathlineto{\pgfqpoint{5.477546in}{2.081585in}}%
\pgfpathlineto{\pgfqpoint{5.501237in}{2.072636in}}%
\pgfpathlineto{\pgfqpoint{5.524928in}{2.061372in}}%
\pgfpathlineto{\pgfqpoint{5.548619in}{2.047919in}}%
\pgfpathlineto{\pgfqpoint{5.578233in}{2.028255in}}%
\pgfpathlineto{\pgfqpoint{5.607847in}{2.005749in}}%
\pgfpathlineto{\pgfqpoint{5.643383in}{1.975546in}}%
\pgfpathlineto{\pgfqpoint{5.684842in}{1.936868in}}%
\pgfpathlineto{\pgfqpoint{5.738147in}{1.883739in}}%
\pgfpathlineto{\pgfqpoint{5.832911in}{1.788654in}}%
\pgfpathlineto{\pgfqpoint{5.856602in}{1.766362in}}%
\pgfpathlineto{\pgfqpoint{5.856602in}{1.766362in}}%
\pgfusepath{stroke}%
\end{pgfscope}%
\begin{pgfscope}%
\pgfpathrectangle{\pgfqpoint{0.573726in}{0.456901in}}{\pgfqpoint{5.578717in}{2.816433in}}%
\pgfusepath{clip}%
\pgfsetrectcap%
\pgfsetroundjoin%
\pgfsetlinewidth{1.003750pt}%
\definecolor{currentstroke}{rgb}{0.750000,0.750000,0.000000}%
\pgfsetstrokecolor{currentstroke}%
\pgfsetdash{}{0pt}%
\pgfpathmoveto{\pgfqpoint{0.563726in}{1.738161in}}%
\pgfpathlineto{\pgfqpoint{0.579438in}{1.733002in}}%
\pgfpathlineto{\pgfqpoint{0.597206in}{1.729372in}}%
\pgfpathlineto{\pgfqpoint{0.614974in}{1.728106in}}%
\pgfpathlineto{\pgfqpoint{0.632742in}{1.729227in}}%
\pgfpathlineto{\pgfqpoint{0.650511in}{1.732715in}}%
\pgfpathlineto{\pgfqpoint{0.668279in}{1.738510in}}%
\pgfpathlineto{\pgfqpoint{0.686047in}{1.746510in}}%
\pgfpathlineto{\pgfqpoint{0.703815in}{1.756575in}}%
\pgfpathlineto{\pgfqpoint{0.727506in}{1.772903in}}%
\pgfpathlineto{\pgfqpoint{0.751197in}{1.792083in}}%
\pgfpathlineto{\pgfqpoint{0.780811in}{1.819158in}}%
\pgfpathlineto{\pgfqpoint{0.822270in}{1.860475in}}%
\pgfpathlineto{\pgfqpoint{0.881498in}{1.919275in}}%
\pgfpathlineto{\pgfqpoint{0.911111in}{1.945450in}}%
\pgfpathlineto{\pgfqpoint{0.934802in}{1.963652in}}%
\pgfpathlineto{\pgfqpoint{0.958493in}{1.978807in}}%
\pgfpathlineto{\pgfqpoint{0.976261in}{1.987887in}}%
\pgfpathlineto{\pgfqpoint{0.994030in}{1.994828in}}%
\pgfpathlineto{\pgfqpoint{1.011798in}{1.999510in}}%
\pgfpathlineto{\pgfqpoint{1.029566in}{2.001851in}}%
\pgfpathlineto{\pgfqpoint{1.047334in}{2.001809in}}%
\pgfpathlineto{\pgfqpoint{1.065103in}{1.999387in}}%
\pgfpathlineto{\pgfqpoint{1.082871in}{1.994626in}}%
\pgfpathlineto{\pgfqpoint{1.100639in}{1.987609in}}%
\pgfpathlineto{\pgfqpoint{1.118407in}{1.978458in}}%
\pgfpathlineto{\pgfqpoint{1.142098in}{1.963218in}}%
\pgfpathlineto{\pgfqpoint{1.165789in}{1.944944in}}%
\pgfpathlineto{\pgfqpoint{1.195403in}{1.918702in}}%
\pgfpathlineto{\pgfqpoint{1.230939in}{1.883919in}}%
\pgfpathlineto{\pgfqpoint{1.307935in}{1.807419in}}%
\pgfpathlineto{\pgfqpoint{1.337549in}{1.781676in}}%
\pgfpathlineto{\pgfqpoint{1.361240in}{1.763927in}}%
\pgfpathlineto{\pgfqpoint{1.384931in}{1.749309in}}%
\pgfpathlineto{\pgfqpoint{1.402699in}{1.740675in}}%
\pgfpathlineto{\pgfqpoint{1.420467in}{1.734209in}}%
\pgfpathlineto{\pgfqpoint{1.438235in}{1.730024in}}%
\pgfpathlineto{\pgfqpoint{1.456004in}{1.728191in}}%
\pgfpathlineto{\pgfqpoint{1.473772in}{1.728744in}}%
\pgfpathlineto{\pgfqpoint{1.491540in}{1.731673in}}%
\pgfpathlineto{\pgfqpoint{1.509308in}{1.736926in}}%
\pgfpathlineto{\pgfqpoint{1.527077in}{1.744412in}}%
\pgfpathlineto{\pgfqpoint{1.544845in}{1.754000in}}%
\pgfpathlineto{\pgfqpoint{1.568536in}{1.769762in}}%
\pgfpathlineto{\pgfqpoint{1.592227in}{1.788474in}}%
\pgfpathlineto{\pgfqpoint{1.621840in}{1.815122in}}%
\pgfpathlineto{\pgfqpoint{1.663300in}{1.856172in}}%
\pgfpathlineto{\pgfqpoint{1.722527in}{1.915291in}}%
\pgfpathlineto{\pgfqpoint{1.752141in}{1.941920in}}%
\pgfpathlineto{\pgfqpoint{1.775832in}{1.960610in}}%
\pgfpathlineto{\pgfqpoint{1.799523in}{1.976347in}}%
\pgfpathlineto{\pgfqpoint{1.817291in}{1.985914in}}%
\pgfpathlineto{\pgfqpoint{1.835059in}{1.993377in}}%
\pgfpathlineto{\pgfqpoint{1.852827in}{1.998605in}}%
\pgfpathlineto{\pgfqpoint{1.870596in}{2.001508in}}%
\pgfpathlineto{\pgfqpoint{1.888364in}{2.002036in}}%
\pgfpathlineto{\pgfqpoint{1.906132in}{2.000178in}}%
\pgfpathlineto{\pgfqpoint{1.923900in}{1.995968in}}%
\pgfpathlineto{\pgfqpoint{1.941669in}{1.989478in}}%
\pgfpathlineto{\pgfqpoint{1.959437in}{1.980822in}}%
\pgfpathlineto{\pgfqpoint{1.977205in}{1.970151in}}%
\pgfpathlineto{\pgfqpoint{2.000896in}{1.953113in}}%
\pgfpathlineto{\pgfqpoint{2.024587in}{1.933353in}}%
\pgfpathlineto{\pgfqpoint{2.054201in}{1.905763in}}%
\pgfpathlineto{\pgfqpoint{2.107505in}{1.852114in}}%
\pgfpathlineto{\pgfqpoint{2.148965in}{1.811355in}}%
\pgfpathlineto{\pgfqpoint{2.178578in}{1.785134in}}%
\pgfpathlineto{\pgfqpoint{2.202269in}{1.766882in}}%
\pgfpathlineto{\pgfqpoint{2.225960in}{1.751668in}}%
\pgfpathlineto{\pgfqpoint{2.243729in}{1.742539in}}%
\pgfpathlineto{\pgfqpoint{2.261497in}{1.735545in}}%
\pgfpathlineto{\pgfqpoint{2.279265in}{1.730809in}}%
\pgfpathlineto{\pgfqpoint{2.297033in}{1.728412in}}%
\pgfpathlineto{\pgfqpoint{2.314801in}{1.728396in}}%
\pgfpathlineto{\pgfqpoint{2.332570in}{1.730762in}}%
\pgfpathlineto{\pgfqpoint{2.350338in}{1.735468in}}%
\pgfpathlineto{\pgfqpoint{2.368106in}{1.742433in}}%
\pgfpathlineto{\pgfqpoint{2.385874in}{1.751535in}}%
\pgfpathlineto{\pgfqpoint{2.409565in}{1.766716in}}%
\pgfpathlineto{\pgfqpoint{2.433256in}{1.784941in}}%
\pgfpathlineto{\pgfqpoint{2.462870in}{1.811136in}}%
\pgfpathlineto{\pgfqpoint{2.498406in}{1.845888in}}%
\pgfpathlineto{\pgfqpoint{2.575402in}{1.922424in}}%
\pgfpathlineto{\pgfqpoint{2.605016in}{1.948216in}}%
\pgfpathlineto{\pgfqpoint{2.628707in}{1.966016in}}%
\pgfpathlineto{\pgfqpoint{2.652398in}{1.980695in}}%
\pgfpathlineto{\pgfqpoint{2.670166in}{1.989378in}}%
\pgfpathlineto{\pgfqpoint{2.687934in}{1.995897in}}%
\pgfpathlineto{\pgfqpoint{2.705703in}{2.000138in}}%
\pgfpathlineto{\pgfqpoint{2.723471in}{2.002027in}}%
\pgfpathlineto{\pgfqpoint{2.741239in}{2.001531in}}%
\pgfpathlineto{\pgfqpoint{2.759007in}{1.998659in}}%
\pgfpathlineto{\pgfqpoint{2.776775in}{1.993460in}}%
\pgfpathlineto{\pgfqpoint{2.794544in}{1.986026in}}%
\pgfpathlineto{\pgfqpoint{2.812312in}{1.976486in}}%
\pgfpathlineto{\pgfqpoint{2.836003in}{1.960781in}}%
\pgfpathlineto{\pgfqpoint{2.859694in}{1.942117in}}%
\pgfpathlineto{\pgfqpoint{2.889308in}{1.915513in}}%
\pgfpathlineto{\pgfqpoint{2.930767in}{1.874492in}}%
\pgfpathlineto{\pgfqpoint{2.995917in}{1.809774in}}%
\pgfpathlineto{\pgfqpoint{3.025531in}{1.783741in}}%
\pgfpathlineto{\pgfqpoint{3.049222in}{1.765688in}}%
\pgfpathlineto{\pgfqpoint{3.072913in}{1.750711in}}%
\pgfpathlineto{\pgfqpoint{3.090681in}{1.741779in}}%
\pgfpathlineto{\pgfqpoint{3.108449in}{1.734996in}}%
\pgfpathlineto{\pgfqpoint{3.126217in}{1.730479in}}%
\pgfpathlineto{\pgfqpoint{3.143986in}{1.728307in}}%
\pgfpathlineto{\pgfqpoint{3.161754in}{1.728519in}}%
\pgfpathlineto{\pgfqpoint{3.179522in}{1.731110in}}%
\pgfpathlineto{\pgfqpoint{3.197290in}{1.736036in}}%
\pgfpathlineto{\pgfqpoint{3.215058in}{1.743210in}}%
\pgfpathlineto{\pgfqpoint{3.232827in}{1.752508in}}%
\pgfpathlineto{\pgfqpoint{3.256518in}{1.767923in}}%
\pgfpathlineto{\pgfqpoint{3.280209in}{1.786345in}}%
\pgfpathlineto{\pgfqpoint{3.309822in}{1.812724in}}%
\pgfpathlineto{\pgfqpoint{3.351282in}{1.853594in}}%
\pgfpathlineto{\pgfqpoint{3.416432in}{1.918483in}}%
\pgfpathlineto{\pgfqpoint{3.446045in}{1.944750in}}%
\pgfpathlineto{\pgfqpoint{3.469736in}{1.963052in}}%
\pgfpathlineto{\pgfqpoint{3.493427in}{1.978324in}}%
\pgfpathlineto{\pgfqpoint{3.511196in}{1.987502in}}%
\pgfpathlineto{\pgfqpoint{3.528964in}{1.994548in}}%
\pgfpathlineto{\pgfqpoint{3.546732in}{1.999340in}}%
\pgfpathlineto{\pgfqpoint{3.564500in}{2.001793in}}%
\pgfpathlineto{\pgfqpoint{3.582269in}{2.001866in}}%
\pgfpathlineto{\pgfqpoint{3.600037in}{1.999556in}}%
\pgfpathlineto{\pgfqpoint{3.617805in}{1.994905in}}%
\pgfpathlineto{\pgfqpoint{3.635573in}{1.987992in}}%
\pgfpathlineto{\pgfqpoint{3.653341in}{1.978940in}}%
\pgfpathlineto{\pgfqpoint{3.677032in}{1.963818in}}%
\pgfpathlineto{\pgfqpoint{3.700723in}{1.945642in}}%
\pgfpathlineto{\pgfqpoint{3.730337in}{1.919494in}}%
\pgfpathlineto{\pgfqpoint{3.765874in}{1.884773in}}%
\pgfpathlineto{\pgfqpoint{3.842869in}{1.808202in}}%
\pgfpathlineto{\pgfqpoint{3.872483in}{1.782361in}}%
\pgfpathlineto{\pgfqpoint{3.896174in}{1.764510in}}%
\pgfpathlineto{\pgfqpoint{3.919865in}{1.749772in}}%
\pgfpathlineto{\pgfqpoint{3.937633in}{1.741038in}}%
\pgfpathlineto{\pgfqpoint{3.955401in}{1.734466in}}%
\pgfpathlineto{\pgfqpoint{3.973170in}{1.730170in}}%
\pgfpathlineto{\pgfqpoint{3.990938in}{1.728225in}}%
\pgfpathlineto{\pgfqpoint{4.008706in}{1.728664in}}%
\pgfpathlineto{\pgfqpoint{4.026474in}{1.731480in}}%
\pgfpathlineto{\pgfqpoint{4.044243in}{1.736624in}}%
\pgfpathlineto{\pgfqpoint{4.062011in}{1.744006in}}%
\pgfpathlineto{\pgfqpoint{4.079779in}{1.753498in}}%
\pgfpathlineto{\pgfqpoint{4.103470in}{1.769145in}}%
\pgfpathlineto{\pgfqpoint{4.127161in}{1.787761in}}%
\pgfpathlineto{\pgfqpoint{4.156775in}{1.814321in}}%
\pgfpathlineto{\pgfqpoint{4.198234in}{1.855312in}}%
\pgfpathlineto{\pgfqpoint{4.263384in}{1.920066in}}%
\pgfpathlineto{\pgfqpoint{4.292998in}{1.946146in}}%
\pgfpathlineto{\pgfqpoint{4.316689in}{1.964249in}}%
\pgfpathlineto{\pgfqpoint{4.340380in}{1.979286in}}%
\pgfpathlineto{\pgfqpoint{4.358148in}{1.988267in}}%
\pgfpathlineto{\pgfqpoint{4.375916in}{1.995103in}}%
\pgfpathlineto{\pgfqpoint{4.393684in}{1.999675in}}%
\pgfpathlineto{\pgfqpoint{4.411453in}{2.001903in}}%
\pgfpathlineto{\pgfqpoint{4.429221in}{2.001748in}}%
\pgfpathlineto{\pgfqpoint{4.446989in}{1.999213in}}%
\pgfpathlineto{\pgfqpoint{4.464757in}{1.994342in}}%
\pgfpathlineto{\pgfqpoint{4.482526in}{1.987220in}}%
\pgfpathlineto{\pgfqpoint{4.500294in}{1.977971in}}%
\pgfpathlineto{\pgfqpoint{4.523985in}{1.962615in}}%
\pgfpathlineto{\pgfqpoint{4.547676in}{1.944242in}}%
\pgfpathlineto{\pgfqpoint{4.577289in}{1.917908in}}%
\pgfpathlineto{\pgfqpoint{4.618749in}{1.877070in}}%
\pgfpathlineto{\pgfqpoint{4.683899in}{1.812149in}}%
\pgfpathlineto{\pgfqpoint{4.713513in}{1.785835in}}%
\pgfpathlineto{\pgfqpoint{4.737204in}{1.767484in}}%
\pgfpathlineto{\pgfqpoint{4.760894in}{1.752154in}}%
\pgfpathlineto{\pgfqpoint{4.778663in}{1.742927in}}%
\pgfpathlineto{\pgfqpoint{4.796431in}{1.735828in}}%
\pgfpathlineto{\pgfqpoint{4.814199in}{1.730982in}}%
\pgfpathlineto{\pgfqpoint{4.831967in}{1.728472in}}%
\pgfpathlineto{\pgfqpoint{4.849736in}{1.728343in}}%
\pgfpathlineto{\pgfqpoint{4.867504in}{1.730596in}}%
\pgfpathlineto{\pgfqpoint{4.885272in}{1.735192in}}%
\pgfpathlineto{\pgfqpoint{4.903040in}{1.742052in}}%
\pgfpathlineto{\pgfqpoint{4.920809in}{1.751055in}}%
\pgfpathlineto{\pgfqpoint{4.944500in}{1.766118in}}%
\pgfpathlineto{\pgfqpoint{4.968191in}{1.784244in}}%
\pgfpathlineto{\pgfqpoint{4.997804in}{1.810345in}}%
\pgfpathlineto{\pgfqpoint{5.033341in}{1.845035in}}%
\pgfpathlineto{\pgfqpoint{5.110336in}{1.921640in}}%
\pgfpathlineto{\pgfqpoint{5.139950in}{1.947529in}}%
\pgfpathlineto{\pgfqpoint{5.163641in}{1.965431in}}%
\pgfpathlineto{\pgfqpoint{5.187332in}{1.980230in}}%
\pgfpathlineto{\pgfqpoint{5.205100in}{1.989013in}}%
\pgfpathlineto{\pgfqpoint{5.222868in}{1.995638in}}%
\pgfpathlineto{\pgfqpoint{5.240637in}{1.999989in}}%
\pgfpathlineto{\pgfqpoint{5.258405in}{2.001991in}}%
\pgfpathlineto{\pgfqpoint{5.276173in}{2.001609in}}%
\pgfpathlineto{\pgfqpoint{5.293941in}{1.998849in}}%
\pgfpathlineto{\pgfqpoint{5.311710in}{1.993759in}}%
\pgfpathlineto{\pgfqpoint{5.329478in}{1.986429in}}%
\pgfpathlineto{\pgfqpoint{5.347246in}{1.976985in}}%
\pgfpathlineto{\pgfqpoint{5.370937in}{1.961396in}}%
\pgfpathlineto{\pgfqpoint{5.394628in}{1.942828in}}%
\pgfpathlineto{\pgfqpoint{5.424242in}{1.916313in}}%
\pgfpathlineto{\pgfqpoint{5.465701in}{1.875352in}}%
\pgfpathlineto{\pgfqpoint{5.530851in}{1.810563in}}%
\pgfpathlineto{\pgfqpoint{5.560465in}{1.784436in}}%
\pgfpathlineto{\pgfqpoint{5.584156in}{1.766283in}}%
\pgfpathlineto{\pgfqpoint{5.607847in}{1.751187in}}%
\pgfpathlineto{\pgfqpoint{5.625615in}{1.742157in}}%
\pgfpathlineto{\pgfqpoint{5.643383in}{1.735268in}}%
\pgfpathlineto{\pgfqpoint{5.661151in}{1.730641in}}%
\pgfpathlineto{\pgfqpoint{5.678920in}{1.728357in}}%
\pgfpathlineto{\pgfqpoint{5.696688in}{1.728455in}}%
\pgfpathlineto{\pgfqpoint{5.714456in}{1.730933in}}%
\pgfpathlineto{\pgfqpoint{5.732224in}{1.735749in}}%
\pgfpathlineto{\pgfqpoint{5.749993in}{1.742819in}}%
\pgfpathlineto{\pgfqpoint{5.767761in}{1.752019in}}%
\pgfpathlineto{\pgfqpoint{5.791452in}{1.767317in}}%
\pgfpathlineto{\pgfqpoint{5.815143in}{1.785641in}}%
\pgfpathlineto{\pgfqpoint{5.844757in}{1.811929in}}%
\pgfpathlineto{\pgfqpoint{5.856602in}{1.823239in}}%
\pgfpathlineto{\pgfqpoint{5.856602in}{1.823239in}}%
\pgfusepath{stroke}%
\end{pgfscope}%
\begin{pgfscope}%
\pgfsetrectcap%
\pgfsetmiterjoin%
\pgfsetlinewidth{0.602250pt}%
\definecolor{currentstroke}{rgb}{0.000000,0.000000,0.000000}%
\pgfsetstrokecolor{currentstroke}%
\pgfsetdash{}{0pt}%
\pgfpathmoveto{\pgfqpoint{0.573726in}{0.456901in}}%
\pgfpathlineto{\pgfqpoint{0.573726in}{3.273333in}}%
\pgfusepath{stroke}%
\end{pgfscope}%
\begin{pgfscope}%
\pgfsetrectcap%
\pgfsetmiterjoin%
\pgfsetlinewidth{0.602250pt}%
\definecolor{currentstroke}{rgb}{0.000000,0.000000,0.000000}%
\pgfsetstrokecolor{currentstroke}%
\pgfsetdash{}{0pt}%
\pgfpathmoveto{\pgfqpoint{0.573726in}{0.456901in}}%
\pgfpathlineto{\pgfqpoint{6.152443in}{0.456901in}}%
\pgfusepath{stroke}%
\end{pgfscope}%
\begin{pgfscope}%
\pgfsetbuttcap%
\pgfsetroundjoin%
\pgfsetlinewidth{2.007500pt}%
\definecolor{currentstroke}{rgb}{0.000000,0.000000,0.000000}%
\pgfsetstrokecolor{currentstroke}%
\pgfsetdash{{7.400000pt}{3.200000pt}}{0.000000pt}%
\pgfpathmoveto{\pgfqpoint{0.526569in}{3.565000in}}%
\pgfpathlineto{\pgfqpoint{0.776569in}{3.565000in}}%
\pgfpathlineto{\pgfqpoint{1.026569in}{3.565000in}}%
\pgfusepath{stroke}%
\end{pgfscope}%
\begin{pgfscope}%
\definecolor{textcolor}{rgb}{0.000000,0.000000,0.000000}%
\pgfsetstrokecolor{textcolor}%
\pgfsetfillcolor{textcolor}%
\pgftext[x=1.159903in,y=3.506667in,left,base]{\color{textcolor}{\rmfamily\fontsize{12.000000}{14.400000}\selectfont\catcode`\^=\active\def^{\ifmmode\sp\else\^{}\fi}\catcode`\%=\active\def%{\%}Исходный сигнал}}%
\end{pgfscope}%
\begin{pgfscope}%
\pgfsetrectcap%
\pgfsetroundjoin%
\pgfsetlinewidth{2.007500pt}%
\definecolor{currentstroke}{rgb}{0.835294,0.368627,0.000000}%
\pgfsetstrokecolor{currentstroke}%
\pgfsetdash{}{0pt}%
\pgfpathmoveto{\pgfqpoint{2.737051in}{3.565000in}}%
\pgfpathlineto{\pgfqpoint{2.987051in}{3.565000in}}%
\pgfpathlineto{\pgfqpoint{3.237051in}{3.565000in}}%
\pgfusepath{stroke}%
\end{pgfscope}%
\begin{pgfscope}%
\definecolor{textcolor}{rgb}{0.000000,0.000000,0.000000}%
\pgfsetstrokecolor{textcolor}%
\pgfsetfillcolor{textcolor}%
\pgftext[x=3.370385in,y=3.506667in,left,base]{\color{textcolor}{\rmfamily\fontsize{12.000000}{14.400000}\selectfont\catcode`\^=\active\def^{\ifmmode\sp\else\^{}\fi}\catcode`\%=\active\def%{\%}Входной РФ}}%
\end{pgfscope}%
\begin{pgfscope}%
\pgfsetbuttcap%
\pgfsetroundjoin%
\pgfsetlinewidth{2.007500pt}%
\definecolor{currentstroke}{rgb}{0.000000,0.619608,0.450980}%
\pgfsetstrokecolor{currentstroke}%
\pgfsetdash{{12.800000pt}{3.200000pt}{2.000000pt}{3.200000pt}}{0.000000pt}%
\pgfpathmoveto{\pgfqpoint{4.577904in}{3.565000in}}%
\pgfpathlineto{\pgfqpoint{4.827904in}{3.565000in}}%
\pgfpathlineto{\pgfqpoint{5.077904in}{3.565000in}}%
\pgfusepath{stroke}%
\end{pgfscope}%
\begin{pgfscope}%
\definecolor{textcolor}{rgb}{0.000000,0.000000,0.000000}%
\pgfsetstrokecolor{textcolor}%
\pgfsetfillcolor{textcolor}%
\pgftext[x=5.211237in,y=3.506667in,left,base]{\color{textcolor}{\rmfamily\fontsize{12.000000}{14.400000}\selectfont\catcode`\^=\active\def^{\ifmmode\sp\else\^{}\fi}\catcode`\%=\active\def%{\%}Выходной РФ}}%
\end{pgfscope}%
\end{pgfpicture}%
\makeatother%
\endgroup%

    \caption{Ряд Фурье реакции на первый периодический сигнал}
    \label{fig:plot_fourier_approx_out_1}
\end{figure}

\begin{figure}[H]
    \centering
    %% Creator: Matplotlib, PGF backend
%%
%% To include the figure in your LaTeX document, write
%%   \input{<filename>.pgf}
%%
%% Make sure the required packages are loaded in your preamble
%%   \usepackage{pgf}
%%
%% Also ensure that all the required font packages are loaded; for instance,
%% the lmodern package is sometimes necessary when using math font.
%%   \usepackage{lmodern}
%%
%% Figures using additional raster images can only be included by \input if
%% they are in the same directory as the main LaTeX file. For loading figures
%% from other directories you can use the `import` package
%%   \usepackage{import}
%%
%% and then include the figures with
%%   \import{<path to file>}{<filename>.pgf}
%%
%% Matplotlib used the following preamble
%%   \def\mathdefault#1{#1}
%%   \everymath=\expandafter{\the\everymath\displaystyle}
%%   \IfFileExists{scrextend.sty}{
%%     \usepackage[fontsize=12.000000pt]{scrextend}
%%   }{
%%     \renewcommand{\normalsize}{\fontsize{12.000000}{14.400000}\selectfont}
%%     \normalsize
%%   }
%%   \usepackage{fontspec}\setmainfont{Times New Roman}\usepackage[english,russian]{babel}\usepackage{amsmath}
%%   \ifdefined\pdftexversion\else  % non-pdftex case.
%%     \usepackage{fontspec}
%%     \setmainfont{times.ttf}[Path=\detokenize{/usr/local/share/fonts/truetype/times-new-roman/}]
%%     \setsansfont{DejaVuSans.ttf}[Path=\detokenize{/usr/share/matplotlib/mpl-data/fonts/ttf/}]
%%     \setmonofont{DejaVuSansMono.ttf}[Path=\detokenize{/usr/share/matplotlib/mpl-data/fonts/ttf/}]
%%   \fi
%%   \makeatletter\@ifpackageloaded{underscore}{}{\usepackage[strings]{underscore}}\makeatother
%%
\begingroup%
\makeatletter%
\begin{pgfpicture}%
\pgfpathrectangle{\pgfpointorigin}{\pgfqpoint{6.295433in}{3.740000in}}%
\pgfusepath{use as bounding box, clip}%
\begin{pgfscope}%
\pgfsetbuttcap%
\pgfsetmiterjoin%
\definecolor{currentfill}{rgb}{1.000000,1.000000,1.000000}%
\pgfsetfillcolor{currentfill}%
\pgfsetlinewidth{0.000000pt}%
\definecolor{currentstroke}{rgb}{1.000000,1.000000,1.000000}%
\pgfsetstrokecolor{currentstroke}%
\pgfsetdash{}{0pt}%
\pgfpathmoveto{\pgfqpoint{0.000000in}{0.000000in}}%
\pgfpathlineto{\pgfqpoint{6.295433in}{0.000000in}}%
\pgfpathlineto{\pgfqpoint{6.295433in}{3.740000in}}%
\pgfpathlineto{\pgfqpoint{0.000000in}{3.740000in}}%
\pgfpathlineto{\pgfqpoint{0.000000in}{0.000000in}}%
\pgfpathclose%
\pgfusepath{fill}%
\end{pgfscope}%
\begin{pgfscope}%
\pgfsetbuttcap%
\pgfsetmiterjoin%
\definecolor{currentfill}{rgb}{1.000000,1.000000,1.000000}%
\pgfsetfillcolor{currentfill}%
\pgfsetlinewidth{0.000000pt}%
\definecolor{currentstroke}{rgb}{0.000000,0.000000,0.000000}%
\pgfsetstrokecolor{currentstroke}%
\pgfsetstrokeopacity{0.000000}%
\pgfsetdash{}{0pt}%
\pgfpathmoveto{\pgfqpoint{0.532060in}{0.456901in}}%
\pgfpathlineto{\pgfqpoint{6.152443in}{0.456901in}}%
\pgfpathlineto{\pgfqpoint{6.152443in}{3.273333in}}%
\pgfpathlineto{\pgfqpoint{0.532060in}{3.273333in}}%
\pgfpathlineto{\pgfqpoint{0.532060in}{0.456901in}}%
\pgfpathclose%
\pgfusepath{fill}%
\end{pgfscope}%
\begin{pgfscope}%
\pgfsetbuttcap%
\pgfsetroundjoin%
\definecolor{currentfill}{rgb}{1.000000,1.000000,1.000000}%
\pgfsetfillcolor{currentfill}%
\pgfsetlinewidth{0.803000pt}%
\definecolor{currentstroke}{rgb}{0.000000,0.000000,0.000000}%
\pgfsetstrokecolor{currentstroke}%
\pgfsetdash{}{0pt}%
\pgfsys@defobject{currentmarker}{\pgfqpoint{0.000000in}{0.000000in}}{\pgfqpoint{0.000000in}{0.048611in}}{%
\pgfpathmoveto{\pgfqpoint{0.000000in}{0.000000in}}%
\pgfpathlineto{\pgfqpoint{0.000000in}{0.048611in}}%
\pgfusepath{stroke,fill}%
}%
\begin{pgfscope}%
\pgfsys@transformshift{0.744953in}{0.456901in}%
\pgfsys@useobject{currentmarker}{}%
\end{pgfscope}%
\end{pgfscope}%
\begin{pgfscope}%
\definecolor{textcolor}{rgb}{0.000000,0.000000,0.000000}%
\pgfsetstrokecolor{textcolor}%
\pgfsetfillcolor{textcolor}%
\pgftext[x=0.744953in,y=0.408290in,,top]{\color{textcolor}{\rmfamily\fontsize{12.000000}{14.400000}\selectfont\catcode`\^=\active\def^{\ifmmode\sp\else\^{}\fi}\catcode`\%=\active\def%{\%}$\mathdefault{0}$}}%
\end{pgfscope}%
\begin{pgfscope}%
\pgfsetbuttcap%
\pgfsetroundjoin%
\definecolor{currentfill}{rgb}{1.000000,1.000000,1.000000}%
\pgfsetfillcolor{currentfill}%
\pgfsetlinewidth{0.803000pt}%
\definecolor{currentstroke}{rgb}{0.000000,0.000000,0.000000}%
\pgfsetstrokecolor{currentstroke}%
\pgfsetdash{}{0pt}%
\pgfsys@defobject{currentmarker}{\pgfqpoint{0.000000in}{0.000000in}}{\pgfqpoint{0.000000in}{0.048611in}}{%
\pgfpathmoveto{\pgfqpoint{0.000000in}{0.000000in}}%
\pgfpathlineto{\pgfqpoint{0.000000in}{0.048611in}}%
\pgfusepath{stroke,fill}%
}%
\begin{pgfscope}%
\pgfsys@transformshift{1.648959in}{0.456901in}%
\pgfsys@useobject{currentmarker}{}%
\end{pgfscope}%
\end{pgfscope}%
\begin{pgfscope}%
\definecolor{textcolor}{rgb}{0.000000,0.000000,0.000000}%
\pgfsetstrokecolor{textcolor}%
\pgfsetfillcolor{textcolor}%
\pgftext[x=1.648959in,y=0.408290in,,top]{\color{textcolor}{\rmfamily\fontsize{12.000000}{14.400000}\selectfont\catcode`\^=\active\def^{\ifmmode\sp\else\^{}\fi}\catcode`\%=\active\def%{\%}$\mathdefault{2}$}}%
\end{pgfscope}%
\begin{pgfscope}%
\pgfsetbuttcap%
\pgfsetroundjoin%
\definecolor{currentfill}{rgb}{1.000000,1.000000,1.000000}%
\pgfsetfillcolor{currentfill}%
\pgfsetlinewidth{0.803000pt}%
\definecolor{currentstroke}{rgb}{0.000000,0.000000,0.000000}%
\pgfsetstrokecolor{currentstroke}%
\pgfsetdash{}{0pt}%
\pgfsys@defobject{currentmarker}{\pgfqpoint{0.000000in}{0.000000in}}{\pgfqpoint{0.000000in}{0.048611in}}{%
\pgfpathmoveto{\pgfqpoint{0.000000in}{0.000000in}}%
\pgfpathlineto{\pgfqpoint{0.000000in}{0.048611in}}%
\pgfusepath{stroke,fill}%
}%
\begin{pgfscope}%
\pgfsys@transformshift{2.552964in}{0.456901in}%
\pgfsys@useobject{currentmarker}{}%
\end{pgfscope}%
\end{pgfscope}%
\begin{pgfscope}%
\definecolor{textcolor}{rgb}{0.000000,0.000000,0.000000}%
\pgfsetstrokecolor{textcolor}%
\pgfsetfillcolor{textcolor}%
\pgftext[x=2.552964in,y=0.408290in,,top]{\color{textcolor}{\rmfamily\fontsize{12.000000}{14.400000}\selectfont\catcode`\^=\active\def^{\ifmmode\sp\else\^{}\fi}\catcode`\%=\active\def%{\%}$\mathdefault{4}$}}%
\end{pgfscope}%
\begin{pgfscope}%
\pgfsetbuttcap%
\pgfsetroundjoin%
\definecolor{currentfill}{rgb}{1.000000,1.000000,1.000000}%
\pgfsetfillcolor{currentfill}%
\pgfsetlinewidth{0.803000pt}%
\definecolor{currentstroke}{rgb}{0.000000,0.000000,0.000000}%
\pgfsetstrokecolor{currentstroke}%
\pgfsetdash{}{0pt}%
\pgfsys@defobject{currentmarker}{\pgfqpoint{0.000000in}{0.000000in}}{\pgfqpoint{0.000000in}{0.048611in}}{%
\pgfpathmoveto{\pgfqpoint{0.000000in}{0.000000in}}%
\pgfpathlineto{\pgfqpoint{0.000000in}{0.048611in}}%
\pgfusepath{stroke,fill}%
}%
\begin{pgfscope}%
\pgfsys@transformshift{3.456970in}{0.456901in}%
\pgfsys@useobject{currentmarker}{}%
\end{pgfscope}%
\end{pgfscope}%
\begin{pgfscope}%
\definecolor{textcolor}{rgb}{0.000000,0.000000,0.000000}%
\pgfsetstrokecolor{textcolor}%
\pgfsetfillcolor{textcolor}%
\pgftext[x=3.456970in,y=0.408290in,,top]{\color{textcolor}{\rmfamily\fontsize{12.000000}{14.400000}\selectfont\catcode`\^=\active\def^{\ifmmode\sp\else\^{}\fi}\catcode`\%=\active\def%{\%}$\mathdefault{6}$}}%
\end{pgfscope}%
\begin{pgfscope}%
\pgfsetbuttcap%
\pgfsetroundjoin%
\definecolor{currentfill}{rgb}{1.000000,1.000000,1.000000}%
\pgfsetfillcolor{currentfill}%
\pgfsetlinewidth{0.803000pt}%
\definecolor{currentstroke}{rgb}{0.000000,0.000000,0.000000}%
\pgfsetstrokecolor{currentstroke}%
\pgfsetdash{}{0pt}%
\pgfsys@defobject{currentmarker}{\pgfqpoint{0.000000in}{0.000000in}}{\pgfqpoint{0.000000in}{0.048611in}}{%
\pgfpathmoveto{\pgfqpoint{0.000000in}{0.000000in}}%
\pgfpathlineto{\pgfqpoint{0.000000in}{0.048611in}}%
\pgfusepath{stroke,fill}%
}%
\begin{pgfscope}%
\pgfsys@transformshift{4.360975in}{0.456901in}%
\pgfsys@useobject{currentmarker}{}%
\end{pgfscope}%
\end{pgfscope}%
\begin{pgfscope}%
\definecolor{textcolor}{rgb}{0.000000,0.000000,0.000000}%
\pgfsetstrokecolor{textcolor}%
\pgfsetfillcolor{textcolor}%
\pgftext[x=4.360975in,y=0.408290in,,top]{\color{textcolor}{\rmfamily\fontsize{12.000000}{14.400000}\selectfont\catcode`\^=\active\def^{\ifmmode\sp\else\^{}\fi}\catcode`\%=\active\def%{\%}$\mathdefault{8}$}}%
\end{pgfscope}%
\begin{pgfscope}%
\pgfsetbuttcap%
\pgfsetroundjoin%
\definecolor{currentfill}{rgb}{1.000000,1.000000,1.000000}%
\pgfsetfillcolor{currentfill}%
\pgfsetlinewidth{0.803000pt}%
\definecolor{currentstroke}{rgb}{0.000000,0.000000,0.000000}%
\pgfsetstrokecolor{currentstroke}%
\pgfsetdash{}{0pt}%
\pgfsys@defobject{currentmarker}{\pgfqpoint{0.000000in}{0.000000in}}{\pgfqpoint{0.000000in}{0.048611in}}{%
\pgfpathmoveto{\pgfqpoint{0.000000in}{0.000000in}}%
\pgfpathlineto{\pgfqpoint{0.000000in}{0.048611in}}%
\pgfusepath{stroke,fill}%
}%
\begin{pgfscope}%
\pgfsys@transformshift{5.264981in}{0.456901in}%
\pgfsys@useobject{currentmarker}{}%
\end{pgfscope}%
\end{pgfscope}%
\begin{pgfscope}%
\definecolor{textcolor}{rgb}{0.000000,0.000000,0.000000}%
\pgfsetstrokecolor{textcolor}%
\pgfsetfillcolor{textcolor}%
\pgftext[x=5.264981in,y=0.408290in,,top]{\color{textcolor}{\rmfamily\fontsize{12.000000}{14.400000}\selectfont\catcode`\^=\active\def^{\ifmmode\sp\else\^{}\fi}\catcode`\%=\active\def%{\%}$\mathdefault{10}$}}%
\end{pgfscope}%
\begin{pgfscope}%
\pgfpathrectangle{\pgfqpoint{0.532060in}{0.456901in}}{\pgfqpoint{5.620383in}{2.816433in}}%
\pgfusepath{clip}%
\pgfsetbuttcap%
\pgfsetroundjoin%
\pgfsetlinewidth{0.501875pt}%
\definecolor{currentstroke}{rgb}{0.501961,0.501961,0.501961}%
\pgfsetstrokecolor{currentstroke}%
\pgfsetstrokeopacity{0.400000}%
\pgfsetdash{{1.850000pt}{0.800000pt}}{0.000000pt}%
\pgfpathmoveto{\pgfqpoint{0.970954in}{0.456901in}}%
\pgfpathlineto{\pgfqpoint{0.970954in}{3.273333in}}%
\pgfusepath{stroke}%
\end{pgfscope}%
\begin{pgfscope}%
\pgfsetbuttcap%
\pgfsetroundjoin%
\definecolor{currentfill}{rgb}{1.000000,1.000000,1.000000}%
\pgfsetfillcolor{currentfill}%
\pgfsetlinewidth{0.602250pt}%
\definecolor{currentstroke}{rgb}{0.000000,0.000000,0.000000}%
\pgfsetstrokecolor{currentstroke}%
\pgfsetdash{}{0pt}%
\pgfsys@defobject{currentmarker}{\pgfqpoint{0.000000in}{0.000000in}}{\pgfqpoint{0.000000in}{0.027778in}}{%
\pgfpathmoveto{\pgfqpoint{0.000000in}{0.000000in}}%
\pgfpathlineto{\pgfqpoint{0.000000in}{0.027778in}}%
\pgfusepath{stroke,fill}%
}%
\begin{pgfscope}%
\pgfsys@transformshift{0.970954in}{0.456901in}%
\pgfsys@useobject{currentmarker}{}%
\end{pgfscope}%
\end{pgfscope}%
\begin{pgfscope}%
\pgfpathrectangle{\pgfqpoint{0.532060in}{0.456901in}}{\pgfqpoint{5.620383in}{2.816433in}}%
\pgfusepath{clip}%
\pgfsetbuttcap%
\pgfsetroundjoin%
\pgfsetlinewidth{0.501875pt}%
\definecolor{currentstroke}{rgb}{0.501961,0.501961,0.501961}%
\pgfsetstrokecolor{currentstroke}%
\pgfsetstrokeopacity{0.400000}%
\pgfsetdash{{1.850000pt}{0.800000pt}}{0.000000pt}%
\pgfpathmoveto{\pgfqpoint{1.196956in}{0.456901in}}%
\pgfpathlineto{\pgfqpoint{1.196956in}{3.273333in}}%
\pgfusepath{stroke}%
\end{pgfscope}%
\begin{pgfscope}%
\pgfsetbuttcap%
\pgfsetroundjoin%
\definecolor{currentfill}{rgb}{1.000000,1.000000,1.000000}%
\pgfsetfillcolor{currentfill}%
\pgfsetlinewidth{0.602250pt}%
\definecolor{currentstroke}{rgb}{0.000000,0.000000,0.000000}%
\pgfsetstrokecolor{currentstroke}%
\pgfsetdash{}{0pt}%
\pgfsys@defobject{currentmarker}{\pgfqpoint{0.000000in}{0.000000in}}{\pgfqpoint{0.000000in}{0.027778in}}{%
\pgfpathmoveto{\pgfqpoint{0.000000in}{0.000000in}}%
\pgfpathlineto{\pgfqpoint{0.000000in}{0.027778in}}%
\pgfusepath{stroke,fill}%
}%
\begin{pgfscope}%
\pgfsys@transformshift{1.196956in}{0.456901in}%
\pgfsys@useobject{currentmarker}{}%
\end{pgfscope}%
\end{pgfscope}%
\begin{pgfscope}%
\pgfpathrectangle{\pgfqpoint{0.532060in}{0.456901in}}{\pgfqpoint{5.620383in}{2.816433in}}%
\pgfusepath{clip}%
\pgfsetbuttcap%
\pgfsetroundjoin%
\pgfsetlinewidth{0.501875pt}%
\definecolor{currentstroke}{rgb}{0.501961,0.501961,0.501961}%
\pgfsetstrokecolor{currentstroke}%
\pgfsetstrokeopacity{0.400000}%
\pgfsetdash{{1.850000pt}{0.800000pt}}{0.000000pt}%
\pgfpathmoveto{\pgfqpoint{1.422957in}{0.456901in}}%
\pgfpathlineto{\pgfqpoint{1.422957in}{3.273333in}}%
\pgfusepath{stroke}%
\end{pgfscope}%
\begin{pgfscope}%
\pgfsetbuttcap%
\pgfsetroundjoin%
\definecolor{currentfill}{rgb}{1.000000,1.000000,1.000000}%
\pgfsetfillcolor{currentfill}%
\pgfsetlinewidth{0.602250pt}%
\definecolor{currentstroke}{rgb}{0.000000,0.000000,0.000000}%
\pgfsetstrokecolor{currentstroke}%
\pgfsetdash{}{0pt}%
\pgfsys@defobject{currentmarker}{\pgfqpoint{0.000000in}{0.000000in}}{\pgfqpoint{0.000000in}{0.027778in}}{%
\pgfpathmoveto{\pgfqpoint{0.000000in}{0.000000in}}%
\pgfpathlineto{\pgfqpoint{0.000000in}{0.027778in}}%
\pgfusepath{stroke,fill}%
}%
\begin{pgfscope}%
\pgfsys@transformshift{1.422957in}{0.456901in}%
\pgfsys@useobject{currentmarker}{}%
\end{pgfscope}%
\end{pgfscope}%
\begin{pgfscope}%
\pgfpathrectangle{\pgfqpoint{0.532060in}{0.456901in}}{\pgfqpoint{5.620383in}{2.816433in}}%
\pgfusepath{clip}%
\pgfsetbuttcap%
\pgfsetroundjoin%
\pgfsetlinewidth{0.501875pt}%
\definecolor{currentstroke}{rgb}{0.501961,0.501961,0.501961}%
\pgfsetstrokecolor{currentstroke}%
\pgfsetstrokeopacity{0.400000}%
\pgfsetdash{{1.850000pt}{0.800000pt}}{0.000000pt}%
\pgfpathmoveto{\pgfqpoint{1.874960in}{0.456901in}}%
\pgfpathlineto{\pgfqpoint{1.874960in}{3.273333in}}%
\pgfusepath{stroke}%
\end{pgfscope}%
\begin{pgfscope}%
\pgfsetbuttcap%
\pgfsetroundjoin%
\definecolor{currentfill}{rgb}{1.000000,1.000000,1.000000}%
\pgfsetfillcolor{currentfill}%
\pgfsetlinewidth{0.602250pt}%
\definecolor{currentstroke}{rgb}{0.000000,0.000000,0.000000}%
\pgfsetstrokecolor{currentstroke}%
\pgfsetdash{}{0pt}%
\pgfsys@defobject{currentmarker}{\pgfqpoint{0.000000in}{0.000000in}}{\pgfqpoint{0.000000in}{0.027778in}}{%
\pgfpathmoveto{\pgfqpoint{0.000000in}{0.000000in}}%
\pgfpathlineto{\pgfqpoint{0.000000in}{0.027778in}}%
\pgfusepath{stroke,fill}%
}%
\begin{pgfscope}%
\pgfsys@transformshift{1.874960in}{0.456901in}%
\pgfsys@useobject{currentmarker}{}%
\end{pgfscope}%
\end{pgfscope}%
\begin{pgfscope}%
\pgfpathrectangle{\pgfqpoint{0.532060in}{0.456901in}}{\pgfqpoint{5.620383in}{2.816433in}}%
\pgfusepath{clip}%
\pgfsetbuttcap%
\pgfsetroundjoin%
\pgfsetlinewidth{0.501875pt}%
\definecolor{currentstroke}{rgb}{0.501961,0.501961,0.501961}%
\pgfsetstrokecolor{currentstroke}%
\pgfsetstrokeopacity{0.400000}%
\pgfsetdash{{1.850000pt}{0.800000pt}}{0.000000pt}%
\pgfpathmoveto{\pgfqpoint{2.100961in}{0.456901in}}%
\pgfpathlineto{\pgfqpoint{2.100961in}{3.273333in}}%
\pgfusepath{stroke}%
\end{pgfscope}%
\begin{pgfscope}%
\pgfsetbuttcap%
\pgfsetroundjoin%
\definecolor{currentfill}{rgb}{1.000000,1.000000,1.000000}%
\pgfsetfillcolor{currentfill}%
\pgfsetlinewidth{0.602250pt}%
\definecolor{currentstroke}{rgb}{0.000000,0.000000,0.000000}%
\pgfsetstrokecolor{currentstroke}%
\pgfsetdash{}{0pt}%
\pgfsys@defobject{currentmarker}{\pgfqpoint{0.000000in}{0.000000in}}{\pgfqpoint{0.000000in}{0.027778in}}{%
\pgfpathmoveto{\pgfqpoint{0.000000in}{0.000000in}}%
\pgfpathlineto{\pgfqpoint{0.000000in}{0.027778in}}%
\pgfusepath{stroke,fill}%
}%
\begin{pgfscope}%
\pgfsys@transformshift{2.100961in}{0.456901in}%
\pgfsys@useobject{currentmarker}{}%
\end{pgfscope}%
\end{pgfscope}%
\begin{pgfscope}%
\pgfpathrectangle{\pgfqpoint{0.532060in}{0.456901in}}{\pgfqpoint{5.620383in}{2.816433in}}%
\pgfusepath{clip}%
\pgfsetbuttcap%
\pgfsetroundjoin%
\pgfsetlinewidth{0.501875pt}%
\definecolor{currentstroke}{rgb}{0.501961,0.501961,0.501961}%
\pgfsetstrokecolor{currentstroke}%
\pgfsetstrokeopacity{0.400000}%
\pgfsetdash{{1.850000pt}{0.800000pt}}{0.000000pt}%
\pgfpathmoveto{\pgfqpoint{2.326963in}{0.456901in}}%
\pgfpathlineto{\pgfqpoint{2.326963in}{3.273333in}}%
\pgfusepath{stroke}%
\end{pgfscope}%
\begin{pgfscope}%
\pgfsetbuttcap%
\pgfsetroundjoin%
\definecolor{currentfill}{rgb}{1.000000,1.000000,1.000000}%
\pgfsetfillcolor{currentfill}%
\pgfsetlinewidth{0.602250pt}%
\definecolor{currentstroke}{rgb}{0.000000,0.000000,0.000000}%
\pgfsetstrokecolor{currentstroke}%
\pgfsetdash{}{0pt}%
\pgfsys@defobject{currentmarker}{\pgfqpoint{0.000000in}{0.000000in}}{\pgfqpoint{0.000000in}{0.027778in}}{%
\pgfpathmoveto{\pgfqpoint{0.000000in}{0.000000in}}%
\pgfpathlineto{\pgfqpoint{0.000000in}{0.027778in}}%
\pgfusepath{stroke,fill}%
}%
\begin{pgfscope}%
\pgfsys@transformshift{2.326963in}{0.456901in}%
\pgfsys@useobject{currentmarker}{}%
\end{pgfscope}%
\end{pgfscope}%
\begin{pgfscope}%
\pgfpathrectangle{\pgfqpoint{0.532060in}{0.456901in}}{\pgfqpoint{5.620383in}{2.816433in}}%
\pgfusepath{clip}%
\pgfsetbuttcap%
\pgfsetroundjoin%
\pgfsetlinewidth{0.501875pt}%
\definecolor{currentstroke}{rgb}{0.501961,0.501961,0.501961}%
\pgfsetstrokecolor{currentstroke}%
\pgfsetstrokeopacity{0.400000}%
\pgfsetdash{{1.850000pt}{0.800000pt}}{0.000000pt}%
\pgfpathmoveto{\pgfqpoint{2.778966in}{0.456901in}}%
\pgfpathlineto{\pgfqpoint{2.778966in}{3.273333in}}%
\pgfusepath{stroke}%
\end{pgfscope}%
\begin{pgfscope}%
\pgfsetbuttcap%
\pgfsetroundjoin%
\definecolor{currentfill}{rgb}{1.000000,1.000000,1.000000}%
\pgfsetfillcolor{currentfill}%
\pgfsetlinewidth{0.602250pt}%
\definecolor{currentstroke}{rgb}{0.000000,0.000000,0.000000}%
\pgfsetstrokecolor{currentstroke}%
\pgfsetdash{}{0pt}%
\pgfsys@defobject{currentmarker}{\pgfqpoint{0.000000in}{0.000000in}}{\pgfqpoint{0.000000in}{0.027778in}}{%
\pgfpathmoveto{\pgfqpoint{0.000000in}{0.000000in}}%
\pgfpathlineto{\pgfqpoint{0.000000in}{0.027778in}}%
\pgfusepath{stroke,fill}%
}%
\begin{pgfscope}%
\pgfsys@transformshift{2.778966in}{0.456901in}%
\pgfsys@useobject{currentmarker}{}%
\end{pgfscope}%
\end{pgfscope}%
\begin{pgfscope}%
\pgfpathrectangle{\pgfqpoint{0.532060in}{0.456901in}}{\pgfqpoint{5.620383in}{2.816433in}}%
\pgfusepath{clip}%
\pgfsetbuttcap%
\pgfsetroundjoin%
\pgfsetlinewidth{0.501875pt}%
\definecolor{currentstroke}{rgb}{0.501961,0.501961,0.501961}%
\pgfsetstrokecolor{currentstroke}%
\pgfsetstrokeopacity{0.400000}%
\pgfsetdash{{1.850000pt}{0.800000pt}}{0.000000pt}%
\pgfpathmoveto{\pgfqpoint{3.004967in}{0.456901in}}%
\pgfpathlineto{\pgfqpoint{3.004967in}{3.273333in}}%
\pgfusepath{stroke}%
\end{pgfscope}%
\begin{pgfscope}%
\pgfsetbuttcap%
\pgfsetroundjoin%
\definecolor{currentfill}{rgb}{1.000000,1.000000,1.000000}%
\pgfsetfillcolor{currentfill}%
\pgfsetlinewidth{0.602250pt}%
\definecolor{currentstroke}{rgb}{0.000000,0.000000,0.000000}%
\pgfsetstrokecolor{currentstroke}%
\pgfsetdash{}{0pt}%
\pgfsys@defobject{currentmarker}{\pgfqpoint{0.000000in}{0.000000in}}{\pgfqpoint{0.000000in}{0.027778in}}{%
\pgfpathmoveto{\pgfqpoint{0.000000in}{0.000000in}}%
\pgfpathlineto{\pgfqpoint{0.000000in}{0.027778in}}%
\pgfusepath{stroke,fill}%
}%
\begin{pgfscope}%
\pgfsys@transformshift{3.004967in}{0.456901in}%
\pgfsys@useobject{currentmarker}{}%
\end{pgfscope}%
\end{pgfscope}%
\begin{pgfscope}%
\pgfpathrectangle{\pgfqpoint{0.532060in}{0.456901in}}{\pgfqpoint{5.620383in}{2.816433in}}%
\pgfusepath{clip}%
\pgfsetbuttcap%
\pgfsetroundjoin%
\pgfsetlinewidth{0.501875pt}%
\definecolor{currentstroke}{rgb}{0.501961,0.501961,0.501961}%
\pgfsetstrokecolor{currentstroke}%
\pgfsetstrokeopacity{0.400000}%
\pgfsetdash{{1.850000pt}{0.800000pt}}{0.000000pt}%
\pgfpathmoveto{\pgfqpoint{3.230968in}{0.456901in}}%
\pgfpathlineto{\pgfqpoint{3.230968in}{3.273333in}}%
\pgfusepath{stroke}%
\end{pgfscope}%
\begin{pgfscope}%
\pgfsetbuttcap%
\pgfsetroundjoin%
\definecolor{currentfill}{rgb}{1.000000,1.000000,1.000000}%
\pgfsetfillcolor{currentfill}%
\pgfsetlinewidth{0.602250pt}%
\definecolor{currentstroke}{rgb}{0.000000,0.000000,0.000000}%
\pgfsetstrokecolor{currentstroke}%
\pgfsetdash{}{0pt}%
\pgfsys@defobject{currentmarker}{\pgfqpoint{0.000000in}{0.000000in}}{\pgfqpoint{0.000000in}{0.027778in}}{%
\pgfpathmoveto{\pgfqpoint{0.000000in}{0.000000in}}%
\pgfpathlineto{\pgfqpoint{0.000000in}{0.027778in}}%
\pgfusepath{stroke,fill}%
}%
\begin{pgfscope}%
\pgfsys@transformshift{3.230968in}{0.456901in}%
\pgfsys@useobject{currentmarker}{}%
\end{pgfscope}%
\end{pgfscope}%
\begin{pgfscope}%
\pgfpathrectangle{\pgfqpoint{0.532060in}{0.456901in}}{\pgfqpoint{5.620383in}{2.816433in}}%
\pgfusepath{clip}%
\pgfsetbuttcap%
\pgfsetroundjoin%
\pgfsetlinewidth{0.501875pt}%
\definecolor{currentstroke}{rgb}{0.501961,0.501961,0.501961}%
\pgfsetstrokecolor{currentstroke}%
\pgfsetstrokeopacity{0.400000}%
\pgfsetdash{{1.850000pt}{0.800000pt}}{0.000000pt}%
\pgfpathmoveto{\pgfqpoint{3.682971in}{0.456901in}}%
\pgfpathlineto{\pgfqpoint{3.682971in}{3.273333in}}%
\pgfusepath{stroke}%
\end{pgfscope}%
\begin{pgfscope}%
\pgfsetbuttcap%
\pgfsetroundjoin%
\definecolor{currentfill}{rgb}{1.000000,1.000000,1.000000}%
\pgfsetfillcolor{currentfill}%
\pgfsetlinewidth{0.602250pt}%
\definecolor{currentstroke}{rgb}{0.000000,0.000000,0.000000}%
\pgfsetstrokecolor{currentstroke}%
\pgfsetdash{}{0pt}%
\pgfsys@defobject{currentmarker}{\pgfqpoint{0.000000in}{0.000000in}}{\pgfqpoint{0.000000in}{0.027778in}}{%
\pgfpathmoveto{\pgfqpoint{0.000000in}{0.000000in}}%
\pgfpathlineto{\pgfqpoint{0.000000in}{0.027778in}}%
\pgfusepath{stroke,fill}%
}%
\begin{pgfscope}%
\pgfsys@transformshift{3.682971in}{0.456901in}%
\pgfsys@useobject{currentmarker}{}%
\end{pgfscope}%
\end{pgfscope}%
\begin{pgfscope}%
\pgfpathrectangle{\pgfqpoint{0.532060in}{0.456901in}}{\pgfqpoint{5.620383in}{2.816433in}}%
\pgfusepath{clip}%
\pgfsetbuttcap%
\pgfsetroundjoin%
\pgfsetlinewidth{0.501875pt}%
\definecolor{currentstroke}{rgb}{0.501961,0.501961,0.501961}%
\pgfsetstrokecolor{currentstroke}%
\pgfsetstrokeopacity{0.400000}%
\pgfsetdash{{1.850000pt}{0.800000pt}}{0.000000pt}%
\pgfpathmoveto{\pgfqpoint{3.908972in}{0.456901in}}%
\pgfpathlineto{\pgfqpoint{3.908972in}{3.273333in}}%
\pgfusepath{stroke}%
\end{pgfscope}%
\begin{pgfscope}%
\pgfsetbuttcap%
\pgfsetroundjoin%
\definecolor{currentfill}{rgb}{1.000000,1.000000,1.000000}%
\pgfsetfillcolor{currentfill}%
\pgfsetlinewidth{0.602250pt}%
\definecolor{currentstroke}{rgb}{0.000000,0.000000,0.000000}%
\pgfsetstrokecolor{currentstroke}%
\pgfsetdash{}{0pt}%
\pgfsys@defobject{currentmarker}{\pgfqpoint{0.000000in}{0.000000in}}{\pgfqpoint{0.000000in}{0.027778in}}{%
\pgfpathmoveto{\pgfqpoint{0.000000in}{0.000000in}}%
\pgfpathlineto{\pgfqpoint{0.000000in}{0.027778in}}%
\pgfusepath{stroke,fill}%
}%
\begin{pgfscope}%
\pgfsys@transformshift{3.908972in}{0.456901in}%
\pgfsys@useobject{currentmarker}{}%
\end{pgfscope}%
\end{pgfscope}%
\begin{pgfscope}%
\pgfpathrectangle{\pgfqpoint{0.532060in}{0.456901in}}{\pgfqpoint{5.620383in}{2.816433in}}%
\pgfusepath{clip}%
\pgfsetbuttcap%
\pgfsetroundjoin%
\pgfsetlinewidth{0.501875pt}%
\definecolor{currentstroke}{rgb}{0.501961,0.501961,0.501961}%
\pgfsetstrokecolor{currentstroke}%
\pgfsetstrokeopacity{0.400000}%
\pgfsetdash{{1.850000pt}{0.800000pt}}{0.000000pt}%
\pgfpathmoveto{\pgfqpoint{4.134974in}{0.456901in}}%
\pgfpathlineto{\pgfqpoint{4.134974in}{3.273333in}}%
\pgfusepath{stroke}%
\end{pgfscope}%
\begin{pgfscope}%
\pgfsetbuttcap%
\pgfsetroundjoin%
\definecolor{currentfill}{rgb}{1.000000,1.000000,1.000000}%
\pgfsetfillcolor{currentfill}%
\pgfsetlinewidth{0.602250pt}%
\definecolor{currentstroke}{rgb}{0.000000,0.000000,0.000000}%
\pgfsetstrokecolor{currentstroke}%
\pgfsetdash{}{0pt}%
\pgfsys@defobject{currentmarker}{\pgfqpoint{0.000000in}{0.000000in}}{\pgfqpoint{0.000000in}{0.027778in}}{%
\pgfpathmoveto{\pgfqpoint{0.000000in}{0.000000in}}%
\pgfpathlineto{\pgfqpoint{0.000000in}{0.027778in}}%
\pgfusepath{stroke,fill}%
}%
\begin{pgfscope}%
\pgfsys@transformshift{4.134974in}{0.456901in}%
\pgfsys@useobject{currentmarker}{}%
\end{pgfscope}%
\end{pgfscope}%
\begin{pgfscope}%
\pgfpathrectangle{\pgfqpoint{0.532060in}{0.456901in}}{\pgfqpoint{5.620383in}{2.816433in}}%
\pgfusepath{clip}%
\pgfsetbuttcap%
\pgfsetroundjoin%
\pgfsetlinewidth{0.501875pt}%
\definecolor{currentstroke}{rgb}{0.501961,0.501961,0.501961}%
\pgfsetstrokecolor{currentstroke}%
\pgfsetstrokeopacity{0.400000}%
\pgfsetdash{{1.850000pt}{0.800000pt}}{0.000000pt}%
\pgfpathmoveto{\pgfqpoint{4.586977in}{0.456901in}}%
\pgfpathlineto{\pgfqpoint{4.586977in}{3.273333in}}%
\pgfusepath{stroke}%
\end{pgfscope}%
\begin{pgfscope}%
\pgfsetbuttcap%
\pgfsetroundjoin%
\definecolor{currentfill}{rgb}{1.000000,1.000000,1.000000}%
\pgfsetfillcolor{currentfill}%
\pgfsetlinewidth{0.602250pt}%
\definecolor{currentstroke}{rgb}{0.000000,0.000000,0.000000}%
\pgfsetstrokecolor{currentstroke}%
\pgfsetdash{}{0pt}%
\pgfsys@defobject{currentmarker}{\pgfqpoint{0.000000in}{0.000000in}}{\pgfqpoint{0.000000in}{0.027778in}}{%
\pgfpathmoveto{\pgfqpoint{0.000000in}{0.000000in}}%
\pgfpathlineto{\pgfqpoint{0.000000in}{0.027778in}}%
\pgfusepath{stroke,fill}%
}%
\begin{pgfscope}%
\pgfsys@transformshift{4.586977in}{0.456901in}%
\pgfsys@useobject{currentmarker}{}%
\end{pgfscope}%
\end{pgfscope}%
\begin{pgfscope}%
\pgfpathrectangle{\pgfqpoint{0.532060in}{0.456901in}}{\pgfqpoint{5.620383in}{2.816433in}}%
\pgfusepath{clip}%
\pgfsetbuttcap%
\pgfsetroundjoin%
\pgfsetlinewidth{0.501875pt}%
\definecolor{currentstroke}{rgb}{0.501961,0.501961,0.501961}%
\pgfsetstrokecolor{currentstroke}%
\pgfsetstrokeopacity{0.400000}%
\pgfsetdash{{1.850000pt}{0.800000pt}}{0.000000pt}%
\pgfpathmoveto{\pgfqpoint{4.812978in}{0.456901in}}%
\pgfpathlineto{\pgfqpoint{4.812978in}{3.273333in}}%
\pgfusepath{stroke}%
\end{pgfscope}%
\begin{pgfscope}%
\pgfsetbuttcap%
\pgfsetroundjoin%
\definecolor{currentfill}{rgb}{1.000000,1.000000,1.000000}%
\pgfsetfillcolor{currentfill}%
\pgfsetlinewidth{0.602250pt}%
\definecolor{currentstroke}{rgb}{0.000000,0.000000,0.000000}%
\pgfsetstrokecolor{currentstroke}%
\pgfsetdash{}{0pt}%
\pgfsys@defobject{currentmarker}{\pgfqpoint{0.000000in}{0.000000in}}{\pgfqpoint{0.000000in}{0.027778in}}{%
\pgfpathmoveto{\pgfqpoint{0.000000in}{0.000000in}}%
\pgfpathlineto{\pgfqpoint{0.000000in}{0.027778in}}%
\pgfusepath{stroke,fill}%
}%
\begin{pgfscope}%
\pgfsys@transformshift{4.812978in}{0.456901in}%
\pgfsys@useobject{currentmarker}{}%
\end{pgfscope}%
\end{pgfscope}%
\begin{pgfscope}%
\pgfpathrectangle{\pgfqpoint{0.532060in}{0.456901in}}{\pgfqpoint{5.620383in}{2.816433in}}%
\pgfusepath{clip}%
\pgfsetbuttcap%
\pgfsetroundjoin%
\pgfsetlinewidth{0.501875pt}%
\definecolor{currentstroke}{rgb}{0.501961,0.501961,0.501961}%
\pgfsetstrokecolor{currentstroke}%
\pgfsetstrokeopacity{0.400000}%
\pgfsetdash{{1.850000pt}{0.800000pt}}{0.000000pt}%
\pgfpathmoveto{\pgfqpoint{5.038979in}{0.456901in}}%
\pgfpathlineto{\pgfqpoint{5.038979in}{3.273333in}}%
\pgfusepath{stroke}%
\end{pgfscope}%
\begin{pgfscope}%
\pgfsetbuttcap%
\pgfsetroundjoin%
\definecolor{currentfill}{rgb}{1.000000,1.000000,1.000000}%
\pgfsetfillcolor{currentfill}%
\pgfsetlinewidth{0.602250pt}%
\definecolor{currentstroke}{rgb}{0.000000,0.000000,0.000000}%
\pgfsetstrokecolor{currentstroke}%
\pgfsetdash{}{0pt}%
\pgfsys@defobject{currentmarker}{\pgfqpoint{0.000000in}{0.000000in}}{\pgfqpoint{0.000000in}{0.027778in}}{%
\pgfpathmoveto{\pgfqpoint{0.000000in}{0.000000in}}%
\pgfpathlineto{\pgfqpoint{0.000000in}{0.027778in}}%
\pgfusepath{stroke,fill}%
}%
\begin{pgfscope}%
\pgfsys@transformshift{5.038979in}{0.456901in}%
\pgfsys@useobject{currentmarker}{}%
\end{pgfscope}%
\end{pgfscope}%
\begin{pgfscope}%
\pgfpathrectangle{\pgfqpoint{0.532060in}{0.456901in}}{\pgfqpoint{5.620383in}{2.816433in}}%
\pgfusepath{clip}%
\pgfsetbuttcap%
\pgfsetroundjoin%
\pgfsetlinewidth{0.501875pt}%
\definecolor{currentstroke}{rgb}{0.501961,0.501961,0.501961}%
\pgfsetstrokecolor{currentstroke}%
\pgfsetstrokeopacity{0.400000}%
\pgfsetdash{{1.850000pt}{0.800000pt}}{0.000000pt}%
\pgfpathmoveto{\pgfqpoint{5.490982in}{0.456901in}}%
\pgfpathlineto{\pgfqpoint{5.490982in}{3.273333in}}%
\pgfusepath{stroke}%
\end{pgfscope}%
\begin{pgfscope}%
\pgfsetbuttcap%
\pgfsetroundjoin%
\definecolor{currentfill}{rgb}{1.000000,1.000000,1.000000}%
\pgfsetfillcolor{currentfill}%
\pgfsetlinewidth{0.602250pt}%
\definecolor{currentstroke}{rgb}{0.000000,0.000000,0.000000}%
\pgfsetstrokecolor{currentstroke}%
\pgfsetdash{}{0pt}%
\pgfsys@defobject{currentmarker}{\pgfqpoint{0.000000in}{0.000000in}}{\pgfqpoint{0.000000in}{0.027778in}}{%
\pgfpathmoveto{\pgfqpoint{0.000000in}{0.000000in}}%
\pgfpathlineto{\pgfqpoint{0.000000in}{0.027778in}}%
\pgfusepath{stroke,fill}%
}%
\begin{pgfscope}%
\pgfsys@transformshift{5.490982in}{0.456901in}%
\pgfsys@useobject{currentmarker}{}%
\end{pgfscope}%
\end{pgfscope}%
\begin{pgfscope}%
\pgfpathrectangle{\pgfqpoint{0.532060in}{0.456901in}}{\pgfqpoint{5.620383in}{2.816433in}}%
\pgfusepath{clip}%
\pgfsetbuttcap%
\pgfsetroundjoin%
\pgfsetlinewidth{0.501875pt}%
\definecolor{currentstroke}{rgb}{0.501961,0.501961,0.501961}%
\pgfsetstrokecolor{currentstroke}%
\pgfsetstrokeopacity{0.400000}%
\pgfsetdash{{1.850000pt}{0.800000pt}}{0.000000pt}%
\pgfpathmoveto{\pgfqpoint{5.716984in}{0.456901in}}%
\pgfpathlineto{\pgfqpoint{5.716984in}{3.273333in}}%
\pgfusepath{stroke}%
\end{pgfscope}%
\begin{pgfscope}%
\pgfsetbuttcap%
\pgfsetroundjoin%
\definecolor{currentfill}{rgb}{1.000000,1.000000,1.000000}%
\pgfsetfillcolor{currentfill}%
\pgfsetlinewidth{0.602250pt}%
\definecolor{currentstroke}{rgb}{0.000000,0.000000,0.000000}%
\pgfsetstrokecolor{currentstroke}%
\pgfsetdash{}{0pt}%
\pgfsys@defobject{currentmarker}{\pgfqpoint{0.000000in}{0.000000in}}{\pgfqpoint{0.000000in}{0.027778in}}{%
\pgfpathmoveto{\pgfqpoint{0.000000in}{0.000000in}}%
\pgfpathlineto{\pgfqpoint{0.000000in}{0.027778in}}%
\pgfusepath{stroke,fill}%
}%
\begin{pgfscope}%
\pgfsys@transformshift{5.716984in}{0.456901in}%
\pgfsys@useobject{currentmarker}{}%
\end{pgfscope}%
\end{pgfscope}%
\begin{pgfscope}%
\pgfpathrectangle{\pgfqpoint{0.532060in}{0.456901in}}{\pgfqpoint{5.620383in}{2.816433in}}%
\pgfusepath{clip}%
\pgfsetbuttcap%
\pgfsetroundjoin%
\pgfsetlinewidth{0.501875pt}%
\definecolor{currentstroke}{rgb}{0.501961,0.501961,0.501961}%
\pgfsetstrokecolor{currentstroke}%
\pgfsetstrokeopacity{0.400000}%
\pgfsetdash{{1.850000pt}{0.800000pt}}{0.000000pt}%
\pgfpathmoveto{\pgfqpoint{5.942985in}{0.456901in}}%
\pgfpathlineto{\pgfqpoint{5.942985in}{3.273333in}}%
\pgfusepath{stroke}%
\end{pgfscope}%
\begin{pgfscope}%
\pgfsetbuttcap%
\pgfsetroundjoin%
\definecolor{currentfill}{rgb}{1.000000,1.000000,1.000000}%
\pgfsetfillcolor{currentfill}%
\pgfsetlinewidth{0.602250pt}%
\definecolor{currentstroke}{rgb}{0.000000,0.000000,0.000000}%
\pgfsetstrokecolor{currentstroke}%
\pgfsetdash{}{0pt}%
\pgfsys@defobject{currentmarker}{\pgfqpoint{0.000000in}{0.000000in}}{\pgfqpoint{0.000000in}{0.027778in}}{%
\pgfpathmoveto{\pgfqpoint{0.000000in}{0.000000in}}%
\pgfpathlineto{\pgfqpoint{0.000000in}{0.027778in}}%
\pgfusepath{stroke,fill}%
}%
\begin{pgfscope}%
\pgfsys@transformshift{5.942985in}{0.456901in}%
\pgfsys@useobject{currentmarker}{}%
\end{pgfscope}%
\end{pgfscope}%
\begin{pgfscope}%
\definecolor{textcolor}{rgb}{0.000000,0.000000,0.000000}%
\pgfsetstrokecolor{textcolor}%
\pgfsetfillcolor{textcolor}%
\pgftext[x=3.342251in,y=0.201367in,,top]{\color{textcolor}{\rmfamily\fontsize{12.000000}{14.400000}\selectfont\catcode`\^=\active\def^{\ifmmode\sp\else\^{}\fi}\catcode`\%=\active\def%{\%}$t$}}%
\end{pgfscope}%
\begin{pgfscope}%
\pgfsetbuttcap%
\pgfsetroundjoin%
\definecolor{currentfill}{rgb}{1.000000,1.000000,1.000000}%
\pgfsetfillcolor{currentfill}%
\pgfsetlinewidth{0.803000pt}%
\definecolor{currentstroke}{rgb}{0.000000,0.000000,0.000000}%
\pgfsetstrokecolor{currentstroke}%
\pgfsetdash{}{0pt}%
\pgfsys@defobject{currentmarker}{\pgfqpoint{0.000000in}{0.000000in}}{\pgfqpoint{0.048611in}{0.000000in}}{%
\pgfpathmoveto{\pgfqpoint{0.000000in}{0.000000in}}%
\pgfpathlineto{\pgfqpoint{0.048611in}{0.000000in}}%
\pgfusepath{stroke,fill}%
}%
\begin{pgfscope}%
\pgfsys@transformshift{0.532060in}{0.787834in}%
\pgfsys@useobject{currentmarker}{}%
\end{pgfscope}%
\end{pgfscope}%
\begin{pgfscope}%
\definecolor{textcolor}{rgb}{0.000000,0.000000,0.000000}%
\pgfsetstrokecolor{textcolor}%
\pgfsetfillcolor{textcolor}%
\pgftext[x=0.272222in, y=0.729973in, left, base]{\color{textcolor}{\rmfamily\fontsize{12.000000}{14.400000}\selectfont\catcode`\^=\active\def^{\ifmmode\sp\else\^{}\fi}\catcode`\%=\active\def%{\%}$\mathdefault{\ensuremath{-}2}$}}%
\end{pgfscope}%
\begin{pgfscope}%
\pgfsetbuttcap%
\pgfsetroundjoin%
\definecolor{currentfill}{rgb}{1.000000,1.000000,1.000000}%
\pgfsetfillcolor{currentfill}%
\pgfsetlinewidth{0.803000pt}%
\definecolor{currentstroke}{rgb}{0.000000,0.000000,0.000000}%
\pgfsetstrokecolor{currentstroke}%
\pgfsetdash{}{0pt}%
\pgfsys@defobject{currentmarker}{\pgfqpoint{0.000000in}{0.000000in}}{\pgfqpoint{0.048611in}{0.000000in}}{%
\pgfpathmoveto{\pgfqpoint{0.000000in}{0.000000in}}%
\pgfpathlineto{\pgfqpoint{0.048611in}{0.000000in}}%
\pgfusepath{stroke,fill}%
}%
\begin{pgfscope}%
\pgfsys@transformshift{0.532060in}{1.326475in}%
\pgfsys@useobject{currentmarker}{}%
\end{pgfscope}%
\end{pgfscope}%
\begin{pgfscope}%
\definecolor{textcolor}{rgb}{0.000000,0.000000,0.000000}%
\pgfsetstrokecolor{textcolor}%
\pgfsetfillcolor{textcolor}%
\pgftext[x=0.272222in, y=1.268614in, left, base]{\color{textcolor}{\rmfamily\fontsize{12.000000}{14.400000}\selectfont\catcode`\^=\active\def^{\ifmmode\sp\else\^{}\fi}\catcode`\%=\active\def%{\%}$\mathdefault{\ensuremath{-}1}$}}%
\end{pgfscope}%
\begin{pgfscope}%
\pgfsetbuttcap%
\pgfsetroundjoin%
\definecolor{currentfill}{rgb}{1.000000,1.000000,1.000000}%
\pgfsetfillcolor{currentfill}%
\pgfsetlinewidth{0.803000pt}%
\definecolor{currentstroke}{rgb}{0.000000,0.000000,0.000000}%
\pgfsetstrokecolor{currentstroke}%
\pgfsetdash{}{0pt}%
\pgfsys@defobject{currentmarker}{\pgfqpoint{0.000000in}{0.000000in}}{\pgfqpoint{0.048611in}{0.000000in}}{%
\pgfpathmoveto{\pgfqpoint{0.000000in}{0.000000in}}%
\pgfpathlineto{\pgfqpoint{0.048611in}{0.000000in}}%
\pgfusepath{stroke,fill}%
}%
\begin{pgfscope}%
\pgfsys@transformshift{0.532060in}{1.865117in}%
\pgfsys@useobject{currentmarker}{}%
\end{pgfscope}%
\end{pgfscope}%
\begin{pgfscope}%
\definecolor{textcolor}{rgb}{0.000000,0.000000,0.000000}%
\pgfsetstrokecolor{textcolor}%
\pgfsetfillcolor{textcolor}%
\pgftext[x=0.401852in, y=1.807256in, left, base]{\color{textcolor}{\rmfamily\fontsize{12.000000}{14.400000}\selectfont\catcode`\^=\active\def^{\ifmmode\sp\else\^{}\fi}\catcode`\%=\active\def%{\%}$\mathdefault{0}$}}%
\end{pgfscope}%
\begin{pgfscope}%
\pgfsetbuttcap%
\pgfsetroundjoin%
\definecolor{currentfill}{rgb}{1.000000,1.000000,1.000000}%
\pgfsetfillcolor{currentfill}%
\pgfsetlinewidth{0.803000pt}%
\definecolor{currentstroke}{rgb}{0.000000,0.000000,0.000000}%
\pgfsetstrokecolor{currentstroke}%
\pgfsetdash{}{0pt}%
\pgfsys@defobject{currentmarker}{\pgfqpoint{0.000000in}{0.000000in}}{\pgfqpoint{0.048611in}{0.000000in}}{%
\pgfpathmoveto{\pgfqpoint{0.000000in}{0.000000in}}%
\pgfpathlineto{\pgfqpoint{0.048611in}{0.000000in}}%
\pgfusepath{stroke,fill}%
}%
\begin{pgfscope}%
\pgfsys@transformshift{0.532060in}{2.403759in}%
\pgfsys@useobject{currentmarker}{}%
\end{pgfscope}%
\end{pgfscope}%
\begin{pgfscope}%
\definecolor{textcolor}{rgb}{0.000000,0.000000,0.000000}%
\pgfsetstrokecolor{textcolor}%
\pgfsetfillcolor{textcolor}%
\pgftext[x=0.401852in, y=2.345897in, left, base]{\color{textcolor}{\rmfamily\fontsize{12.000000}{14.400000}\selectfont\catcode`\^=\active\def^{\ifmmode\sp\else\^{}\fi}\catcode`\%=\active\def%{\%}$\mathdefault{1}$}}%
\end{pgfscope}%
\begin{pgfscope}%
\pgfsetbuttcap%
\pgfsetroundjoin%
\definecolor{currentfill}{rgb}{1.000000,1.000000,1.000000}%
\pgfsetfillcolor{currentfill}%
\pgfsetlinewidth{0.803000pt}%
\definecolor{currentstroke}{rgb}{0.000000,0.000000,0.000000}%
\pgfsetstrokecolor{currentstroke}%
\pgfsetdash{}{0pt}%
\pgfsys@defobject{currentmarker}{\pgfqpoint{0.000000in}{0.000000in}}{\pgfqpoint{0.048611in}{0.000000in}}{%
\pgfpathmoveto{\pgfqpoint{0.000000in}{0.000000in}}%
\pgfpathlineto{\pgfqpoint{0.048611in}{0.000000in}}%
\pgfusepath{stroke,fill}%
}%
\begin{pgfscope}%
\pgfsys@transformshift{0.532060in}{2.942400in}%
\pgfsys@useobject{currentmarker}{}%
\end{pgfscope}%
\end{pgfscope}%
\begin{pgfscope}%
\definecolor{textcolor}{rgb}{0.000000,0.000000,0.000000}%
\pgfsetstrokecolor{textcolor}%
\pgfsetfillcolor{textcolor}%
\pgftext[x=0.401852in, y=2.884539in, left, base]{\color{textcolor}{\rmfamily\fontsize{12.000000}{14.400000}\selectfont\catcode`\^=\active\def^{\ifmmode\sp\else\^{}\fi}\catcode`\%=\active\def%{\%}$\mathdefault{2}$}}%
\end{pgfscope}%
\begin{pgfscope}%
\definecolor{textcolor}{rgb}{0.000000,0.000000,0.000000}%
\pgfsetstrokecolor{textcolor}%
\pgfsetfillcolor{textcolor}%
\pgftext[x=0.216667in,y=1.865117in,,bottom,rotate=90.000000]{\color{textcolor}{\rmfamily\fontsize{12.000000}{14.400000}\selectfont\catcode`\^=\active\def^{\ifmmode\sp\else\^{}\fi}\catcode`\%=\active\def%{\%}$i_{\mathrm{н}}(t)$}}%
\end{pgfscope}%
\begin{pgfscope}%
\pgfpathrectangle{\pgfqpoint{0.532060in}{0.456901in}}{\pgfqpoint{5.620383in}{2.816433in}}%
\pgfusepath{clip}%
\pgfsetbuttcap%
\pgfsetroundjoin%
\pgfsetlinewidth{2.007500pt}%
\definecolor{currentstroke}{rgb}{0.000000,0.000000,0.000000}%
\pgfsetstrokecolor{currentstroke}%
\pgfsetdash{{7.400000pt}{3.200000pt}}{0.000000pt}%
\pgfpathmoveto{\pgfqpoint{0.522060in}{0.787834in}}%
\pgfpathlineto{\pgfqpoint{0.740691in}{0.787834in}}%
\pgfpathlineto{\pgfqpoint{0.746658in}{2.942400in}}%
\pgfpathlineto{\pgfqpoint{2.870903in}{2.942400in}}%
\pgfpathlineto{\pgfqpoint{2.876870in}{0.787834in}}%
\pgfpathlineto{\pgfqpoint{5.001114in}{0.787834in}}%
\pgfpathlineto{\pgfqpoint{5.007081in}{2.942400in}}%
\pgfpathlineto{\pgfqpoint{5.854392in}{2.942400in}}%
\pgfpathlineto{\pgfqpoint{5.854392in}{2.942400in}}%
\pgfusepath{stroke}%
\end{pgfscope}%
\begin{pgfscope}%
\pgfpathrectangle{\pgfqpoint{0.532060in}{0.456901in}}{\pgfqpoint{5.620383in}{2.816433in}}%
\pgfusepath{clip}%
\pgfsetrectcap%
\pgfsetroundjoin%
\pgfsetlinewidth{2.007500pt}%
\definecolor{currentstroke}{rgb}{0.835294,0.368627,0.000000}%
\pgfsetstrokecolor{currentstroke}%
\pgfsetdash{}{0pt}%
\pgfpathmoveto{\pgfqpoint{0.522060in}{0.767144in}}%
\pgfpathlineto{\pgfqpoint{0.537814in}{0.815339in}}%
\pgfpathlineto{\pgfqpoint{0.555715in}{0.877106in}}%
\pgfpathlineto{\pgfqpoint{0.573615in}{0.945975in}}%
\pgfpathlineto{\pgfqpoint{0.597483in}{1.048260in}}%
\pgfpathlineto{\pgfqpoint{0.621351in}{1.161526in}}%
\pgfpathlineto{\pgfqpoint{0.651186in}{1.316533in}}%
\pgfpathlineto{\pgfqpoint{0.681021in}{1.483565in}}%
\pgfpathlineto{\pgfqpoint{0.722790in}{1.730817in}}%
\pgfpathlineto{\pgfqpoint{0.812295in}{2.266258in}}%
\pgfpathlineto{\pgfqpoint{0.848097in}{2.463868in}}%
\pgfpathlineto{\pgfqpoint{0.877931in}{2.614437in}}%
\pgfpathlineto{\pgfqpoint{0.901799in}{2.723531in}}%
\pgfpathlineto{\pgfqpoint{0.925667in}{2.821201in}}%
\pgfpathlineto{\pgfqpoint{0.943568in}{2.886385in}}%
\pgfpathlineto{\pgfqpoint{0.961469in}{2.944330in}}%
\pgfpathlineto{\pgfqpoint{0.979370in}{2.994850in}}%
\pgfpathlineto{\pgfqpoint{0.997271in}{3.037854in}}%
\pgfpathlineto{\pgfqpoint{1.015172in}{3.073353in}}%
\pgfpathlineto{\pgfqpoint{1.027106in}{3.092897in}}%
\pgfpathlineto{\pgfqpoint{1.039040in}{3.109205in}}%
\pgfpathlineto{\pgfqpoint{1.050974in}{3.122349in}}%
\pgfpathlineto{\pgfqpoint{1.062908in}{3.132418in}}%
\pgfpathlineto{\pgfqpoint{1.074842in}{3.139517in}}%
\pgfpathlineto{\pgfqpoint{1.086776in}{3.143767in}}%
\pgfpathlineto{\pgfqpoint{1.098710in}{3.145304in}}%
\pgfpathlineto{\pgfqpoint{1.110644in}{3.144276in}}%
\pgfpathlineto{\pgfqpoint{1.122578in}{3.140842in}}%
\pgfpathlineto{\pgfqpoint{1.134512in}{3.135175in}}%
\pgfpathlineto{\pgfqpoint{1.146446in}{3.127455in}}%
\pgfpathlineto{\pgfqpoint{1.164346in}{3.112437in}}%
\pgfpathlineto{\pgfqpoint{1.182247in}{3.093886in}}%
\pgfpathlineto{\pgfqpoint{1.206115in}{3.064835in}}%
\pgfpathlineto{\pgfqpoint{1.235950in}{3.023896in}}%
\pgfpathlineto{\pgfqpoint{1.313521in}{2.914510in}}%
\pgfpathlineto{\pgfqpoint{1.337389in}{2.885462in}}%
\pgfpathlineto{\pgfqpoint{1.355290in}{2.866358in}}%
\pgfpathlineto{\pgfqpoint{1.373191in}{2.849958in}}%
\pgfpathlineto{\pgfqpoint{1.391092in}{2.836545in}}%
\pgfpathlineto{\pgfqpoint{1.408993in}{2.826326in}}%
\pgfpathlineto{\pgfqpoint{1.426894in}{2.819433in}}%
\pgfpathlineto{\pgfqpoint{1.444794in}{2.815917in}}%
\pgfpathlineto{\pgfqpoint{1.462695in}{2.815754in}}%
\pgfpathlineto{\pgfqpoint{1.480596in}{2.818849in}}%
\pgfpathlineto{\pgfqpoint{1.498497in}{2.825034in}}%
\pgfpathlineto{\pgfqpoint{1.516398in}{2.834077in}}%
\pgfpathlineto{\pgfqpoint{1.534299in}{2.845691in}}%
\pgfpathlineto{\pgfqpoint{1.558167in}{2.864580in}}%
\pgfpathlineto{\pgfqpoint{1.582035in}{2.886510in}}%
\pgfpathlineto{\pgfqpoint{1.617837in}{2.922870in}}%
\pgfpathlineto{\pgfqpoint{1.671540in}{2.978012in}}%
\pgfpathlineto{\pgfqpoint{1.701375in}{3.005137in}}%
\pgfpathlineto{\pgfqpoint{1.725243in}{3.023452in}}%
\pgfpathlineto{\pgfqpoint{1.743143in}{3.034705in}}%
\pgfpathlineto{\pgfqpoint{1.761044in}{3.043533in}}%
\pgfpathlineto{\pgfqpoint{1.778945in}{3.049734in}}%
\pgfpathlineto{\pgfqpoint{1.796846in}{3.053164in}}%
\pgfpathlineto{\pgfqpoint{1.814747in}{3.053744in}}%
\pgfpathlineto{\pgfqpoint{1.832648in}{3.051462in}}%
\pgfpathlineto{\pgfqpoint{1.850549in}{3.046369in}}%
\pgfpathlineto{\pgfqpoint{1.868450in}{3.038583in}}%
\pgfpathlineto{\pgfqpoint{1.886351in}{3.028283in}}%
\pgfpathlineto{\pgfqpoint{1.904252in}{3.015708in}}%
\pgfpathlineto{\pgfqpoint{1.928120in}{2.995915in}}%
\pgfpathlineto{\pgfqpoint{1.957955in}{2.967488in}}%
\pgfpathlineto{\pgfqpoint{2.059393in}{2.865686in}}%
\pgfpathlineto{\pgfqpoint{2.083261in}{2.846610in}}%
\pgfpathlineto{\pgfqpoint{2.101162in}{2.834825in}}%
\pgfpathlineto{\pgfqpoint{2.119063in}{2.825588in}}%
\pgfpathlineto{\pgfqpoint{2.136964in}{2.819191in}}%
\pgfpathlineto{\pgfqpoint{2.154865in}{2.815869in}}%
\pgfpathlineto{\pgfqpoint{2.172766in}{2.815794in}}%
\pgfpathlineto{\pgfqpoint{2.184700in}{2.817602in}}%
\pgfpathlineto{\pgfqpoint{2.202601in}{2.823131in}}%
\pgfpathlineto{\pgfqpoint{2.220502in}{2.832015in}}%
\pgfpathlineto{\pgfqpoint{2.238403in}{2.844156in}}%
\pgfpathlineto{\pgfqpoint{2.256304in}{2.859377in}}%
\pgfpathlineto{\pgfqpoint{2.274205in}{2.877425in}}%
\pgfpathlineto{\pgfqpoint{2.298073in}{2.905310in}}%
\pgfpathlineto{\pgfqpoint{2.321940in}{2.936658in}}%
\pgfpathlineto{\pgfqpoint{2.363709in}{2.996226in}}%
\pgfpathlineto{\pgfqpoint{2.405478in}{3.055263in}}%
\pgfpathlineto{\pgfqpoint{2.429346in}{3.085527in}}%
\pgfpathlineto{\pgfqpoint{2.447247in}{3.105316in}}%
\pgfpathlineto{\pgfqpoint{2.465148in}{3.121845in}}%
\pgfpathlineto{\pgfqpoint{2.483049in}{3.134443in}}%
\pgfpathlineto{\pgfqpoint{2.494983in}{3.140339in}}%
\pgfpathlineto{\pgfqpoint{2.506917in}{3.144020in}}%
\pgfpathlineto{\pgfqpoint{2.518851in}{3.145314in}}%
\pgfpathlineto{\pgfqpoint{2.530785in}{3.144059in}}%
\pgfpathlineto{\pgfqpoint{2.542719in}{3.140106in}}%
\pgfpathlineto{\pgfqpoint{2.554653in}{3.133318in}}%
\pgfpathlineto{\pgfqpoint{2.566587in}{3.123573in}}%
\pgfpathlineto{\pgfqpoint{2.578521in}{3.110764in}}%
\pgfpathlineto{\pgfqpoint{2.590455in}{3.094798in}}%
\pgfpathlineto{\pgfqpoint{2.602389in}{3.075604in}}%
\pgfpathlineto{\pgfqpoint{2.614322in}{3.053123in}}%
\pgfpathlineto{\pgfqpoint{2.632223in}{3.013163in}}%
\pgfpathlineto{\pgfqpoint{2.650124in}{2.965681in}}%
\pgfpathlineto{\pgfqpoint{2.668025in}{2.910724in}}%
\pgfpathlineto{\pgfqpoint{2.685926in}{2.848441in}}%
\pgfpathlineto{\pgfqpoint{2.703827in}{2.779076in}}%
\pgfpathlineto{\pgfqpoint{2.727695in}{2.676174in}}%
\pgfpathlineto{\pgfqpoint{2.751563in}{2.562354in}}%
\pgfpathlineto{\pgfqpoint{2.781398in}{2.406765in}}%
\pgfpathlineto{\pgfqpoint{2.817200in}{2.204680in}}%
\pgfpathlineto{\pgfqpoint{2.864936in}{1.919448in}}%
\pgfpathlineto{\pgfqpoint{2.936539in}{1.490940in}}%
\pgfpathlineto{\pgfqpoint{2.972341in}{1.291294in}}%
\pgfpathlineto{\pgfqpoint{3.002176in}{1.138469in}}%
\pgfpathlineto{\pgfqpoint{3.026044in}{1.027265in}}%
\pgfpathlineto{\pgfqpoint{3.049912in}{0.927269in}}%
\pgfpathlineto{\pgfqpoint{3.067813in}{0.860231in}}%
\pgfpathlineto{\pgfqpoint{3.085714in}{0.800368in}}%
\pgfpathlineto{\pgfqpoint{3.103615in}{0.747891in}}%
\pgfpathlineto{\pgfqpoint{3.121516in}{0.702916in}}%
\pgfpathlineto{\pgfqpoint{3.139417in}{0.665457in}}%
\pgfpathlineto{\pgfqpoint{3.151351in}{0.644626in}}%
\pgfpathlineto{\pgfqpoint{3.163285in}{0.627054in}}%
\pgfpathlineto{\pgfqpoint{3.175219in}{0.612677in}}%
\pgfpathlineto{\pgfqpoint{3.187153in}{0.601412in}}%
\pgfpathlineto{\pgfqpoint{3.199086in}{0.593160in}}%
\pgfpathlineto{\pgfqpoint{3.211020in}{0.587806in}}%
\pgfpathlineto{\pgfqpoint{3.222954in}{0.585220in}}%
\pgfpathlineto{\pgfqpoint{3.234888in}{0.585259in}}%
\pgfpathlineto{\pgfqpoint{3.246822in}{0.587767in}}%
\pgfpathlineto{\pgfqpoint{3.258756in}{0.592577in}}%
\pgfpathlineto{\pgfqpoint{3.270690in}{0.599513in}}%
\pgfpathlineto{\pgfqpoint{3.288591in}{0.613495in}}%
\pgfpathlineto{\pgfqpoint{3.306492in}{0.631186in}}%
\pgfpathlineto{\pgfqpoint{3.324393in}{0.651909in}}%
\pgfpathlineto{\pgfqpoint{3.348261in}{0.683059in}}%
\pgfpathlineto{\pgfqpoint{3.384063in}{0.734008in}}%
\pgfpathlineto{\pgfqpoint{3.431799in}{0.801677in}}%
\pgfpathlineto{\pgfqpoint{3.455667in}{0.832262in}}%
\pgfpathlineto{\pgfqpoint{3.479535in}{0.859121in}}%
\pgfpathlineto{\pgfqpoint{3.497435in}{0.876259in}}%
\pgfpathlineto{\pgfqpoint{3.515336in}{0.890478in}}%
\pgfpathlineto{\pgfqpoint{3.533237in}{0.901549in}}%
\pgfpathlineto{\pgfqpoint{3.551138in}{0.909322in}}%
\pgfpathlineto{\pgfqpoint{3.569039in}{0.913723in}}%
\pgfpathlineto{\pgfqpoint{3.586940in}{0.914757in}}%
\pgfpathlineto{\pgfqpoint{3.604841in}{0.912502in}}%
\pgfpathlineto{\pgfqpoint{3.622742in}{0.907107in}}%
\pgfpathlineto{\pgfqpoint{3.640643in}{0.898786in}}%
\pgfpathlineto{\pgfqpoint{3.658544in}{0.887815in}}%
\pgfpathlineto{\pgfqpoint{3.676445in}{0.874520in}}%
\pgfpathlineto{\pgfqpoint{3.700313in}{0.853827in}}%
\pgfpathlineto{\pgfqpoint{3.730148in}{0.824583in}}%
\pgfpathlineto{\pgfqpoint{3.813685in}{0.739724in}}%
\pgfpathlineto{\pgfqpoint{3.837553in}{0.719174in}}%
\pgfpathlineto{\pgfqpoint{3.861421in}{0.701951in}}%
\pgfpathlineto{\pgfqpoint{3.879322in}{0.691651in}}%
\pgfpathlineto{\pgfqpoint{3.897223in}{0.683865in}}%
\pgfpathlineto{\pgfqpoint{3.915124in}{0.678772in}}%
\pgfpathlineto{\pgfqpoint{3.933025in}{0.676490in}}%
\pgfpathlineto{\pgfqpoint{3.950926in}{0.677070in}}%
\pgfpathlineto{\pgfqpoint{3.968827in}{0.680500in}}%
\pgfpathlineto{\pgfqpoint{3.986728in}{0.686701in}}%
\pgfpathlineto{\pgfqpoint{4.004629in}{0.695530in}}%
\pgfpathlineto{\pgfqpoint{4.022530in}{0.706782in}}%
\pgfpathlineto{\pgfqpoint{4.046398in}{0.725097in}}%
\pgfpathlineto{\pgfqpoint{4.070266in}{0.746494in}}%
\pgfpathlineto{\pgfqpoint{4.106067in}{0.782384in}}%
\pgfpathlineto{\pgfqpoint{4.171704in}{0.849438in}}%
\pgfpathlineto{\pgfqpoint{4.195572in}{0.870700in}}%
\pgfpathlineto{\pgfqpoint{4.219440in}{0.888680in}}%
\pgfpathlineto{\pgfqpoint{4.237341in}{0.899472in}}%
\pgfpathlineto{\pgfqpoint{4.255242in}{0.907592in}}%
\pgfpathlineto{\pgfqpoint{4.273143in}{0.912770in}}%
\pgfpathlineto{\pgfqpoint{4.291044in}{0.914794in}}%
\pgfpathlineto{\pgfqpoint{4.302978in}{0.914318in}}%
\pgfpathlineto{\pgfqpoint{4.314912in}{0.912349in}}%
\pgfpathlineto{\pgfqpoint{4.332813in}{0.906579in}}%
\pgfpathlineto{\pgfqpoint{4.350714in}{0.897459in}}%
\pgfpathlineto{\pgfqpoint{4.368615in}{0.885091in}}%
\pgfpathlineto{\pgfqpoint{4.386515in}{0.869658in}}%
\pgfpathlineto{\pgfqpoint{4.404416in}{0.851420in}}%
\pgfpathlineto{\pgfqpoint{4.428284in}{0.823323in}}%
\pgfpathlineto{\pgfqpoint{4.458119in}{0.783548in}}%
\pgfpathlineto{\pgfqpoint{4.553591in}{0.650343in}}%
\pgfpathlineto{\pgfqpoint{4.571492in}{0.629814in}}%
\pgfpathlineto{\pgfqpoint{4.589393in}{0.612365in}}%
\pgfpathlineto{\pgfqpoint{4.607294in}{0.598674in}}%
\pgfpathlineto{\pgfqpoint{4.619228in}{0.591957in}}%
\pgfpathlineto{\pgfqpoint{4.631162in}{0.587385in}}%
\pgfpathlineto{\pgfqpoint{4.643096in}{0.585134in}}%
\pgfpathlineto{\pgfqpoint{4.655030in}{0.585368in}}%
\pgfpathlineto{\pgfqpoint{4.666963in}{0.588244in}}%
\pgfpathlineto{\pgfqpoint{4.678897in}{0.593903in}}%
\pgfpathlineto{\pgfqpoint{4.690831in}{0.602473in}}%
\pgfpathlineto{\pgfqpoint{4.702765in}{0.614067in}}%
\pgfpathlineto{\pgfqpoint{4.714699in}{0.628782in}}%
\pgfpathlineto{\pgfqpoint{4.726633in}{0.646701in}}%
\pgfpathlineto{\pgfqpoint{4.738567in}{0.667885in}}%
\pgfpathlineto{\pgfqpoint{4.750501in}{0.692380in}}%
\pgfpathlineto{\pgfqpoint{4.768402in}{0.735385in}}%
\pgfpathlineto{\pgfqpoint{4.786303in}{0.785904in}}%
\pgfpathlineto{\pgfqpoint{4.804204in}{0.843849in}}%
\pgfpathlineto{\pgfqpoint{4.822105in}{0.909033in}}%
\pgfpathlineto{\pgfqpoint{4.840006in}{0.981173in}}%
\pgfpathlineto{\pgfqpoint{4.863874in}{1.087514in}}%
\pgfpathlineto{\pgfqpoint{4.887742in}{1.204398in}}%
\pgfpathlineto{\pgfqpoint{4.917577in}{1.363177in}}%
\pgfpathlineto{\pgfqpoint{4.953379in}{1.567992in}}%
\pgfpathlineto{\pgfqpoint{5.007081in}{1.890996in}}%
\pgfpathlineto{\pgfqpoint{5.066751in}{2.246669in}}%
\pgfpathlineto{\pgfqpoint{5.102553in}{2.445770in}}%
\pgfpathlineto{\pgfqpoint{5.132388in}{2.597987in}}%
\pgfpathlineto{\pgfqpoint{5.156256in}{2.708620in}}%
\pgfpathlineto{\pgfqpoint{5.180124in}{2.807986in}}%
\pgfpathlineto{\pgfqpoint{5.198025in}{2.874520in}}%
\pgfpathlineto{\pgfqpoint{5.215926in}{2.933861in}}%
\pgfpathlineto{\pgfqpoint{5.233827in}{2.985805in}}%
\pgfpathlineto{\pgfqpoint{5.251727in}{3.030243in}}%
\pgfpathlineto{\pgfqpoint{5.269628in}{3.067166in}}%
\pgfpathlineto{\pgfqpoint{5.281562in}{3.087646in}}%
\pgfpathlineto{\pgfqpoint{5.293496in}{3.104872in}}%
\pgfpathlineto{\pgfqpoint{5.305430in}{3.118912in}}%
\pgfpathlineto{\pgfqpoint{5.317364in}{3.129849in}}%
\pgfpathlineto{\pgfqpoint{5.329298in}{3.137785in}}%
\pgfpathlineto{\pgfqpoint{5.341232in}{3.142836in}}%
\pgfpathlineto{\pgfqpoint{5.353166in}{3.145133in}}%
\pgfpathlineto{\pgfqpoint{5.365100in}{3.144822in}}%
\pgfpathlineto{\pgfqpoint{5.377034in}{3.142059in}}%
\pgfpathlineto{\pgfqpoint{5.388968in}{3.137012in}}%
\pgfpathlineto{\pgfqpoint{5.400902in}{3.129859in}}%
\pgfpathlineto{\pgfqpoint{5.418803in}{3.115591in}}%
\pgfpathlineto{\pgfqpoint{5.436704in}{3.097660in}}%
\pgfpathlineto{\pgfqpoint{5.454605in}{3.076747in}}%
\pgfpathlineto{\pgfqpoint{5.478473in}{3.045421in}}%
\pgfpathlineto{\pgfqpoint{5.520242in}{2.985705in}}%
\pgfpathlineto{\pgfqpoint{5.562010in}{2.926842in}}%
\pgfpathlineto{\pgfqpoint{5.585878in}{2.896430in}}%
\pgfpathlineto{\pgfqpoint{5.609746in}{2.869799in}}%
\pgfpathlineto{\pgfqpoint{5.627647in}{2.852859in}}%
\pgfpathlineto{\pgfqpoint{5.645548in}{2.838859in}}%
\pgfpathlineto{\pgfqpoint{5.663449in}{2.828019in}}%
\pgfpathlineto{\pgfqpoint{5.681350in}{2.820486in}}%
\pgfpathlineto{\pgfqpoint{5.699251in}{2.816326in}}%
\pgfpathlineto{\pgfqpoint{5.717152in}{2.815530in}}%
\pgfpathlineto{\pgfqpoint{5.735053in}{2.818015in}}%
\pgfpathlineto{\pgfqpoint{5.752954in}{2.823627in}}%
\pgfpathlineto{\pgfqpoint{5.770855in}{2.832147in}}%
\pgfpathlineto{\pgfqpoint{5.788756in}{2.843296in}}%
\pgfpathlineto{\pgfqpoint{5.806657in}{2.856743in}}%
\pgfpathlineto{\pgfqpoint{5.830525in}{2.877595in}}%
\pgfpathlineto{\pgfqpoint{5.854392in}{2.900881in}}%
\pgfpathlineto{\pgfqpoint{5.854392in}{2.900881in}}%
\pgfusepath{stroke}%
\end{pgfscope}%
\begin{pgfscope}%
\pgfpathrectangle{\pgfqpoint{0.532060in}{0.456901in}}{\pgfqpoint{5.620383in}{2.816433in}}%
\pgfusepath{clip}%
\pgfsetbuttcap%
\pgfsetroundjoin%
\pgfsetlinewidth{2.007500pt}%
\definecolor{currentstroke}{rgb}{0.000000,0.619608,0.450980}%
\pgfsetstrokecolor{currentstroke}%
\pgfsetdash{{12.800000pt}{3.200000pt}{2.000000pt}{3.200000pt}}{0.000000pt}%
\pgfpathmoveto{\pgfqpoint{0.522060in}{1.182260in}}%
\pgfpathlineto{\pgfqpoint{0.543781in}{1.193846in}}%
\pgfpathlineto{\pgfqpoint{0.567648in}{1.209397in}}%
\pgfpathlineto{\pgfqpoint{0.591516in}{1.227888in}}%
\pgfpathlineto{\pgfqpoint{0.615384in}{1.249276in}}%
\pgfpathlineto{\pgfqpoint{0.639252in}{1.273455in}}%
\pgfpathlineto{\pgfqpoint{0.669087in}{1.307349in}}%
\pgfpathlineto{\pgfqpoint{0.698922in}{1.344900in}}%
\pgfpathlineto{\pgfqpoint{0.734724in}{1.393985in}}%
\pgfpathlineto{\pgfqpoint{0.776493in}{1.455304in}}%
\pgfpathlineto{\pgfqpoint{0.919700in}{1.670026in}}%
\pgfpathlineto{\pgfqpoint{0.955502in}{1.718104in}}%
\pgfpathlineto{\pgfqpoint{0.985337in}{1.754901in}}%
\pgfpathlineto{\pgfqpoint{1.015172in}{1.788331in}}%
\pgfpathlineto{\pgfqpoint{1.045007in}{1.818174in}}%
\pgfpathlineto{\pgfqpoint{1.074842in}{1.844365in}}%
\pgfpathlineto{\pgfqpoint{1.104677in}{1.866991in}}%
\pgfpathlineto{\pgfqpoint{1.134512in}{1.886284in}}%
\pgfpathlineto{\pgfqpoint{1.164346in}{1.902609in}}%
\pgfpathlineto{\pgfqpoint{1.194181in}{1.916442in}}%
\pgfpathlineto{\pgfqpoint{1.229983in}{1.930552in}}%
\pgfpathlineto{\pgfqpoint{1.283686in}{1.948912in}}%
\pgfpathlineto{\pgfqpoint{1.349323in}{1.971738in}}%
\pgfpathlineto{\pgfqpoint{1.385125in}{1.986469in}}%
\pgfpathlineto{\pgfqpoint{1.420927in}{2.003855in}}%
\pgfpathlineto{\pgfqpoint{1.450761in}{2.020775in}}%
\pgfpathlineto{\pgfqpoint{1.480596in}{2.040101in}}%
\pgfpathlineto{\pgfqpoint{1.510431in}{2.061879in}}%
\pgfpathlineto{\pgfqpoint{1.546233in}{2.091116in}}%
\pgfpathlineto{\pgfqpoint{1.582035in}{2.123361in}}%
\pgfpathlineto{\pgfqpoint{1.623804in}{2.163953in}}%
\pgfpathlineto{\pgfqpoint{1.695408in}{2.237322in}}%
\pgfpathlineto{\pgfqpoint{1.755077in}{2.297457in}}%
\pgfpathlineto{\pgfqpoint{1.796846in}{2.336557in}}%
\pgfpathlineto{\pgfqpoint{1.832648in}{2.367048in}}%
\pgfpathlineto{\pgfqpoint{1.862483in}{2.389898in}}%
\pgfpathlineto{\pgfqpoint{1.892318in}{2.410202in}}%
\pgfpathlineto{\pgfqpoint{1.922153in}{2.427874in}}%
\pgfpathlineto{\pgfqpoint{1.951988in}{2.442940in}}%
\pgfpathlineto{\pgfqpoint{1.981823in}{2.455534in}}%
\pgfpathlineto{\pgfqpoint{2.011658in}{2.465887in}}%
\pgfpathlineto{\pgfqpoint{2.047459in}{2.475797in}}%
\pgfpathlineto{\pgfqpoint{2.089228in}{2.484725in}}%
\pgfpathlineto{\pgfqpoint{2.148898in}{2.494727in}}%
\pgfpathlineto{\pgfqpoint{2.238403in}{2.509714in}}%
\pgfpathlineto{\pgfqpoint{2.292106in}{2.521131in}}%
\pgfpathlineto{\pgfqpoint{2.351775in}{2.536230in}}%
\pgfpathlineto{\pgfqpoint{2.447247in}{2.560902in}}%
\pgfpathlineto{\pgfqpoint{2.483049in}{2.567837in}}%
\pgfpathlineto{\pgfqpoint{2.512884in}{2.571553in}}%
\pgfpathlineto{\pgfqpoint{2.542719in}{2.572771in}}%
\pgfpathlineto{\pgfqpoint{2.566587in}{2.571568in}}%
\pgfpathlineto{\pgfqpoint{2.590455in}{2.568142in}}%
\pgfpathlineto{\pgfqpoint{2.614322in}{2.562258in}}%
\pgfpathlineto{\pgfqpoint{2.638190in}{2.553721in}}%
\pgfpathlineto{\pgfqpoint{2.662058in}{2.542382in}}%
\pgfpathlineto{\pgfqpoint{2.685926in}{2.528143in}}%
\pgfpathlineto{\pgfqpoint{2.709794in}{2.510963in}}%
\pgfpathlineto{\pgfqpoint{2.733662in}{2.490859in}}%
\pgfpathlineto{\pgfqpoint{2.757530in}{2.467909in}}%
\pgfpathlineto{\pgfqpoint{2.781398in}{2.442251in}}%
\pgfpathlineto{\pgfqpoint{2.811233in}{2.406676in}}%
\pgfpathlineto{\pgfqpoint{2.841068in}{2.367679in}}%
\pgfpathlineto{\pgfqpoint{2.876870in}{2.317239in}}%
\pgfpathlineto{\pgfqpoint{2.924605in}{2.245849in}}%
\pgfpathlineto{\pgfqpoint{3.032011in}{2.083759in}}%
\pgfpathlineto{\pgfqpoint{3.067813in}{2.033963in}}%
\pgfpathlineto{\pgfqpoint{3.103615in}{1.988104in}}%
\pgfpathlineto{\pgfqpoint{3.133450in}{1.953439in}}%
\pgfpathlineto{\pgfqpoint{3.163285in}{1.922299in}}%
\pgfpathlineto{\pgfqpoint{3.193120in}{1.894806in}}%
\pgfpathlineto{\pgfqpoint{3.222954in}{1.870927in}}%
\pgfpathlineto{\pgfqpoint{3.252789in}{1.850479in}}%
\pgfpathlineto{\pgfqpoint{3.282624in}{1.833143in}}%
\pgfpathlineto{\pgfqpoint{3.312459in}{1.818481in}}%
\pgfpathlineto{\pgfqpoint{3.348261in}{1.803655in}}%
\pgfpathlineto{\pgfqpoint{3.395997in}{1.786831in}}%
\pgfpathlineto{\pgfqpoint{3.479535in}{1.758008in}}%
\pgfpathlineto{\pgfqpoint{3.515336in}{1.743194in}}%
\pgfpathlineto{\pgfqpoint{3.545171in}{1.728829in}}%
\pgfpathlineto{\pgfqpoint{3.575006in}{1.712275in}}%
\pgfpathlineto{\pgfqpoint{3.604841in}{1.693333in}}%
\pgfpathlineto{\pgfqpoint{3.634676in}{1.671938in}}%
\pgfpathlineto{\pgfqpoint{3.670478in}{1.643129in}}%
\pgfpathlineto{\pgfqpoint{3.706280in}{1.611247in}}%
\pgfpathlineto{\pgfqpoint{3.748049in}{1.570961in}}%
\pgfpathlineto{\pgfqpoint{3.813685in}{1.503907in}}%
\pgfpathlineto{\pgfqpoint{3.885289in}{1.431534in}}%
\pgfpathlineto{\pgfqpoint{3.927058in}{1.392536in}}%
\pgfpathlineto{\pgfqpoint{3.962860in}{1.362157in}}%
\pgfpathlineto{\pgfqpoint{3.992695in}{1.339413in}}%
\pgfpathlineto{\pgfqpoint{4.022530in}{1.319220in}}%
\pgfpathlineto{\pgfqpoint{4.052365in}{1.301662in}}%
\pgfpathlineto{\pgfqpoint{4.082199in}{1.286705in}}%
\pgfpathlineto{\pgfqpoint{4.112034in}{1.274213in}}%
\pgfpathlineto{\pgfqpoint{4.141869in}{1.263949in}}%
\pgfpathlineto{\pgfqpoint{4.177671in}{1.254127in}}%
\pgfpathlineto{\pgfqpoint{4.219440in}{1.245269in}}%
\pgfpathlineto{\pgfqpoint{4.279110in}{1.235308in}}%
\pgfpathlineto{\pgfqpoint{4.368615in}{1.220275in}}%
\pgfpathlineto{\pgfqpoint{4.422317in}{1.208804in}}%
\pgfpathlineto{\pgfqpoint{4.481987in}{1.193662in}}%
\pgfpathlineto{\pgfqpoint{4.577459in}{1.169048in}}%
\pgfpathlineto{\pgfqpoint{4.613261in}{1.162195in}}%
\pgfpathlineto{\pgfqpoint{4.643096in}{1.158573in}}%
\pgfpathlineto{\pgfqpoint{4.672930in}{1.157475in}}%
\pgfpathlineto{\pgfqpoint{4.696798in}{1.158791in}}%
\pgfpathlineto{\pgfqpoint{4.720666in}{1.162343in}}%
\pgfpathlineto{\pgfqpoint{4.744534in}{1.168365in}}%
\pgfpathlineto{\pgfqpoint{4.768402in}{1.177048in}}%
\pgfpathlineto{\pgfqpoint{4.792270in}{1.188541in}}%
\pgfpathlineto{\pgfqpoint{4.816138in}{1.202936in}}%
\pgfpathlineto{\pgfqpoint{4.840006in}{1.220274in}}%
\pgfpathlineto{\pgfqpoint{4.863874in}{1.240534in}}%
\pgfpathlineto{\pgfqpoint{4.887742in}{1.263633in}}%
\pgfpathlineto{\pgfqpoint{4.911610in}{1.289430in}}%
\pgfpathlineto{\pgfqpoint{4.941445in}{1.325164in}}%
\pgfpathlineto{\pgfqpoint{4.971279in}{1.364294in}}%
\pgfpathlineto{\pgfqpoint{5.007081in}{1.414856in}}%
\pgfpathlineto{\pgfqpoint{5.054817in}{1.486336in}}%
\pgfpathlineto{\pgfqpoint{5.156256in}{1.639713in}}%
\pgfpathlineto{\pgfqpoint{5.198025in}{1.697982in}}%
\pgfpathlineto{\pgfqpoint{5.233827in}{1.743685in}}%
\pgfpathlineto{\pgfqpoint{5.263661in}{1.778204in}}%
\pgfpathlineto{\pgfqpoint{5.293496in}{1.809188in}}%
\pgfpathlineto{\pgfqpoint{5.323331in}{1.836525in}}%
\pgfpathlineto{\pgfqpoint{5.353166in}{1.860252in}}%
\pgfpathlineto{\pgfqpoint{5.383001in}{1.880560in}}%
\pgfpathlineto{\pgfqpoint{5.412836in}{1.897771in}}%
\pgfpathlineto{\pgfqpoint{5.442671in}{1.912330in}}%
\pgfpathlineto{\pgfqpoint{5.478473in}{1.927065in}}%
\pgfpathlineto{\pgfqpoint{5.526209in}{1.943829in}}%
\pgfpathlineto{\pgfqpoint{5.609746in}{1.972717in}}%
\pgfpathlineto{\pgfqpoint{5.645548in}{1.987615in}}%
\pgfpathlineto{\pgfqpoint{5.675383in}{2.002068in}}%
\pgfpathlineto{\pgfqpoint{5.705218in}{2.018721in}}%
\pgfpathlineto{\pgfqpoint{5.735053in}{2.037768in}}%
\pgfpathlineto{\pgfqpoint{5.764888in}{2.059267in}}%
\pgfpathlineto{\pgfqpoint{5.800690in}{2.088194in}}%
\pgfpathlineto{\pgfqpoint{5.836491in}{2.120176in}}%
\pgfpathlineto{\pgfqpoint{5.854392in}{2.137135in}}%
\pgfpathlineto{\pgfqpoint{5.854392in}{2.137135in}}%
\pgfusepath{stroke}%
\end{pgfscope}%
\begin{pgfscope}%
\pgfpathrectangle{\pgfqpoint{0.532060in}{0.456901in}}{\pgfqpoint{5.620383in}{2.816433in}}%
\pgfusepath{clip}%
\pgfsetrectcap%
\pgfsetroundjoin%
\pgfsetlinewidth{1.003750pt}%
\definecolor{currentstroke}{rgb}{0.000000,0.750000,0.750000}%
\pgfsetstrokecolor{currentstroke}%
\pgfsetdash{}{0pt}%
\pgfpathmoveto{\pgfqpoint{0.522060in}{1.289207in}}%
\pgfpathlineto{\pgfqpoint{0.567648in}{1.313264in}}%
\pgfpathlineto{\pgfqpoint{0.615384in}{1.341133in}}%
\pgfpathlineto{\pgfqpoint{0.663120in}{1.371600in}}%
\pgfpathlineto{\pgfqpoint{0.710856in}{1.404516in}}%
\pgfpathlineto{\pgfqpoint{0.764559in}{1.444268in}}%
\pgfpathlineto{\pgfqpoint{0.818262in}{1.486662in}}%
\pgfpathlineto{\pgfqpoint{0.877931in}{1.536540in}}%
\pgfpathlineto{\pgfqpoint{0.943568in}{1.594331in}}%
\pgfpathlineto{\pgfqpoint{1.021139in}{1.665859in}}%
\pgfpathlineto{\pgfqpoint{1.110644in}{1.751569in}}%
\pgfpathlineto{\pgfqpoint{1.265785in}{1.903936in}}%
\pgfpathlineto{\pgfqpoint{1.397059in}{2.031598in}}%
\pgfpathlineto{\pgfqpoint{1.480596in}{2.109889in}}%
\pgfpathlineto{\pgfqpoint{1.552200in}{2.174082in}}%
\pgfpathlineto{\pgfqpoint{1.617837in}{2.229915in}}%
\pgfpathlineto{\pgfqpoint{1.677507in}{2.277718in}}%
\pgfpathlineto{\pgfqpoint{1.731210in}{2.318017in}}%
\pgfpathlineto{\pgfqpoint{1.784912in}{2.355473in}}%
\pgfpathlineto{\pgfqpoint{1.832648in}{2.386191in}}%
\pgfpathlineto{\pgfqpoint{1.880384in}{2.414324in}}%
\pgfpathlineto{\pgfqpoint{1.928120in}{2.439733in}}%
\pgfpathlineto{\pgfqpoint{1.969889in}{2.459632in}}%
\pgfpathlineto{\pgfqpoint{2.011658in}{2.477273in}}%
\pgfpathlineto{\pgfqpoint{2.053426in}{2.492589in}}%
\pgfpathlineto{\pgfqpoint{2.095195in}{2.505522in}}%
\pgfpathlineto{\pgfqpoint{2.136964in}{2.516023in}}%
\pgfpathlineto{\pgfqpoint{2.178733in}{2.524051in}}%
\pgfpathlineto{\pgfqpoint{2.220502in}{2.529577in}}%
\pgfpathlineto{\pgfqpoint{2.262271in}{2.532580in}}%
\pgfpathlineto{\pgfqpoint{2.304040in}{2.533047in}}%
\pgfpathlineto{\pgfqpoint{2.345808in}{2.530978in}}%
\pgfpathlineto{\pgfqpoint{2.387577in}{2.526380in}}%
\pgfpathlineto{\pgfqpoint{2.429346in}{2.519271in}}%
\pgfpathlineto{\pgfqpoint{2.471115in}{2.509677in}}%
\pgfpathlineto{\pgfqpoint{2.512884in}{2.497635in}}%
\pgfpathlineto{\pgfqpoint{2.554653in}{2.483191in}}%
\pgfpathlineto{\pgfqpoint{2.596422in}{2.466400in}}%
\pgfpathlineto{\pgfqpoint{2.638190in}{2.447325in}}%
\pgfpathlineto{\pgfqpoint{2.679959in}{2.426039in}}%
\pgfpathlineto{\pgfqpoint{2.727695in}{2.399108in}}%
\pgfpathlineto{\pgfqpoint{2.775431in}{2.369529in}}%
\pgfpathlineto{\pgfqpoint{2.823167in}{2.337448in}}%
\pgfpathlineto{\pgfqpoint{2.876870in}{2.298565in}}%
\pgfpathlineto{\pgfqpoint{2.930572in}{2.256961in}}%
\pgfpathlineto{\pgfqpoint{2.990242in}{2.207861in}}%
\pgfpathlineto{\pgfqpoint{3.055879in}{2.150800in}}%
\pgfpathlineto{\pgfqpoint{3.127483in}{2.085516in}}%
\pgfpathlineto{\pgfqpoint{3.211020in}{2.006298in}}%
\pgfpathlineto{\pgfqpoint{3.330360in}{1.889716in}}%
\pgfpathlineto{\pgfqpoint{3.521303in}{1.703116in}}%
\pgfpathlineto{\pgfqpoint{3.604841in}{1.624652in}}%
\pgfpathlineto{\pgfqpoint{3.676445in}{1.560258in}}%
\pgfpathlineto{\pgfqpoint{3.742082in}{1.504200in}}%
\pgfpathlineto{\pgfqpoint{3.801751in}{1.456161in}}%
\pgfpathlineto{\pgfqpoint{3.855454in}{1.415626in}}%
\pgfpathlineto{\pgfqpoint{3.909157in}{1.377911in}}%
\pgfpathlineto{\pgfqpoint{3.956893in}{1.346948in}}%
\pgfpathlineto{\pgfqpoint{4.004629in}{1.318555in}}%
\pgfpathlineto{\pgfqpoint{4.052365in}{1.292872in}}%
\pgfpathlineto{\pgfqpoint{4.094133in}{1.272725in}}%
\pgfpathlineto{\pgfqpoint{4.135902in}{1.254827in}}%
\pgfpathlineto{\pgfqpoint{4.177671in}{1.239247in}}%
\pgfpathlineto{\pgfqpoint{4.219440in}{1.226044in}}%
\pgfpathlineto{\pgfqpoint{4.261209in}{1.215268in}}%
\pgfpathlineto{\pgfqpoint{4.302978in}{1.206960in}}%
\pgfpathlineto{\pgfqpoint{4.344747in}{1.201152in}}%
\pgfpathlineto{\pgfqpoint{4.386515in}{1.197865in}}%
\pgfpathlineto{\pgfqpoint{4.428284in}{1.197113in}}%
\pgfpathlineto{\pgfqpoint{4.470053in}{1.198897in}}%
\pgfpathlineto{\pgfqpoint{4.511822in}{1.203212in}}%
\pgfpathlineto{\pgfqpoint{4.553591in}{1.210041in}}%
\pgfpathlineto{\pgfqpoint{4.595360in}{1.219357in}}%
\pgfpathlineto{\pgfqpoint{4.637129in}{1.231127in}}%
\pgfpathlineto{\pgfqpoint{4.678897in}{1.245304in}}%
\pgfpathlineto{\pgfqpoint{4.720666in}{1.261835in}}%
\pgfpathlineto{\pgfqpoint{4.762435in}{1.280657in}}%
\pgfpathlineto{\pgfqpoint{4.804204in}{1.301699in}}%
\pgfpathlineto{\pgfqpoint{4.851940in}{1.328362in}}%
\pgfpathlineto{\pgfqpoint{4.899676in}{1.357687in}}%
\pgfpathlineto{\pgfqpoint{4.947412in}{1.389529in}}%
\pgfpathlineto{\pgfqpoint{5.001114in}{1.428163in}}%
\pgfpathlineto{\pgfqpoint{5.054817in}{1.469539in}}%
\pgfpathlineto{\pgfqpoint{5.114487in}{1.518414in}}%
\pgfpathlineto{\pgfqpoint{5.180124in}{1.575263in}}%
\pgfpathlineto{\pgfqpoint{5.251727in}{1.640360in}}%
\pgfpathlineto{\pgfqpoint{5.335265in}{1.719422in}}%
\pgfpathlineto{\pgfqpoint{5.448638in}{1.830024in}}%
\pgfpathlineto{\pgfqpoint{5.651515in}{2.028341in}}%
\pgfpathlineto{\pgfqpoint{5.735053in}{2.106758in}}%
\pgfpathlineto{\pgfqpoint{5.806657in}{2.171098in}}%
\pgfpathlineto{\pgfqpoint{5.854392in}{2.212137in}}%
\pgfpathlineto{\pgfqpoint{5.854392in}{2.212137in}}%
\pgfusepath{stroke}%
\end{pgfscope}%
\begin{pgfscope}%
\pgfpathrectangle{\pgfqpoint{0.532060in}{0.456901in}}{\pgfqpoint{5.620383in}{2.816433in}}%
\pgfusepath{clip}%
\pgfsetrectcap%
\pgfsetroundjoin%
\pgfsetlinewidth{1.003750pt}%
\definecolor{currentstroke}{rgb}{0.750000,0.000000,0.750000}%
\pgfsetstrokecolor{currentstroke}%
\pgfsetdash{}{0pt}%
\pgfpathmoveto{\pgfqpoint{0.522060in}{1.806390in}}%
\pgfpathlineto{\pgfqpoint{0.567648in}{1.816764in}}%
\pgfpathlineto{\pgfqpoint{0.621351in}{1.831471in}}%
\pgfpathlineto{\pgfqpoint{0.686988in}{1.851924in}}%
\pgfpathlineto{\pgfqpoint{0.842130in}{1.901406in}}%
\pgfpathlineto{\pgfqpoint{0.895832in}{1.915708in}}%
\pgfpathlineto{\pgfqpoint{0.943568in}{1.926053in}}%
\pgfpathlineto{\pgfqpoint{0.985337in}{1.932894in}}%
\pgfpathlineto{\pgfqpoint{1.027106in}{1.937424in}}%
\pgfpathlineto{\pgfqpoint{1.068875in}{1.939489in}}%
\pgfpathlineto{\pgfqpoint{1.110644in}{1.939018in}}%
\pgfpathlineto{\pgfqpoint{1.152412in}{1.936027in}}%
\pgfpathlineto{\pgfqpoint{1.194181in}{1.930619in}}%
\pgfpathlineto{\pgfqpoint{1.235950in}{1.922977in}}%
\pgfpathlineto{\pgfqpoint{1.283686in}{1.911847in}}%
\pgfpathlineto{\pgfqpoint{1.337389in}{1.896869in}}%
\pgfpathlineto{\pgfqpoint{1.408993in}{1.874280in}}%
\pgfpathlineto{\pgfqpoint{1.546233in}{1.830435in}}%
\pgfpathlineto{\pgfqpoint{1.599936in}{1.815883in}}%
\pgfpathlineto{\pgfqpoint{1.647672in}{1.805250in}}%
\pgfpathlineto{\pgfqpoint{1.689441in}{1.798119in}}%
\pgfpathlineto{\pgfqpoint{1.731210in}{1.793272in}}%
\pgfpathlineto{\pgfqpoint{1.772978in}{1.790874in}}%
\pgfpathlineto{\pgfqpoint{1.814747in}{1.791008in}}%
\pgfpathlineto{\pgfqpoint{1.856516in}{1.793668in}}%
\pgfpathlineto{\pgfqpoint{1.898285in}{1.798765in}}%
\pgfpathlineto{\pgfqpoint{1.940054in}{1.806123in}}%
\pgfpathlineto{\pgfqpoint{1.987790in}{1.816978in}}%
\pgfpathlineto{\pgfqpoint{2.041492in}{1.831722in}}%
\pgfpathlineto{\pgfqpoint{2.107129in}{1.852201in}}%
\pgfpathlineto{\pgfqpoint{2.262271in}{1.901651in}}%
\pgfpathlineto{\pgfqpoint{2.315974in}{1.915914in}}%
\pgfpathlineto{\pgfqpoint{2.363709in}{1.926214in}}%
\pgfpathlineto{\pgfqpoint{2.405478in}{1.933010in}}%
\pgfpathlineto{\pgfqpoint{2.447247in}{1.937491in}}%
\pgfpathlineto{\pgfqpoint{2.489016in}{1.939505in}}%
\pgfpathlineto{\pgfqpoint{2.530785in}{1.938982in}}%
\pgfpathlineto{\pgfqpoint{2.572554in}{1.935940in}}%
\pgfpathlineto{\pgfqpoint{2.614322in}{1.930485in}}%
\pgfpathlineto{\pgfqpoint{2.656091in}{1.922800in}}%
\pgfpathlineto{\pgfqpoint{2.703827in}{1.911628in}}%
\pgfpathlineto{\pgfqpoint{2.757530in}{1.896614in}}%
\pgfpathlineto{\pgfqpoint{2.829134in}{1.874001in}}%
\pgfpathlineto{\pgfqpoint{2.960407in}{1.831936in}}%
\pgfpathlineto{\pgfqpoint{3.014110in}{1.817161in}}%
\pgfpathlineto{\pgfqpoint{3.061846in}{1.806270in}}%
\pgfpathlineto{\pgfqpoint{3.103615in}{1.798874in}}%
\pgfpathlineto{\pgfqpoint{3.145384in}{1.793736in}}%
\pgfpathlineto{\pgfqpoint{3.187153in}{1.791033in}}%
\pgfpathlineto{\pgfqpoint{3.228921in}{1.790855in}}%
\pgfpathlineto{\pgfqpoint{3.270690in}{1.793209in}}%
\pgfpathlineto{\pgfqpoint{3.312459in}{1.798015in}}%
\pgfpathlineto{\pgfqpoint{3.354228in}{1.805108in}}%
\pgfpathlineto{\pgfqpoint{3.401964in}{1.815703in}}%
\pgfpathlineto{\pgfqpoint{3.455667in}{1.830223in}}%
\pgfpathlineto{\pgfqpoint{3.521303in}{1.850543in}}%
\pgfpathlineto{\pgfqpoint{3.688379in}{1.903595in}}%
\pgfpathlineto{\pgfqpoint{3.742082in}{1.917536in}}%
\pgfpathlineto{\pgfqpoint{3.789817in}{1.927473in}}%
\pgfpathlineto{\pgfqpoint{3.831586in}{1.933905in}}%
\pgfpathlineto{\pgfqpoint{3.873355in}{1.937991in}}%
\pgfpathlineto{\pgfqpoint{3.915124in}{1.939593in}}%
\pgfpathlineto{\pgfqpoint{3.956893in}{1.938655in}}%
\pgfpathlineto{\pgfqpoint{3.998662in}{1.935211in}}%
\pgfpathlineto{\pgfqpoint{4.040431in}{1.929376in}}%
\pgfpathlineto{\pgfqpoint{4.082199in}{1.921351in}}%
\pgfpathlineto{\pgfqpoint{4.129935in}{1.909850in}}%
\pgfpathlineto{\pgfqpoint{4.183638in}{1.894562in}}%
\pgfpathlineto{\pgfqpoint{4.261209in}{1.869802in}}%
\pgfpathlineto{\pgfqpoint{4.374581in}{1.833455in}}%
\pgfpathlineto{\pgfqpoint{4.428284in}{1.818464in}}%
\pgfpathlineto{\pgfqpoint{4.476020in}{1.807319in}}%
\pgfpathlineto{\pgfqpoint{4.517789in}{1.799662in}}%
\pgfpathlineto{\pgfqpoint{4.559558in}{1.794237in}}%
\pgfpathlineto{\pgfqpoint{4.601327in}{1.791229in}}%
\pgfpathlineto{\pgfqpoint{4.643096in}{1.790740in}}%
\pgfpathlineto{\pgfqpoint{4.684864in}{1.792786in}}%
\pgfpathlineto{\pgfqpoint{4.726633in}{1.797299in}}%
\pgfpathlineto{\pgfqpoint{4.768402in}{1.804124in}}%
\pgfpathlineto{\pgfqpoint{4.816138in}{1.814454in}}%
\pgfpathlineto{\pgfqpoint{4.863874in}{1.827037in}}%
\pgfpathlineto{\pgfqpoint{4.923544in}{1.845076in}}%
\pgfpathlineto{\pgfqpoint{5.126421in}{1.908752in}}%
\pgfpathlineto{\pgfqpoint{5.174157in}{1.920447in}}%
\pgfpathlineto{\pgfqpoint{5.215926in}{1.928675in}}%
\pgfpathlineto{\pgfqpoint{5.257694in}{1.934737in}}%
\pgfpathlineto{\pgfqpoint{5.299463in}{1.938425in}}%
\pgfpathlineto{\pgfqpoint{5.341232in}{1.939614in}}%
\pgfpathlineto{\pgfqpoint{5.383001in}{1.938262in}}%
\pgfpathlineto{\pgfqpoint{5.424770in}{1.934417in}}%
\pgfpathlineto{\pgfqpoint{5.466539in}{1.928209in}}%
\pgfpathlineto{\pgfqpoint{5.508308in}{1.919850in}}%
\pgfpathlineto{\pgfqpoint{5.556043in}{1.908031in}}%
\pgfpathlineto{\pgfqpoint{5.609746in}{1.892483in}}%
\pgfpathlineto{\pgfqpoint{5.693284in}{1.865588in}}%
\pgfpathlineto{\pgfqpoint{5.794723in}{1.833201in}}%
\pgfpathlineto{\pgfqpoint{5.848425in}{1.818245in}}%
\pgfpathlineto{\pgfqpoint{5.854392in}{1.816732in}}%
\pgfpathlineto{\pgfqpoint{5.854392in}{1.816732in}}%
\pgfusepath{stroke}%
\end{pgfscope}%
\begin{pgfscope}%
\pgfpathrectangle{\pgfqpoint{0.532060in}{0.456901in}}{\pgfqpoint{5.620383in}{2.816433in}}%
\pgfusepath{clip}%
\pgfsetrectcap%
\pgfsetroundjoin%
\pgfsetlinewidth{1.003750pt}%
\definecolor{currentstroke}{rgb}{0.750000,0.750000,0.000000}%
\pgfsetstrokecolor{currentstroke}%
\pgfsetdash{}{0pt}%
\pgfpathmoveto{\pgfqpoint{0.522060in}{1.816897in}}%
\pgfpathlineto{\pgfqpoint{0.549748in}{1.811751in}}%
\pgfpathlineto{\pgfqpoint{0.579582in}{1.808706in}}%
\pgfpathlineto{\pgfqpoint{0.609417in}{1.808384in}}%
\pgfpathlineto{\pgfqpoint{0.639252in}{1.810800in}}%
\pgfpathlineto{\pgfqpoint{0.669087in}{1.815838in}}%
\pgfpathlineto{\pgfqpoint{0.698922in}{1.823253in}}%
\pgfpathlineto{\pgfqpoint{0.734724in}{1.834781in}}%
\pgfpathlineto{\pgfqpoint{0.776493in}{1.850824in}}%
\pgfpathlineto{\pgfqpoint{0.895832in}{1.898511in}}%
\pgfpathlineto{\pgfqpoint{0.931634in}{1.909393in}}%
\pgfpathlineto{\pgfqpoint{0.961469in}{1.916141in}}%
\pgfpathlineto{\pgfqpoint{0.991304in}{1.920426in}}%
\pgfpathlineto{\pgfqpoint{1.021139in}{1.922041in}}%
\pgfpathlineto{\pgfqpoint{1.050974in}{1.920910in}}%
\pgfpathlineto{\pgfqpoint{1.080809in}{1.917085in}}%
\pgfpathlineto{\pgfqpoint{1.110644in}{1.910753in}}%
\pgfpathlineto{\pgfqpoint{1.140479in}{1.902218in}}%
\pgfpathlineto{\pgfqpoint{1.176280in}{1.889655in}}%
\pgfpathlineto{\pgfqpoint{1.229983in}{1.867925in}}%
\pgfpathlineto{\pgfqpoint{1.301587in}{1.838878in}}%
\pgfpathlineto{\pgfqpoint{1.337389in}{1.826595in}}%
\pgfpathlineto{\pgfqpoint{1.367224in}{1.818370in}}%
\pgfpathlineto{\pgfqpoint{1.397059in}{1.812401in}}%
\pgfpathlineto{\pgfqpoint{1.426894in}{1.808977in}}%
\pgfpathlineto{\pgfqpoint{1.456728in}{1.808262in}}%
\pgfpathlineto{\pgfqpoint{1.486563in}{1.810290in}}%
\pgfpathlineto{\pgfqpoint{1.516398in}{1.814965in}}%
\pgfpathlineto{\pgfqpoint{1.546233in}{1.822061in}}%
\pgfpathlineto{\pgfqpoint{1.582035in}{1.833281in}}%
\pgfpathlineto{\pgfqpoint{1.623804in}{1.849098in}}%
\pgfpathlineto{\pgfqpoint{1.749110in}{1.899089in}}%
\pgfpathlineto{\pgfqpoint{1.784912in}{1.909840in}}%
\pgfpathlineto{\pgfqpoint{1.814747in}{1.916454in}}%
\pgfpathlineto{\pgfqpoint{1.844582in}{1.920591in}}%
\pgfpathlineto{\pgfqpoint{1.874417in}{1.922051in}}%
\pgfpathlineto{\pgfqpoint{1.904252in}{1.920763in}}%
\pgfpathlineto{\pgfqpoint{1.934087in}{1.916789in}}%
\pgfpathlineto{\pgfqpoint{1.963922in}{1.910321in}}%
\pgfpathlineto{\pgfqpoint{1.993757in}{1.901672in}}%
\pgfpathlineto{\pgfqpoint{2.029558in}{1.889007in}}%
\pgfpathlineto{\pgfqpoint{2.083261in}{1.867210in}}%
\pgfpathlineto{\pgfqpoint{2.154865in}{1.838244in}}%
\pgfpathlineto{\pgfqpoint{2.190667in}{1.826071in}}%
\pgfpathlineto{\pgfqpoint{2.220502in}{1.817965in}}%
\pgfpathlineto{\pgfqpoint{2.250337in}{1.812135in}}%
\pgfpathlineto{\pgfqpoint{2.280172in}{1.808862in}}%
\pgfpathlineto{\pgfqpoint{2.310007in}{1.808304in}}%
\pgfpathlineto{\pgfqpoint{2.339841in}{1.810488in}}%
\pgfpathlineto{\pgfqpoint{2.369676in}{1.815308in}}%
\pgfpathlineto{\pgfqpoint{2.399511in}{1.822533in}}%
\pgfpathlineto{\pgfqpoint{2.435313in}{1.833877in}}%
\pgfpathlineto{\pgfqpoint{2.477082in}{1.849787in}}%
\pgfpathlineto{\pgfqpoint{2.602389in}{1.899661in}}%
\pgfpathlineto{\pgfqpoint{2.638190in}{1.910279in}}%
\pgfpathlineto{\pgfqpoint{2.668025in}{1.916760in}}%
\pgfpathlineto{\pgfqpoint{2.697860in}{1.920748in}}%
\pgfpathlineto{\pgfqpoint{2.727695in}{1.922051in}}%
\pgfpathlineto{\pgfqpoint{2.757530in}{1.920607in}}%
\pgfpathlineto{\pgfqpoint{2.787365in}{1.916484in}}%
\pgfpathlineto{\pgfqpoint{2.817200in}{1.909882in}}%
\pgfpathlineto{\pgfqpoint{2.847035in}{1.901120in}}%
\pgfpathlineto{\pgfqpoint{2.882837in}{1.888355in}}%
\pgfpathlineto{\pgfqpoint{2.936539in}{1.866494in}}%
\pgfpathlineto{\pgfqpoint{3.008143in}{1.837615in}}%
\pgfpathlineto{\pgfqpoint{3.043945in}{1.825553in}}%
\pgfpathlineto{\pgfqpoint{3.073780in}{1.817567in}}%
\pgfpathlineto{\pgfqpoint{3.103615in}{1.811877in}}%
\pgfpathlineto{\pgfqpoint{3.133450in}{1.808756in}}%
\pgfpathlineto{\pgfqpoint{3.163285in}{1.808355in}}%
\pgfpathlineto{\pgfqpoint{3.193120in}{1.810694in}}%
\pgfpathlineto{\pgfqpoint{3.222954in}{1.815659in}}%
\pgfpathlineto{\pgfqpoint{3.252789in}{1.823011in}}%
\pgfpathlineto{\pgfqpoint{3.288591in}{1.834478in}}%
\pgfpathlineto{\pgfqpoint{3.330360in}{1.850478in}}%
\pgfpathlineto{\pgfqpoint{3.449700in}{1.898221in}}%
\pgfpathlineto{\pgfqpoint{3.485502in}{1.909167in}}%
\pgfpathlineto{\pgfqpoint{3.515336in}{1.915981in}}%
\pgfpathlineto{\pgfqpoint{3.545171in}{1.920339in}}%
\pgfpathlineto{\pgfqpoint{3.575006in}{1.922033in}}%
\pgfpathlineto{\pgfqpoint{3.604841in}{1.920980in}}%
\pgfpathlineto{\pgfqpoint{3.634676in}{1.917231in}}%
\pgfpathlineto{\pgfqpoint{3.664511in}{1.910966in}}%
\pgfpathlineto{\pgfqpoint{3.694346in}{1.902489in}}%
\pgfpathlineto{\pgfqpoint{3.730148in}{1.889978in}}%
\pgfpathlineto{\pgfqpoint{3.783851in}{1.868283in}}%
\pgfpathlineto{\pgfqpoint{3.855454in}{1.839196in}}%
\pgfpathlineto{\pgfqpoint{3.891256in}{1.826860in}}%
\pgfpathlineto{\pgfqpoint{3.921091in}{1.818576in}}%
\pgfpathlineto{\pgfqpoint{3.950926in}{1.812538in}}%
\pgfpathlineto{\pgfqpoint{3.980761in}{1.809038in}}%
\pgfpathlineto{\pgfqpoint{4.010596in}{1.808244in}}%
\pgfpathlineto{\pgfqpoint{4.040431in}{1.810195in}}%
\pgfpathlineto{\pgfqpoint{4.070266in}{1.814797in}}%
\pgfpathlineto{\pgfqpoint{4.100100in}{1.821827in}}%
\pgfpathlineto{\pgfqpoint{4.135902in}{1.832984in}}%
\pgfpathlineto{\pgfqpoint{4.177671in}{1.848755in}}%
\pgfpathlineto{\pgfqpoint{4.308945in}{1.900788in}}%
\pgfpathlineto{\pgfqpoint{4.344747in}{1.911137in}}%
\pgfpathlineto{\pgfqpoint{4.374581in}{1.917347in}}%
\pgfpathlineto{\pgfqpoint{4.404416in}{1.921035in}}%
\pgfpathlineto{\pgfqpoint{4.434251in}{1.922025in}}%
\pgfpathlineto{\pgfqpoint{4.464086in}{1.920268in}}%
\pgfpathlineto{\pgfqpoint{4.493921in}{1.915850in}}%
\pgfpathlineto{\pgfqpoint{4.523756in}{1.908983in}}%
\pgfpathlineto{\pgfqpoint{4.553591in}{1.899999in}}%
\pgfpathlineto{\pgfqpoint{4.589393in}{1.887040in}}%
\pgfpathlineto{\pgfqpoint{4.649063in}{1.862556in}}%
\pgfpathlineto{\pgfqpoint{4.708732in}{1.838560in}}%
\pgfpathlineto{\pgfqpoint{4.744534in}{1.826332in}}%
\pgfpathlineto{\pgfqpoint{4.774369in}{1.818167in}}%
\pgfpathlineto{\pgfqpoint{4.804204in}{1.812267in}}%
\pgfpathlineto{\pgfqpoint{4.834039in}{1.808918in}}%
\pgfpathlineto{\pgfqpoint{4.863874in}{1.808282in}}%
\pgfpathlineto{\pgfqpoint{4.893709in}{1.810388in}}%
\pgfpathlineto{\pgfqpoint{4.923544in}{1.815136in}}%
\pgfpathlineto{\pgfqpoint{4.953379in}{1.822296in}}%
\pgfpathlineto{\pgfqpoint{4.989180in}{1.833578in}}%
\pgfpathlineto{\pgfqpoint{5.030949in}{1.849442in}}%
\pgfpathlineto{\pgfqpoint{5.156256in}{1.899375in}}%
\pgfpathlineto{\pgfqpoint{5.192058in}{1.910060in}}%
\pgfpathlineto{\pgfqpoint{5.221893in}{1.916608in}}%
\pgfpathlineto{\pgfqpoint{5.251727in}{1.920671in}}%
\pgfpathlineto{\pgfqpoint{5.281562in}{1.922052in}}%
\pgfpathlineto{\pgfqpoint{5.311397in}{1.920686in}}%
\pgfpathlineto{\pgfqpoint{5.341232in}{1.916637in}}%
\pgfpathlineto{\pgfqpoint{5.371067in}{1.910103in}}%
\pgfpathlineto{\pgfqpoint{5.400902in}{1.901397in}}%
\pgfpathlineto{\pgfqpoint{5.436704in}{1.888681in}}%
\pgfpathlineto{\pgfqpoint{5.490407in}{1.866852in}}%
\pgfpathlineto{\pgfqpoint{5.562010in}{1.837929in}}%
\pgfpathlineto{\pgfqpoint{5.597812in}{1.825811in}}%
\pgfpathlineto{\pgfqpoint{5.627647in}{1.817765in}}%
\pgfpathlineto{\pgfqpoint{5.657482in}{1.812005in}}%
\pgfpathlineto{\pgfqpoint{5.687317in}{1.808808in}}%
\pgfpathlineto{\pgfqpoint{5.717152in}{1.808328in}}%
\pgfpathlineto{\pgfqpoint{5.746987in}{1.810590in}}%
\pgfpathlineto{\pgfqpoint{5.776822in}{1.815483in}}%
\pgfpathlineto{\pgfqpoint{5.806657in}{1.822771in}}%
\pgfpathlineto{\pgfqpoint{5.842458in}{1.834177in}}%
\pgfpathlineto{\pgfqpoint{5.854392in}{1.838500in}}%
\pgfpathlineto{\pgfqpoint{5.854392in}{1.838500in}}%
\pgfusepath{stroke}%
\end{pgfscope}%
\begin{pgfscope}%
\pgfsetrectcap%
\pgfsetmiterjoin%
\pgfsetlinewidth{0.602250pt}%
\definecolor{currentstroke}{rgb}{0.000000,0.000000,0.000000}%
\pgfsetstrokecolor{currentstroke}%
\pgfsetdash{}{0pt}%
\pgfpathmoveto{\pgfqpoint{0.532060in}{0.456901in}}%
\pgfpathlineto{\pgfqpoint{0.532060in}{3.273333in}}%
\pgfusepath{stroke}%
\end{pgfscope}%
\begin{pgfscope}%
\pgfsetrectcap%
\pgfsetmiterjoin%
\pgfsetlinewidth{0.602250pt}%
\definecolor{currentstroke}{rgb}{0.000000,0.000000,0.000000}%
\pgfsetstrokecolor{currentstroke}%
\pgfsetdash{}{0pt}%
\pgfpathmoveto{\pgfqpoint{0.532060in}{0.456901in}}%
\pgfpathlineto{\pgfqpoint{6.152443in}{0.456901in}}%
\pgfusepath{stroke}%
\end{pgfscope}%
\begin{pgfscope}%
\pgfsetbuttcap%
\pgfsetroundjoin%
\pgfsetlinewidth{2.007500pt}%
\definecolor{currentstroke}{rgb}{0.000000,0.000000,0.000000}%
\pgfsetstrokecolor{currentstroke}%
\pgfsetdash{{7.400000pt}{3.200000pt}}{0.000000pt}%
\pgfpathmoveto{\pgfqpoint{0.555736in}{3.565000in}}%
\pgfpathlineto{\pgfqpoint{0.955736in}{3.565000in}}%
\pgfusepath{stroke}%
\end{pgfscope}%
\begin{pgfscope}%
\definecolor{textcolor}{rgb}{0.000000,0.000000,0.000000}%
\pgfsetstrokecolor{textcolor}%
\pgfsetfillcolor{textcolor}%
\pgftext[x=1.139069in,y=3.506667in,left,base]{\color{textcolor}{\rmfamily\fontsize{12.000000}{14.400000}\selectfont\catcode`\^=\active\def^{\ifmmode\sp\else\^{}\fi}\catcode`\%=\active\def%{\%}Исходный сигнал}}%
\end{pgfscope}%
\begin{pgfscope}%
\pgfsetrectcap%
\pgfsetroundjoin%
\pgfsetlinewidth{2.007500pt}%
\definecolor{currentstroke}{rgb}{0.835294,0.368627,0.000000}%
\pgfsetstrokecolor{currentstroke}%
\pgfsetdash{}{0pt}%
\pgfpathmoveto{\pgfqpoint{2.766218in}{3.565000in}}%
\pgfpathlineto{\pgfqpoint{3.166218in}{3.565000in}}%
\pgfusepath{stroke}%
\end{pgfscope}%
\begin{pgfscope}%
\definecolor{textcolor}{rgb}{0.000000,0.000000,0.000000}%
\pgfsetstrokecolor{textcolor}%
\pgfsetfillcolor{textcolor}%
\pgftext[x=3.349551in,y=3.506667in,left,base]{\color{textcolor}{\rmfamily\fontsize{12.000000}{14.400000}\selectfont\catcode`\^=\active\def^{\ifmmode\sp\else\^{}\fi}\catcode`\%=\active\def%{\%}Входной РФ}}%
\end{pgfscope}%
\begin{pgfscope}%
\pgfsetbuttcap%
\pgfsetroundjoin%
\pgfsetlinewidth{2.007500pt}%
\definecolor{currentstroke}{rgb}{0.000000,0.619608,0.450980}%
\pgfsetstrokecolor{currentstroke}%
\pgfsetdash{{12.800000pt}{3.200000pt}{2.000000pt}{3.200000pt}}{0.000000pt}%
\pgfpathmoveto{\pgfqpoint{4.607071in}{3.565000in}}%
\pgfpathlineto{\pgfqpoint{5.007071in}{3.565000in}}%
\pgfusepath{stroke}%
\end{pgfscope}%
\begin{pgfscope}%
\definecolor{textcolor}{rgb}{0.000000,0.000000,0.000000}%
\pgfsetstrokecolor{textcolor}%
\pgfsetfillcolor{textcolor}%
\pgftext[x=5.190404in,y=3.506667in,left,base]{\color{textcolor}{\rmfamily\fontsize{12.000000}{14.400000}\selectfont\catcode`\^=\active\def^{\ifmmode\sp\else\^{}\fi}\catcode`\%=\active\def%{\%}Выходной РФ}}%
\end{pgfscope}%
\end{pgfpicture}%
\makeatother%
\endgroup%

    \caption{Ряд Фурье реакции на второй периодический сигнал}
    \label{fig:plot_fourier_approx_out_2}
\end{figure}